\documentclass[11pt,twoside]{book}

%  Double-sided layout: chapters open on a recto (right) page, and the inner
%  (spine) margin is wider than the outer to leave room for binding.
\usepackage[a4paper,twoside,inner=1.25in,outer=0.85in,top=1in,bottom=1in]{geometry}
\usepackage[T1]{fontenc}
\usepackage{textcomp}
%  Body text + math in a Times New Roman equivalent (URW Nimbus Roman, the
%  metric-compatible Times clone); works under pdflatex.  Code listings keep a
%  monospace font.
\usepackage{mathptmx}
\usepackage{graphicx}
\usepackage{listings}
\usepackage{xcolor}
\usepackage{tocloft}
%  Two-digit chapter numbers make section numbers like "15.10" wider than the
%  default ToC number box, so they collide with the heading.  Widen the number
%  boxes (and shift the matching indents) to leave a clear gap.
\setlength{\cftchapnumwidth}{2.0em}
\setlength{\cftsecindent}{2.0em}
\setlength{\cftsecnumwidth}{3.2em}
\setlength{\cftsubsecindent}{5.2em}
\setlength{\cftsubsecnumwidth}{4.0em}
\usepackage{hyperref}
\usepackage{titlesec}
\usepackage{parskip}
\usepackage{enumitem}
\usepackage{longtable}
\usepackage{tikz}
\usetikzlibrary{arrows.meta, positioning, fit, calc}
\usepackage{makeidx}
\makeindex

%  Long inline-code wrapping (so fully-qualified Ada names do not run into the
%  margin).  (3) microtype tightens spacing and emergencystretch lets TeX relax
%  a line rather than overrun.  (1) \texttt is taught to break after the natural
%  separators of a qualified name --- '.', '_' and '/'.  '_' is the macro \_ in
%  our source, handled by redefinition; '.' and '/' are ordinary characters, so
%  \texttt activates them before reading its argument (this works for the normal
%  \texttt{...} call; where the text was already tokenised, e.g. a section title,
%  it simply adds no extra break point).
\usepackage{microtype}
\setlength{\emergencystretch}{3em}
\makeatletter
\newcommand*\bm@dbrk{\discretionary{}{}{}}
\begingroup
  \catcode`\.\active \catcode`\/\active
  \gdef\bm@ttbreaks{\catcode`\.\active \catcode`\/\active}%
  \gdef.{\char`\.\bm@dbrk}%
  \gdef/{\char`\/\bm@dbrk}%
\endgroup
\let\bm@origtexttt\texttt
\protected\def\texttt{\begingroup\bm@ttbreaks\bm@grabtt}
\def\bm@grabtt#1{\bm@origtexttt{\def\_{\textunderscore\bm@dbrk}#1}\endgroup}
\makeatother

\hypersetup{colorlinks=true, linkcolor=blue!50!black, urlcolor=blue!60!black,
            pdftitle={Bare-Metal Ada on the ESP32-S3},
            pdfauthor={The ada_esp32s3 project}}

\definecolor{kw}{rgb}{0.10,0.10,0.55}
\definecolor{cmt}{rgb}{0.30,0.45,0.30}
\definecolor{str}{rgb}{0.55,0.10,0.10}
\definecolor{bg}{gray}{0.97}

\lstset{
  language=Ada,
  basicstyle=\ttfamily\small,
  keywordstyle=\color{kw}\bfseries,
  commentstyle=\color{cmt}\itshape,
  stringstyle=\color{str},
  showstringspaces=false,
  frame=single,
  framesep=4pt,
  rulecolor=\color{gray!40},
  backgroundcolor=\color{bg},
  breaklines=true,
  columns=fullflexible,
  keepspaces=true,
  numbers=none,
  morekeywords={overriding,aliased,synchronized,interface,limited,protected,
                requeue,tagged,until,with}
}
% A non-Ada style for shell snippets.
\lstdefinestyle{sh}{language=bash, basicstyle=\ttfamily\small,
  keywordstyle=\color{kw}, frame=single, backgroundcolor=\color{bg},
  breaklines=true, showstringspaces=false}

\titleformat{\chapter}[display]{\normalfont\huge\bfseries}{\chaptertitlename\ \thechapter}{16pt}{\Huge}

%  A short "if you come from C / ESP-IDF" aside, set off in a light frame.
\newcommand{\cnote}[1]{%
  \par\smallskip\noindent
  \fbox{\parbox{\dimexpr\linewidth-2\fboxsep-2\fboxrule\relax}{%
    \small\textbf{From C / ESP-IDF:}\ #1}}%
  \par\smallskip}

%  Common TikZ box styles for the diagrams.
\tikzset{
  blk/.style   = {draw, rounded corners=2pt, align=center, font=\small,
                  minimum height=7mm, inner sep=4pt, fill=blue!4},
  layer/.style = {draw, align=center, font=\small, minimum height=8mm,
                  text width=10.5cm, fill=blue!4},
  flow/.style  = {-{Stealth[length=2.4mm]}, thick},
}

\begin{document}

\frontmatter
\begin{titlepage}
  \centering
  \vspace*{1.5cm}
  \includegraphics[height=4.5cm]{Ada.pdf}\par
  \vspace{1.2cm}
  {\Huge\bfseries Bare-Metal Ada on the ESP32-S3\par}
  \vspace{1cm}
  {\Large A practical guide to the runtime, the peripheral HAL,\\
   and a pure-Ada ext4 filesystem\par}
  \vspace{2cm}
  {\large Written in Ada, with a little Xtensa assembly.\par}
  \vfill
  {\large\ttfamily github.com/rowsail/ada\_esp32s3\par}
\end{titlepage}

\tableofcontents

\mainmatter

\chapter{Why Ada for Microcontrollers}

Microcontroller firmware lives in a hostile place: a few hundred kilobytes of
RAM, no operating system to catch mistakes, interrupts that fire between any two
instructions, and hardware registers whose exact bit layout is law. A bug is not
a stack trace on a screen. It is a board that silently locks up in the field.
Ada was designed for exactly this kind of software, and the language gives the
firmware author tools that C simply does not have.

\section{The case for Ada}\index{type safety}\index{representation clause}\index{memory-mapped I/O}\index{private child}\index{contracts}

\begin{description}[style=nextline]
\item[Strong, range-checked types.] In Ada a ``GPIO pin'' need not be an
  \texttt{int}. It can be a type that admits \emph{only} the pins the silicon
  actually exposes; passing a flash-mapped pad becomes a compile error, not a
  hung chip. Numeric types carry ranges, and the compiler (or a run-time check)
  enforces them.
\item[Representation control.] Ada lets you state, in the language, the precise
  bit and byte layout of a record and the exact address of an object. A hardware
  register is just an Ada record at a fixed address: no casts, no macros, no
  guessing about padding (Sections~\ref{ch:registers} and~\ref{ch:bitfields} of
  the Talking to Hardware chapter).
\item[Packages and strong interfaces.] Drivers are packages with a visible
  specification and a hidden body. The unsafe register pokes can be sealed inside
  a \emph{private child} that application code literally cannot name.
\item[Tasking and real-time in the language.] Tasks, protected objects (data
  guarded by a monitor), \texttt{delay until}, and priorities are part of Ada
  itself, not a bolted-on RTOS API. On this port they are implemented by a
  native Ada runtime that runs directly on the silicon.
\item[Restriction profiles.] With one pragma you can forbid features that do not
  belong in certified or tiny firmware (no heap, no recursion in the runtime, no
  unbounded constructs) and have the compiler \emph{prove} you stayed within the
  subset -- the Ravenscar and Jorvik profiles (Chapter~\ref{ch:profiles}).
\item[Contracts.] Preconditions, postconditions and type predicates let you
  state intent (``this key must be 128 or 256 bits'') and have it checked.
\end{description}

The result is firmware where a large class of mistakes (wrong-width register
writes, out-of-range pins, aliasing of a peripheral by two tasks, forgetting to
release a DMA channel) are caught by the compiler or by a cheap run-time
check, long before they reach silicon.

\section{What this book covers}

This book documents a complete bare-metal stack for the dual-core ESP32-S3
(Xtensa LX7) written almost entirely in Ada, running directly on the silicon
with its own runtime, boot path and tooling:

\begin{enumerate}
\item the Ada language techniques that make register and bit-field programming
  safe and pleasant (the Ada Language Techniques and Talking to Hardware
  chapters, Chapters~\ref{ch:embedded} and~\ref{ch:hardware});
\item the anatomy of an Ada application, from \texttt{Main} down through the
  runtime to the boot ROM, and the three selectable runtime profiles
  (Chapters~\ref{ch:anatomy}--\ref{ch:build});
\item the reusable peripheral HAL (twenty-plus task-safe drivers) with the
  API and worked examples for each (Chapters~\ref{ch:hal}--\ref{ch:hal3});
\item mass storage: the SD-card block drivers and the pure-Ada ext2/3/4
  filesystem (Chapter~\ref{ch:storage}).
\end{enumerate}

New to Ada? Chapter~\ref{ch:ada-primer}, \emph{Ada for C Programmers}, is a
one-sitting crash course in the slice of the language this book uses
(declarations, subprograms and parameter modes, packages, control flow), framed
against the C you already know; the rest of the book assumes that much. Individual
Ada terms that may still be unfamiliar to a reader coming from C
(\emph{library-level}, \emph{elaboration}, \emph{aspect}, \emph{protected object}
and the like) are also collected in the Terminology section of the Ada Language
Techniques chapter (\S\ref{sec:terms}).


\chapter{Getting Started}
\label{ch:getting-started}

This chapter takes you from a fresh clone to an LED blinking under your own Ada,
in a few commands. The later chapters explain the runtime, the HAL and the
language techniques; this one just gets you running.

\section{What you need}\index{Alire}

\begin{description}[style=nextline]
\item[An ESP32-S3 board.] Any devkit with the built-in USB-Serial-JTAG port (most
  of them). One USB cable is enough -- it carries the console \emph{and} the
  debugger.
\item[Alire (\texttt{alr}).] The Ada package manager fetches the toolchains the
  build needs: the \texttt{gnat\_xtensa\_esp32\_elf} cross-compiler, a native
  \texttt{gnat\_native} (for the host tools), and \texttt{gprbuild}. Install Alire,
  then make sure its toolchains are on your \texttt{PATH} for the shell you build
  from.
\item[Serial-port permission.] On Linux, add yourself to the \texttt{dialout}
  group (and install the udev rule for the USB-Serial-JTAG, see the debugging
  chapter) so you can open \texttt{/dev/ttyACM0} without \texttt{sudo}.
\end{description}

\section{Get the code}

The runtime depends on two git submodules (the \texttt{bb-runtimes} fork that
carries the ESP32-S3 bareboard runtime, and \texttt{xtensa-dynconfig}), so clone
\emph{recursively}:

\begin{lstlisting}[style=sh]
git clone --recurse-submodules \
    https://github.com/rowsail/ada_esp32s3.git
cd ada_esp32s3
#  (if you already cloned without submodules:)
git submodule update --init --recursive
\end{lstlisting}

\section{The \texttt{./x} dispatcher}\index{x helper script}

The repository ships with over sixty ready-to-run examples, and one script at
its root -- \texttt{./x} -- drives all of them \emph{from inside the clone}. (For
your \emph{own} application, kept outside the repo, the counterpart command is
\texttt{esp32-ada}, \S\ref{sec:esp32-ada}; the two share the same verbs.) List the
examples, then build, flash and watch one:

\begin{lstlisting}[style=sh]
./x list                                   # the available examples + their profiles
./x run esp32s3_gpio0_blink -p /dev/ttyACM0   # build + flash + serial monitor
\end{lstlisting}

\texttt{./x run} builds the example, packages it (the Ada \texttt{esp\_elf2image}
host tool), writes it to the board (the Ada \texttt{esp\_flash} tool, over the
chip's ROM bootloader), resets, and streams the console. The first build also
generates the runtime for the chosen profile, so it takes a little longer; later
ones are incremental. The blink example toggles an LED and prints over the
console -- your first Ada running bare-metal on the S3.

The pieces are also separate commands when you want them:

\begin{lstlisting}[style=sh]
./x build   esp32s3_gpio0_blink            # cross-compile + link + package app.bin
./x flash   esp32s3_gpio0_blink -p /dev/ttyACM0
./x monitor -p /dev/ttyACM0                # just the serial console (115200)
./x clean   esp32s3_gpio0_blink
\end{lstlisting}

The examples are meant to be \emph{read}, not just run: each opens with a header
that says what it demonstrates, how to build and run it, what the console should
print, and what hardware (if any) it needs; magic numbers are named, and the
``why'' is in the code rather than only in this book. A short house style at
\texttt{examples/STYLE.md} records the bar, with \texttt{esp32s3\_gpio0\_blink}
and \texttt{esp32s3\_gdma\_copy} as the models to copy; \texttt{./x new} (and
\texttt{esp32-ada init}, \S\ref{sec:esp32-ada}) scaffold a new project already in
that shape. So once a chapter here points you at an example, the example itself
carries the rest of the explanation.

\section{Choosing a profile}\index{runtime profile}\index{light-tasking profile}\index{embedded profile}\index{full profile}

Most examples build under more than one runtime profile (Chapter~\ref{ch:profiles}).
Pass \texttt{--profile} to pick one; \texttt{auto} (the default) uses the
example's own:

\begin{lstlisting}[style=sh]
./x run esp32s3_heartbeat --profile embedded -p /dev/ttyACM0
\end{lstlisting}

\texttt{light-tasking} is the smallest, \texttt{embedded} adds exceptions, RAII
handles, a heap and the full HAL (the usual choice), and \texttt{full} adds the
complete Ada tasking model.

\section{Board configuration}\index{board.ads}

Each project owns a \texttt{board.ads} (flash and PSRAM sizes) that sizes the image
and the second-stage bootloader. Edit it directly, or:

\begin{lstlisting}[style=sh]
./x config esp32s3_gpio0_blink show
./x config esp32s3_gpio0_blink flash-size 4MB
\end{lstlisting}

\section{Debugging}

On-chip debugging needs the pinned OpenOCD + S3 GDB, fetched once; then
\texttt{./x debug} (or \texttt{F5} in VS Code) halts at the boot entry
\texttt{app\_main} (the shared boot glue; these images have no C \texttt{main}):

\begin{lstlisting}[style=sh]
./x get-debug-tools
./x debug esp32s3_gpio0_blink              # build, flash, OpenOCD + GDB, break at app_main
\end{lstlisting}

The Debugging chapter (Chapter~\ref{ch:debug}) covers GDB, decoding a Guru
Meditation, and the editor integration in full.

\section{Your own project, out of tree: \texttt{esp32-ada}}\index{esp32-ada launcher}\index{scaffolding}
\label{sec:esp32-ada}

\texttt{./x} is for the examples that ship \emph{inside} the clone. For your own
application you do not edit an example or copy the runtime into your tree: you
treat this repository as an \emph{SDK} and scaffold a standalone project anywhere
on disk. The launcher that does this -- the out-of-tree counterpart of
\texttt{./x}, with the same verbs -- is \texttt{esp32-ada}.

It becomes available when you \emph{source} the SDK's \texttt{export.sh} once per
shell. That one step puts \texttt{esp32-ada} on your \texttt{PATH}, exports
\texttt{ESP32S3\_ADA\_SDK} (so a project folder anywhere can find the SDK), and
adds the runtime and HAL projects to \texttt{GPR\_PROJECT\_PATH} -- so both
\texttt{gprbuild} \emph{and} the Ada Language Server resolve
\texttt{with "esp32s3\_rts.gpr"} with no path baked into your project:

\begin{lstlisting}[style=sh]
. /path/to/ada_esp32s3/export.sh    # once per shell (add to ~/.bashrc to persist)
mkdir ~/myblink && cd ~/myblink
esp32-ada init                      # scaffold app.gpr, board.ads, src/main.adb
#  ... edit src/main.adb ...
esp32-ada run -p /dev/ttyACM0       # build + flash + monitor
\end{lstlisting}

Every \texttt{esp32-ada} verb mirrors the matching \texttt{./x} one; the only
difference is the target. \texttt{./x build esp32s3\_gpio0\_blink} names an example
inside the clone, while \texttt{esp32-ada build} operates on the project folder you
are standing in:

\begin{longtable}{p{0.46\linewidth} p{0.48\linewidth}}
\toprule
\textbf{Command} & \textbf{What it does} \\ \midrule
\endhead
\texttt{esp32-ada init [\textit{dir}]} & Scaffold a project in \textit{dir} (or the current folder). \\
\texttt{esp32-ada build [-P \textit{prof}]} & Build to \texttt{app.bin}. \textit{prof} = \texttt{light-tasking} (default), \texttt{embedded}, or \texttt{full}. \\
\texttt{esp32-ada flash [-p \textit{port}]} & Build if needed, then flash over the USB ROM bootloader. \\
\texttt{esp32-ada run [-p \textit{port}] [-P \textit{prof}]} & Build, flash, and open the serial monitor. \\
\texttt{esp32-ada monitor [-p \textit{port}]} & Just the serial console (115200). \\
\texttt{esp32-ada clean} & Remove this project's build artifacts. \\
\texttt{esp32-ada config [show | flash-size \textit{S} | psram-size \textit{S}]} & Show or set this project's flash / PSRAM size (its \texttt{board.ads}). \\
\texttt{esp32-ada debug [-p \textit{port}]} & On-chip debug (OpenOCD + GDB), halting at \texttt{Main}. \\
\texttt{esp32-ada kill-openocd} & Kill every OpenOCD (release captured USB-JTAG ports). \\
\texttt{esp32-ada install-ide} & Install the VS Code extension (committed \texttt{.vsix}). \\
\texttt{esp32-ada install-vim} & Symlink the Vim/Neovim plugin. \\ \midrule
\bottomrule
\end{longtable}

\noindent\texttt{-p}/\texttt{--port} defaults to \texttt{\$ESPPORT} or
\texttt{/dev/ttyACM0}; \texttt{-C \textit{dir}} runs as if from \textit{dir}.
\texttt{init} writes a complete, self-contained project:

\begin{lstlisting}[style=sh]
myblink/
  app.gpr            # withs esp32s3_rts.gpr + esp32s3_hal.gpr (resolved by name)
  board.ads          # this project's flash / PSRAM sizes
  src/main.adb       # your code (procedure Main; boots both cores, then idles)
  build.sh flash.sh  # thin shims into the SDK's shared bare-boot
  .vscode/           # ALS project + build/flash/run tasks driving esp32-ada
\end{lstlisting}

Nothing from the SDK is copied in -- \texttt{app.gpr} references the runtime and
HAL \emph{by name} through \texttt{GPR\_PROJECT\_PATH}, and the build shims call the
SDK's shared bare-boot through \texttt{\$ESP32S3\_ADA\_SDK}. Add another SDK library
the same way: \texttt{with "<name>.gpr";} in \texttt{app.gpr} for anything under the
SDK's \texttt{libs/}. The default profile is \texttt{light-tasking}; pass
\texttt{-P embedded} (or \texttt{full}) to build under another
(Chapter~\ref{ch:profiles}).

\subsection{Taking an example out of tree to edit}\index{examples}

The in-repo examples are wired to their location -- an Alire path-pin in
\texttt{alire.toml}, a \texttt{\textit{name}.gpr}, and build shims relative to
\texttt{examples/common/bare/} -- so copying a whole example directory elsewhere
will not build. Instead, scaffold a fresh project and bring the \emph{sources}
across; the scaffold already supplies the \texttt{app.gpr}, board config and build
shims that make them build:

\begin{lstlisting}[style=sh]
. /path/to/ada_esp32s3/export.sh
esp32-ada init ~/myblink && cd ~/myblink
SDK=$ESP32S3_ADA_SDK
cp "$SDK"/examples/esp32s3_gpio0_blink/src/*.ad? src/   # the Ada sources (incl. main.adb)
esp32-ada build                                         # gpio0_blink uses the default light-tasking
esp32-ada run -p /dev/ttyACM0
\end{lstlisting}

A few things make the copy build cleanly:

\begin{itemize}
\item \textbf{Only \texttt{src/} is yours.} \texttt{esp32-ada init} already wrote
  \texttt{app.gpr} (\texttt{Source\_Dirs => "src"}, \texttt{Main => "main.adb"}),
  so the example's \texttt{src/main.adb} (a \texttt{procedure Main}) and its other
  units drop straight in. Overwrite the scaffold's \texttt{src/main.adb}.
\item \textbf{Match the profile.} Most examples assume \texttt{embedded} (the HAL's
  RAII handles need it); the scaffold defaults to \texttt{light-tasking}. The
  example's profile is shown by \texttt{./x list} (and its \texttt{README}); pass
  it with \texttt{-P}.
\item \textbf{Bring board config and C natives if customized.} If the example sets
  a non-default flash/PSRAM size, copy its \texttt{board.ads} (or run
  \texttt{esp32-ada config}); if it has C natives, copy its \texttt{glue.c}.
\item \textbf{Extra libraries.} If the example \texttt{with}s more of the SDK's
  \texttt{libs/} (beyond the HAL), add the matching \texttt{with "<name>.gpr";} to
  \texttt{app.gpr}.
\end{itemize}

\cnote{There is no fixed ``project location'' and no \texttt{idf.py}: the SDK is
found through \texttt{\$ESP32S3\_ADA\_SDK} (set by \texttt{export.sh}), so your
project folder can live anywhere and stays portable -- check it into its own git
repo, with no SDK paths baked into any file.}

\section{Where to go next}

\begin{itemize}
\item The anatomy of an application and the runtime profiles --
  Chapters~\ref{ch:anatomy}--\ref{ch:build}.
\item The peripheral HAL -- start at Chapter~\ref{ch:hal}; each driver has worked
  examples you can run with \texttt{./x run}.
\item The Ada techniques behind it all -- the Ada Language Techniques and
  Talking to Hardware chapters (Chapters~\ref{ch:embedded}
  and~\ref{ch:hardware}), including a Terminology section (\S\ref{sec:terms})
  for unfamiliar terms.
\end{itemize}

To start your \emph{own} project rather than run an in-repo example, use the
\texttt{esp32-ada} launcher (\S\ref{sec:esp32-ada}), including the recipe for
lifting an example's sources out of the tree to build on as your own.

\chapter{Anatomy of an Ada Application}
\label{ch:anatomy}

An Ada program for the ESP32-S3 is not a flat blob. It is a stack of layers, each
written in Ada (with a thin sliver of Xtensa assembly at the very bottom). Knowing
the layers explains what links into your binary, what runs before \texttt{Main},
and where each feature comes from.

\section{The layers, top to bottom}

\begin{description}[style=nextline]
\item[Your application.] One or more Ada packages and a main procedure
  (\texttt{Main}). It \texttt{with}s the peripheral HAL (\texttt{ESP32S3.*}) and,
  if it uses tasks or protected objects, those language features directly.
\item[The HAL (\texttt{libs/esp32s3\_hal}).] Hand-written driver packages over a
  generated register layer (\texttt{ESP32S3\_Registers.*}). Optional --- a program
  can poke registers itself --- but it is where all the device knowledge lives.
\item[The GNAT binder output.] \texttt{gnatbind} generates a package
  (\texttt{b\_\_main}) that elaborates every unit in the right order and calls
  \texttt{Main}. Its \texttt{adainit}/\texttt{adafinal} bracket your program.
\item[The Ada runtime (RTS).] The \texttt{System.*} and \texttt{Ada.*} bodies:
  secondary stack, exceptions (if enabled), \texttt{Ada.Text\_IO} over the
  USB-serial console, and --- crucially --- the tasking kernel (GNARL/GNARLI):
  tasks, protected objects, delays, scheduling.
\item[Board support.] The silicon-specific bottom of the runtime: the Xtensa
  windowed \emph{context switch}, \emph{interrupt entry/exit} vectors, the
  \emph{clock/alarm tick}, SMP scheduling and the inter-core interrupt. This is
  Ada plus a little assembly (\texttt{start.S}, \texttt{highint5.S},
  \texttt{context.S}). This is the layer an off-the-shelf RTOS would otherwise
  provide --- here it is ours.
\item[The 2nd-stage bootloader.] Before any of the above runs, our own minimal
  loader (a ZFP-Ada \texttt{boot\_main} + an assembly prologue + a thin C PSRAM
  shim) maps the flash cache/MMU and brings up the octal PSRAM, then jumps to the
  application image. It replaces the chip vendor's stock second-stage bootloader.
\item[The mask ROM.] The chip's on-chip ROM loads the 2nd-stage bootloader from
  flash at reset.
\end{description}

\begin{figure}[ht]
\centering
\begin{tikzpicture}[node distance=2mm]
\node[layer, fill=blue!9] (app)  {Your application --- \texttt{Main} + packages};
\node[layer, below=of app]  (hal)  {Peripheral HAL: \texttt{ESP32S3.*} over generated \texttt{ESP32S3\_Registers.*}};
\node[layer, below=of hal]  (bind) {GNAT binder --- \texttt{adainit} / \texttt{adafinal}, calls \texttt{Main}};
\node[layer, below=of bind] (rts)  {Ada runtime: \texttt{System.*}, \texttt{Ada.*}, tasking kernel (GNARL)};
\node[layer, below=of rts]  (bsp)  {Board support: context switch, interrupt vectors, tick, SMP\, (\texttt{*.S})};
\node[layer, below=of bsp]  (boot) {2nd-stage bootloader --- cache/MMU, octal PSRAM};
\node[layer, below=of boot, fill=gray!12] (rom) {Mask ROM --- loads the bootloader at reset};
\end{tikzpicture}
\caption{The layers of an Ada application, from your code down to the chip's ROM.
Everything above the ROM is Ada (with a little Xtensa assembly in board support).}
\label{fig:layers}
\end{figure}

\section{What runs before \texttt{Main}}

At power-on: ROM $\rightarrow$ our 2nd-stage bootloader (cache/MMU + PSRAM)
$\rightarrow$ \texttt{start.S} (selects the 240\,MHz PLL, sets up the stack and
the interrupt vectors, and on the SMP examples releases the second core)
$\rightarrow$ the runtime's startup, which calls \texttt{adainit} (elaboration)
$\rightarrow$ your \texttt{Main}. There is no vendor startup shim and no
third-party RTOS scheduler --- the Ada runtime is the only scheduler, and it is
linked from source.

\begin{figure}[ht]
\centering
\resizebox{\linewidth}{!}{%
\begin{tikzpicture}[node distance=7mm]
\node[blk] (rom)   {ROM};
\node[blk, right=of rom]   (boot)  {2nd-stage\\bootloader\\\scriptsize cache, PSRAM};
\node[blk, right=of boot]  (start) {\texttt{start.S}\\\scriptsize PLL, vectors,\\\scriptsize core 1};
\node[blk, right=of start] (init)  {\texttt{adainit}\\\scriptsize elaboration};
\node[blk, right=of init]  (main)  {\texttt{Main}};
\draw[flow] (rom)--(boot);   \draw[flow] (boot)--(start);
\draw[flow] (start)--(init); \draw[flow] (init)--(main);
\end{tikzpicture}}
\caption{Reset to \texttt{Main}.}
\label{fig:boot}
\end{figure}

\section{A minimal application}

\begin{lstlisting}
with Ada.Real_Time; use Ada.Real_Time;
with ESP32S3.GPIO;

procedure Main is
   Led : constant ESP32S3.GPIO.Pin_Id := 2;
begin
   ESP32S3.GPIO.Configure (Led, Mode => ESP32S3.GPIO.Output);
   loop
      ESP32S3.GPIO.Toggle (Led);
      delay until Clock + Milliseconds (250);
   end loop;
end Main;
\end{lstlisting}

\texttt{delay until} is served by the board-support tick; \texttt{ESP32S3.GPIO} is
the HAL; everything links against the runtime crate.

\chapter{The Runtime Profiles}
\label{ch:profiles}

The runtime is offered in three \emph{profiles}, selected by the environment
variable \texttt{ESP32S3\_RTS\_PROFILE} when the crate is built. They trade
language features against footprint and certifiability. The same application
source can often build under more than one; the HAL drivers that use finalization
are restricted to the richer two.

\begin{longtable}{p{2.6cm} p{11.3cm}}
\textbf{Profile} & \textbf{What it gives you} \\ \hline
\endhead
\texttt{light-tasking} &
  The Ravenscar/Jorvik subset (\texttt{pragma Profile (Jorvik)}). Tasks,
  protected objects, \texttt{delay until}, priorities --- but with restrictions
  the compiler enforces: \texttt{No\_Finalization}, \texttt{No\_Exception\_%
  Propagation}, \texttt{No\_Secondary\_Stack}, and a heap-less bump allocator.
  Smallest and most analysable; a raised exception is fatal (resets the board).
  HAL drivers that need none of those (GPIO, RNG, Temperature, SDMMC, Touch,
  RTC, crypto) build here. \\
\texttt{embedded} &
  Still Jorvik, but exception-capable: full exception propagation with messages,
  controlled-type \emph{finalization} (RAII), a secondary stack, and a real
  heap (the O(1) TLSF allocator, \S\ref{ch:heap}). This is the default target for most of the HAL and for the
  ext4 filesystem. The RAII driver handles (SPI/I2C/UART sessions, GDMA/LEDC/RMT/
  \dots\ channels) live here. \\
\texttt{full} &
  Complete Ada tasking with no Jorvik restrictions: dynamic task creation, task
  hierarchies, selective \texttt{accept}/rendezvous, \texttt{abort} and
  asynchronous transfer of control. Built on a fully ported GNARL kernel.
  Experimental (work in progress), but real rendezvous and \texttt{abort} run on
  the silicon. \\
\end{longtable}

\section{Choosing a profile}

\begin{itemize}
\item Certifiable, deterministic, minimal-footprint firmware, or code that must
  never allocate: \texttt{light-tasking}.
\item Most applications --- you want exceptions, RAII handles, the secondary
  stack, a heap, and the full HAL including the filesystem: \texttt{embedded}.
\item Code that genuinely needs dynamic tasks, rendezvous or \texttt{abort}:
  \texttt{full}.
\end{itemize}

The selection happens at build time. An example's \texttt{build.sh} sets it:

\begin{lstlisting}[style=sh]
export ESP32S3_RTS_PROFILE=embedded   # or light-tasking / full
\end{lstlisting}

Because the HAL keys its object directory by profile, the same shared driver tree
is safely reused across profiles without stale objects.

\chapter{Building and Running on the ESP32-S3}
\label{ch:build}

\section{Prerequisites}

\begin{itemize}
\item The Alire-managed GNAT cross-toolchain for the target
  (\texttt{gnat\_xtensa\_esp32\_elf}) and \texttt{gprbuild}; these are fetched by
  Alire and discovered automatically by the build scripts.
\item A USB cable to the board's built-in USB-serial-JTAG port (it appears as
  \texttt{/dev/ttyACM0}).
\item Nothing else: no external SDK or build system. An Ada \texttt{esp\_flash}
  host tool (built with the Alire \texttt{gnat\_native} toolchain) packages and
  writes the image.
\end{itemize}

\section{The \texttt{./x} helper}

From the repository root, the \texttt{./x} script builds, flashes, monitors and
debugs any example by name:

\begin{lstlisting}[style=sh]
./x build  esp32s3_gpio0_blink            # cross-compile + link + package app.bin
./x flash  esp32s3_gpio0_blink -p /dev/ttyACM0
./x run    esp32s3_gpio0_blink -p /dev/ttyACM0   # build + flash + serial monitor
\end{lstlisting}

\section{What a build does}

\texttt{build.sh} sets the profile and heap size, then \texttt{bare\_build.sh}:

\begin{enumerate}
\item regenerates/selects the runtime crate for the chosen profile
  (\texttt{gen\_runtime.sh}; how that works, and how to port it, is
  Chapter~\ref{ch:runtime-build});
\item compiles the Ada application against the pinned runtime with
  \texttt{gprbuild}, producing a relocatable \texttt{ada\_app.o}
  (\texttt{Main} renamed to \texttt{\_ada\_main});
\item compiles the shared bare boot (\texttt{start.S}, the interrupt vectors,
  the heap and a freestanding libc) and the example's C console glue;
\item links everything against the linker scripts into \texttt{app.elf}; and
\item packages \texttt{app.bin} (image header + segments).
\end{enumerate}

\texttt{flash.sh} writes the 2nd-stage bootloader, the partition table and
\texttt{app.bin} over USB-serial.

The last two steps --- turning the linked \texttt{app.elf} into a flashable
\texttt{app.bin}, and writing it to the chip --- are handled by two small
\emph{host} tools written in native Ada (built with the Alire
\texttt{gnat\_native} toolchain), so the whole build/flash path needs only the
Alire toolchain and the committed vendored blobs.

\section{Packaging and flashing: two Ada host tools}

\subsection{\texttt{esp\_elf2image} --- ELF to flash image}

\texttt{examples/common/bare/elf2image} converts \texttt{app.elf} into the
ESP32-S3 flash image \texttt{app.bin} --- the exact on-flash format the chip's
ROM and our second-stage bootloader expect (flash mode \texttt{dio}, freq
\texttt{80m}, size \texttt{2MB}). The generated image boots unmodified on
hardware across every example. The packaging is more than a flat copy:

\begin{itemize}
  \item \textbf{Segments} are the ELF's allocated \texttt{PROGBITS} sections,
        merged where contiguous and each padded to a multiple of four.
  \item A \textbf{24-byte header} --- the 8-byte common part (\texttt{magic 16\#E9\#},
        segment count, the \texttt{dio}/\texttt{2MB}/\texttt{80m} flash byte, the
        entry point) plus a 16-byte extended part (\texttt{chip\_id=9},
        \texttt{wp\_pin=16\#EE\#}, \texttt{hash\_appended=1}).
  \item \textbf{Flash (IROM/DROM) segments aligned to 64\,KB}, with the RAM
        segments written \emph{interleaved as the alignment padding} (then a zero
        \texttt{PADDING} segment for the remainder) --- the exact scheme the
        bootloader's MMU mapping relies on.
  \item An \textbf{XOR checksum} (seed \texttt{16\#EF\#}) as the last byte of a
        16-aligned block, and finally a \textbf{SHA-256} of the whole image
        appended (computed by an in-tree \texttt{sha256.adb}).
\end{itemize}

The flash parameters (\texttt{dio}/\texttt{80m}/\texttt{2MB}, chip = ESP32-S3) are
fixed to match the bare boot; they are constants at the top of the source.

\subsection{\texttt{esp\_flash} --- writing it to the chip}

\texttt{examples/common/bare/espflash} writes the image over the chip's
\emph{ROM} serial bootloader.
It is 100\,\% Ada: the OS serial interface (open, termios raw mode, the
\texttt{TIOCM*} DTR/RTS modem lines, poll, read, write) is bound directly to libc
through \texttt{Interfaces.C} --- there is no C source. (\texttt{GNAT.Serial\_-
Communications} would cover the byte I/O but not the DTR/RTS control the reset
needs, so the tool binds \texttt{ioctl(TIOCMBIS/TIOCMBIC)} itself.) The sequence
is the standard ROM protocol, with no RAM stub and no compression:

\begin{enumerate}
  \item \textbf{Reset into download mode} --- the USB-JTAG DTR/RTS pulse.
  \item \textbf{SYNC} --- the \texttt{16\#07\# 16\#07\# 16\#12\# 16\#20\#} + 32$\times$\texttt{16\#55\#}
        handshake until the ROM answers.
  \item \textbf{SPI attach} and \textbf{set-params} (the flash geometry).
  \item Per file, \textbf{flash\_begin} (erase) then \textbf{flash\_data} in
        \texttt{16\#400\#}-byte blocks (the last padded with \texttt{16\#FF\#}), each
        carrying the ROM XOR checksum (seed \texttt{16\#EF\#}).
  \item \textbf{flash\_end}, then a \textbf{hard reset to run} (suppressible with
        \texttt{--no-reset} for capture workflows like the ACATS sweep, which
        does its own reset).
\end{enumerate}

Every command is a \texttt{<direction, op, length, checksum>} packet in SLIP
framing (\texttt{16\#C0\#} delimiters, \texttt{16\#DB\#} escaping) over the serial port.
A typical invocation writes the three images at their flash offsets:

\begin{lstlisting}[style=sh]
esp_flash -p /dev/ttyACM0 --flash-size 2MB \
  0x0     bootloader.bin \
  0x8000  partition-table.bin \
  0x10000 app.bin
\end{lstlisting}

Note that \texttt{esp\_flash} writes \emph{flash} only; the external PSRAM is not
flashed --- the 2nd-stage bootloader maps it at runtime. On hardware a $\sim$225\,KB
image flashes in about three seconds (ROM speed, no stub).

\section{Starting your own project}

The \texttt{esp32-ada} launcher scaffolds a fresh application in any empty
directory (no runtime sources are copied --- it pins the crate):

\begin{lstlisting}[style=sh]
. export.sh                 # put esp32-ada on PATH
mkdir my_app && cd my_app
esp32-ada new               # scaffold build.sh, the .gpr, src/main.adb, board.ads
# edit src/main.adb ...
esp32-ada run -p /dev/ttyACM0
\end{lstlisting}

Each project owns a \texttt{config/board.ads} that states the flash and PSRAM
sizes (used to size the image header and build/select the bootloader):

\begin{lstlisting}
package Board is
   Flash_Size : constant := 2 * 1024 * 1024;   --  2 MB
   PSRAM_Size : constant := 2 * 1024 * 1024;   --  mapped at 16#3D00_0000#
end Board;
\end{lstlisting}

\section{Console output}

The runtime wires \texttt{Ada.Text\_IO} to the built-in USB-serial-JTAG console,
so \texttt{Put\_Line} reaches \texttt{/dev/ttyACM0}. The examples additionally use
a tiny C glue layer calling the ROM \texttt{esp\_rom\_printf} for formatted status
lines. Either way, \texttt{./x run} opens a monitor so you see the output live.

\chapter{Debugging on the Chip}
\label{ch:debug}

When something misbehaves on the silicon there are three complementary tools: an
interactive debugger over JTAG, the post-mortem backtrace, and the runtime's own
fault diagnostics. Each works against the bare application image directly, using
only the pinned cross-toolchain.

\section{Interactive debugging: OpenOCD + GDB}\index{debugging}\index{OpenOCD}\index{GDB}\index{JTAG}\index{USB-JTAG}\index{breakpoint}\index{backtrace}

The ESP32-S3's built-in USB-Serial-JTAG port is both the console \emph{and} a JTAG
debug interface, so a single USB cable is enough. The project pins its own copies
of the debug tools; one command fetches them:

\begin{lstlisting}[style=sh]
./x get-debug-tools      # downloads OpenOCD + xtensa-esp32s3-elf-gdb into tools/
./x debug blink          # build, flash, start OpenOCD, attach GDB, rest at Main
\end{lstlisting}

One subtlety is worth knowing: the S3 needs the \emph{S3-specific}
\texttt{xtensa-esp32s3-elf-gdb}; the Alire crate's esp32 (LX6) GDB silently fails
on it. \texttt{./x debug} selects the right one. From there it is ordinary GDB on
\texttt{app.elf}: breakpoints, \texttt{step}, \texttt{print},
\texttt{backtrace}. Two options matter on this target:

\begin{itemize}
  \item \texttt{--smp} exposes \emph{both} LX7 cores as GDB threads
        (\texttt{info threads}), so you can see what each core is doing,
        essential for a cross-core hang.
  \item \texttt{--attach} halts the board \emph{in place} without resetting it, so
        you can examine a live hang or crash exactly where it stopped (a
        post-mortem).
\end{itemize}

The command line and both editors (next section) drive this same OpenOCD + GDB
path; when finished, \texttt{./x kill-openocd} releases the USB-JTAG port.

\section{Editor integration: VS Code and Vim}\index{editor integration}\index{Ada Language Server}

The same build, flash and debug actions are wired into both VS Code and Vim, so
you need never leave the editor. Both drive the top-level \texttt{./x} dispatcher
and both use the project's pinned OpenOCD + s3 GDB for on-chip debugging.

Two prerequisites are shared by either editor. On-chip debugging first needs the
tools fetched once:

\begin{lstlisting}[style=sh]
./x get-debug-tools     # pinned, SHA-256-verified OpenOCD + xtensa-esp32s3-elf-gdb
\end{lstlisting}

And \emph{code intelligence} (completion, go-to-definition, diagnostics, hover)
comes from the \textbf{Ada Language Server}, which ships inside the AdaCore
toolchain (nothing to install separately). It only needs to be pointed at the
example's \texttt{.gpr} project file, shown for each editor below.

\subsection{VS Code}\index{VS Code}

There are two complementary pieces. The first is the committed workspace
configuration in \texttt{.vscode/} (\texttt{tasks.json}, \texttt{launch.json},
\texttt{settings.json}), which wires the editor to \texttt{./x} and needs two
extensions:

\begin{description}[style=nextline]
\item[Ada \& SPARK (\texttt{AdaCore.ada}).] The Ada Language Server: required,
  for all code intelligence.
\item[Native Debug (\texttt{webfreak.debug}).] A GDB frontend, needed only for
  \texttt{F5} debugging. \texttt{cppdbg} cannot walk the Xtensa register-window
  stack (it connects but shows no call stack or source); Native Debug simply
  relays GDB's own frames, which unwind correctly.
\end{description}

Install both from a terminal (or from the Extensions view, which also offers them
as workspace recommendations):

\begin{lstlisting}[style=sh]
code --install-extension AdaCore.ada
code --install-extension webfreak.debug
\end{lstlisting}

Then point the language server at the example you are working on by editing
\texttt{ada.projectFile} in \texttt{.vscode/settings.json} and reloading the
window (\texttt{Ctrl-Shift-P} $\rightarrow$ \emph{Developer: Reload Window}):

\begin{lstlisting}[style=sh]
"ada.projectFile": "examples/esp32s3_gpio0_blink/gpio_blink.gpr"
\end{lstlisting}

Build, flash and monitor are tasks: \texttt{Ctrl-Shift-B} runs the default
\texttt{Ada: build}, and \emph{Tasks: Run Task} offers flash, monitor, run and
clean. Each prompts for the example, serial port and runtime profile and routes
to \texttt{./x}; build errors land in the \textbf{Problems} panel via a GNAT
problem matcher.

To \textbf{debug} (\texttt{F5}, ``Ada: debug (OpenOCD + GDB)''): build \emph{and
flash} the example first, so the chip runs the same image whose \texttt{app.elf}
you debug. The launch only \emph{attaches}, so a stale or different-profile
image leaves the breakpoint addresses mismatched. It starts OpenOCD and halts at
\texttt{app\_main} (the shared boot-glue entry; these GNAT images have no C
\texttt{main}); set breakpoints in the Ada source and continue/step.

The second piece is an optional \textbf{status-bar extension}
(\texttt{ide/vscode-ada-esp32}) that adds clickable build / flash / monitor /
debug icons plus persistent port and profile selectors. A prebuilt
\texttt{.vsix} is committed to the repository, so installing needs only the
\texttt{code} CLI on \texttt{PATH}:

\begin{lstlisting}[style=sh]
./x install-ide         # installs the committed vscode-ada-esp32.vsix
\end{lstlisting}

Reload the window afterwards; re-run it to update after a \texttt{git pull} (a
repo-hosted \texttt{.vsix} does not auto-update). Its debug icon uses the same
OpenOCD + GDB launch described above.

\subsection{Vim and Neovim}\index{Vim}\index{termdebug}

The \texttt{ide/vim-ada-esp32} plugin gives the same actions as Vim commands. It
is pure Vimscript (nothing to compile) and needs \textbf{Vim 8.1+ with
\texttt{+terminal}} (which bundles the \texttt{termdebug} GDB frontend); Neovim
works too. Install it with one command, which symlinks the plugin into Vim's and
Neovim's native package directories; because it is a symlink \emph{into} the repo,
a \texttt{git pull} updates it in place:

\begin{lstlisting}[style=sh]
./x install-vim
\end{lstlisting}

Or wire it up by hand; it has the standard \texttt{plugin/} layout:

\begin{lstlisting}[style=sh]
" vim-plug
Plug '/path/to/ada_esp32s3/ide/vim-ada-esp32'
" or, with no plugin manager:
set runtimepath^=/path/to/ada_esp32s3/ide/vim-ada-esp32
\end{lstlisting}

The commands auto-detect context by searching upward for \texttt{./x}, so they
work from anywhere inside the repo:

\begin{description}[style=nextline]
\item[{\texttt{:AdaEsp32Build [example]}}] build via \texttt{./x build}.
\item[{\texttt{:AdaEsp32Flash [example]}}] flash to the selected port.
\item[{\texttt{:AdaEsp32BuildFlash [example]}}] build then flash.
\item[{\texttt{:AdaEsp32Run [example]}}] build $+$ flash $+$ monitor.
\item[{\texttt{:AdaEsp32Monitor}}] serial console at 115200.
\item[{\texttt{:AdaEsp32Debug [example]}}] OpenOCD $+$ \texttt{termdebug} (below).
\item[{\texttt{:AdaEsp32Port [port]} / \texttt{:AdaEsp32Profile [prof]}}] set the
  serial port / runtime profile; both persist across later commands.
\end{description}

The example argument accepts the short name (\texttt{gpio0\_blink}) or the full
directory (\texttt{esp32s3\_gpio0\_blink}), with tab-completion; the chosen
example, port and profile persist (in \texttt{g:ada\_esp32\_example} and friends)
so later commands reuse them without re-prompting.

\textbf{Debugging} uses Vim's built-in \texttt{termdebug}, the Vim analog of
the Native Debug setup, relaying GDB's own frames so it handles Xtensa correctly.
After a build:

\begin{lstlisting}[style=sh]
:AdaEsp32Build gpio0_blink
:AdaEsp32Debug gpio0_blink
\end{lstlisting}

This starts OpenOCD in the background, waits for its GDB server, then opens
\texttt{termdebug} reset and halted with a breakpoint armed at \texttt{app\_main}.
From there \texttt{:Continue} runs to it, \texttt{:Break src/foo.adb:53} (or
\texttt{:Break} on the cursor line) sets an Ada breakpoint, and
\texttt{:Step}/\texttt{:Over}/\texttt{:Finish} are the step controls; OpenOCD is
stopped automatically when you quit the session. One inherent caveat: a
\texttt{:Over} across a \emph{blocking} statement (a \texttt{delay}, an entry
call) can step into the hand-written Xtensa assembly (the interrupt vectors or
context switch), which carries no debug bounds. GDB then reports ``cannot find
bounds of current function.'' Set a breakpoint past the statement and
\texttt{:Continue} instead.

Code intelligence again comes from the Ada Language Server through your LSP
client, pointed at the example's \texttt{.gpr}. For example, with
\texttt{coc.nvim} in \texttt{coc-settings.json}:

\begin{lstlisting}[style=sh]
{ "ada.projectFile": "examples/esp32s3_gpio0_blink/gpio_blink.gpr" }
\end{lstlisting}

or register \texttt{ada\_language\_server} for the \texttt{ada} filetype with
native LSP / vim-lsp and set the same \texttt{ada.projectFile} setting.

\section{Post-mortem: decoding a Guru Meditation}\index{Guru Meditation}\index{panic}\index{addr2line}\index{post-mortem}

An unhandled fault prints a register dump and a backtrace of program-counter
values over the console (a ``Guru Meditation''). You do not need a live debugger
to read it: turn the addresses back into source with \texttt{addr2line} against
the ELF:

\begin{lstlisting}[style=sh]
xtensa-esp32s3-elf-addr2line -fe build/app.elf 0x4200abcd 0x4200ef01
\end{lstlisting}

In the dump, \texttt{EXCVADDR} is the faulting data address and \texttt{EXCCAUSE}
the reason: a \texttt{LoadStoreError} (a bad data access), or an
\texttt{InstructionFetchError} (a wild jump, often a corrupted stack or a stray
pointer). Flash, reproduce, decode: that loop localises most faults without ever
attaching a debugger.

\section{The runtime's own diagnostics}\index{stack overflow}\index{last-chance handler}\index{safe state}

The runtime helps before you reach for a tool at all:

\begin{itemize}
  \item \textbf{Precise stack-overflow detection.} Each task arms a hardware
        watchpoint near its stack limit, so an overflow faults \emph{precisely}
        (a debug exception) instead of silently corrupting a neighbour, and on
        the \texttt{full} profile the runtime can turn it into a catchable
        \texttt{Storage\_Error}.
  \item \textbf{A fault safe-state hook.} The \texttt{esp32s3\_exceptions}
        example wraps the panic handler to classify a fault (stack overflow
        versus other) against the running task's stack window, latch a safe
        state, and print an explicit \texttt{*** STACK OVERFLOW ***} verdict
        before the board resets.
  \item \textbf{The last-chance handler.} Any exception that escapes every handler
        ends up in the runtime's last-chance handler; on the leanest profile,
        where exceptions cannot propagate across calls, that is your fault report:
        a message and a controlled reset.
\end{itemize}

\textbf{A note on console output while debugging.} \texttt{Ada.Text\_IO} reaches
\texttt{/dev/ttyACM0}; for terse status from interrupt or early-boot context the
examples call the ROM \texttt{esp\_rom\_printf}, which is always available. When
two cores (or an interrupt handler and a task) might print at once, funnel
output through one task: interleaved single-character writes have themselves been
mistaken for bugs more than once.

\section{Breadcrumbs at the panic park}\index{panic}\index{breadcrumbs}\index{double exception}\index{PRID}

When a fault escalates past every handler \textemdash{} an unhandled
exception, a double exception, a debug exception with no debugger attached
\textemdash{} the vendored vectors funnel into \texttt{\_xt\_panic}, which
used to be a bare \texttt{j .}: the board parked with the evidence locked in
special registers that the next reset destroys.  \texttt{start.S} now dumps
the machine state to a fixed block (\texttt{\_\_gnat\_panic\_dump}) before
parking: EXCCAUSE and EXCVADDR, EPC1..EPC7 with their EPS twins, DEPC, PS,
\texttt{a0}/\texttt{a1}, WINDOWBASE/WINDOWSTART, PRID, and a marker for which
entry point was taken.  The first writer wins \textemdash{} a magic word gates
the block, so when both cores pile into the park the original fault survives
\textemdash{} and the block lives in \texttt{.data}, re-initialised at every
boot, so a stale dump cannot masquerade as a fresh one.  One JTAG
\texttt{mdw} reads the whole story out of a wedged board.

Two decoding traps, both paid for in hours.  The double-exception vector
\emph{overwrites} EXCCAUSE with a panic-reason code before calling
\texttt{\_xt\_panic}, so ``cause 2'' at a double-exception park means
\emph{double exception}, not \emph{instruction fetch error}.  And PRID does
not name the core by its whole value: bit~13 alone does.  On the ESP32-S3,
\texttt{16\#CDCD\#} is core~0 and \texttt{16\#ABAB\#} is core~1
\textemdash{} misread that and every per-core deduction that follows is
backwards.

\section{Hardware watchpoints without a debugger}\index{DBREAK}\index{watchpoint}\index{debug exception}

OpenOCD watchpoints have two failure modes on this chip: they cannot be armed
early enough to catch a corruption that happens inside the first second of
boot, and an armed watchpoint over an address the ROM itself uses turns reset
into a watchdog loop.  The escape is to let the \emph{firmware} arm the
watchpoint at exactly the right moment \textemdash{} three instructions:

\begin{lstlisting}[style=sh]
wsr.dbreaka0  <address>       # the word to watch
wsr.dbreakc0  16#8000_003C#   # store-break, 4-byte mask
isync                          # (16#4000_003C# = load-break)
\end{lstlisting}

If a debugger is attached, the hit halts the core at the offending store with
the writer's PC on screen.  If \emph{no} debugger is attached the hit becomes
a debug exception that funnels into \texttt{\_xt\_panic} \textemdash{} and
the breadcrumb block's EPC6 slot (the debug level is~6) records the writer's
PC anyway.  A store-watch armed after setup, plus one \texttt{mdw} after the
park, identified in one boot a rogue writer that three days of halt-and-inspect
archaeology had missed.

The technique has a sting worth knowing even if you never use it: DBREAK state
\emph{survives soft resets}.  The full runtime's stack-overflow guard arms a
DBREAK in normal operation, so reflashing from a full-profile image to any
other image used to inherit a live watchpoint over what is now unrelated data
\textemdash{} and the new image died at boot with a debuggerless debug
exception.  \texttt{start.S} therefore zeroes \texttt{DBREAKC0/1} and
\texttt{IBREAKENABLE} on both cores before anything else runs; if you port the
boot path elsewhere, keep that.

\section{Decoding the compiler's bare-metal complaints}
\index{No\_Implicit\_Dynamic\_Code}\index{accessibility check}\index{restrictions}

The bare-metal profiles forbid things a hosted program takes for granted, and the
compiler enforces the bans. The messages are precise but terse, and the first time
one stops a build it can be baffling. Here are the ones you are most likely to meet
coming from C, what each really means, and the fix.

\paragraph{\texttt{subprogram must not be deeper than access type}.}
You took \texttt{'Access} of a \emph{nested} subprogram (one declared inside
another subprogram) to use as a callback. On a hosted target the compiler would
build a \emph{trampoline} (a scrap of code on the stack) to carry the enclosing
frame; the \texttt{No\_Implicit\_Dynamic\_Code} restriction forbids writing code to
the stack at all, and on this chip such a trampoline faults outright. \emph{Fix:}
move the callback to library level (the top of a package) and pass any state it
needs through a context parameter rather than capturing it: closure-free. This is
why every callback in this book is a library-level, \texttt{Ctx}-carrying
subprogram (\S\ref{ch:ada-primer} introduces subprograms; the HAL callbacks follow
this rule throughout).

\paragraph{\texttt{accessibility check failed} (a \texttt{Program\_Error} at run
time), or \texttt{non-local pointer cannot point to local object} (at compile
time).} You stored an \texttt{'Access} of an object into something that outlives it
into a library-level variable, a registry, a longer-lived record. Once the object's
scope ends the stored pointer would dangle, so Ada rejects it: at compile time when
it can prove the lifetimes, at run time otherwise. \emph{Fix:} give the object the
lifetime the pointer assumes, which usually means declaring it at library level.
(This is exactly the bug behind the FTP server that stored a \texttt{Main}-local
filesystem mount in a library-level handler, cured by moving the mount into a
library-level package.)

\paragraph{\texttt{violation of restriction No\_X}.}
You used a feature the chosen runtime profile bans: a heap allocator, a second
task entry, an exception handler under a no-propagation profile, and so on. The
profile is a contract, and the compiler holds you to it. \emph{Fix:} either avoid
the feature, or build under a richer profile (light-tasking $\rightarrow$ embedded
$\rightarrow$ full) that permits it; Chapter~\ref{ch:profiles} maps features to
profiles.

\paragraph{\texttt{Constraint\_Error} (a range or validity check).}
A value left the range its (sub)type allows: an index past the end of an array,
an out-of-range assignment, a bad enumeration read from external data. Unlike the
above, this is almost always a genuine bug in your data or arithmetic rather than a
profile rule. \emph{Fix:} find the offending value; the type told you its range, so
the check is doing its job. Build with checks on (\texttt{-gnata}) while developing;
a release build may have them off, turning the same mistake into silent
corruption.

The throughline is that each message names a guarantee the bare-metal target makes
(no stack-written code, no dangling pointers, no banned features, no out-of-range
values) and points at the one place you broke it. They read as obstacles at first
and as a safety net soon after.

%  Language + hardware techniques BEFORE the HAL that uses them, then the
%  encapsulation features the drivers are built from.
\chapter{Ada Language Techniques}
\label{ch:embedded}

A handful of Ada features come up again and again in embedded code: writing
literals clearly, constraining values to only what is legal, validating data that
arrives from outside, reinterpreting bits deliberately, attaching checkable
contracts, and handling (or forbidding) exceptions. This chapter collects the
language-level techniques; the next, \emph{Talking to Hardware}, applies them to
registers and memory. Together they underpin the HAL and the filesystem, so a few
terms are worth learning before the driver chapters that use them.

\section{Terminology}
\label{sec:terms}\index{terminology}

A few Ada words recur throughout the book and can trip up a reader coming from C.
They are defined once here.

\begin{description}[style=nextline]
\item[Access type / \texttt{'Access}.] Ada's typed pointer; \texttt{X'Access}
  yields a pointer to \texttt{X}. Strongly typed, with accessibility rules that
  guard against dangling references.
\item[Activation.] The moment a declared task begins executing: after its master's
  declarations, or, for a library-level task, as part of elaboration before
  \texttt{Main}.
\item[Aspect.] A property attached to a declaration with the word \texttt{with}:
  \texttt{with Volatile}, \texttt{with Address => \ldots}, \texttt{with Pre =>
  \ldots}. Comparable to a GCC \texttt{\_\_attribute\_\_}, but standardised and
  checked.
\item[Attribute.] A built-in query or operation written with an apostrophe:
  \texttt{X'Image}, \texttt{T'Last}, \texttt{Obj'Address}, \texttt{P'Access}. There
  is no C equivalent.
\item[Ceiling priority.] The priority at which a protected object's operations
  run (the priority-ceiling protocol); for an interrupt handler it is the
  interrupt's hardware level.
\item[Constraint check / \texttt{Constraint\_Error}.] The run-time check that a
  value is in range or valid; a failure raises \texttt{Constraint\_Error}. The
  checks can be switched off in a release build.
\item[Elaboration.] The once-only, ordered set-up that runs \emph{before}
  \texttt{Main}: evaluating initial values, executing the statements in package
  bodies, and activating library-level tasks. Like C's static initialisers, but
  ordered and able to run arbitrary code.
\item[Entry / barrier.] A protected (or task) operation guarded by a Boolean
  condition, the \emph{barrier}; a caller blocks until the barrier is open. This is
  how condition synchronisation is expressed.
\item[Generic.] A compile-time template (a package or subprogram parameterised
  by types or values) instantiated to produce a concrete unit, much like a C++
  template.
\item[Library unit / library-level.] A unit (package or subprogram) declared at
  the outermost level, \emph{not} nested inside a subprogram or block. A
  library-level object or task exists for the whole run of the program. C has only
  file scope and block scope; Ada nests arbitrarily, so ``library-level'' names
  the top, program-lifetime level, and the Jorvik profiles require tasks (and
  some other entities) to live there.
\item[Package (specification / body).] A named module. The \emph{specification}
  is the visible interface (roughly a C header) and the \emph{body} is the
  implementation (roughly the \texttt{.c}), but they are one logical unit and
  the body's internals are genuinely hidden by the compiler, not merely by
  convention.
\item[Pragma.] A directive to the compiler: \texttt{pragma Pack},
  \texttt{pragma Attach\_Handler}, \texttt{pragma Restrictions}. Like
  \texttt{\#pragma}, but part of the language with defined meanings.
\item[Profile (Ravenscar / Jorvik).] A named bundle of \texttt{Restrictions}
  defining an analysable tasking subset. Distinct from this project's runtime
  \emph{build} profiles, which choose \emph{which} such subset the runtime
  supports.
\item[Protected object.] Data plus operations with \emph{automatic} mutual
  exclusion: a monitor. Access is race-free without a manual lock; the compiler
  generates the locking.
\item[Rendezvous.] A synchronous call from one task to another task's
  \emph{entry}: the two meet, exchange parameters, the accept body runs, then both
  continue (the \texttt{full} profile).
\item[Representation clause.] A clause that fixes a type's machine layout (bit
  positions, size, byte order): \texttt{for T use record \ldots\ end record}. It is
  how a record is mapped onto a hardware register.
\item[Secondary stack.] A runtime area used to return a result whose size the
  caller does not know in advance (e.g.\ a function returning \texttt{String}). C
  has no equivalent; the leanest profile omits it.
\item[Task.] A thread of control defined in the language. It \emph{activates}
  (starts) automatically. There is no explicit start call.
\item[Type vs.\ subtype.] A \emph{type} is a distinct set of values and
  operations; two distinct types never mix without an explicit conversion. A
  \emph{subtype} is the same type plus a constraint (a range or predicate),
  fully compatible with its base, just narrower.
\item[Unconstrained type.] A type whose size is settled only when an object is
  created: \texttt{String}, \texttt{array (Positive range <>)}, a record with
  defaulted discriminants.
\end{description}

%  The lexical basics (numeric bases, underscores, and case insignificance)
%  are introduced early, with the first program, in \S\ref{sec:bases} (the Anatomy
%  chapter), so the register code in the intervening chapters reads correctly.
%  This chapter picks up where that leaves off: the TYPES those numbers take.

\section{Subtypes and modular types}
\label{sec:subtypes}\index{subtype}

The other way the type carries intent is by constraining which values are legal
at all. Two facilities earn their keep constantly in embedded code: subtypes that
narrow a value to a valid range or set, and modular types that give wrap-around
arithmetic and bit operations.

\subsection{Range subtypes: only the legal values}

A subtype attaches a constraint to an existing type. The everyday embedded use is
to make an out-of-range value impossible:

\begin{lstlisting}
   subtype Channel is Natural range 0 .. 7;     -- a peripheral with 8 channels
   subtype Percent is Natural range 0 .. 100;   -- a duty cycle, say
\end{lstlisting}

A value known at compile time that violates the range is a compile error; a value
computed at run time is caught by a constraint check (wherever checks are
enabled). Either way ``channel 9'' never reaches the hardware, and the range is
documented in the type rather than in a comment.

\subsection{Static-predicate subtypes: a non-contiguous valid set}

When the legal values are not one contiguous range, a \emph{static predicate}
states the exact set. The HAL's pin type is the real example: of the 0\,..\,48 pad
numbers, some are absent on the package and some are bonded to the in-package SPI
flash and octal PSRAM. Driving any of those hangs the chip:

\begin{lstlisting}
   type Pad_Number is range -1 .. 48;            -- -1 is the "no pin" sentinel

   subtype Pin_Id is Pad_Number range 0 .. 48 with
     Static_Predicate => Pin_Id in 0 .. 21 | 38 .. 48;   -- 22..37 reserved
\end{lstlisting}

Because the predicate is \emph{static}, naming a reserved pad as a constant of
type \texttt{Pin\_Id} is rejected at compile time (\emph{static expression fails
static predicate check}); a pin computed at run time is verified by the predicate
under \texttt{-gnata}. Every driver names its pins with \texttt{Pin\_Id}, so a
flash/PSRAM pad cannot be wired by accident.

\cnote{a pin is an \texttt{int} or a \texttt{\#define}, and wiring a flash pad is a
run-time hang you debug with an oscilloscope. \texttt{Pin\_Id} turns a bad
constant into a \emph{compile} error and a bad computed value into a checked one.}

\subsection{Modular types: wrap-around and bits}

A modular type (\texttt{mod 2**n}) holds the values $0\,..\,2^n-1$, and its
arithmetic \emph{wraps} instead of overflowing. That is exactly the behaviour you
want for a register word, a free-running counter, or a checksum; modular types
also support the logical operators \texttt{and}, \texttt{or}, \texttt{xor} and
\texttt{not}:

\begin{lstlisting}
   type Byte is mod 2 ** 8 with Size => 8;       -- 0 .. 255, wraps at 256

   Sum : Byte := 0;
   ...
   Sum   := Sum + B;                       -- no overflow error; wraps (a checksum)
   Flags := Flags or  2#0000_0100#;        -- set a bit
   Flags := Flags and not 2#0000_0100#;    -- clear a bit
\end{lstlisting}

The predefined \texttt{Interfaces.Unsigned\_8/16/32/64} are the usual choice for
register-sized words: they are modular and add \texttt{Shift\_Left} /
\texttt{Shift\_Right} / \texttt{Rotate\_*}. A plain signed integer would instead
raise \texttt{Constraint\_Error} on overflow and offers no bit operators. So the
rule of thumb is: registers, masks and counters are \emph{modular}, while
quantities that must never silently wrap (an index, a length) stay
\emph{range-constrained}.

\subsection{And beyond}

Two more worth keeping in mind:
\begin{itemize}
  \item \textbf{Fixed-point types} (\texttt{type Q8 is delta 1.0 / 256.0 range
        0.0 .. 1.0}) give fractional values with exact, integer-backed
        arithmetic: a duty cycle or a scaled sensor reading without dragging in
        floating point.
  \item \textbf{Derived types} make otherwise-identical numbers
        \emph{incompatible}: with \texttt{type Hz is new Natural;} and
        \texttt{type Millis is new Natural;}, a frequency can never be passed
        where a delay was meant. The mix-up is a compile error, not a bug.
\end{itemize}

\section{Arrays and arbitrary index ranges}
\label{sec:arrays}\index{arrays}\index{array index range}

An Ada array type fixes not just its element type but its \emph{index range}, and
that range is yours to choose: any discrete range at all. The lower bound need
not be 0 or 1; it can be any value (even negative), and the index can be an
enumeration type rather than an integer. The range is part of the type:

\begin{lstlisting}
type Sine_Table is array (0  .. 255) of Integer;     -- 0-based lookup table
type CAN_Data   is array (1  .. 8)   of Unsigned_8;  -- 1-based, eight bytes
type Window     is array (-3 .. 3)   of Float;       -- centred on zero
type Channel    is (Red, Green, Blue);
type Mix        is array (Channel)   of Unsigned_8;  -- indexed by an enum
\end{lstlisting}

The point is not novelty: it lets the index \emph{match the thing it models}, so
the code reads like the datasheet or the protocol, with no mental ``subtract
one'' and none of the off-by-one bugs that come with it.

\textbf{Iterating without knowing the bounds.} Every array value carries its own
bounds, available as attributes: \texttt{A'First}, \texttt{A'Last},
\texttt{A'Length} and \texttt{A'Range} (which is \texttt{A'First\,..\,A'Last}). The
idiomatic loop ranges over \texttt{'Range}, so it is correct whatever the bounds
are, and stays correct if you later change them:

\begin{lstlisting}
Total : Natural := 0;
for I in Buf'Range loop          -- I takes each actual index value in turn
   Total := Total + Natural (Buf (I));
end loop;
\end{lstlisting}

The same code works for a \texttt{0\,..\,255} buffer and a \texttt{1\,..\,8} one
because the loop names \texttt{Buf'Range}; there is no literal bound to get
wrong. (When you do not need the index itself, \texttt{for E of Buf loop}
iterates the elements directly.)

\textbf{A worked example: J1939 CAN payloads, indexed from 1.} A Classic CAN
frame carries up to eight data bytes. SAE J1939 (the protocol on commercial
vehicles) numbers those bytes \emph{1 through 8} in its specifications
(``byte 1,'' ``bytes 4--5,'' and so on). Declaring the payload with a
\texttt{1\,..\,8} index makes the Ada line up with the document exactly:

\begin{lstlisting}
type J1939_Payload is array (1 .. 8) of Interfaces.Unsigned_8;

--  PGN 61444 (EEC1): engine speed is SPN 190, in BYTES 4-5,
--  little-endian, scaled 0.125 rpm/bit.  The byte numbers here are
--  taken verbatim from the J1939 spec -- no off-by-one translation.
function Engine_RPM (D : J1939_Payload) return Float is
   Raw : constant Interfaces.Unsigned_16 :=
            Interfaces.Unsigned_16 (D (4)) +              -- byte 4: low
            Interfaces.Unsigned_16 (D (5)) * 256;         -- byte 5: high
begin
   return Float (Raw) * 0.125;
end Engine_RPM;
\end{lstlisting}

\texttt{D (4)} \emph{is} ``byte 4'' (the same number a J1939 reference uses),
so the code can be checked against the standard line by line, and the loop idiom
still applies: \texttt{for I in D'Range loop} walks bytes \texttt{1\,..\,8}.

\cnote{In C the eight bytes are \texttt{data[0..7]}, so every J1939 byte number
from the spec has to be mentally decremented as you read or write the code (spec
``byte 4'' becomes \texttt{data[3]}). An arbitrary lower bound removes that whole
translation step.}

\section{The \texttt{others} choice: filling in ``everything else''}
\label{sec:others}\index{others}

\texttt{others} is a catch-all \emph{choice}: it stands for ``every case not
named so far,'' and it must always come \emph{last}. Ada uses it wherever you
enumerate possibilities: array and record aggregates, \texttt{case} statements,
and exception handlers. It is everywhere in embedded Ada (it fills the reserved
fields of the register writes in Chapter~\ref{ch:hardware}), so it is worth a
close look.

\textbf{In array aggregates: fill the rest.} An aggregate must give a value for
\emph{every} element; \texttt{others} supplies the ones you did not name.

\begin{lstlisting}
Zeros  : Buffer := (others => 0);                            -- every element 0
Header : Buffer := (0 => 16#AA#, 1 => 16#55#, others => 0);  -- first two, rest 0
\end{lstlisting}

\textbf{In record aggregates, and the \texttt{<>} ``box''.} A record aggregate
must also set every component. Writing \texttt{others => <>} means ``give each
remaining component its \emph{default} value'': the default from the type's
own declaration. This is the central idiom for writing a hardware register: set
the one or two fields you mean, and let every reserved or don't-care field take
its declared default.

\begin{lstlisting}
--  A register record whose fields mostly default to 0 / False:
R.TX_CONF := (TX_START => True, others => <>);   -- set one bit, default the rest
\end{lstlisting}

Use \texttt{others => <>} rather than \texttt{others => 0} for a record whose
remaining components are of \emph{different} types: each takes its own default, so
the types need not agree. \texttt{others => 0} only works when everything it
covers is the same type (e.g.\ an all-zero byte buffer).

\textbf{In \texttt{case}: cover the remaining values.} A \texttt{case} must be
exhaustive; \texttt{when others} handles whatever you did not list.

\begin{lstlisting}
case Mode is
   when Standard => ...;
   when others   => ...;     -- every other value of the type
end case;
\end{lstlisting}

\textbf{In exception handlers: catch the rest.} \texttt{when others} catches
any exception not named earlier in the handler.

\begin{lstlisting}
begin
   ...
exception
   when Constraint_Error => ...;   -- a specific one
   when others           => ...;   -- anything else
end;
\end{lstlisting}

Two properties tie these uses together: \texttt{others} is always the \emph{last}
choice, and it covers \emph{exactly} the cases not already named, so adding a
component, value or exception elsewhere silently changes what ``everything else''
means. For a register write that is precisely what you want (a newly-defined
reserved bit still defaults). For a \texttt{case} over an enumeration it is often
the opposite: omitting \texttt{others} is safer because then \emph{adding} a value
to the type forces the compiler to make you handle it.

\section{Fixed-point arithmetic}
\label{sec:fixed}\index{fixed-point types}

Microcontroller code is full of fractional quantities (a duty cycle, a gain, a
scaled sensor reading) that do not justify (or cannot afford) floating point.
A \emph{fixed-point} type gives fractions with exact, integer-backed arithmetic:
you state the resolution with \texttt{delta}, and the compiler represents values
as integers scaled by that \texttt{'Small}.

\begin{lstlisting}
   type Duty is delta 1.0 / 256.0 range 0.0 .. 1.0;   -- 8-bit fraction, 0 .. 1
   D : Duty := 0.75;                 -- stored as the integer 192 (= 0.75 * 256)

   type Volts is delta 0.001 range 0.0 .. 3.3;        -- millivolt resolution
\end{lstlisting}

Arithmetic on \texttt{Duty} is integer arithmetic underneath (no FPU, no
rounding surprises), yet the source reads in real-world units. The PWM and ADC
drivers scale this way: a duty cycle or a voltage is a fixed-point value the
driver converts to the register's raw counts.

\section{Trusting external data: the \texttt{'Valid} attribute}
\label{sec:valid}\index{Valid attribute}

Anything that enters from outside the program (a value read from a register, a
field unpacked from an SD sector, a message off the wire) might not be a legal
value of its type. Reading it into a constrained type and then using it is
\emph{erroneous} if the bits are out of range. \texttt{X'Valid} asks ``is the
object actually a valid value of its (sub)type?'' and answers with a plain
Boolean (no exception, safe to ask even of garbage):

\begin{lstlisting}
   Raw : Parity with Address => Reg'Address, Volatile;  -- enum from hardware
begin
   if Raw'Valid then
      Use_Parity (Raw);             -- known to be None/Even/Odd
   else
      Handle_Corrupt;               -- some other bit pattern
   end if;
\end{lstlisting}

This is the safe way to consume an enumeration with reserved codes, a range field
that hardware might leave uninitialised, or a structure read from possibly-corrupt
storage: validate first, trust after.

\section{\texttt{Unchecked\_Conversion}: deliberate reinterpretation}
\label{sec:uncheck}\index{Unchecked\_Conversion}

Occasionally you must view the same bits as two different types: serialise a
record to bytes, decode a fat pointer, read a float's bit pattern.
\texttt{Ada.Unchecked\_Conversion} is the explicit, greppable way to do it; unlike
a C cast it is a named instantiation, so every reinterpretation in the codebase is
searchable.

\begin{lstlisting}
   function To_Fat is new Ada.Unchecked_Conversion (Handler_Access, Fat_Ptr);
\end{lstlisting}

The runtime uses exactly this to inspect an \texttt{access protected procedure}
as a \texttt{(object, code)} pair when matching a registered interrupt handler.
The discipline: prefer a representation clause or an address overlay when the two
views have a defined relationship; keep \texttt{Unchecked\_Conversion} for the
genuine bit-reinterpretation cases, where its name documents the hazard.

\section{Access types: where you need a pointer, and where you don't}
\label{sec:access}\index{access types}

Ada has pointers (\emph{access types}), but a great many of the places a C
programmer reaches for one do not need it. The embedded habit is to lead with the
pointer-free option and use an access type only where the problem genuinely
demands it.

\subsection{Where you can avoid them}

\textbf{Parameters are already passed efficiently.} You do not pass
\texttt{\&thing} and dereference. A parameter mode (\texttt{in},
\texttt{out}, \texttt{in out}) says how it is used, and the compiler passes a
large object by reference automatically:

\begin{lstlisting}
   procedure Process (Frame : in out Packet);   -- no pointer; Frame is the object
\end{lstlisting}

\textbf{Whole arrays and records are values.} You can assign, return and pass
entire structures; there is no need for a pointer just to ``return a struct.''

\begin{lstlisting}
   R := S;                          -- copies a whole record
   function Identity return Matrix; -- returns a whole array by value
\end{lstlisting}

\textbf{Unconstrained types size themselves.} A \texttt{String} or an
\texttt{array (Positive range <>)} takes its bounds from its value, so a function
can return one without \texttt{malloc} and without a pointer (the secondary stack
carries it):

\begin{lstlisting}
   function Pin_Name (P : Pin_Id) return String;   -- no heap, no pointer
\end{lstlisting}

\textbf{A register is an address overlay, not a pointer.} As
\S\ref{ch:registers} showed, you place a typed object \emph{at} the register's
address; you never hold a \texttt{volatile uint32\_t *} and dereference it.

\textbf{A resource is a handle, not a raw pointer.} The limited-controlled handles
of Chapter~\ref{ch:encap} own a peripheral by scope, not by a pointer you must
remember to free.

\subsection{Where you genuinely need one}

\begin{itemize}
  \item \textbf{Dynamic or linked data built at run time:} a list, a tree, a
        variable-length buffer pool: \texttt{type Node\_Ref is access Node;} and
        \texttt{new Node}.
  \item \textbf{Callbacks and interrupt handlers:} an access-to-subprogram
        value lets you register code to be called later. This is the
        \texttt{On\_Edge'Access} from the GPIO counter (\S\ref{sec:tasks}) and the
        handler registration in the Interrupts chapter.
  \item \textbf{Hardware structures that reference a buffer:} a DMA descriptor
        holds the \emph{address} of the data to move, so the buffer is
        \texttt{aliased} and the descriptor stores its \texttt{'Access}
        (\S\ref{sec:placement}).
  \item \textbf{Talking to C:} C interfaces are expressed in pointers, so
        \texttt{System.Address} and access types cross that boundary
        (\S\ref{sec:interfacing}).
\end{itemize}

\subsection{Why the pointer you do use is safer}

Even when you need one, an Ada access type is far less hazardous than a C pointer:

\begin{itemize}
  \item it is \textbf{named and typed:} a \texttt{Node\_Ref} is not a
        \texttt{Byte\_Ref} and neither is a \texttt{void *}; the compiler keeps
        them apart;
  \item there is \textbf{no pointer arithmetic:} you index an array, you do not
        increment a pointer past its end;
  \item \textbf{accessibility rules} reject \texttt{Local'Access} returned from the
        subprogram that owns \texttt{Local}, so a dangling reference is a compile
        error, not a crash;
  \item \texttt{\textbf{not null}} access subtypes exclude the null case at the
        type level;
  \item and \texttt{\textbf{for T'Storage\_Size use 0;}} (\S\ref{sec:memory})
        forbids \texttt{new} on an access type outright: a compile-time promise
        that this pointer type never touches a heap, which is how the
        \texttt{light-tasking} profile stays heap-less.
\end{itemize}

So the rule of thumb: prefer value semantics and address overlays; reach for an
access type for the genuinely dynamic structure, the callback, or the
hardware-referenced buffer, and even then it is typed, arithmetic-free, and, if
you wish, provably heap-free.

\section{Contracts for hardware-safe APIs}
\label{sec:contracts}\index{contracts}

Finally, a driver's \emph{interface} can carry its rules as checkable contracts,
so misuse is caught at the call site rather than producing a wrong register write.
Preconditions and postconditions state what must hold on entry and exit; type
invariants and predicates state what must always hold of a value. The AES driver,
for instance, accepts only a legal key length:

\begin{lstlisting}
   function Supported_Key (Key : Byte_Array) return Boolean is
     (Key'Length in 16 | 32);            -- AES-128 or AES-256 only

   procedure Set_Key (Key : Byte_Array)
     with Pre => Supported_Key (Key);    -- a 24-byte key is rejected at the call
\end{lstlisting}

The cost is tunable: \texttt{pragma Assertion\_Policy} (and the \texttt{-gnata}
switch) turns the run-time checks on for development and \emph{off} for a release
build, where the contracts remain as machine-checked documentation. Combined with
the range and predicate subtypes of \S\ref{sec:subtypes}, the effect is that a
large class of driver-misuse bugs (a bad pin, a wrong-width key, an
out-of-range channel) cannot reach the hardware.

\section{Exceptions and error handling}
\label{sec:exceptions}\index{exceptions}

An Ada \emph{exception} is a named error condition: \texttt{raise} signals it and
a handler catches it. The mechanics are familiar:

\begin{lstlisting}
   begin
      Read_Sector (N, Buf);
   exception
      when Storage_Error =>           -- out of memory / stack
         Recover;
      when E : others =>              -- anything else; E names it
         Log (Ada.Exceptions.Exception_Message (E));
   end;
\end{lstlisting}

A raise propagates up the call chain to the nearest matching handler, running
finalizers (Chapter~\ref{ch:encap}) on the way, which is why a controlled
handle is released even when an exception unwinds past it. If no handler matches,
it keeps propagating; if it escapes the whole program, the runtime's
\emph{last-chance handler} runs (below).

\textbf{Predefined and user-defined exceptions.} The language defines a handful
the runtime raises for you: \texttt{Constraint\_Error} (a failed range, index,
\texttt{null} or overflow check, what a bad \texttt{Pin\_Id} or an out-of-range
index produces), \texttt{Program\_Error} (an illegal runtime situation),
\texttt{Storage\_Error} (out of stack or heap), and \texttt{Tasking\_Error}. You
declare your own just as easily, and raise it with a message:

\begin{lstlisting}
   Not_Owned : exception;          -- a named error condition of your own
   ...
   if not Holding then
      raise Not_Owned with "use the bus without holding it";
   end if;
\end{lstlisting}

A handler can inspect the exception through \texttt{Ada.Exceptions}
(\texttt{Exception\_Name}, \texttt{Exception\_Message},
\texttt{Exception\_Information}), and a bare \texttt{raise;} inside a handler
\emph{re-raises} the current one, so you can clean up locally and still let an
outer scope deal with it.

\textbf{How this project uses them.} The ownership-checked HAL gateways raise
\texttt{Not\_Owned} / \texttt{Not\_Initialized} when a session is used without
being held or set up; the AES driver's \texttt{Supported\_Key} precondition
(\S\ref{sec:contracts}) rejects a wrong-length key; the range and predicate
subtypes (\S\ref{sec:subtypes}) raise \texttt{Constraint\_Error} on a bad value
under \texttt{-gnata}; and the ext4 filesystem reports I/O and corruption errors
by raising. The lock-free drivers instead report failure through a \texttt{Status}
enum or an \texttt{out Boolean}, precisely so they work on \emph{every} profile,
which is the next point.

\textbf{The cost (and even the availability) depends on the profile.} This
is the embedded-specific part, and it differs across all three:

\begin{itemize}
  \item \texttt{light-tasking} is \texttt{No\_Exception\_Propagation}: there is no
        unwinder, so a \texttt{raise} not handled \emph{in the same subprogram}
        cannot propagate. In practice it reaches the last-chance handler and
        resets the board. The range and predicate checks still \emph{detect}
        errors; you simply cannot catch and recover across calls. Design for
        ``detect, log, reset,'' and report an expected failure with a status
        result, not an exception.
  \item \texttt{embedded} adds \emph{zero-cost} exceptions (ZCX): full propagation,
        handlers, \texttt{Ada.Exceptions} names and messages, and finalization on
        the way out (so an RAII handle releases during an unwind). A handler that
        is not entered costs nothing on the happy path: the cost lives in static
        tables, not in instructions. This is the minimum profile for the RAII
        drivers and the filesystem; here you catch and recover normally.
  \item \texttt{full} has the \emph{same} exception model as \texttt{embedded},
        plus the complete tasking runtime, so exceptions interoperate with dynamic
        tasks, rendezvous and \texttt{abort}, and a stack overflow caught by the
        runtime's guard can be delivered as a catchable \texttt{Storage\_Error}
        instead of a reset (see the Debugging chapter).
\end{itemize}

When an exception escapes everything (always on light-tasking, and as a final
backstop elsewhere) the runtime's \emph{last-chance handler} runs; on this port
it reports the fault and resets, and the \texttt{esp32s3\_safe\_state} example
shows a fault being classified there (see the Debugging chapter).

\textbf{Strategy on a microcontroller.} Prefer the cheap, local mechanisms to
\emph{prevent} bad values: range and predicate subtypes (\S\ref{sec:subtypes}),
\texttt{'Valid} (\S\ref{sec:valid}), contracts (\S\ref{sec:contracts}); use a
status result for an expected, recoverable condition (a NACK, a CRC mismatch); and
reserve exceptions for the genuinely exceptional, remembering that on the leanest
profile ``handling'' one means a controlled reset, not a resumption.

\textbf{A runnable demonstration.} The \texttt{esp32s3\_exceptions} example
exercises all four outcomes on the chip: local handling, propagation across a
call, a re-raise to an outer handler, and an unhandled exception reaching the
last-chance handler (it installs a custom one, exported as
\texttt{\_\_gnat\_last\_chance\_handler}, that prints the exception and halts).

\section{Restricting the language: \texttt{pragma Restrictions} and profiles}
\label{sec:restrictions}\index{Restrictions pragma}

The flip side of choosing features is forbidding them. \texttt{pragma
Restrictions} lets you ban constructs that have no place in tiny or certifiable
firmware and have the compiler \emph{prove} you stayed within the subset:

\begin{lstlisting}
   pragma Restrictions (No_Dependence => Ada.Unchecked_Deallocation);
   pragma Restrictions (No_Recursion, No_Secondary_Stack);
   pragma Restrictions (No_Exceptions);          -- if a raise would be fatal anyway
\end{lstlisting}

A named bundle of these is a \emph{profile}: \texttt{pragma Profile (Ravenscar)}
(and its superset \texttt{Jorvik}) selects the analysable tasking subset
(no dynamic task creation, no entry queues longer than one, a single delay form),
which is what the \texttt{light-tasking} and \texttt{embedded} runtimes are built
against. The Runtime Profiles chapter covers the runtime side; the point here is
that the restriction is a \emph{language} statement the compiler enforces unit by
unit, so a violation is a build error, not a surprise on silicon.

\section{Controlling elaboration: \texttt{Pure}, \texttt{Preelaborate}, \texttt{No\_Elaboration\_Code}}
\label{sec:elab}\index{elaboration}

Elaboration (\S\ref{sec:terms}) is the code that runs before \texttt{Main} to
bring declarations to life. On a microcontroller you usually want \emph{less} of
it: less start-up code, and data that can live in flash rather than be built in
RAM at boot. Category pragmas let you state (and have the compiler enforce)
how little a unit needs:

\begin{itemize}
  \item \texttt{pragma Pure:} the unit has no state and no elaboration at all
        (only pure functions and constants); the strongest promise, and its data
        can sit in flash.
  \item \texttt{pragma Preelaborate:} the unit may be elaborated without
        running any code: all its data has static initial values, so there is no
        boot-time cost and nothing is built at start-up.
  \item \texttt{pragma Elaborate\_Body:} forces the body to be elaborated
        immediately after the spec, a sequencing guarantee that keeps elaboration
        order predictable.
  \item \texttt{pragma Restrictions (No\_Elaboration\_Code):} the compiler
        \emph{rejects} any unit that would emit elaboration code: a hard guarantee
        that start-up is free.
\end{itemize}

\begin{lstlisting}
package Sensor_Tables is
   pragma Pure;                          -- pure: a flash-resident lookup table
   Lut : constant array (0 .. 255) of Integer := ( ... );
end Sensor_Tables;

package Ring is
   pragma Preelaborate;                  -- no code runs to elaborate this unit
   ...
end Ring;
\end{lstlisting}

\textbf{Embedded payoff.} A \texttt{Pure} or \texttt{Preelaborate} unit
contributes nothing to boot time and keeps its constants out of RAM; the runtime
itself builds its lowest layers with \texttt{No\_Elaboration\_Code} so the kernel
has no hidden start-up. The pragmas turn ``this should be cheap to start'' into a
checked property.


\chapter{Talking to Hardware}
\label{ch:hardware}

A peripheral register is a fixed memory location with side effects; a structure
read off storage or the wire is a foreign byte layout. Ada lets the \emph{type}
carry both -- the exact address, the bit positions, the volatility, the
endianness -- so once a register or structure is described correctly it is safe at
every use, with no pointer casts or shift-and-mask macros. This chapter is those
techniques; they build directly on the language features of
Chapter~\ref{ch:embedded}.

\section{Register Addressing in Ada}
\label{ch:registers}\index{register addressing}

A peripheral register is a fixed memory location with side effects. To program it
in Ada you describe \emph{three} things: its layout (a type), its location (an
address), and the fact that touching it is observable (volatility). You never
need a pointer cast.

\subsection{An object at a fixed address}

The simplest register is a 32-bit word. Place an Ada object at its address and
mark it \texttt{Volatile} so the compiler never caches, reorders, or elides
accesses:

\begin{lstlisting}
with System;
with Interfaces; use Interfaces;

procedure Toggle_Raw is
   --  The GPIO "output set" and "output clear" W1TS/W1TC registers.
   GPIO_OUT_W1TS : Unsigned_32
     with Volatile, Address => System'To_Address (16#6000_4008#);
   GPIO_OUT_W1TC : Unsigned_32
     with Volatile, Address => System'To_Address (16#6000_400C#);
begin
   GPIO_OUT_W1TS := 2 ** 2;   --  set GPIO2  (writing a 1 sets that bit)
   GPIO_OUT_W1TC := 2 ** 2;   --  clear GPIO2 (writing a 1 clears it)
end Toggle_Raw;
\end{lstlisting}

\texttt{System'To\_Address} turns a numeric literal into an address; the
\texttt{Address} aspect binds the object there. Because the type is exactly
\texttt{Unsigned\_32}, the write is exactly one 32-bit store -- the compiler
cannot ``optimise'' it into a byte write, which a peripheral might not accept.

\cnote{this replaces \texttt{*(volatile uint32\_t *)0x60004008 = x;}\ -- but with
no cast, and with the access width fixed by the type, so the compiler cannot
silently split or merge the access.}

\subsection{Generated register packages}

Writing every register by hand is error-prone, so the whole ESP32-S3 register map
is generated from the vendor SVD file into typed packages under
\texttt{ESP32S3\_Registers.*} (one package per peripheral). Most registers are
records of named bit-fields with a representation clause; a few -- like the
write-one-to-set ports -- are bare 32-bit words. Each peripheral is one big record
placed at its base address, and a driver writes them by name:

\begin{lstlisting}
with ESP32S3_Registers.GPIO; use ESP32S3_Registers.GPIO;

--  GPIO_Periph is the peripheral record, already placed at 16#6000_4000#.
GPIO_Periph.OUT_W1TS := 2 ** 2;          --  set GPIO2 high (write-1-to-set)
--  ... and a register that IS a field record is written field by field:
GPIO_Periph.PIN (2).INT_ENA := 1;        --  GPIO2's interrupt enable, by name
\end{lstlisting}

The application almost never touches these directly -- it calls the hand-written
HAL (\texttt{ESP32S3.GPIO} and friends) -- but the same techniques underpin every
driver.

\subsection{Array overlays}

Sometimes the SVD flattens a regular bank of registers (a FIFO, a descriptor
table) into separate fields. You can overlay your own array on top of them with
\texttt{Address}, recovering clean indexing:

\begin{lstlisting}
type Reg_Bank is array (0 .. 7) of Interfaces.Unsigned_32 with Volatile;
Bank : Reg_Bank with Import, Volatile, Address => Some_Periph.First_Reg'Address;
--  Bank (3) is now the fourth 32-bit register, addressed normally.
\end{lstlisting}

The \texttt{Import} aspect says ``this object already exists at that address; do
not allocate or initialise it.'' This idiom appears throughout the HAL where the
register layout is regular but the generator did not capture it as an array.

\section{Bit Fields and Record Representation}
\label{ch:bitfields}\index{representation clause}

Hardware registers pack many fields into one word: a 3-bit mode here, a single
enable bit there, a 12-bit divider somewhere else. In C this means shift-and-mask
macros that are easy to get wrong. In Ada you declare a record whose components
have the right types and then a \emph{representation clause} that nails each
component to specific bits. The compiler does the shifting and masking for you,
correctly, every time.

\subsection{A register as a record}

\begin{lstlisting}
type Mode_Bits is mod 2 ** 3;      --  a 3-bit field
type Divider   is mod 2 ** 12;     --  a 12-bit field

type Control_Register is record
   Enable   : Boolean;             --  1 bit
   Mode     : Mode_Bits;           --  3 bits
   Reserved : Boolean := False;
   Div      : Divider;             --  12 bits
end record
  with Volatile_Full_Access, Size => 32, Bit_Order => System.Low_Order_First;

for Control_Register use record
   Enable   at 0 range 0  .. 0;
   Mode     at 0 range 1  .. 3;
   Reserved at 0 range 4  .. 4;
   Div      at 0 range 5  .. 16;
end record;
\end{lstlisting}

Now you assign fields by name and the layout is guaranteed:

\begin{lstlisting}
Ctrl := (Enable => True, Mode => 2, Div => 1023, others => <>);
\end{lstlisting}

\texttt{at 0 range L\,..\,H} means ``bits $L$ through $H$ of the byte at offset
0.'' \texttt{Volatile\_Full\_Access} guarantees the whole 32-bit word is read or
written in one access -- essential for registers that latch on a full-word
write. This is precisely how the generated \texttt{ESP32S3\_Registers} packages
describe every register, and how the HAL builds register values:

\begin{lstlisting}
--  From the real SPI driver: build a USER register value by field, in one write.
Regs.USER := (DOUTDIN => True, USR_MOSI => True, USR_MISO => True,
              CK_OUT_EDGE => Out_Edge, others => <>);
\end{lstlisting}

\subsection{Enumerations with explicit codes}

When a field's values are an enumeration, give the enumeration an explicit
representation so the symbolic names map to the exact hardware codes:

\begin{lstlisting}
type Parity is (None, Even, Odd);
for Parity use (None => 0, Even => 2, Odd => 3);
\end{lstlisting}

\subsection{Scalar storage order: decoding on-disk and on-wire data}

The same machinery decodes \emph{external} byte layouts, not just registers. A
record can declare its \texttt{Scalar\_Storage\_Order}, and GNAT will read each
field with the correct endianness automatically -- invaluable for filesystem
and protocol structures. The ext4 filesystem in Chapter~\ref{ch:storage} relies
on exactly this idea (though it ultimately uses explicit little-/big-endian
readers for clarity), and JBD2's big-endian journal records show why being able
to say ``this structure is big-endian'' in the type is so useful.

\section{Volatile, Atomic, and full-word access}
\label{sec:volatile}\index{Volatile}

Placing an object at a register's address is only half the story; you must also
tell the compiler how that location behaves, because the optimiser otherwise
assumes ordinary memory -- caching reads, dropping ``dead'' writes, reordering.
A small family of aspects states the truth:

\begin{itemize}
  \item \texttt{Volatile} -- every read and write happens, in program order, and
        none is cached or elided. The baseline for any register.
  \item \texttt{Atomic} -- additionally, each access is indivisible (a single
        load or store). Use it for a flag shared with an interrupt handler or the
        other core, where a torn half-update would be a bug.
  \item \texttt{Volatile\_Full\_Access} -- the whole object is read or written in
        \emph{one} access of its size. Essential for a register that latches on a
        full-word write or clears on read: updating one bit-field still emits a
        single 32-bit access, not a byte poke.
  \item \texttt{Volatile\_Components} / \texttt{Atomic\_Components} -- the same,
        applied to each element of an array. The runtime's per-core
        \texttt{Switch\_Pending} and \texttt{In\_Native\_Int} flags are
        \texttt{Volatile\_Components} arrays for exactly this reason.
\end{itemize}

\begin{lstlisting}
   Ready : Boolean with Atomic, Address => Status_Reg'Address;  -- ISR-shared flag
   --  Control_Register from the bit-field section is Volatile_Full_Access:
   --  Ctrl.Enable := True  emits one 32-bit store, not a read-modify-byte-write.
\end{lstlisting}

Pick the weakest that is correct: \texttt{Volatile} for most registers,
\texttt{Volatile\_Full\_Access} for latch-on-write ones, \texttt{Atomic} only
where indivisibility actually matters (it can be more expensive).

\section{Packing: \texttt{pragma Pack} and packed arrays}
\label{sec:pack}\index{packed arrays}

The bit-field record nails \emph{named} fields to bits. When you instead have a
homogeneous run of small components -- a bit vector, a table of 2-bit codes --
\texttt{pragma Pack} (or a \texttt{Component\_Size} clause) tells the compiler to
store them at their natural bit width instead of rounding each up to a byte or
word:

\begin{lstlisting}
   type Pin_Mask is array (0 .. 31) of Boolean;
   pragma Pack (Pin_Mask);          -- 32 bits total, one bit per pin
   --  or, explicitly:
   type Codes is array (0 .. 15) of Mode_Bits with Component_Size => 2;  -- 32 bits
\end{lstlisting}

A packed \texttt{array (...) of Boolean} is the idiomatic Ada bitmap: you index
it by position (\texttt{Mask (Pin) := True}) and the compiler does the
shift-and-mask, while the whole thing still occupies one word you can overlay on
a register. Reach for the named record when the fields differ; reach for a packed
array when they are uniform and indexed.

\section{Placing objects in memory}
\label{sec:placement}\index{Linker\_Section}

On a microcontroller \emph{where} an object lives is part of its meaning: some
data must sit in fast IRAM, some in external PSRAM, some in RTC memory that
survives deep sleep, and some must \emph{not} be cleared at startup. Two aspects
cover it:

\begin{itemize}
  \item \texttt{Linker\_Section => "..."} places an object in a named output
        section the linker script positions (e.g.\ an IRAM or \texttt{.noinit}
        section).
  \item an \texttt{Address} overlay (\S\ref{ch:registers}) pins it at a fixed
        location -- which is how the RTC-retained boot counter in the deep-sleep
        example lives in RTC slow memory and survives a reset:
\end{itemize}

\begin{lstlisting}
   Boot_Count : Unsigned_32 with Volatile,
     Address => System'To_Address (16#5000_0000#);   -- RTC slow RAM, retained

   IRAM_Buf : Buffer with Linker_Section => ".iram1.data";
\end{lstlisting}

Alignment matters too: \texttt{with Alignment => 64} forces DMA buffers onto a
cache-line boundary, and the GDMA driver relies on aligned, internal-RAM (not
PSRAM) descriptors -- placement and alignment together.

\begin{table}[ht]
\centering
\begin{tabular}{@{}l l l@{}}
\toprule
\textbf{Base} & \textbf{Top} & \textbf{Region} \\ \midrule
\texttt{16\#3C00\_0000\#} & \texttt{16\#3D00\_0000\#} & Flash $\to$ DROM -- cached read-only data \\
\texttt{16\#3D00\_0000\#} & \texttt{16\#3E00\_0000\#} & External PSRAM -- mapped by the bootloader \\
\texttt{16\#3FC8\_8000\#} & \texttt{16\#3FD0\_0000\#} & Internal SRAM (DRAM view) -- data, stacks, heap \\
\texttt{16\#4037\_0000\#} & \texttt{16\#403E\_0000\#} & Internal SRAM (IRAM view) -- code, vectors, DMA \\
\texttt{16\#4200\_0000\#} & \texttt{16\#4400\_0000\#} & Flash $\to$ IROM -- cached code \\
\texttt{16\#5000\_0000\#} & \texttt{16\#5000\_2000\#} & RTC slow RAM -- retained across deep sleep \\
\bottomrule
\end{tabular}
\caption{Principal address-space regions, by ascending base address. Flash
(DROM) and external PSRAM are not separate windows: they are two sub-ranges of
one 32\,MB external data-bus window (\texttt{16\#3C00\_0000\#}--\texttt{16\#3E00%
\_0000\#}) that the cache MMU partitions -- this project maps flash from the
bottom and PSRAM at \texttt{16\#3D00\_0000\#}, the PSRAM's actual extent set by
the fitted device. Placement pins objects where they belong: RTC memory for
state that survives sleep, internal SRAM for DMA buffers (PSRAM cannot be a DMA
source), flash for constants.}
\label{tab:memmap}
\end{table}

\section{Controlling memory: bounded allocation}
\label{sec:memory}\index{Storage\_Size}

Embedded firmware usually wants allocation to be \emph{bounded} or absent. Ada
makes the size explicit rather than leaving it to a global heap:

\begin{itemize}
  \item \texttt{for T'Storage\_Size use 0;} on an access type forbids
        \texttt{new} of \texttt{T} -- a compile-time guarantee that pointer type
        never heap-allocates.
  \item A task can cap its own stack: \texttt{task body Worker} with
        \texttt{pragma Storage\_Size} (or the \texttt{Storage\_Size} aspect) sizes
        the primary stack precisely.
  \item The binder caps the secondary stacks (\texttt{-D<n>}, \texttt{-Q<n>}), and
        the \texttt{light-tasking} profile removes the heap and secondary stack
        entirely (a bump allocator only).
\end{itemize}

\begin{lstlisting}
   type Node_Ref is access Node with Storage_Size => 0;  -- new Node => compile error
   task Sampler with Storage_Size => 4 * 1024;           -- a 4 KB primary stack
\end{lstlisting}

Sizing memory in the source -- not discovering it when the heap is exhausted at
3 a.m.\ in the field -- is the embedded mindset the language supports directly.

\section{Interfacing with C, ROM, and assembly}
\label{sec:interfacing}\index{interfacing!with C}

A bare-metal stack still meets non-Ada code: a ROM routine, a C glue function, a
few instructions only assembly can express. Ada's interfacing pragmas make those
boundaries explicit and typed.

\textbf{Calling out} -- \texttt{pragma Import} (or the \texttt{Import} aspect)
with \texttt{Convention => C} binds an Ada declaration to an external symbol; no
header, no wrapper:

\begin{lstlisting}
   procedure Rom_Printf (Fmt : System.Address)
     with Import, Convention => C, External_Name => "esp_rom_printf";
\end{lstlisting}

\textbf{Being called} -- \texttt{pragma Export} publishes an Ada subprogram or
object under a C name so the boot assembly or glue can reach it (the runtime
exports its interrupt-dispatch entries and the per-core flags this way).
\texttt{Interfaces.C} (and \texttt{.Strings}, \texttt{.Pointers}) provides the
C-compatible scalar and array types for the data that crosses.

\textbf{Inline assembly} -- for the handful of operations with no Ada spelling
(a special register read, a cache or barrier instruction),
\texttt{System.Machine\_Code.Asm} embeds Xtensa instructions with typed operands;
the GDMA driver uses it for a memory barrier around descriptor hand-off. It is a
scalpel, not a habit: everything that can be Ada, is.

\section{The thread through it all}

One idea runs through this chapter and the last: in Ada the \emph{type} -- with
its layout, its range, its volatility, its contracts -- carries the intent, so the
compiler (or a cheap, switch-off-able run-time check) catches the mistake. Describe
a register, a value, or an interface correctly once, and every use of it afterwards
is safe by construction.

\chapter{Encapsulation, Ownership and Concurrency}
\label{ch:encap}

Four Ada features let an interface control, in the type system itself, what client
code is allowed to do, and together they make a peripheral driver that is
misuse-resistant by construction. \texttt{private} hides a type's representation;
\texttt{limited} forbids copying it; \texttt{controlled} gives it automatic setup
and teardown; \texttt{protected} serialises access to it. None costs more at run
time than the discipline you would otherwise have to write and remember by hand,
and each turns a class of ``be careful to\ldots'' rules into a compiler-enforced
guarantee. This chapter explains each keyword with small examples and its embedded
payoff; the \emph{Writing a Task-Safe Peripheral Driver} chapter shows them
combined in the real HAL.

\section{\texttt{private} types: hiding the representation}\index{private types}

A private type is declared in two halves. The visible part of a package gives the
type a name and the operations on it; the \emph{private part} (everything after
the word \texttt{private}) gives its real definition. Client code may use the
name and call the operations, but cannot see or touch the components.

\begin{lstlisting}
package Counters is
   type Counter is private;             -- clients see only this and the ops
   procedure Bump  (C : in out Counter);
   function  Value (C : Counter) return Natural;
private
   type Counter is record               -- the real definition, hidden
      N : Natural := 0;
   end record;
end Counters;
\end{lstlisting}

A user can declare a \texttt{Counter}, \texttt{Bump} it and read its
\texttt{Value}, but cannot write \texttt{C.N}, cannot rely on it being a record,
and will not break if the representation later changes. This is the difference
between a \texttt{struct} in a public C header (whose fields anyone may scribble
on) and a genuinely opaque handle.

\textbf{Embedded payoff.} The representation of a driver handle (which host it
owns, a DMA-descriptor pointer, a hardware index) is hidden, so application code
cannot corrupt it and the driver author stays free to change it.

\section{\texttt{private} child packages: sealing the unsafe layer}

Visibility extends to whole packages. A child package declared \texttt{private} is
visible only inside its own family's bodies; application code cannot even
\emph{name} it in a \texttt{with} clause. The HAL uses this to seal the only code
that touches registers:

\begin{lstlisting}
private package ESP32S3.SPI.Engine is   -- un-with-able from an application
   --  the raw, lock-free register driver: the ONLY code that pokes SPI.
   ...
end ESP32S3.SPI.Engine;
\end{lstlisting}

The public \texttt{ESP32S3.SPI} package \texttt{with}s the engine in its body and
mediates every operation. \textbf{Embedded payoff.} ``Poke the registers
directly'' is not merely discouraged: it is \emph{unwritable}. A
\texttt{with ESP32S3.SPI.Engine;} in an application fails to compile, so the unsafe
core is reachable only through the safe fa\c{c}ade.

\section{\texttt{limited} types: values that cannot be copied}\index{limited type}

A limited type has no assignment (\texttt{:=}) and no predefined equality
(\texttt{=}). You can declare one, pass it (by reference) and operate on it, but
you cannot copy it. Make a type limited when a value stands for a \emph{unique}
thing that must not be duplicated: a handle to a single hardware resource, a
lock, an open session:

\begin{lstlisting}
   type Session is limited private;     -- a SPI bus session
   ...
   A, B : Session;
   B := A;       -- COMPILE ERROR: Session is limited, it cannot be copied
\end{lstlisting}

If you \emph{could} copy a handle, two copies would each believe they own the bus;
releasing one would leave the other dangling. \texttt{limited} makes that bug
impossible to write. It is the real shape of the SPI \texttt{Session}, the GDMA
\texttt{Channel}, the ADC \texttt{Reader}, the LEDC/PCNT \texttt{Channel} and the
rest of the HAL's ownership handles.

\section{\texttt{controlled} types and RAII: deterministic teardown}\index{controlled type}\index{RAII}

\texttt{Ada.Finalization.Controlled} (and its non-copyable sibling
\texttt{Limited\_Controlled}) give a type three hooks the runtime calls
\emph{automatically}: \texttt{Initialize} when an object comes into being,
\texttt{Finalize} when it leaves scope (\emph{including} when an exception unwinds
the scope), and \texttt{Adjust} after a copy (for the non-limited kind). That is
RAII: a resource's lifetime is tied to an object's scope.

\begin{lstlisting}
   type Session is new Ada.Finalization.Limited_Controlled with record
      Host : ...;
   end record;
   overriding procedure Finalize (S : in out Session);   -- auto-release the bus
\end{lstlisting}

\begin{lstlisting}
   declare
      S : Session;                  -- Initialize
   begin
      Acquire (S, Host => SPI2);    -- take the bus
      Transfer (S, ...);
   end;                             -- Finalize runs HERE: the bus is released,
                                    -- even if Transfer raised an exception
\end{lstlisting}

\textbf{Embedded payoff.} You \emph{cannot forget} to release a DMA channel,
unlock a bus, or power down a peripheral: the release is bound to the scope and
happens on every exit path, with no leak, deterministic timing, and no
\texttt{goto cleanup} ladder. \texttt{limited} + \texttt{controlled} is the
project's standard handle: limited stops you copying the owner, controlled
guarantees it is released. (The real SPI \texttt{Session}'s overriding
\texttt{Finalize} is commented ``auto-release on scope exit'' for exactly this.)

\cnote{ESP-IDF hands you an opaque handle and a matching \texttt{*\_delete} /
\texttt{*\_free} you must remember to call on every path (including the error
ones). A controlled \texttt{Session} releases itself when it leaves scope (even
when an exception unwinds past it) so the leak and the missed-cleanup bug
cannot happen.}

\section{\texttt{protected} types and objects: race-free sharing}\index{protected object}

A protected object is data plus operations with mutual exclusion built in: a
monitor. Only one task at a time runs a protected \emph{procedure} or
\emph{entry}; the compiler generates the locking, and on this runtime the lock is
a \emph{priority ceiling}, so worst-case blocking is bounded (real-time
friendly). There are three kinds of operation:

\begin{itemize}
  \item a protected \textbf{function} reads the state and may run concurrently
        with other functions (a shared read);
  \item a protected \textbf{procedure} reads and writes with exclusive access
        (one caller at a time);
  \item a protected \textbf{entry} is a procedure with a Boolean \emph{barrier};
        a caller blocks until the barrier is open (condition synchronisation).
\end{itemize}

\begin{lstlisting}
   protected Counter is
      procedure Bump;                  -- exclusive write
      function  Value return Natural;  -- shared read
   private
      N : Natural := 0;
   end Counter;

   protected body Counter is
      procedure Bump is begin N := N + 1; end Bump;
      function  Value return Natural is (N);
   end Counter;
\end{lstlisting}

\texttt{Counter.Bump} from one task and \texttt{Counter.Value} from another never
race: no \texttt{Volatile}, no hand-rolled critical section. The HAL uses
exactly this to arbitrate a shared peripheral: a per-port \texttt{protected Guard}
(or \texttt{Pool}, or \texttt{Lock}) hands out ownership so two tasks cannot drive
the same controller at once, and an interrupt handler is itself a protected
procedure (the Interrupts chapter). \textbf{Embedded payoff.} Shared state across
tasks, cores and interrupt handlers is correct by construction, with bounded,
analysable blocking.

\section{The four together: a misuse-proof driver}

Combined, these keywords make a peripheral that is hard to misuse even on purpose:

\begin{itemize}
  \item the registers live in a \texttt{private} child: you cannot see or poke
        them;
  \item the public handle is \texttt{limited}: you cannot copy ownership;
  \item it is \texttt{controlled}: you cannot forget to release it;
  \item access is \texttt{protected}: you cannot race it.
\end{itemize}

That is precisely the HAL's task-safe driver skeleton; the \emph{Writing a
Task-Safe Peripheral Driver} chapter builds one end to end (Engine $+$ Session
$+$ Guard). The payoff is that a whole class of embedded bugs (a dangling
handle, a forgotten unlock, a double-driven bus, a torn shared counter) becomes
a compile error or a structural impossibility instead of a field failure.

\section{Related keywords}

A few neighbours round out the family. Each is shown with a small example.

\paragraph{\texttt{constant}: a value that never changes.} A constant is given
its value once and can never be assigned again; the compiler may place it in
read-only memory (flash), saving RAM. A constant with no type (a \emph{named
number}) is a pure compile-time value.

\begin{lstlisting}
   Max_Pins : constant := 48;                 -- named number (no type, no storage)
   Sine     : constant array (0 .. 255) of Integer := ( ... );  -- table, in flash
   --  Sine (K) := 0;  -- COMPILE ERROR: Sine is constant
\end{lstlisting}

In embedded code lookup tables, calibration data and configuration are
\texttt{constant} so they live in flash, not in scarce RAM.

\paragraph{\texttt{aliased}: an object you may point at.} By default you cannot
take \texttt{'Access} of a declared object (the accessibility rules guard against
dangling pointers). Marking it \texttt{aliased} permits it, needed when, say, a
DMA buffer must be handed to a descriptor by reference.

\begin{lstlisting}
   Buf  : aliased DMA_Buffer;          -- 'Access may now be taken
   Desc.Src := Buf'Access;             -- give the DMA engine a pointer to it
\end{lstlisting}

\paragraph{\texttt{renames}: a new name for an existing thing.} A renaming
declaration introduces an \emph{alias}: a second name that denotes the \emph{same}
entity, with no copy and no new object made. It works for objects, components,
subprograms, packages and exceptions.

\begin{lstlisting}
   --  alias a deeply-nested register field: short name, same storage
   Ctrl : Unsigned_32 renames Periph.SPI2.User.Control;
   Ctrl := 16#1#;                      -- writes Periph.SPI2.User.Control

   --  alias a package to shorten a long path (just this one name, unlike `use`)
   package GE renames ESP32S3.GPIO;
   GE.Toggle (2);

   --  alias a subprogram -- give an operation a local or operator name
   function "+" (L, R : Vector) return Vector renames Vectors.Add;
\end{lstlisting}

For an object or component the alias shares the original's storage (writing
through \texttt{Ctrl} above writes the real register field), so \texttt{renames}
is the way to name one specific deeply-qualified thing once and use the short name
afterwards. Note how it differs from \texttt{use} (\S\ref{ch:registers}):
\texttt{use} makes \emph{all} of a package's names directly visible, whereas
\texttt{renames} introduces exactly \emph{one} new name and says precisely what it
stands for, so it shortens a single long reference without the wholesale
visibility (and possible name clashes) that \texttt{use} brings.

\paragraph{\texttt{tagged}: records with inheritance and dispatch.} A tagged
type can be \emph{extended} (a derived type adds components) and its primitive
operations can \emph{dispatch} at run time: a call through a class-wide value
\texttt{T'Class} selects the right override for the actual object.

\begin{lstlisting}
   type Sensor is tagged record Id : Natural; end record;
   function Read (S : Sensor) return Integer;

   type Temp_Sensor is new Sensor with record Pin : Pin_Id; end record;  -- extends
   overriding function Read (S : Temp_Sensor) return Integer;            -- override
\end{lstlisting}

A call \texttt{Read (S)} where \texttt{S} is \texttt{Sensor'Class} dispatches to
\texttt{Temp\_Sensor}'s \texttt{Read} if that is the actual type. It is used
sparingly in embedded code (the \texttt{embedded}-profile demo shows library-level
tagged dispatch) because a dispatching call is an indirect call.

\paragraph{\texttt{abstract}: a type or operation with no implementation.} An
abstract type cannot be instantiated, and an abstract operation has no body:
every concrete descendant \emph{must} provide one. This is how Ada expresses an
interface: a contract that says ``every sensor shall have a \texttt{Read}.''

\begin{lstlisting}
   type Sensor is abstract tagged null record;
   function Read (S : Sensor) return Integer is abstract;  -- no body; must override
   --  S : Sensor;  -- COMPILE ERROR: Sensor is abstract
\end{lstlisting}

These are reached for less often than \texttt{private} /\texttt{limited} /
\texttt{protected} /\texttt{controlled}, but knowing them completes the picture of
how Ada lets the \emph{declaration} of a type state what may -- and may not -- be
done with it.

\chapter{The Peripheral HAL}
\label{ch:hal}

The hand-written drivers live under \texttt{libs/esp32s3\_hal} as children of a
single root package, \texttt{ESP32S3}: \texttt{ESP32S3.GPIO},
\texttt{ESP32S3.SPI}, and so on. A consumer adds one line to its project file,

\begin{lstlisting}
with "esp32s3_hal.gpr";   --  or a relative path in-repo
\end{lstlisting}

and then \texttt{with}s the drivers it needs. Every driver in this book has been
exercised on real silicon.

\section{Conventions used by every driver}

\begin{description}[style=nextline]
\item[Valid-pin types.] A pin is named with \texttt{ESP32S3.GPIO.Pin\_Id}, a
  subtype whose \emph{static predicate} admits only the pads the ESP32-S3 exposes
  (0--21 and 38--48); the flash/PSRAM pads are excluded. A bad constant pin is a
  compile-time error.
\item[Errors.] Simple drivers report failure through a \texttt{Status} enum or an
  \texttt{out Boolean}; the filesystem and richer layers raise exceptions. Either
  way, failure is explicit.
\item[Task-safety, two shapes.] A peripheral that is a single shared controller
  (SPI host, I2C host, UART port, I2S port, the TWAI/LCD controller) is mediated
  by a protected object plus a limited, non-copyable \texttt{Session} handle: you
  \texttt{Acquire} it, use it, and it is released automatically when it leaves
  scope. A peripheral that is a \emph{pool} of identical units (GDMA channels,
  LEDC/RMT/PCNT/SDM channels, timers, ADC units, MCPWM channels) hands out a
  limited, controlled handle you \texttt{Claim}; ownership makes per-handle
  operations safe without a global lock, and the unit is reclaimed on scope exit.
\item[RAII.] Those \texttt{Session}/\texttt{Channel}/\texttt{Reader} handles are
  controlled types -- they release on scope exit, including during exception
  unwinding -- so they require the \texttt{embedded} or \texttt{full} profile.
  The finalization-free drivers (GPIO, RNG, Temperature, RTC, RTC\_IO, Touch, SHA,
  AES, SDMMC) use no controlled types -- a few serialise internally with a protected
  object instead -- so they work under every profile.
\item[Init-before-use guard.] For the drivers that configure a controller once
  at startup and only later hand out ownership (SPI, I2C, UART, I2S, TWAI, LCD,
  and MCPWM), \texttt{Acquire}/\texttt{Claim} raises \texttt{Not\_Initialized} if
  the port/host/unit was never \texttt{Setup}. Configuration must precede
  ownership; the contract is enforced, not merely documented.
\item[Ownership gateway.] Across all six \texttt{Session} drivers (SPI, I2C,
  UART, I2S, TWAI, LCD), all register access lives in a private \texttt{Engine}
  child, and the configured handle it returns is hidden in a nested package body.
  The only way to reach it is one ownership-checked gateway, \texttt{Owned (S)},
  which raises \texttt{Not\_Owned} unless the caller holds the port. Every
  operation that touches a live controller -- transfers \emph{and} every
  post-\texttt{Setup} reconfiguration (pins, baud, inversion, loopback) -- goes
  through it. So a transfer or a setting change without ownership fails loudly,
  and a new operation cannot be added that skips the check: there is no other
  door to the hardware. This is a project rule for any exclusive driver that can
  be reconfigured after setup (see the UART section for the worked pattern).
\end{description}

\paragraph{Recurring parameter types.} Each driver below lists its parameters
with their allowed values and defaults. A few types recur throughout and are
described once here rather than on every driver:
\begin{description}[style=nextline]
\item[\texttt{Pin\_Id}] a usable GPIO pad -- \texttt{0\,..\,21} or
  \texttt{38\,..\,48} (the gaps are flash/PSRAM/absent pads). A bad constant is a
  compile error.
\item[\texttt{Optional\_Pin}] a \texttt{Pin\_Id} \emph{or} \texttt{No\_Pin} (the
  value $-1$) to leave that line unrouted.
\item[Ownership handles (\texttt{Session}, \texttt{Channel}, \texttt{Reader},
  \texttt{Unit}, \texttt{Timer})] the limited, controlled handles described above:
  you declare one, \texttt{Acquire}/\texttt{Claim} it, pass it to each operation,
  and it releases on scope exit.
\item[\texttt{...\_Hz}, \texttt{Freq}, \texttt{Baud}, \texttt{Bit\_Rate}]
  frequencies / rates in hertz (or bits per second); \texttt{Positive} unless a
  narrower subtype is noted.
\item[\texttt{...\_Percent}] a \texttt{Float} in \texttt{0.0\,..\,100.0}.
\item[Buffers] a transfer takes a buffer \texttt{'Address} and a byte
  \texttt{Length}; the DMA-backed drivers cap a single transfer at \texttt{4095}
  bytes.
\end{description}

\section{GPIO -- digital input and output}\index{GPIO driver}

\texttt{ESP32S3.GPIO} configures a pad's direction, pull and drive strength and
then drives or reads it. \texttt{Set}/\texttt{Clear}/\texttt{Write}/\texttt{Read}
are hardware-atomic (lock-free); \texttt{Configure}/\texttt{Toggle} are serialised
by a protected object.

\begin{lstlisting}
procedure Configure (Pin : Pin_Id; Mode : Pin_Mode;
                     Pull  : Pull_Mode      := Floating;
                     Drive : Drive_Strength := Drive_Medium);
procedure Set    (Pin : Pin_Id);   procedure Clear (Pin : Pin_Id);
procedure Toggle (Pin : Pin_Id);   procedure Write (Pin : Pin_Id; On : Boolean);
function  Read   (Pin : Pin_Id) return Boolean;
\end{lstlisting}

\textbf{Parameters.}
\begin{description}[style=nextline]
\item[\texttt{Pin}] the pad to act on (\texttt{Pin\_Id}).
\item[\texttt{Mode}] direction: \texttt{Input} or \texttt{Output}.
\item[\texttt{Pull}] internal resistor: \texttt{Floating} (default), \texttt{Pull\_Up}
  or \texttt{Pull\_Down}.
\item[\texttt{Drive}] output current: \texttt{Drive\_Weak}, \texttt{Drive\_Medium}
  (default), \texttt{Drive\_Strong}, \texttt{Drive\_Strongest} (roughly
  5 / 10 / 20 / 40\,mA).
\item[\texttt{On}] (\texttt{Write}) the level to drive -- \texttt{True} = high,
  \texttt{False} = low.
\end{description}

\texttt{Configure} always sets the pad up as a software-controlled GPIO (routed
through the GPIO matrix). Handing a pad to a \emph{peripheral} instead is the job
of that peripheral's own \texttt{Configure\_Pins}, which programs the matrix
directly -- the ESP32-S3 GPIO matrix is a full crossbar, so almost any signal
can reach almost any pad.

\textbf{Example 1 -- blink an LED.}
\begin{lstlisting}
with Ada.Real_Time; use Ada.Real_Time;
with ESP32S3.GPIO;  use ESP32S3.GPIO;
procedure Blink is
   Led : constant Pin_Id := 2;
begin
   Configure (Led, Mode => Output);
   loop
      Toggle (Led);
      delay until Clock + Milliseconds (250);
   end loop;
end Blink;
\end{lstlisting}

\textbf{Example 2 -- a button drives an LED.}
\begin{lstlisting}
with ESP32S3.GPIO; use ESP32S3.GPIO;
procedure Button is
   Button_Pin : constant Pin_Id := 9;    --  active-low, internal pull-up
   Led        : constant Pin_Id := 2;
begin
   Configure (Button_Pin, Mode => Input,  Pull => Pull_Up);
   Configure (Led,        Mode => Output);
   loop
      Write (Led, not Read (Button_Pin)); --  pressed (low) -> LED on
   end loop;
end Button;
\end{lstlisting}

\section{RNG -- hardware random numbers}\index{RNG driver}

\begin{lstlisting}
function Read return Word;                 --  a 32-bit random word
procedure Fill (Buffer : out Byte_Array);  --  fill any length
\end{lstlisting}

\textbf{Parameters.} \texttt{Read} takes none and returns a \texttt{Word} (a 32-bit
\texttt{Interfaces.Unsigned\_32}). \texttt{Fill}'s \texttt{Buffer} is an
\texttt{out Byte\_Array} of \emph{any} length -- the call fills every element with
random bytes.

\textbf{Example 1 -- a random word.}
\begin{lstlisting}
with Ada.Text_IO; use Ada.Text_IO;
with ESP32S3.RNG;
procedure RNG_Word is
   X : constant ESP32S3.RNG.Word := ESP32S3.RNG.Read;
begin
   Put_Line (ESP32S3.RNG.Word'Image (X));
end RNG_Word;
\end{lstlisting}

\textbf{Example 2 -- a random nonce.}
\begin{lstlisting}
with ESP32S3.RNG;
procedure RNG_Nonce is
   Nonce : ESP32S3.RNG.Byte_Array (0 .. 15);
begin
   ESP32S3.RNG.Fill (Nonce);   --  16 random bytes
end RNG_Nonce;
\end{lstlisting}

(For true entropy the radio/ADC clocks should be running; see the package spec.)

\section{Temperature -- the on-die sensor}\index{Temperature driver}

\begin{lstlisting}
procedure Initialize (Span : Measure_Range := Range_Minus10_80);
function Read_Celsius return Integer;        --  whole degrees
function Read_Centi_Celsius return Integer;  --  hundredths of a degree
\end{lstlisting}

\textbf{Parameters.} \texttt{Initialize}'s \texttt{Span} (a \texttt{Measure\_Range})
picks the calibrated window -- accuracy is best when it brackets the real
temperature. The five values are \texttt{Range\_Minus10\_80} (default, $-10\,..\,80$\,\textdegree C),
\texttt{Range\_20\_100}, \texttt{Range\_50\_125}, \texttt{Range\_Minus30\_50}
($-30\,..\,50$) and \texttt{Range\_Minus40\_20} ($-40\,..\,20$). The two
\texttt{Read\_*} functions take no parameters; one returns whole degrees, the
other hundredths.

\textbf{Example 1 -- read degrees.}
\begin{lstlisting}
with Ada.Text_IO;        use Ada.Text_IO;
with ESP32S3.Temperature; use ESP32S3.Temperature;
procedure Temp_Degrees is
begin
   Initialize;
   Put_Line ("die temp =" & Integer'Image (Read_Celsius) & " C");
end Temp_Degrees;
\end{lstlisting}

\textbf{Example 2 -- higher resolution and a different span.}
\begin{lstlisting}
with Ada.Text_IO;        use Ada.Text_IO;
with ESP32S3.Temperature; use ESP32S3.Temperature;
procedure Temp_Centi is
begin
   Initialize (Span => Range_50_125);          --  a hotter measurement range
   declare
      Centi : constant Integer := Read_Centi_Celsius;
   begin
      Put_Line ("temp =" & Integer'Image (Centi / 100) & "." &
                Integer'Image (Centi mod 100) & " C");
   end;
end Temp_Centi;
\end{lstlisting}

\section{SPI -- full-duplex master}\index{SPI driver}

\texttt{ESP32S3.SPI} drives either general-purpose host (\texttt{SPI2}/\texttt{SPI3})
over GDMA. Bring the host up and route its shared wires once, then \texttt{Acquire}
a \texttt{Session} per transaction -- which is where each device's own SPI mode and
bit clock are applied, so two devices on one host can differ. \texttt{Transfer}
shifts \texttt{Tx} out while capturing \texttt{Rx}.

\begin{lstlisting}
procedure Setup (Host : SPI_Host);                          --  controller + GDMA
procedure Configure_Pins (Host : SPI_Host; Sclk, Mosi, Miso : Optional_Pin;
                          Cs : Optional_Pin := No_Pin);      --  shared wiring
procedure Acquire  (S : in out Session; Host : SPI_Host;
                    Mode      : SPI_Mode := 0;               --  per-device
                    Clock_Hz  : Positive := 1_000_000;       --  per-device
                    Sclk, Mosi, Miso : Optional_Pin := No_Pin; --  pin override
                    CS_Pin    : Optional_Pin := No_Pin;      --  built-in sw select
                    Select_CB : CS_Select    := null;        --  or app chip select
                    Ctx : System.Address := System.Null_Address);
procedure Select_Device (S : in out Session; On : Boolean);
procedure Transfer (S : Session; Tx, Rx : System.Address; Length : Natural);
procedure Release  (S : in out Session);
\end{lstlisting}

The three host calls split by lifetime and scope, the same shape as the other
buses: \texttt{Setup} (once, controller + GDMA), \texttt{Configure\_Pins} (once,
the shared \texttt{Sclk}/\texttt{Mosi}/\texttt{Miso} every device uses), and
\texttt{Acquire} (per transaction, this device's mode/clock/select). SPI's
\texttt{Setup} carries no mode or clock -- unlike, say, \texttt{I2C.Setup} --
precisely because on SPI those are per-\emph{device}, not per-host.

\textbf{Parameters.}
\begin{description}[style=nextline]
\item[\texttt{Host}] which general-purpose host: \texttt{SPI2} or \texttt{SPI3}.
\item[\texttt{Sclk}, \texttt{Mosi}, \texttt{Miso}, \texttt{Cs}] (\texttt{Configure\_Pins})
  the shared bus lines as \texttt{Optional\_Pin}; \texttt{Cs} is the host's single
  hardware chip-select and defaults to \texttt{No\_Pin} -- with more than one device
  on the bus, give each its own select via \texttt{CS\_Pin} at \texttt{Acquire}.
\item[\texttt{Mode}] (\texttt{Acquire}) this device's SPI clock mode \texttt{0\,..\,3}
  (the CPOL/CPHA combinations); default \texttt{0}, applied under the hold.
\item[\texttt{Clock\_Hz}] (\texttt{Acquire}) this device's bit clock in Hz; default
  \texttt{1\_000\_000}, clamped to roughly 80\,kHz\,..\,80\,MHz.
\item[\texttt{Sclk}, \texttt{Mosi}, \texttt{Miso}] (\texttt{Acquire}) usually
  \texttt{No\_Pin} = keep the \texttt{Configure\_Pins} routing; set them only for a
  device wired to a \emph{different} set of pads on the same controller (rare).
\item[\texttt{CS\_Pin}] a built-in software chip-select: name one GPIO and the
  driver owns it, driving it active-low and holding it asserted across the whole
  transaction. The common case for a shared bus; defaults to \texttt{No\_Pin}.
\item[\texttt{Select\_CB}, \texttt{Ctx}] an application chip-select callback and its
  opaque context, for a select that is \emph{not} one plain GPIO (a decoder, an
  I/O-expander line); leave all three out for the hardware \texttt{CS0}.
\item[\texttt{Tx}, \texttt{Rx}, \texttt{Length}] transmit/receive buffer
  \texttt{'Address}es and a byte count, \texttt{1\,..\,4095}; full-duplex, so both
  move \texttt{Length} bytes together.
\item[\texttt{Pad}] (\texttt{Enable\_Loopback}) one pad that folds MOSI back to MISO
  for the wiring-free self-test.
\end{description}

\textbf{Example 1 -- write a command to a device.}
\begin{lstlisting}
with Interfaces;
with ESP32S3.SPI; use ESP32S3.SPI;
procedure SPI_Write is
   S  : Session;
   Tx : aliased array (0 .. 2) of Interfaces.Unsigned_8 := (16#9F#, 0, 0);
   Rx : aliased array (0 .. 2) of Interfaces.Unsigned_8;
begin
   Setup (SPI2);
   Configure_Pins (SPI2, Sclk => 12, Mosi => 11, Miso => 13, Cs => 10);
   Acquire (S, SPI2, Mode => 0, Clock_Hz => 8_000_000);  --  this device's mode/clock
   Transfer (S, Tx'Address, Rx'Address, 3);   --  Rx now holds the response
end SPI_Write;                                 --  S releases the host here
\end{lstlisting}

\textbf{Example 2 -- wiring-free self-test via internal loopback.}
\begin{lstlisting}
with Interfaces;
with ESP32S3.SPI; use ESP32S3.SPI;
procedure SPI_Loopback is
   S  : Session;
   Tx : aliased array (0 .. 31) of Interfaces.Unsigned_8 := (others => 16#A5#);
   Rx : aliased array (0 .. 31) of Interfaces.Unsigned_8 := (others => 0);
begin
   Setup (SPI2);
   Enable_Loopback (SPI2, Pad => 5);   --  MOSI->MISO through one pad
   Acquire (S, SPI2);
   Transfer (S, Tx'Address, Rx'Address, 32);
   --  Rx = Tx
end SPI_Loopback;
\end{lstlisting}

\textbf{Per-device chip select.}\index{chip-select}\index{shared SPI bus}
The host has one hardware chip-select (\texttt{CS0}); to put several devices on one
bus, each device brings its \emph{own} select instead. The simple case is one
GPIO: pass its pin as \texttt{CS\_Pin} to \texttt{Acquire} and the driver owns it
-- it suppresses the hardware \texttt{CS0} for that session, parks the GPIO high,
and \texttt{Select\_Device} drives it low to assert and high to release. Crucially
the line is toggled only at \texttt{Select\_Device}, never per \texttt{Transfer}, so
it stays asserted across a whole multi-transfer transaction (a streamed flash read
keeps \texttt{CS} low the entire time) -- which the hardware \texttt{CS0}, pulsed
once per transfer, cannot do. This is how the W25Q flash shares the SPI2 bus with
the W5500 in Chapter~\ref{ch:storage}.

\begin{lstlisting}
   Setup (SPI2);
   Configure_Pins (SPI2, Sclk => 12, Mosi => 11, Miso => 13);  --  no hardware Cs
   Acquire (S, SPI2, Clock_Hz => 8_000_000, CS_Pin => 10);  --  drive IO10 active-low
   Select_Device (S, On => True);      --  assert IO10
   Transfer (S, Tx'Address, Rx'Address, 2);
   Select_Device (S, On => False);     --  deassert IO10
\end{lstlisting}

\textbf{When the select is not one GPIO.} A select behind a 3:8 decoder or an
I/O-expander cannot be expressed as a single pin. For those, leave
\texttt{CS\_Pin} at \texttt{No\_Pin} and pass a \texttt{CS\_Select} callback (with
an opaque \texttt{Ctx}) instead; \texttt{Select\_Device} then calls back to assert
and deassert, with exactly the same held-across-the-transaction semantics. The
callback must be \emph{library-level and closure-free} (this HAL builds under
\texttt{No\_Implicit\_Dynamic\_Code}, so a nested closure would emit a faulting
trampoline; per-device state travels in \texttt{Ctx}).

\begin{lstlisting}
--  Library-level, closure-free: drive IO10 low to select, high to release.
procedure CS_IO10 (Ctx : System.Address; Active : Boolean) is
begin
   if Active then ESP32S3.GPIO.Clear (10); else ESP32S3.GPIO.Set (10); end if;
end CS_IO10;

   Acquire (S, SPI2, Select_CB => CS_IO10'Access);   --  bring our own select
\end{lstlisting}

\textbf{Per-device mode and clock.}\index{SPI!per-device clock} Because
\texttt{Acquire} sets the mode and clock under the exclusive hold -- a few register
writes while no transfer is in flight -- two devices on one host need not agree on
them. A flash and a display can share \texttt{SPI2}, each clocked its own way:

\begin{lstlisting}
   Setup (SPI2);
   Configure_Pins (SPI2, Sclk => 12, Mosi => 11, Miso => 13);  --  shared wires
   --  ... later, in the flash driver:
   Acquire (S, SPI2, Mode => 0, Clock_Hz =>  8_000_000, CS_Pin => 21);
   --  ... and in the display driver, on the same host:
   Acquire (S, SPI2, Mode => 3, Clock_Hz => 40_000_000, CS_Pin => 10);
\end{lstlisting}

\noindent Each \texttt{Acquire} reprograms the controller for its caller, so the
flash always sees mode 0 at 8\,MHz and the display mode 3 at 40\,MHz -- the SD-card
driver uses exactly this to run its init handshake at 400\,kHz and switch to several
MHz for data. A device wired to \emph{different} pads on the same controller can
likewise pass \texttt{Sclk}/\texttt{Mosi}/\texttt{Miso} to \texttt{Acquire}, which
re-routes the GPIO matrix for the hold; left at \texttt{No\_Pin} they keep the shared
\texttt{Configure\_Pins} routing.

\section{I2C -- master}\index{I2C driver}

\begin{lstlisting}
procedure Setup (Host : I2C_Host; Clock_Hz : Positive := 100_000);
procedure Configure_Pins (Host : I2C_Host; Scl, Sda : Pin_Id);
procedure Acquire (S : in out Session; Host : I2C_Host);
procedure Write (S : Session; Addr : Slave_Address; Data : Byte_Array;
                 Success : out Boolean; Check_Ack : Boolean := True);
procedure Read  (S : Session; Addr : Slave_Address; Data : out Byte_Array;
                 Success : out Boolean);
\end{lstlisting}

\textbf{Parameters.}
\begin{description}[style=nextline]
\item[\texttt{Host}] \texttt{I2C0} or \texttt{I2C1}.
\item[\texttt{Clock\_Hz}] SCL frequency; default \texttt{100\_000} (standard mode),
  \texttt{400\_000} for fast mode.
\item[\texttt{Scl}, \texttt{Sda}] the two bus lines (\texttt{Pin\_Id}); the bus is
  open-drain, so fit external pull-ups.
\item[\texttt{Addr}] 7-bit slave address, \texttt{0\,..\,16\#7F\#}
  (\texttt{Slave\_Address}).
\item[\texttt{Data}] the bytes to write, or the \texttt{out} buffer to read into
  (\texttt{Byte\_Array} -- its length sets the count).
\item[\texttt{Success}] \texttt{out Boolean} -- \texttt{True} if the slave
  acknowledged.
\item[\texttt{Check\_Ack}] (\texttt{Write}) require an ACK on each byte; default
  \texttt{True}.
\end{description}

\textbf{Example 1 -- write a register on a sensor at address \texttt{16\#68\#}.}
\begin{lstlisting}
with ESP32S3.I2C; use ESP32S3.I2C;
procedure I2C_Write_Reg is
   S  : Session;
   Ok : Boolean;
begin
   Setup (I2C0, Clock_Hz => 400_000);
   Configure_Pins (I2C0, Scl => 9, Sda => 8);
   Acquire (S, I2C0);
   Write (S, Addr => 16#68#, Data => (16#6B#, 16#00#), Success => Ok);  --  reg<=0
end I2C_Write_Reg;
\end{lstlisting}

\textbf{Example 2 -- register read (write the address, then read).}
\begin{lstlisting}
with ESP32S3.I2C; use ESP32S3.I2C;
procedure I2C_Read_Reg is
   S    : Session;
   Ok   : Boolean;
   Who  : Byte_Array (0 .. 0);
begin
   Setup (I2C0, Clock_Hz => 400_000);
   Configure_Pins (I2C0, Scl => 9, Sda => 8);
   Acquire (S, I2C0);
   Write (S, 16#68#, (1 => 16#75#), Ok);     --  point at WHO_AM_I register
   Read  (S, 16#68#, Who, Ok);               --  read it back
end I2C_Read_Reg;
\end{lstlisting}

\section{UART -- asynchronous serial}\index{UART driver}

\begin{lstlisting}
--  Startup (single-threaded): the ONLY port-based configuration call.
procedure Setup (Port : UART_Port; Baud : Baud_Rate := 115_200;
                 Bits : Data_Bits := 8; Parity : Parity_Mode := None;
                 Stop : Stop_Bits := One;
                 Tx, Rx, Rts, Cts : Optional_Pin := No_Pin;  --  routes the pads too
                 Rx_Flow_Threshold : Natural := 100);
procedure Acquire (S : in out Session; Port : UART_Port);
procedure Write (S : Session; Data : Byte_Array);
procedure Read  (S : Session; Data : out Byte_Array; Count : out Natural);
function  Available (S : Session) return Natural;
--  Post-Setup reconfiguration -- each takes the HELD Session, not a Port:
procedure Set_Baud      (S : Session; Baud   : Baud_Rate);
procedure Set_Data_Bits (S : Session; Bits   : Data_Bits);
procedure Set_Parity    (S : Session; Parity : Parity_Mode);
procedure Set_Stop_Bits (S : Session; Stop   : Stop_Bits);
procedure Configure_Pins  (S : Session; ...);   --  re-route pads
procedure Enable_Loopback (S : Session; On : Boolean := True);
\end{lstlisting}

\textbf{Parameters.}
\begin{description}[style=nextline]
\item[\texttt{Port}] \texttt{UART0}, \texttt{UART1} or \texttt{UART2}.
\item[\texttt{Baud}] bit rate (\texttt{Baud\_Rate}, \texttt{300\,..\,5\_000\_000});
  default \texttt{115\_200}.
\item[\texttt{Bits}] data bits per frame (\texttt{Data\_Bits}, \texttt{5\,..\,8});
  default \texttt{8}.
\item[\texttt{Parity}] \texttt{None} (default), \texttt{Even} or \texttt{Odd}.
\item[\texttt{Stop}] stop bits: \texttt{One} (default) or \texttt{Two}.
\item[\texttt{Tx}, \texttt{Rx}, \texttt{Rts}, \texttt{Cts}] the four pads as
  \texttt{Optional\_Pin}; all default to \texttt{No\_Pin} (route only the lines the
  link uses).
\item[\texttt{Rx\_Flow\_Threshold}] RX FIFO fill, in bytes, at which RTS
  de-asserts; default \texttt{100} (the FIFO holds 128).
\item[\texttt{On}] (\texttt{Enable\_Loopback}) \texttt{True} (default) folds TX back
  to RX internally.
\item[the \texttt{Set\_*} setters] each take the held \texttt{Session} and one new
  value of the matching \texttt{Setup} type (e.g. \texttt{Set\_Parity} takes a
  \texttt{Parity\_Mode}); \texttt{Configure\_Pins} additionally accepts per-line
  \texttt{*\_Invert} booleans (default \texttt{False}).
\end{description}

\texttt{Setup} both configures the frame format \emph{and} routes the four pins
(all optional -- pass only what the link uses, or none for an internal loopback
self-test). It is the one port-based call, meant to run once at startup before
any task contends for the port. \emph{Every} later change -- baud, data bits,
parity, stop bits, pin re-routing, loopback -- takes a \texttt{Session} instead
of a \texttt{Port}, so you must \texttt{Acquire} the port before you can
reconfigure it. Changing a parameter therefore requires owning the port and can
never race another task's transfer. Each frame setter is an independent
read-modify-write that leaves the other attributes and the pin routing untouched.

This rule is enforced \emph{structurally}, not by convention. Inside the body,
the per-port register handle lives in a nested package whose body hides it; the
only way to obtain it is one ownership-checked gateway that raises
\texttt{Not\_Owned} unless the passed \texttt{Session} holds the port. Every
transfer and every reconfiguration goes through that gateway, so a call without
ownership fails loudly -- and, just as importantly, a \emph{new} operation
added to the driver cannot reach the registers without passing a \texttt{Session}
through the same check. The guard is impossible to forget because there is no
other door to the hardware.

\textbf{Example 1 -- transmit a line.}
\begin{lstlisting}
with ESP32S3.UART; use ESP32S3.UART;
procedure UART_Tx is
   S : Session;
begin
   Setup (UART1, Baud => 115_200, Tx => 17, Rx => 16);   --  format + pins in one
   Acquire (S, UART1);
   Write (S, (Character'Pos ('h'), Character'Pos ('i'), 10));
end UART_Tx;
\end{lstlisting}

\textbf{Example 2 -- echo received bytes (wiring-free, internal loopback).}
\begin{lstlisting}
with ESP32S3.UART; use ESP32S3.UART;
procedure UART_Echo is
   S   : Session;
   Buf : Byte_Array (0 .. 63);
   N   : Natural;
begin
   Setup (UART1);
   Acquire (S, UART1);
   Enable_Loopback (S);              --  TX folded back to RX (held port)
   Write (S, (1 => 16#42#));
   loop
      exit when Available (S) > 0;
   end loop;
   Read (S, Buf, N);                 --  Buf (0 .. N-1) holds the echoed bytes
end UART_Echo;
\end{lstlisting}

\textbf{Example 2, step by step.} This example is worth reading line by line,
because it exercises every part of the \texttt{Session} pattern.

\emph{The declarations (before \texttt{begin}).}
\begin{itemize}
\item \texttt{S : Session} -- \texttt{Session} is a \emph{limited controlled}
  type. \emph{Limited} means it cannot be copied or assigned (two tasks can
  never share one hold); \emph{controlled} means it has a \texttt{Finalize} that
  runs automatically when \texttt{S} leaves scope. It starts inactive -- it
  owns nothing yet.
\item \texttt{Buf : Byte\_Array (0 .. 63)} -- a 64-element array of
  \texttt{Byte} (which is \texttt{new Interfaces.Unsigned\_8}); the destination
  for the bytes we read back.
\item \texttt{N : Natural} -- receives the count of bytes actually read.
\end{itemize}

\emph{The statements, in execution order.}
\begin{enumerate}
\item \texttt{Setup (UART1)} brings UART1 up with the default frame format
  (115200 baud, 8 data bits, no parity, one stop bit -- the defaults on
  \texttt{Setup}). It opens the controller and records that UART1 is configured.
  No pins are routed: we pass no pin arguments to \texttt{Setup} (they all
  default to \texttt{No\_Pin}) because this is a wiring-free self-test.
\item \texttt{Acquire (S, UART1)} takes exclusive ownership of UART1. It first
  checks the init-before-use guard: had UART1 never been \texttt{Setup}, this
  would raise \texttt{Not\_Initialized}. It then blocks on UART1's protected
  guard until the port is free (here, immediately), marks \texttt{S} active, and
  records that it holds UART1. Until \texttt{S} is released no other task can
  acquire UART1 -- and, because every reconfiguration call below takes the held
  \texttt{Session}, no other task can change UART1's settings either.
\item \texttt{Enable\_Loopback (S)} sets the controller's internal TX$\rightarrow$RX
  loopback bit \emph{on the held port}. Everything the transmitter shifts out is
  fed straight back into the receiver \emph{inside the chip}. Because UART is
  push-pull (unlike I2C's open-drain bus) this loopback is faithful and needs
  \emph{no external wiring and no GPIO pads} -- the whole point of the example.
  It takes \texttt{S} (not \texttt{UART1}): you can only flip loopback on a port
  you own.
\item \texttt{Write (S, (1 => 16\#42\#))} -- \texttt{(1 => 16\#42\#)} is a
  one-element array aggregate, a \texttt{Byte\_Array} of length 1 holding the
  byte \texttt{16\#42\#} (66, ASCII \texttt{'B'}). \texttt{Write} pushes it into
  the TX FIFO and returns once it is queued. This runs \emph{outside} the
  protected lock -- the lock only arbitrated \emph{ownership} in the previous
  step; the blocking I/O is never done while holding it. With loopback on, that
  byte will reappear in the RX FIFO shortly.
\item The poll loop \texttt{loop exit when Available (S) > 0; end loop;} spins,
  re-reading how many bytes sit in the RX FIFO, until at least one has arrived.
  The looped-back byte takes roughly one character-time ($\sim$87\,\textmu s at
  115200 baud), so the loop spins briefly then exits. (A self-test idiom;
  production code would use a timeout or an interrupt instead of a bare spin.)
\item \texttt{Read (S, Buf, N)} reads up to \texttt{Buf'Length} (64) bytes into
  \texttt{Buf} and sets \texttt{N} to how many it actually got -- here
  \texttt{N = 1}, with \texttt{Buf (0) = 16\#42\#}. Elements
  \texttt{Buf (N .. 63)} are left untouched. \texttt{Read} waits only briefly
  per byte and does a short read on timeout, so it never hangs on bytes that
  never come.
\item At \texttt{end UART\_Echo}, \texttt{S} goes out of scope. Being controlled,
  its \texttt{Finalize} runs automatically and calls \texttt{Release}, handing
  UART1 back to its guard. This happens \emph{even if an exception had been
  raised} between \texttt{Acquire} and here -- that is why the RAII handle
  exists: the lock cannot leak.
\end{enumerate}

The two phases are deliberately separate. \texttt{Setup} (step 1) is the one
port-based call, meant to run single-threaded at startup before any task
contends for the port. \texttt{Acquire} (step 2) is the per-transaction
ownership gate that hands you a \texttt{Session} -- and everything after it,
including \texttt{Enable\_Loopback} here, acts through that \texttt{Session}, so
a reconfiguration is only possible while you hold the port. Two guards tie the
phases together: you cannot \texttt{Acquire} a port you never \texttt{Setup}, and
you cannot change a port's parameters without holding it.

\textbf{Example 3 -- change baud and frame format independently.} Each frame
attribute has its own setter, so a protocol that renegotiates the link (or a
bootloader that bumps the speed after handshaking) can change just one without
re-running \texttt{Setup} or touching the pin routing. The setters take the held
\texttt{Session}, so you \texttt{Acquire} the port first -- the reconfiguration
runs under the same ownership as the transfers:
\begin{lstlisting}
with ESP32S3.UART; use ESP32S3.UART;
procedure UART_Reconfig is
   S : Session;
begin
   Setup (UART1, Tx => 17, Rx => 16);   --  115200 8-N-1 to start
   Acquire (S, UART1);                  --  own the port before reconfiguring
   Set_Baud      (S, 9_600);            --  then change each attribute alone
   Set_Data_Bits (S, 7);
   Set_Parity    (S, Even);
   Set_Stop_Bits (S, Two);              --  now 9600 7-E-2
   Write (S, (1 => 16#55#));
end UART_Reconfig;
\end{lstlisting}

\chapter{DMA, PWM, Audio, Pulse and Analog Peripherals}
\label{ch:hal2}

These drivers follow the \emph{claim-a-handle} pattern: a peripheral with several
identical units hands out a limited, controlled handle through a protected pool.
You \texttt{Claim} a unit, use it through the handle, and it is reclaimed when the
handle leaves scope.

\section{GDMA: general DMA}\index{GDMA driver}

\begin{lstlisting}
procedure Claim (C : in out Channel; Peri : Peripheral);
procedure Copy  (C : Channel; Dst, Src : System.Address; Length : Natural);
procedure Start (C : Channel; Dir : Direction; ...);  --  peripheral transfers
function  Done  (C : Channel; Dir : Direction) return Boolean;
procedure Wait  (C : Channel; Dir : Direction);
--  Gapless double-buffered streaming, one ring per direction (see below).
procedure Start_Stream    (C : Channel; Buffer : System.Address; Half_Length : Natural);
function  Await_Half      (C : Channel) return Natural;   --  half to refill (0/1)
procedure Start_In_Stream (C : Channel; Buffer : System.Address; Half_Length : Natural);
function  Await_In_Half   (C : Channel) return Natural;   --  half to read (0/1)
procedure Stop   (C : Channel; Dir : Direction);
procedure Unbind (C : Channel; Dir : Direction);  --  detach an unused direction
\end{lstlisting}

\textbf{Parameters.}
\begin{description}[style=nextline]
\item[\texttt{Peri}] which peripheral the channel serves: \texttt{Mem2Mem} (a
  plain memory copy) or a real peripheral: \texttt{SPI2}, \texttt{SPI3},
  \texttt{UHCI0}, \texttt{I2S0}, \texttt{I2S1}, \texttt{LCD\_CAM}, \texttt{AES},
  \texttt{SHA}, \texttt{ADC\_DAC}, \texttt{RMT}. The pool has five channels.
\item[\texttt{Dir}] transfer direction: \texttt{Mem\_To\_Periph} or
  \texttt{Periph\_To\_Mem}.
\item[\texttt{Dst}, \texttt{Src}, \texttt{Length}] destination/source
  \texttt{'Address}es and a byte count, \texttt{1\,..\,4095}.
\end{description}

\lstexample{Memory-to-memory copy.}
\begin{lstlisting}
with Interfaces;
with ESP32S3.GDMA; use ESP32S3.GDMA;
procedure GDMA_Copy is
   C   : Channel;
   Src : aliased array (0 .. 255) of Interfaces.Unsigned_8 := (others => 16#5A#);
   Dst : aliased array (0 .. 255) of Interfaces.Unsigned_8 := (others => 0);
begin
   Claim (C, Peri => Mem2Mem);                --  a free channel for mem2mem
   Copy  (C, Dst'Address, Src'Address, 256);  --  blocking DMA copy; Dst = Src
end GDMA_Copy;                                 --  C is released on scope exit
\end{lstlisting}

\textbf{Example 2: the RAII guarantee.} Claiming all five channels and then
one more raises (the pool is exhausted); each goes back automatically:
\begin{lstlisting}
with ESP32S3.GDMA; use ESP32S3.GDMA;
procedure GDMA_RAII is
   A, B, C2, D, E : Channel;
begin
   Claim (A, Mem2Mem); Claim (B, Mem2Mem); Claim (C2, Mem2Mem);
   Claim (D, Mem2Mem); Claim (E, Mem2Mem);          --  all five claimed
end GDMA_RAII;                                       --  all five reclaimed here
\end{lstlisting}
(The SPI and I2S drivers use the peripheral \texttt{Start}/\texttt{Wait} path
internally, so most applications use GDMA only for memory copies.)

\texttt{Wait} (and \texttt{Copy}) is interrupt-driven: after a brief spin it
suspends the calling task, and the channel's end-of-transfer interrupt wakes it,
so a transfer hands the core back rather than busy-waiting. The driver owns the
\texttt{Device\_L2\_0} interrupt (CPU\_INT 19) for this.

That suspend is \emph{bounded by a deadline}. A stuck or misconfigured transfer
whose end-of-transfer interrupt never arrives would otherwise block the waiting task
forever, so before it suspends, \texttt{Wait} arms a per-channel
\texttt{Ada.Real\_Time.Timing\_Events} timer a few seconds out; whichever fires
first --- the real EOF interrupt or the timer --- releases the \emph{same}
suspension object, and \texttt{Wait} cancels the other on wake. \texttt{Set\_True}
is idempotent, so the two racing near the deadline is harmless. The normal path
keeps its exact, zero-latency interrupt wake (the timer is cancelled before it ever
runs); only a genuine fault trips the deadline, and then \texttt{Wait} returns
rather than deadlocking the system. A timed entry call would read more directly, but
the Ravenscar/Jorvik runtime forbids \texttt{select}, so a timing-event wake is the
profile-legal way to put a ceiling on a \texttt{Suspend\_Until\_True}.

\textbf{Streaming rings.} \texttt{Start\_Stream} services a continuous
\emph{producer}: it links two descriptors, one per half of the caller's buffer,
into a ring the OUT engine walks forever, and enables the per-descriptor
completion interrupt. \texttt{Await\_Half} suspends until a half has drained
and returns which one (0 or 1) is now safe to refill. The refill is therefore
paced by the DMA itself, not by a timer, so live-generated audio never drifts
against the peripheral clock. \texttt{Start\_In\_Stream} and
\texttt{Await\_In\_Half} are the receive mirror for a continuous
\emph{consumer} (a demodulator, a recorder): the IN engine fills the two halves
forever, and the caller reads each half as it completes, missing no samples.
Run one ring in each direction on two channels and you have true full duplex.
\texttt{Stop} with the matching direction ends a ring.

\textbf{One channel per direction.} \texttt{Claim} binds \emph{both}
directions of a channel to the peripheral, but the hardware routes each
peripheral request line to only one channel per direction. When two channels
serve the same peripheral (one streaming TX, another streaming RX), the idle
binding on one channel \emph{shadows} the active binding on the other, and that
direction silently transfers nothing. (Measured: a capture stream beside a
running I2S transmit stream delivered zero bytes until the transmit channel's
unused IN binding was detached.) \texttt{Unbind} detaches one direction; every
claimant that uses a channel one-way should unbind the other, and the I2S driver
now does.

\section{MCPWM: motor-control PWM}\index{MCPWM driver}

\begin{lstlisting}
--  The unit's clock comes up on the first Claim of any of its channels.
procedure Claim (C : in out Channel; Unit : MCPWM_Unit; Index : Channel_Index);
procedure Configure_Channel (C : Channel; Freq : Positive; Pin : Pin_Id;
                             Complement_Pin : Optional_Pin := No_Pin;
                             Dead_Time_Ns   : Natural := 0);
procedure Start (C : Channel);   procedure Set_Duty (C : Channel; Percent : Duty_Percent);
\end{lstlisting}

\textbf{Parameters.}
\begin{description}[style=nextline]
\item[\texttt{Unit}] \texttt{MCPWM0} or \texttt{MCPWM1}.
\item[\texttt{Index}] generator channel within the unit: \texttt{Ch0},
  \texttt{Ch1} or \texttt{Ch2}.
\item[\texttt{Freq}] PWM frequency in Hz.
\item[\texttt{Pin}] the output pad.
\item[\texttt{Complement\_Pin}] optional inverted output for a half-bridge
  (\texttt{Optional\_Pin}, default \texttt{No\_Pin}).
\item[\texttt{Dead\_Time\_Ns}] gap inserted between the complementary edges, in
  nanoseconds; default \texttt{0}.
\item[\texttt{Percent}] duty cycle, \texttt{0.0\,..\,100.0}.
\end{description}

\lstexample{20\,kHz at 25\,\% on one pin.}
\begin{lstlisting}
with ESP32S3.MCPWM; use ESP32S3.MCPWM;
procedure MCPWM_Simple is
   C : Channel;
begin
   Claim (C, MCPWM0, Ch0);   --  first Claim brings up the MCPWM0 clock
   Configure_Channel (C, Freq => 20_000, Pin => 4);
   Set_Duty (C, 25.0);
   Start (C);
end MCPWM_Simple;
\end{lstlisting}

\lstexample{A complementary pair with 500\,ns dead-time (half-bridge).}
\begin{lstlisting}
with ESP32S3.MCPWM; use ESP32S3.MCPWM;
procedure MCPWM_Pair is
   C : Channel;
begin
   Claim (C, MCPWM0, Ch0);   --  first Claim brings up the MCPWM0 clock
   Configure_Channel (C, Freq => 20_000, Pin => 4,
                      Complement_Pin => 5, Dead_Time_Ns => 500);
   Set_Duty (C, 60.0);
   Start (C);
end MCPWM_Pair;
\end{lstlisting}
The driver also offers a carrier (\texttt{Set\_Carrier}: \texttt{Prescale}
\texttt{0\,..\,15}, \texttt{Duty\_Eighths} \texttt{1\,..\,7}, \texttt{First\_Pulse}
\texttt{0\,..\,15}) and fault/trip-zone protection: \texttt{Configure\_Fault}
binds a \texttt{Fault\_Input} (\texttt{Fault0\,..\,2}) to a pad, and
\texttt{Protect\_Channel} sets the \texttt{Fault\_Mode} (\texttt{One\_Shot} or
\texttt{Cycle\_By\_Cycle}) and \texttt{Trip\_Action} (\texttt{No\_Change},
\texttt{Force\_Low} or \texttt{Force\_High}).

\section{I2S: digital audio over GDMA}\index{I2S driver}

\begin{lstlisting}
type I2S_Mode is (Standard, PDM);
type PCM_8  is array (Natural range <>) of Interfaces.Integer_8;   -- the typed,
type PCM_16 is array (Natural range <>) of Interfaces.Integer_16;  --  signed PCM
type PCM_32 is array (Natural range <>) of Interfaces.Integer_32;  --  buffers
--  Acquire claims a port AND brings it up (audio format + pin routing).  The
--  first Acquire of a port opens it; a later Acquire reuses it (Reconfigure
--  re-opens to change the format), so a codec's MCLK stays running.
procedure Acquire (S : in out Session; Port : I2S_Port;
                   Sample_Rate : Positive := 16_000; Bits : Sample_Bits := Bits_16;
                   Mode : I2S_Mode := Standard;
                   Bclk, Ws, Dout, Din, Mclk : Optional_Pin := No_Pin);
procedure Reconfigure (S : Session; ...);  -- re-open the held port's format
function  Configured_Bits (S : Session) return Sample_Bits;
--  Typed transfers (preferred): the element type fixes the on-wire width, so the
--  byte count is derived and a precondition checks it against the port's Bits.
procedure Write    (S : Session; Samples : PCM_16)
   with Pre => Configured_Bits (S) = Bits_16;
procedure Read     (S : Session; Samples : out PCM_16)
   with Pre => Configured_Bits (S) = Bits_16;
procedure Transfer (S : Session; Tx : PCM_16; Rx : out PCM_16)
   with Pre => Configured_Bits (S) = Bits_16 and then Tx'Length = Rx'Length;
--  PCM_8 / PCM_32 overloads likewise (also Start_Continuous / Capture).  The
--  *_Raw forms take a 'Address + byte Length for bytes already framed elsewhere:
procedure Write_Raw (S : Session; Tx : System.Address; Length : Natural);
--  Gapless double-buffered streaming, both directions (see below).
procedure Start_Stream (S : Session; Samples : PCM_16);   --  TX ring, 2 halves
function  Await_Half   (S : Session) return Natural;      --  half to refill (0/1)
procedure Start_Capture_Stream (S : Session; Rx : System.Address;
                                Half_Length : Natural);   --  RX ring, own channel
function  Await_Capture_Half   (S : Session) return Natural;  --  half to read
procedure Stop_Capture         (S : Session);
\end{lstlisting}

\textbf{Parameters.}
\begin{description}[style=nextline]
\item[\texttt{Port}] \texttt{I2S0} or \texttt{I2S1}.
\item[\texttt{Sample\_Rate}] frames per second (\texttt{Positive}); default
  \texttt{16\_000}.
\item[\texttt{Bits}] sample width: \texttt{Bits\_8}, \texttt{Bits\_16} (default),
  \texttt{Bits\_24} or \texttt{Bits\_32}.
\item[\texttt{Mode}] data path: \texttt{Standard} (default, plain I2S/TDM) or
  \texttt{PDM} (the sigma-delta converters; see below).
\item[\texttt{Bclk}, \texttt{Ws}, \texttt{Dout}, \texttt{Din}] the bit-clock,
  word-select, data-out and data-in lines (\texttt{Optional\_Pin};
  \texttt{Dout}/\texttt{Din} default \texttt{No\_Pin}).
\item[\texttt{Samples}] a typed \texttt{PCM\_8\,/\,16\,/\,32} buffer; the driver
  derives the byte count and a precondition checks the width against the port.
  The \texttt{*\_Raw} forms instead take a \texttt{'Address} + byte
  \texttt{Length} (\texttt{1\,..\,4095}) for bytes already framed elsewhere.
\end{description}

The element type of a \texttt{PCM\_*} buffer fixes the on-wire width, so a
\texttt{Write\,(S,\,Samples)} needs no caller-side \texttt{'Address} or
\texttt{*\,2}: the driver computes the byte count from \texttt{Samples'Length},
and the precondition \texttt{Configured\_Bits\,(S)\,=\,Bits\_16} turns a
buffer-vs.-port width mismatch from silent noise into a caught contract
violation. Reach for \texttt{Write\_Raw}/\texttt{Transfer\_Raw} only for bytes
that are \emph{already} framed (a DMA buffer built elsewhere, or an opaque
bit-pattern self-test).

A one-shot \texttt{Write} also has to know when it is truly \emph{finished}. The
GDMA completing means only that the samples have reached the I2S TX FIFO --- not
that the serializer has clocked them out onto the wire. Clearing \texttt{TX\_START}
the instant the DMA is done would discard whatever is still in the FIFO, truncating
the tail of every buffer (the end of the sound simply cut off). So after the DMA
wait the driver spins on the \texttt{STATE.TX\_IDLE} status --- FIFO drained,
serializer idle --- before it stops the channel, so the last samples actually leave
the pin.

\lstexample{Play a buffer of samples.}
\begin{lstlisting}
with ESP32S3.I2S; use ESP32S3.I2S;
procedure I2S_Play is
   S    : Session;
   Tone : PCM_16 (0 .. 255) := (others => 0);
begin
   Acquire (S, I2S0, Sample_Rate => 44_100, Bits => Bits_16,
            Bclk => 4, Ws => 5, Dout => 6);   --  claim + configure
   Write (S, Tone);                           --  256 samples x 2 bytes, derived
end I2S_Play;
\end{lstlisting}

\lstexample{Wiring-free full-duplex loopback self-test.}
\begin{lstlisting}
with ESP32S3.I2S; use ESP32S3.I2S;
procedure I2S_Loopback is
   S  : Session;
   Tx : PCM_16 (0 .. 63) := (others => 0);    --  fill with a test pattern
   Rx : PCM_16 (0 .. 63) := (others => 0);
begin
   Acquire (S, I2S0);                          --  claim (defaults: 16k 16-bit)
   Enable_Loopback (S, Pad => 7);              --  TX folded to RX (held port)
   Transfer (S, Tx, Rx);                       --  Rx = Tx (same length checked)
end I2S_Loopback;
\end{lstlisting}

\subsection{Gapless streaming and full duplex}

A one-shot \texttt{Write} stops the clock when it finishes, and back-to-back
writes leave an audible gap. \texttt{Start\_Continuous} closes the gap for a
\emph{static} buffer (it replays one buffer forever), but live-generated audio
needs to refill while playing, and pacing the refill from a CPU timer drifts
against the I2S clock. (Measured: a timer-paced double buffer feeding a codec
produced a beat that returned every few seconds as the two clocks slid past each
other.) \texttt{Start\_Stream} fixes the pacing at the source: the DMA loops
the two halves of the buffer forever, and \texttt{Await\_Half} blocks until a
half has drained, returning which one (0 or 1) to refill. The producer is paced
by the same clock that consumes the samples, so it can never drift.

\texttt{Start\_Capture\_Stream} is the receive mirror. It claims its
\emph{own} DMA channel, so it runs concurrently with a streaming transmit:
play out of one ring while recording into the other, both gapless, with no
clock interruption in either direction (true full duplex).
\texttt{Await\_Capture\_Half} returns each captured half as the DMA fills
it, so a continuous consumer (a recorder, a software modem) misses no samples;
\texttt{Stop\_Capture} ends it.

Two facts make external capture work, and both are invisible until they bite.
First, the RX block runs as a \emph{slave}: it samples on the bit and word
clocks it receives, not on clocks it generates. The driver therefore feeds the
RX clock inputs from the very pads the TX master drives, so capture samples on
the true on-wire clock, phase-exact with the external device. (Without that
routing, slave RX has no clock at all and capture waits forever; the internal
loopback self-test masks the bug, because \texttt{SIG\_LOOPBACK} bypasses the
pads.) Second, the one-channel-per-direction rule from the GDMA section: the
transmit stream's channel unbinds its unused RX direction and the capture
channel unbinds its unused TX direction, or one shadows the other.

\lstexample{Full duplex: synthesize live audio while recording the input.}
\begin{lstlisting}
Half : constant := 500;                        --  31 ms per half at 16 kHz
Play_Ring, Capture_Ring : PCM_16 (0 .. 4 * Half - 1);
S : Session;
begin
   Acquire (S, I2S0, Bclk => 4, Ws => 5, Dout => 6, Din => 7);
   Fill (Play_Ring);                           --  prime both halves
   Start_Stream (S, Play_Ring);
   Start_Capture_Stream (S, Capture_Ring'Address, 2 * Half * 2);
   loop
      Refill  (Play_Ring,    Await_Half (S));          --  DMA-paced, no drift
      Consume (Capture_Ring, Await_Capture_Half (S));  --  misses no samples
   end loop;
\end{lstlisting}

\subsection{PDM mode: the hardware PCM\,$\leftrightarrow$\,PDM converters}

\texttt{Mode => PDM} routes the data path through the sigma-delta converters
instead of the plain serial framer. PDM (pulse-density modulation) is the
sigma-delta sibling of PWM: a 1-bit oversampled stream whose \emph{density} of
ones encodes the analog level. The S3 I2S converts in both directions, and the
DMA still moves ordinary PCM, so only the on-wire format changes:

\begin{itemize}
  \item \textbf{TX} (\texttt{PCM2PDM}) turns each PCM sample into a 1-bit stream:
        feed it to a class-D amplifier or an RC low-pass for analog audio out
        (the S3 has no analog DAC, so this is how you get one).
  \item \textbf{RX} (\texttt{PDM2PCM}) decimates a 1-bit PDM input back to PCM.
        PDM is the format \emph{emitted} by digital \textbf{PDM microphones} and
        some PDM-output \textbf{ADCs}; the converter turns their bitstream into the
        PCM samples your program reads. In PDM the receiver is the clock master:
        it drives the device clock out on \texttt{Ws} and reads the bitstream in on
        \texttt{Din}.
\end{itemize}

\begin{lstlisting}
with ESP32S3.I2S; use ESP32S3.I2S;
procedure PDM_Mic is
   S  : Session;
   Rx : PCM_16 (0 .. 511) := (others => 0);
begin
   Acquire (S, I2S0, Sample_Rate => 16_000, Bits => Bits_16, Mode => PDM,
            Ws => 5, Din => 6);               --  claim + configure (PDM mic)
   Read (S, Rx);                              --  PDM mic -> signed PCM_16
end PDM_Mic;
\end{lstlisting}

\textbf{A note on verifying PDM.} Every other HAL driver self-tests with no
external hardware; PDM is the exception, and the reason is instructive. There is
no on-chip loopback for the converters (the \texttt{SIG\_LOOPBACK} bit only
shares the \emph{standard}-I2S WS and BCK, not the bit timing between
\texttt{PCM2PDM} and \texttt{PDM2PCM}), so an internal
\texttt{PCM}\,$\rightarrow$\,\texttt{PDM}\,$\rightarrow$\,\texttt{PCM} loop
recovers nothing. Nor can a static level
synthesised from a GPIO (a pull resistor or a driven pin) stand in for a mic: the
decimator's mandatory high-pass filter strips DC, so a constant 1 or 0 is rejected.
Genuine PDM verification needs a real device producing a toggling bitstream, so
the \texttt{esp32s3\_i2s\_pdm} example is a microphone-capture demo rather than a
self-test.

\section{LEDC: LED / general PWM}\index{LEDC driver}

\begin{lstlisting}
procedure Claim (C : in out Channel; Index : Channel_Index);
procedure Configure (C : in out Channel; Freq : Positive; Pin : Pin_Id;
                     Bits : Resolution := 10);
procedure Set_Duty (C : Channel; Percent : Duty_Percent);
\end{lstlisting}

\textbf{Parameters.}
\begin{description}[style=nextline]
\item[\texttt{Index}] channel \texttt{0\,..\,7} (\texttt{Channel\_Index}).
\item[\texttt{Freq}] PWM frequency in Hz.
\item[\texttt{Pin}] the output pad.
\item[\texttt{Bits}] duty resolution in bits, \texttt{1\,..\,14}
  (\texttt{Resolution}); default \texttt{10}. More bits give finer duty steps but a
  lower maximum frequency.
\item[\texttt{Percent}] duty cycle, \texttt{0.0\,..\,100.0}.
\end{description}

\lstexample{5\,kHz at 50\,\%.}
\begin{lstlisting}
with ESP32S3.LEDC; use ESP32S3.LEDC;
procedure LEDC_Simple is
   C : Channel;
begin
   Claim (C, 0);
   Configure (C, Freq => 5_000, Pin => 2);
   Set_Duty (C, 50.0);
end LEDC_Simple;
\end{lstlisting}

\lstexample{Fade an LED.}
\begin{lstlisting}
with Ada.Real_Time; use Ada.Real_Time;
with ESP32S3.LEDC;   use ESP32S3.LEDC;
procedure LEDC_Fade is
   C : Channel;
begin
   Claim (C, 0);
   Configure (C, Freq => 5_000, Pin => 2);
   for D in 0 .. 100 loop
      Set_Duty (C, Duty_Percent (D));
      delay until Clock + Milliseconds (10);
   end loop;
end LEDC_Fade;
\end{lstlisting}

\section{RMT: remote-control pulse TX/RX (IR, WS2812)}\index{RMT driver}

\begin{lstlisting}
type RMT_Symbol is record
   Level0, Level1 : Boolean := False; Duration0, Duration1 : Tick_Count := 0;
end record;
procedure Configure (C : in out TX_Channel; Resolution_Hz : Positive; Pin : Pin_Id);
procedure Transmit  (C : TX_Channel; Symbols : Symbol_Array);
procedure Configure (C : in out RX_Channel; Resolution_Hz : Positive; Pin : Pin_Id);
procedure Start (C : RX_Channel);
procedure Receive (C : RX_Channel; Into : out Symbol_Array; Count : out Natural);
\end{lstlisting}

\textbf{Parameters.}
\begin{description}[style=nextline]
\item[\texttt{Index}] channel \texttt{0\,..\,3}: four TX channels
  (\texttt{TX\_Index}) and four RX (\texttt{RX\_Index}), separate pools.
\item[\texttt{Resolution\_Hz}] the tick clock: one \texttt{Tick\_Count} unit lasts
  \texttt{1/Resolution\_Hz} seconds.
\item[\texttt{Pin}] the carrier pad.
\item[\texttt{Idle\_Ticks}] (RX) the silence that ends a frame,
  \texttt{0\,..\,32\_767}; default \texttt{2\_000}.
\item[\texttt{Symbols} / \texttt{Into}] a \texttt{Symbol\_Array} of
  \texttt{RMT\_Symbol}s, each holding two level+duration halves;
  \texttt{Duration0}/\allowbreak\texttt{Duration1} are \texttt{Tick\_Count}
  (\texttt{0\,..\,32\_767}). One symbol-RAM block holds 48 symbols; a longer burst
  is streamed by refilling the block in halves (wrap mode).
\item[\texttt{Count}] (\texttt{Receive}) \texttt{out} count of symbols captured.
\end{description}

\lstexample{Transmit a pulse burst (1\,\textmu s resolution).}
\begin{lstlisting}
with ESP32S3.RMT; use ESP32S3.RMT;
procedure RMT_Tx is
   Tx    : TX_Channel;
   Burst : constant Symbol_Array :=
     ((Level0 => True,  Duration0 => 10, Level1 => False, Duration1 => 10),
      (Level0 => True,  Duration0 => 20, Level1 => False, Duration1 => 20));
begin
   Claim (Tx, 0);
   Configure (Tx, Resolution_Hz => 1_000_000, Pin => 4);
   Transmit (Tx, Burst);
end RMT_Tx;
\end{lstlisting}

\lstexample{Receive into a buffer (e.g. an IR remote).}
\begin{lstlisting}
with ESP32S3.RMT; use ESP32S3.RMT;
procedure RMT_Rx is
   Rx   : RX_Channel;
   Syms : Symbol_Array (0 .. 63);
   N    : Natural;
begin
   Claim (Rx, 0);
   Configure (Rx, Resolution_Hz => 1_000_000, Pin => 5);
   Start (Rx);
   Receive (Rx, Syms, N);   --  Syms (0 .. N-1) are the captured pulses
end RMT_Rx;
\end{lstlisting}

\section{PCNT: pulse / edge counter}\index{PCNT driver}

\begin{lstlisting}
procedure Configure (U : in out Unit; Pin : Pin_Id; Both_Edges : Boolean := False);
function  Count (U : Unit) return Integer;   procedure Clear (U : Unit);
\end{lstlisting}

\textbf{Parameters.}
\begin{description}[style=nextline]
\item[\texttt{Index}] unit \texttt{0\,..\,3} (\texttt{Unit\_Index}).
\item[\texttt{Pin}] the input pad whose edges are counted.
\item[\texttt{Both\_Edges}] \texttt{False} (default) counts rising edges only;
  \texttt{True} counts both (double resolution).
\item[\texttt{Count}] returns the signed counter (16-bit, wraps at $\pm$32768).
\end{description}

\lstexample{Count rising edges on a pin.}
\begin{lstlisting}
with Ada.Text_IO;  use Ada.Text_IO;
with ESP32S3.PCNT; use ESP32S3.PCNT;
procedure PCNT_Count is
   U : Unit;
begin
   Claim (U, 0);
   Configure (U, Pin => 6);
   Clear (U);
   --  ... pulses arrive ...
   Put_Line ("edges =" & Integer'Image (Count (U)));
end PCNT_Count;
\end{lstlisting}

\lstexample{Count both edges (double the resolution).}
\begin{lstlisting}
with ESP32S3.PCNT; use ESP32S3.PCNT;
procedure PCNT_Both is
   U : Unit;
begin
   Claim (U, 0);
   Configure (U, Pin => 6, Both_Edges => True);
end PCNT_Both;
\end{lstlisting}

\section{SDM: sigma-delta density output}\index{SDM driver}

\begin{lstlisting}
procedure Configure (C : in out Channel; Pin : Pin_Id;
                     Carrier_Hz : Positive := 1_000_000);
procedure Set_Density (C : Channel; Percent : Density_Percent);
\end{lstlisting}

\textbf{Parameters.}
\begin{description}[style=nextline]
\item[\texttt{Index}] channel \texttt{0\,..\,7} (\texttt{Channel\_Index}).
\item[\texttt{Pin}] the output pad.
\item[\texttt{Carrier\_Hz}] the desired pulse-stream (carrier) frequency in Hz;
  default \texttt{1\_000\_000}. The modulator divides the APB clock
  ($\approx$80\,MHz) by an integer \texttt{1\,..\,256}, so the achieved rate is
  the nearest \texttt{APB/N} (roughly \texttt{312\_500\,..\,80\_000\_000}\,Hz),
  rounded for you. A higher carrier sits further above your RC low-pass cutoff and
  smooths to a cleaner level; the exact value rarely matters.
\item[\texttt{Percent}] target density, \texttt{0.0\,..\,100.0}.
\end{description}

\lstexample{A 25\,\% density bitstream (analog-ish level via an RC filter).}
\begin{lstlisting}
with ESP32S3.SDM; use ESP32S3.SDM;
procedure SDM_Density is
   C : Channel;
begin
   Claim (C, 0);
   Configure (C, Pin => 7, Carrier_Hz => 1_000_000);
   Set_Density (C, 25.0);
end SDM_Density;
\end{lstlisting}

\lstexample{Breathe an LED by sweeping density.}
\begin{lstlisting}
with Ada.Real_Time; use Ada.Real_Time;
with ESP32S3.SDM;    use ESP32S3.SDM;
procedure SDM_Breathe is
   C : Channel;
begin
   Claim (C, 0);
   Configure (C, Pin => 7);
   for P in 0 .. 100 loop
      Set_Density (C, Density_Percent (P));
      delay until Clock + Milliseconds (10);
   end loop;
end SDM_Breathe;
\end{lstlisting}

\section{TWAI: CAN 2.0}\index{TWAI driver}

\begin{lstlisting}
subtype Standard_Id is Interfaces.Unsigned_32 range 0 .. 16#7FF#;        -- 11-bit
subtype Extended_Id is Interfaces.Unsigned_32 range 0 .. 16#1FFF_FFFF#;  -- 29-bit
type Standard_Frame is record   -- 11-bit id (CAN 2.0A)
   Id : Standard_Id := 0; Remote : Boolean := False;   -- Remote => RTR request
   Length : Data_Length := 0; Data : Data_Bytes := (others => 0);
end record;
type Extended_Frame is record   -- 29-bit id (CAN 2.0B)
   Id : Extended_Id := 0; Remote : Boolean := False;
   Length : Data_Length := 0; Data : Data_Bytes := (others => 0);
end record;
procedure Acquire (S : in out Session;    -- claim AND configure the controller
                   Mode : Bus_Mode := Normal; Bit_Rate : Positive := 125_000);
procedure Reconfigure (S : Session;       -- re-apply mode + bit rate, held
                       Mode : Bus_Mode := Normal; Bit_Rate : Positive := 125_000);
procedure Send (S : Session; F : Standard_Frame);   -- overloaded on the frame type
procedure Send (S : Session; F : Extended_Frame);
function  Available   (S : Session) return Boolean; -- a frame is waiting
function  Is_Extended (S : Session) return Boolean; -- ... and it is 29-bit
procedure Receive (S : Session; F : out Standard_Frame; Got : out Boolean);
procedure Receive (S : Session; F : out Extended_Frame; Got : out Boolean);
\end{lstlisting}

\textbf{Parameters.}
\begin{description}[style=nextline]
\item[\texttt{Mode}] \texttt{Normal} (default), \texttt{Listen\_Only} (receive
  without ACKing) or \texttt{Self\_Test} (ACK its own frames, the wiring-free
  loopback).
\item[\texttt{Bit\_Rate}] bus speed in bits/s; default \texttt{125\_000} (e.g.
  \texttt{250\_000}, \texttt{500\_000}, \texttt{1\_000\_000}).
\item[\texttt{Tx}, \texttt{Rx}] the bus pads (\texttt{Optional\_Pin}).
\item[\texttt{F} (\texttt{Standard\_Frame} / \texttt{Extended\_Frame})] the frame:
  \texttt{Id} (range-checked to its width: \texttt{Standard\_Id}
  \texttt{0\,..\,16\#7FF\#} or \texttt{Extended\_Id}
  \texttt{0\,..\,16\#1FFF\_FFFF\#}), \texttt{Length} (\texttt{0\,..\,8} bytes),
  \texttt{Data} (an 8-byte array), and \texttt{Remote}: \texttt{False} (default)
  for a data frame, \texttt{True} for a remote-transmission request (RTR), which
  carries \texttt{Id} and \texttt{Length} but no \texttt{Data}. \texttt{Send} is
  overloaded on the frame type.
\item[\texttt{Got}] (\texttt{Receive}) \texttt{out Boolean}: \texttt{True} if a
  frame of the requested width was read.
\item[\texttt{Pad}] (\texttt{Enable\_Loopback}) the single self-test pad.
\end{description}

\textbf{Standard vs.\ extended frames.} A CAN identifier is either 11-bit
(standard) or 29-bit (extended, CAN 2.0B). The two are distinct frame types, so
the id is range-checked to its width and \texttt{Send} is overloaded: you
cannot put a 29-bit id in a standard frame. On receive the \emph{sender} picks
the width, so you peek first: \texttt{Available} reports a waiting frame,
\texttt{Is\_Extended} reports its width, and you call the matching
\texttt{Receive} overload (Example 2). Orthogonal to the width, \texttt{Remote}
marks a \emph{remote-transmission request} (RTR): a frame that carries an id and a
length but no data, asking whichever node owns that id to reply with the data.
Set it on either frame type; a received frame reports it.

\lstexample{Send a standard and an extended frame.}
\begin{lstlisting}
with ESP32S3.TWAI; use ESP32S3.TWAI;
procedure TWAI_Send is
   S : Session;
begin
   Acquire (S, Mode => Normal, Bit_Rate => 500_000);   --  own + configure
   Configure_Pins (S, Tx => 4, Rx => 5);               --  then route the pins
   Send (S, Standard_Frame'(Id => 16#123#, Remote => False, Length => 2,
                            Data => (16#DE#, 16#AD#, others => 0)));
   Send (S, Extended_Frame'(Id => 16#14AB_CDE#, Remote => False, Length => 1,
                            Data => (0 => 16#42#, others => 0)));  -- type picks Send
   Send (S, Standard_Frame'(Id => 16#7A5#, Remote => True,        -- RTR: request
                            Length => 8, Data => (others => 0)));  -- 8 bytes, none sent
end TWAI_Send;
\end{lstlisting}

\lstexample{Self-test loopback (no transceiver, no wiring).}
\begin{lstlisting}
with ESP32S3.TWAI; use ESP32S3.TWAI;
procedure TWAI_Selftest is
   S   : Session;
   RE  : Extended_Frame;
   Got : Boolean := False;
begin
   Acquire (S, Mode => Self_Test);
   Enable_Loopback (S, Pad => 4);
   Send (S, Extended_Frame'(Id => 16#14AB_CDE#, Remote => False, Length => 1,
                            Data => (0 => 16#42#, others => 0)));
   if Available (S) and then Is_Extended (S) then
      Receive (S, RE, Got);   --  Got, RE.Id = 16#14AB_CDE#, RE.Remote = False
   end if;
end TWAI_Selftest;
\end{lstlisting}

\section{Timer: general-purpose timers}\index{Timer driver}

\begin{lstlisting}
procedure Configure (T : in out Timer; Tick_Hz : Positive := 1_000_000);
procedure Start (T : Timer);  function Value (T : Timer) return Ticks;
procedure Set_Alarm (T : Timer; At_Ticks : Ticks);
function  Alarm_Fired (T : Timer) return Boolean;  procedure Clear_Alarm (T : Timer);
\end{lstlisting}

\textbf{Parameters.}
\begin{description}[style=nextline]
\item[\texttt{Index}] \texttt{0} or \texttt{1} (\texttt{Timer\_Index}, the two
  timer groups TIMG0/TIMG1).
\item[\texttt{Tick\_Hz}] counter rate in Hz; default \texttt{1\_000\_000}
  (1 tick = 1\,\textmu s).
\item[\texttt{At\_Ticks} / \texttt{Value}] a 54-bit count (\texttt{Ticks}, an
  \texttt{Unsigned\_64}); \texttt{Set\_Alarm} fires when the counter reaches
  \texttt{At\_Ticks}.
\end{description}

\lstexample{Measure an elapsed interval in microseconds.}
\begin{lstlisting}
with Ada.Text_IO;   use Ada.Text_IO;
with ESP32S3.Timer; use ESP32S3.Timer;
procedure Timer_Measure is
   T       : Timer;
   Start_T : Ticks;
begin
   Claim (T, 0);
   Configure (T, Tick_Hz => 1_000_000);   --  1 tick = 1 us
   Start (T);
   Start_T := Value (T);
   --  ... work ...
   Put_Line ("elapsed us =" & Ticks'Image (Value (T) - Start_T));
end Timer_Measure;
\end{lstlisting}

\lstexample{A one-shot alarm 50\,ms out.}
\begin{lstlisting}
with ESP32S3.Timer; use ESP32S3.Timer;
procedure Timer_Alarm is
   T : Timer;
begin
   Claim (T, 0);
   Configure (T, Tick_Hz => 1_000_000);
   Start (T);
   Set_Alarm (T, Value (T) + 50_000);
   loop exit when Alarm_Fired (T); end loop;
   Clear_Alarm (T);
end Timer_Alarm;
\end{lstlisting}

\section{LCD: 8-bit i80 parallel master}\index{LCD driver}

\begin{lstlisting}
--  Acquire claims the controller AND brings it up (pixel clock + pin routing):
procedure Acquire (S : in out Session; Pclk_Hz : Positive := 1_000_000;
                   Data : Data_Pins := (others => No_Pin);
                   Pclk : Optional_Pin := No_Pin);
procedure Reconfigure (S : Session; ...);  -- re-apply clock + pins, held
procedure Transmit (S : Session; Tx : System.Address; Length : Natural;
                    Ok : out Boolean);
\end{lstlisting}

\textbf{Parameters.}
\begin{description}[style=nextline]
\item[\texttt{Pclk\_Hz}] pixel clock in Hz; default \texttt{1\_000\_000} (the
  achieved value is \texttt{20\,MHz} divided by an integer).
\item[\texttt{Data}] a \texttt{Data\_Pins} value: an 8-element array of
  \texttt{Optional\_Pin} for the data lines D0\,..\,D7.
\item[\texttt{Pclk} / \texttt{Pclk\_Pad}] the pixel-clock pad.
\item[\texttt{Tx}, \texttt{Length}] frame-buffer \texttt{'Address} and a byte
  count, \texttt{1\,..\,4095}.
\item[\texttt{Ok}] \texttt{out Boolean}: \texttt{True} on a completed transfer.
\end{description}

\lstexample{Push a frame buffer to a parallel display.}
\begin{lstlisting}
with Interfaces;
with ESP32S3.LCD; use ESP32S3.LCD;
procedure LCD_Frame is
   S  : Session;
   Ok : Boolean;
   FB : aliased array (0 .. 3999) of Interfaces.Unsigned_8 := (others => 16#FF#);
begin
   Acquire (S, Pclk_Hz => 200_000,
            Data => (1, 2, 3, 4, 5, 6, 7, 8), Pclk => 9);  --  claim + configure
   Transmit (S, FB'Address, FB'Length, Ok);
end LCD_Frame;
\end{lstlisting}

\lstexample{Expose the pixel clock on a pad (for a scope).}
\begin{lstlisting}
with ESP32S3.LCD; use ESP32S3.LCD;
procedure LCD_Clk is
   Lcd : Session;
begin
   Acquire (Lcd, Pclk_Hz => 1_000_000);    -- claim + bring up the clock
   Enable_Clock_Out (Lcd, Pclk_Pad => 9);  -- expose PCLK on IO9
end LCD_Clk;
\end{lstlisting}

\section{ADC: SAR analog input}\index{ADC driver}

\begin{lstlisting}
procedure Claim (R : in out Reader; Unit : ADC_Unit);
function Read (R : Reader; Ch : Channel_Index; Atten : Attenuation := Db_12)
   return Raw_Value;                              --  0 .. 4095 (12-bit)
\end{lstlisting}

\textbf{Parameters.}
\begin{description}[style=nextline]
\item[\texttt{Unit}] \texttt{ADC1} or \texttt{ADC2}.
\item[\texttt{Ch}] channel \texttt{0\,..\,9} (\texttt{Channel\_Index}); each maps to
  a fixed pad.
\item[\texttt{Atten}] input attenuation, trading range for resolution:
  \texttt{Db\_0} ($\sim$1.1\,V full scale), \texttt{Db\_2\_5} ($\sim$1.5\,V),
  \texttt{Db\_6} ($\sim$2.2\,V) or \texttt{Db\_12} ($\sim$3.3\,V, default).
\item[returns] a \texttt{Raw\_Value}, \texttt{0\,..\,4095} (12-bit).
\end{description}

\lstexample{Read a 12-bit sample.}
\begin{lstlisting}
with Ada.Text_IO; use Ada.Text_IO;
with ESP32S3.ADC; use ESP32S3.ADC;
procedure ADC_Read_Sample is
   R : Reader;
   V : Raw_Value;
begin
   Claim (R, ADC1);
   V := Read (R, Ch => 0);     --  0 .. 4095 at full (12 dB) attenuation
   Put_Line (Raw_Value'Image (V));
end ADC_Read_Sample;
\end{lstlisting}

\lstexample{A lower attenuation for a small signal.}
\begin{lstlisting}
with Ada.Text_IO; use Ada.Text_IO;
with ESP32S3.ADC; use ESP32S3.ADC;
procedure ADC_Atten is
   R : Reader;
   V : Raw_Value;
begin
   Claim (R, ADC1);
   V := Read (R, Ch => 3, Atten => Db_0);   --  ~0..1.1 V full scale
   Put_Line (Raw_Value'Image (V));
end ADC_Atten;
\end{lstlisting}

\textbf{Per-attenuation calibration.} The SAR ADC needs a self-calibration
``init code'' --- found once at start-up by grounding the input and binary-searching
the code that reads zero --- and that code is \emph{attenuation-dependent}: each
\texttt{Atten} setting switches a different analog gain in front of the converter.
So the driver calibrates every unit at \emph{every} attenuation once, stores the
resulting code per \texttt{(unit, attenuation)}, and programs the matching one into
the SAR immediately before each \texttt{Read}. Calibrating at a single attenuation
and reusing that one code at the others --- the easy trap --- skews every reading
not taken at the calibrated setting; keying the code to the attenuation the caller
actually asked for keeps all four accurate.

\chapter{Low-Power, Touch and Crypto}
\label{ch:hal3}

These drivers need no finalization, so they build under \emph{every} runtime
profile.

\section{RTC: deep sleep and retained memory}\index{RTC driver}

The RTC domain keeps a small slow-memory block alive across deep sleep, and
drives the wake timer. \texttt{Read}/\texttt{Write} access retained 32-bit words;
\texttt{Deep\_Sleep\_For}/\texttt{Until} enter deep sleep (the chip resets on
wake, so execution resumes at the top of \texttt{Main}); \texttt{Last\_Wake}
reports why it woke.

\begin{lstlisting}
function  Read  (Index : Word_Index) return Interfaces.Unsigned_32;
procedure Write (Index : Word_Index; Value : Interfaces.Unsigned_32);
function  Last_Wake return Wake_Cause;   --  Power_On / Deep_Sleep_Timer / ...
procedure Deep_Sleep_For   (Wake_After : Duration);
procedure Deep_Sleep_Until (Pin : Pin_Id; High : Boolean := True);
\end{lstlisting}

\textbf{Parameters.}
\begin{description}[style=nextline]
\item[\texttt{Index}] retained-word slot, \texttt{0\,..\,2047}
  (\texttt{Word\_Index}, the 8\,KB RTC slow memory as 32-bit words).
\item[\texttt{Value}] the 32-bit word to store (\texttt{Unsigned\_32}).
\item[\texttt{Wake\_After}] (\texttt{Deep\_Sleep\_For}) the sleep length as an Ada
  \texttt{Duration} in \emph{seconds}; approximate, since the RTC slow clock is
  uncalibrated.
\item[\texttt{Pin}, \texttt{High}] (\texttt{Deep\_Sleep\_Until}) the wake pad and
  the level to wake on; \texttt{High => True} (default) wakes on a high level.
\item[\texttt{Last\_Wake}] returns \texttt{Power\_On}, \texttt{Deep\_Sleep\_Timer},
  \texttt{Deep\_Sleep\_GPIO} or \texttt{Other\_Reset}.
\end{description}

\lstcap{Example 1: a boot counter that survives deep sleep.}
\begin{lstlisting}
with Interfaces;     use Interfaces;
with Ada.Text_IO;    use Ada.Text_IO;
with ESP32S3.RTC;    use ESP32S3.RTC;
procedure RTC_Boot is
   N : Interfaces.Unsigned_32;
begin
   if Last_Wake = Deep_Sleep_Timer then
      N := Read (0) + 1;     --  woke from our own timer
   else
      N := 1;                --  cold boot
   end if;
   Write (0, N);
   Put_Line ("boot #" & N'Image);
   Deep_Sleep_For (5.0);     --  Duration in seconds; wakes and re-enters Main
end RTC_Boot;
\end{lstlisting}

\lstcap{Example 2: wake on a button.}
\begin{lstlisting}
with ESP32S3.RTC; use ESP32S3.RTC;
procedure RTC_Wake is
begin
   Deep_Sleep_Until (Pin => 0, High => False);  --  sleep until GPIO0 goes low
end RTC_Wake;
\end{lstlisting}
(After a deep-sleep wake the driver also disarms the super-watchdog
(\texttt{Disable\_Super\_Watchdog}) so a staying-awake app is not reset.)

\section{RTC\_IO: pad hold and RTC pulls}\index{RTC\_IO driver}

Latch a pad's level so it survives \texttt{GPIO} writes and deep sleep, and apply
RTC-domain pulls.

\begin{lstlisting}
procedure Hold (Pin : RTC_Pin);   procedure Release (Pin : RTC_Pin);
function  Is_Held (Pin : RTC_Pin) return Boolean;
procedure Set_Pull (Pin : RTC_Pin; Mode : Pull_Mode);   --  No_Pull / Up / Down
\end{lstlisting}

\textbf{Parameters.}
\begin{description}[style=nextline]
\item[\texttt{Pin}] an RTC-capable pad (\texttt{RTC\_Pin}), restricted to
  \texttt{0\,..\,21} (only those reach the RTC domain).
\item[\texttt{Mode}] (\texttt{Set\_Pull}) the RTC-domain pull: \texttt{No\_Pull},
  \texttt{Up} or \texttt{Down}.
\end{description}

\lstcap{Example 1: hold a relay line high through sleep.}
\begin{lstlisting}
with ESP32S3.GPIO;
with ESP32S3.RTC_IO; use ESP32S3.RTC_IO;
procedure RTCIO_Hold is
begin
   ESP32S3.GPIO.Configure (5, Mode => ESP32S3.GPIO.Output);
   ESP32S3.GPIO.Set (5);
   Hold (5);                 --  GPIO writes now ignored; level frozen
   --  ... deep sleep ...    --  pad stays high across sleep
   Release (5);              --  GPIO control restored
end RTCIO_Hold;
\end{lstlisting}

\lstcap{Example 2: an RTC pull-up on a wake pin.}
\begin{lstlisting}
with ESP32S3.RTC_IO; use ESP32S3.RTC_IO;
procedure RTCIO_Pull is
begin
   Set_Pull (0, Up);
end RTCIO_Pull;
\end{lstlisting}

\section{Touch: capacitive sensing}\index{Touch driver}

\begin{lstlisting}
procedure Setup;                       --  bring up the touch FSM (call once)
procedure Enable (Ch : Channel);       --  channel n is on GPIO n (n = 1 .. 14)
function  Read (Ch : Channel) return Natural;   --  raw self-capacitance count
function  Touched (Ch : Channel; Reference : Natural; Margin : Natural := 20_000)
   return Boolean;
\end{lstlisting}

\textbf{Parameters.}
\begin{description}[style=nextline]
\item[\texttt{Ch}] touch channel \texttt{1\,..\,14} (\texttt{Channel}); channel
  \emph{n} is wired to GPIO \emph{n}.
\item[\texttt{Read}] returns a raw self-capacitance count (\texttt{Natural};
  larger means more capacitance; \texttt{0} means not measured yet).
\item[\texttt{Reference}] (\texttt{Touched}) the untouched baseline count you
  captured earlier with \texttt{Read}.
\item[\texttt{Margin}] how far the live count must deviate from \texttt{Reference}
  to register a touch; default \texttt{20\_000}.
\end{description}

\lstcap{Example 1: read two pads' baselines.}
\begin{lstlisting}
with Ada.Text_IO;   use Ada.Text_IO;
with ESP32S3.Touch; use ESP32S3.Touch;
procedure Touch_Read_Base is
begin
   Setup;  Enable (1);  Enable (3);
   Put_Line ("ch1 =" & Natural'Image (Read (1)) &
             "  ch3 =" & Natural'Image (Read (3)));
end Touch_Read_Base;
\end{lstlisting}

\lstcap{Example 2: detect a finger against a captured baseline.}
\begin{lstlisting}
with Ada.Text_IO;   use Ada.Text_IO;
with ESP32S3.Touch; use ESP32S3.Touch;
procedure Touch_Detect is
   Base : constant Natural := Read (1);   --  with nothing touching
begin
   loop
      if Touched (1, Base, Margin => 50_000) then
         Put_Line ("touched!");
      end if;
   end loop;
end Touch_Detect;
\end{lstlisting}

\section{SHA: hardware hashing}\index{SHA driver}

\begin{lstlisting}
function Hash_1   (Data : Byte_Array) return SHA1_Digest;    --  20 bytes
function Hash_224 (Data : Byte_Array) return SHA224_Digest;  --  28 bytes
function Hash_256 (Data : Byte_Array) return SHA256_Digest;  --  32 bytes
\end{lstlisting}

\textbf{Parameters.} Each \texttt{Hash\_*} takes one \texttt{Data : Byte\_Array} of
\emph{any} length and returns a fixed-size digest: \texttt{SHA1\_Digest}
(20 bytes), \texttt{SHA224\_Digest} (28) or \texttt{SHA256\_Digest} (32).

\lstcap{Example 1: SHA-256 of a message.}
\begin{lstlisting}
with ESP32S3.SHA; use ESP32S3.SHA;
procedure SHA_256_Ex is
   Msg    : constant Byte_Array := (16#61#, 16#62#, 16#63#);  --  "abc"
   Digest : constant SHA256_Digest := Hash_256 (Msg);
begin
   null;   --  Digest is the FIPS-180 vector for "abc"
end SHA_256_Ex;
\end{lstlisting}

\lstcap{Example 2: SHA-1 (e.g. for a legacy checksum).}
\begin{lstlisting}
with ESP32S3.SHA; use ESP32S3.SHA;
procedure SHA_1_Ex is
   Msg : constant Byte_Array := (16#61#, 16#62#, 16#63#);
   D1  : constant SHA1_Digest := Hash_1 (Msg);
begin
   null;
end SHA_1_Ex;
\end{lstlisting}

The accelerator is a shared resource serialised by a protected object, so
concurrent \texttt{Hash} calls from different tasks are safe.

\section{AES: hardware block cipher}\index{AES driver}

\begin{lstlisting}
subtype Key_128 is Key_Bytes (0 .. 15);   subtype Key_256 is Key_Bytes (0 .. 31);
function Encrypt_ECB (Key : Key_Bytes; Plain  : Block) return Block
   with Pre => Supported_Key (Key);        --  16- or 32-byte keys only
function Decrypt_ECB (Key : Key_Bytes; Cipher : Block) return Block
   with Pre => Supported_Key (Key);
\end{lstlisting}

\textbf{Parameters.}
\begin{description}[style=nextline]
\item[\texttt{Key}] the cipher key (\texttt{Key\_Bytes}); its \emph{length} selects
  the cipher: 16 bytes (\texttt{Key\_128}) for AES-128, 32 bytes
  (\texttt{Key\_256}) for AES-256. The \texttt{Supported\_Key} precondition rejects
  any other length (there is no AES-192 on this silicon).
\item[\texttt{Plain} / \texttt{Cipher}] one 16-byte \texttt{Block} to transform; the
  result is a \texttt{Block}.
\end{description}

\lstcap{Example 1: AES-128 encrypt then decrypt one block.}
\begin{lstlisting}
with ESP32S3.AES; use ESP32S3.AES;
procedure AES_128_Ex is
   Key    : constant Key_128 := (16#00#, 16#01#, 16#02#, 16#03#, 16#04#, 16#05#,
                                 16#06#, 16#07#, 16#08#, 16#09#, 16#0A#, 16#0B#,
                                 16#0C#, 16#0D#, 16#0E#, 16#0F#);
   Plain  : constant Block := (others => 16#11#);
   Cipher : constant Block := Encrypt_ECB (Key, Plain);
   Back   : constant Block := Decrypt_ECB (Key, Cipher);   --  Back = Plain
begin
   null;
end AES_128_Ex;
\end{lstlisting}

\lstcap{Example 2: AES-256 (a 32-byte key selects the cipher automatically).}
\begin{lstlisting}
with ESP32S3.AES; use ESP32S3.AES;
procedure AES_256_Ex is
   Key256 : constant Key_256 := (others => 16#2B#);
   Plain  : constant Block := (others => 16#11#);
   C      : constant Block := Encrypt_ECB (Key256, Plain);
begin
   null;
end AES_256_Ex;
\end{lstlisting}

The key parameter is one unconstrained type; its \emph{length} (16 or 32 bytes)
chooses AES-128 or AES-256, and the \texttt{Supported\_Key} precondition rejects
anything else: the contract makes the unsupported AES-192 (absent on this
silicon) a compile/run-time error rather than a silent fallback.

\chapter{Writing a Task-Safe Peripheral Driver}
\label{ch:driver-design}\index{driver design}\index{task-safe driver}

Chapters~\ref{ch:hal}--\ref{ch:hal3} \emph{used} the drivers. This chapter shows
how to \emph{write} one, and -- just as important -- why it is built the way it
is. The shape described here is a project rule: every driver in the HAL for a
peripheral that is \textbf{exclusive} (cannot be shared between tasks) and can be
\textbf{reconfigured after setup} follows it exactly. By the end you should be
able to add a new one and get the same guarantees for free.

\section{The problem}

Consider a typical controller -- a SPI host, a UART, a CAN controller. Three
facts make it dangerous to share:

\begin{enumerate}
\item \textbf{It is one physical resource.} Two tasks driving the same FIFO or
  the same configuration registers at once corrupt each other's transfers.
\item \textbf{It is configured in stages.} You bring it up once (clocks, frame
  format), then later you may re-route pins, change the baud rate, flip a line's
  polarity -- often interleaved with real traffic.
\item \textbf{The registers are global.} On a bare-metal target the peripheral is
  a fixed memory-mapped block. Any code that can name it can poke it.
\end{enumerate}

A naive driver exports \texttt{Configure}, \texttt{Write}, \texttt{Read} as plain
procedures over a global. Nothing stops task B from calling \texttt{Set\_Baud}
while task A is mid-\texttt{Write}, or from starting a transfer it has no business
starting. The corruption is intermittent and miserable to debug. We want a design
where those mistakes are \emph{impossible to write}, not merely discouraged.

\section{The ownership model}\index{ownership gateway}\index{session}\index{capability}

The cure is to make \emph{ownership} a value you hold, and to make every operation
that touches the live controller require it.

\begin{description}[style=nextline]
\item[A \texttt{Session} is a capability.] \texttt{Acquire} hands a task a
  \texttt{Session} -- a limited, controlled handle that represents exclusive
  ownership of the controller. ``Limited'' means it cannot be copied (two tasks
  can never share one), and ``controlled'' means it releases automatically when it
  leaves scope, even on an exception. While you hold it, no other task can
  acquire the controller.
\item[Three things become impossible.] You cannot acquire a controller that was
  never set up (\texttt{Not\_Initialized}); you cannot transfer or reconfigure
  without holding it (\texttt{Not\_Owned}); and you cannot leak the lock, because
  the handle's finalizer always releases.
\end{description}

The \texttt{Session} is the public face. The interesting part is the machinery
behind it that makes ``you cannot reconfigure without holding it'' a structural
fact rather than a rule a future maintainer must remember.

\section{Three layers}\index{Engine pattern}\index{private child}

A compliant driver is three packages:

\begin{enumerate}
\item \textbf{The Engine} (\texttt{ESP32S3.<Drv>.Engine}, a \emph{private child}).
  It owns \emph{all} register access and exposes a \texttt{Bus} handle. It is the
  only unit in the whole program that names the peripheral's registers.
\item \textbf{The State gateway} (a package nested in the parent's \emph{body}). It
  holds the one configured \texttt{Bus} handle and hands it out through a single
  function, \texttt{Owned (S)}, which raises unless \texttt{S} holds the
  controller.
\item \textbf{The public package} (\texttt{ESP32S3.<Drv>}). It declares the
  \texttt{Session}, the exceptions, \texttt{Setup}, \texttt{Acquire}, and the
  operations -- each of which is a one-liner that reaches the hardware only
  through \texttt{Owned}.
\end{enumerate}

The rest of this chapter builds a small driver, \texttt{ESP32S3.Widget}, one layer
at a time.

\begin{figure}[ht]
\centering
\resizebox{\linewidth}{!}{%
\begin{tikzpicture}[node distance=12mm]
\node[blk, text width=2.4cm] (app) {Application\\task};
\node[blk, right=of app, fill=blue!8, text width=3.4cm] (pub)
  {Public package\\\scriptsize limited \texttt{Session} +\\\scriptsize protected \texttt{Guard}};
\node[blk, right=of pub, fill=blue!8, text width=2.8cm] (gw)
  {\texttt{Owned (S)}\\\scriptsize gateway\\\scriptsize (raises if not held)};
\node[blk, right=of gw, fill=orange!10, text width=2.6cm] (eng)
  {\texttt{Engine}\\\scriptsize private child};
\node[blk, right=of eng, fill=gray!12, text width=2.0cm] (reg) {Registers};
\draw[flow] (app) -- node[above,font=\scriptsize]{Acquire} (pub);
\draw[flow] (pub) -- (gw);
\draw[flow] (gw)  -- (eng);
\draw[flow] (eng) -- (reg);
\end{tikzpicture}}
\caption{The task-safe driver skeleton: a transfer reaches the registers only via
the ownership-checked \texttt{Owned} gateway, and only the private \texttt{Engine}
names them. Forgetting the check is structurally impossible.}
\label{fig:driver-skeleton}
\end{figure}

\section{Layer 1 -- the Engine (the only code that sees the registers)}

The engine is a \emph{private child}: a unit named \texttt{ESP32S3.Widget.Engine}
can be \texttt{with}-ed only from inside the \texttt{ESP32S3.Widget} subtree.
Application code -- and even other drivers -- physically cannot reach it.

\begin{lstlisting}
private package ESP32S3.Widget.Engine is
   type Bus is private;                       --  an opaque, configured controller

   function  Open (Clock_Hz : Positive) return Bus;
   procedure Configure_Pins (B : Bus; Tx, Rx : ESP32S3.GPIO.Optional_Pin);
   procedure Transfer (B : Bus; Tx, Rx : System.Address; Length : Natural);
private
   type Bus is record
      Regs  : Periph_Ref := null;             --  pointer to the register block
      Valid : Boolean    := False;
   end record;
end ESP32S3.Widget.Engine;
\end{lstlisting}

The engine \emph{body} is the one place that \texttt{with}s the SVD register
package. It does the bit-twiddling -- the same code you would have written
anyway -- but now it is sealed behind the \texttt{Bus} handle. \texttt{Open}
brings the controller up and returns a \texttt{Bus}; every other operation takes a
\texttt{Bus} and works through it. Crucially, \textbf{the parent package never
\texttt{with}s the SVD registers at all} -- so once the engine exists, the global
controller is unreachable from anywhere except through a \texttt{Bus} handle.

\section{Layer 2 -- ownership: the Session and the Guard}\index{protected object}\index{limited type}\index{controlled type}

The public package declares the handle and a per-controller lock. The
\texttt{Session} is limited and controlled:

\begin{lstlisting}
package ESP32S3.Widget is
   type Session is limited private;

   Not_Initialized : exception;   --  Acquire before Setup
   Not_Owned       : exception;   --  use without holding the controller

   procedure Setup (Clock_Hz : Positive := 1_000_000);          --  startup only
   procedure Acquire (S : in out Session);
   procedure Configure_Pins (S : Session; Tx, Rx : ESP32S3.GPIO.Optional_Pin);
   procedure Transfer (S : Session; Tx, Rx : System.Address; Length : Natural);
   procedure Release (S : in out Session);
private
   type Session is new Ada.Finalization.Limited_Controlled with record
      Active : Boolean := False;
   end record;
   overriding procedure Finalize (S : in out Session);          --  auto-release
end ESP32S3.Widget;
\end{lstlisting}

In the body, a protected object arbitrates ownership. Its guarded section is
tiny -- it only flips a flag -- so a task is never blocked inside it while a
slow transfer runs:

\begin{lstlisting}
protected Guard is
   entry    Acquire;
   procedure Release;
private
   Held : Boolean := False;
end Guard;

protected body Guard is
   entry Acquire when not Held is begin Held := True;  end Acquire;
   procedure Release         is begin Held := False; end Release;
end Guard;
\end{lstlisting}

\texttt{entry Acquire when not Held} is the whole mutual-exclusion story: a second
task calling \texttt{Acquire} \emph{suspends} on the barrier until the first
releases. (For a peripheral with several identical instances -- e.g. two SPI
hosts -- this is an \texttt{array (Host) of Guard} and a matching array of
handles; the singleton case is shown here for clarity.)

\section{Layer 3 -- the gateway that cannot be bypassed}\index{Owned gateway}

Here is the heart of the design. The configured \texttt{Bus} lives inside a
package nested in the \emph{body}, so no other part of the driver can name it. The
only way out is \texttt{Owned}:

\begin{lstlisting}
package State is
   procedure Open (Clock_Hz : Positive);       --  the Setup path
   function  Ready return Boolean;
   function  Owned (S : Session) return E.Bus; --  THE gateway
end State;

package body State is
   The_Bus : E.Bus;                            --  hidden here; unreachable elsewhere
   Is_Set  : Boolean := False;

   procedure Open (Clock_Hz : Positive) is
   begin
      The_Bus := E.Open (Clock_Hz);
      Is_Set  := True;
   end Open;

   function Ready return Boolean is (Is_Set);

   function Owned (S : Session) return E.Bus is
   begin
      if not S.Active then
         raise Not_Owned with "Widget used without holding it -- Acquire first";
      end if;
      return The_Bus;
   end Owned;
end State;
\end{lstlisting}

Now every public operation is a single line that funnels through \texttt{Owned},
and the ownership check happens in exactly one place:

\begin{lstlisting}
procedure Setup (Clock_Hz : Positive := 1_000_000) is
begin
   State.Open (Clock_Hz);
end Setup;

procedure Acquire (S : in out Session) is
begin
   if not State.Ready then
      raise Not_Initialized with "Widget acquired before Setup";
   end if;
   Guard.Acquire;                              --  suspends until free
   S.Active := True;
end Acquire;

procedure Configure_Pins (S : Session; Tx, Rx : ESP32S3.GPIO.Optional_Pin) is
begin
   E.Configure_Pins (State.Owned (S), Tx, Rx); --  Owned raises unless held
end Configure_Pins;

procedure Transfer (S : Session; Tx, Rx : System.Address; Length : Natural) is
begin
   E.Transfer (State.Owned (S), Tx, Rx, Length);
end Transfer;

procedure Release (S : in out Session) is
begin
   if S.Active then
      S.Active := False;
      Guard.Release;
   end if;
end Release;

overriding procedure Finalize (S : in out Session) is
begin
   Release (S);                                --  scope exit / exception unwind
end Finalize;
\end{lstlisting}

Look at \texttt{Configure\_Pins} and \texttt{Transfer}: there is \emph{no}
\texttt{if S.Active then} in either. They cannot reach the hardware without
calling \texttt{State.Owned (S)}, and \texttt{Owned} is the one place that checks.
A maintainer adding a new operation tomorrow has no choice but to go through the
same door -- there is no \texttt{Bus} in scope to grab any other way.

\section{The limited-handle variant}

If the engine's \texttt{Bus} is a \emph{limited} type -- typically because it
holds something limited, like a finalization-controlled DMA channel -- it cannot
be returned by value. Make the hidden instance \texttt{aliased} and have
\texttt{Owned} return an access to it:

\begin{lstlisting}
The_Bus : aliased E.Bus;
function Owned (S : Session) return access E.Bus is
begin
   if not S.Active then raise Not_Owned with "..."; end if;
   return The_Bus'Access;
end Owned;
--  callers: E.Transfer (State.Owned (S).all, Tx, Rx, Length);
\end{lstlisting}

In the HAL, SPI and I2S (whose engine \texttt{Bus} is limited) and LCD (whose
handle owns a GDMA channel) use this form; I2C, UART and TWAI return the handle by
value. Nothing else about the pattern changes.

\section{Why it is built this way}\index{RAII}

Each layer earns its place by removing a specific class of mistake:

\begin{description}[style=nextline]
\item[Private-child engine $\rightarrow$ the registers have one owner.] Because the
  parent never \texttt{with}s the SVD package, the controller global is
  unreachable outside the engine. The ``any code can poke it'' hazard is gone at
  the visibility level, not by convention.
\item[Hidden handle + single \texttt{Owned} gateway $\rightarrow$ the guard cannot be
  forgotten.] The earlier, weaker approach put an \texttt{if S.Active} check at the
  top of each operation. That works until someone adds an operation and forgets
  the check. Hiding the only \texttt{Bus} behind \texttt{Owned} makes the check
  structural: forgetting it is not a bug you can write, because there is no other
  way to obtain a \texttt{Bus}.
\item[Limited \texttt{Session} $\rightarrow$ ownership cannot be duplicated.] A
  non-copyable handle means two tasks can never hold the same controller, and a
  handle can never outlive its scope as a dangling copy.
\item[Controlled \texttt{Session} (RAII) $\rightarrow$ the lock cannot leak.] The
  finalizer releases on every exit path, including exception unwinding. A fault
  between \texttt{Acquire} and \texttt{Release} hands the controller back instead
  of deadlocking the next task forever.
\item[Tiny protected section $\rightarrow$ ownership without serialising work.] The
  lock is held only long enough to flip a flag; the blocking transfer runs
  outside it. Tasks contend for \emph{ownership}, never for the bus itself.
\item[Two named exceptions $\rightarrow$ misuse fails loudly.] \texttt{Not\_Initialized}
  and \texttt{Not\_Owned} turn ``used in the wrong state'' into an immediate,
  explanatory error instead of a silent no-op or a corrupted transfer.
\end{description}

\section{The benefits}

\begin{itemize}
\item \textbf{Mistakes become unrepresentable.} ``Configure a port you don't
  own,'' ``transmit from the wrong task,'' ``poke the registers directly'' --
  none of these can be written. The compiler and the visibility rules enforce
  the discipline, so code review does not have to.
\item \textbf{One pattern, every peripheral.} Because all six \texttt{Session}
  drivers share this skeleton, learning one teaches all of them, and a new driver
  is a fill-in-the-blanks exercise rather than a fresh concurrency design.
\item \textbf{The unsafe core is testable in isolation.} The engine is plain,
  lock-free register logic; the parent is plain ownership logic. Each can be
  reasoned about -- and the engine bench-tested -- on its own.
\item \textbf{RAII hygiene for free.} Scope-based release means the common
  ``acquire, do work, release'' path needs no explicit \texttt{Release} and is
  exception-safe by construction.
\item \textbf{It costs almost nothing.} The handle is a couple of words, the
  gateway is one comparison, and the protected section is a flag flip. The safety
  is structural, not runtime overhead.
\end{itemize}

\section{When the rule applies (and when it does not)}

Use this pattern for a driver that is \textbf{both} exclusive (cannot be shared
between tasks) \textbf{and} reconfigurable after setup. That is the combination
that makes uncoordinated access dangerous.

It does \emph{not} apply to lock-free or shareable peripherals. \texttt{GPIO}'s
set/clear/read are hardware-atomic; \texttt{RNG} is a stateless read; the
\texttt{Temperature}, \texttt{RTC} and crypto blocks either are internally
serialised by a small protected object or have nothing a second task can corrupt.
Those drivers stay simple on purpose and run under every runtime profile,
including light-tasking, where the controlled \texttt{Session} types (which need
finalization) are not available. Reach for the full ownership machinery only when
the hazard is real; for everything else, the simplest safe thing is better.

\section{Letting the application plug in: dispatch or callback}
\label{driver:plugin}\index{callback}\index{dispatching}\index{closure-free}

A driver often needs the \emph{application} to supply a piece of behaviour: assert
a chip-select line, pin a fresh socket to one interface, decide what a Modbus
register read returns. In most languages the reach-for tool is a \emph{closure}.

A closure is a subprogram bundled with the variables it uses but does not itself
declare -- the surrounding state it ``closes over.'' A nested subprogram that reads
a variable from its enclosing scope is the everyday case, and on a driver boundary
it shows up like this:

\begin{lstlisting}
procedure Configure (Bus : in out E.Bus; CS : ESP32S3.GPIO.Pin_Id) is
   procedure Select_It (Active : Boolean) is
   begin
      ESP32S3.GPIO.Write (CS, Active);     -- CS is captured from Configure
   end Select_It;
begin
   Open (Bus, Chip_Select => Select_It'Access);   -- hand the nested proc out
end Configure;
\end{lstlisting}

\texttt{Select\_It} declares only \texttt{Active} but uses \texttt{CS}, so to be
callable later it has to remember \emph{which} \texttt{CS} -- the one from the
\texttt{Configure} call that produced it. That pairing of code with captured state
is the closure. A subprogram that captures \emph{nothing} -- using only its own
parameters, globals and library-level state -- is not a closure but a plain
function pointer, and that distinction is the whole of what follows.

A closure is exactly what this platform forbids. A subprogram that captures its
environment compiles to a \emph{trampoline} -- a scrap of executable code GCC
writes onto the stack, which is where the captured \texttt{CS} gets bound -- and on
the ESP32-S3 the stack lives in non-executable data memory, so calling such a
trampoline faults outright. The HAL is built under
\texttt{pragma Restrictions (No\_Implicit\_Dynamic\_Code)}, which turns that into a
\emph{compile-time} refusal rather than a runtime crash (the same fault, seen from
the debugger's side, is the war story in Chapter~\ref{ch:debug}).

So a driver's extension point must need \emph{no} captured state. Two shapes
satisfy that, and which one a driver picks comes down to how much per-instance
state the hook carries.

\textbf{A closure-free callback with an explicit context.} The hook is a
library-level access-to-subprogram, and any state it needs travels through an
opaque context parameter the driver hands back to it -- not captured, passed. The
SPI driver's application chip-select is this shape (Chapter~\ref{ch:hal}):

\begin{lstlisting}
type CS_Select is access procedure (Ctx : System.Address; Active : Boolean);
--  the app's procedure is library-level; per-device state rides in Ctx
\end{lstlisting}

It is light and runs under every profile. The cost is that \texttt{Ctx} is
untyped -- the callback converts it back itself -- so it suits a hook that is
essentially stateless or one-shot: ``drive this pin,'' ``call \texttt{Set\_Interface}
on this socket.''

\textbf{Dispatching through a tagged type.} You derive a type, keep your data in
its fields, and override the operations you support. The Modbus slave is this
shape (Chapter~\ref{ch:modbus}):

\begin{lstlisting}
type Device is new Modbus.Slave.Server with record
   Holding : Word_Array (0 .. 63) := (others => 0);   -- your storage, in the object
end record;
overriding procedure On_Read_Holding_Registers (Self : in out Device; ...);
\end{lstlisting}

Here \texttt{Self} \emph{is} the context: no \texttt{Ctx} to thread, no conversion,
and the call dispatches through a static vtable -- a table, not generated code, so
\texttt{No\_Implicit\_Dynamic\_Code} is satisfied. The set of primitives you
override expresses exactly which operations the device implements; the ones you
leave alone take the base type's default (for Modbus, an \texttt{Illegal\_Function}
reply). This suits a hook whose context is rich and structured, or a \emph{family}
of related operations with sensible defaults.

The rule of thumb: a stateless or one-shot hook wants the closure-free callback
(carry any scrap of context in an opaque pointer); a hook with rich per-instance
state, or a group of related operations, wants a tagged type you derive and
override. Both exist for one reason -- to recover the convenience of a capturing
closure without the trampoline this chip cannot run.

\chapter{Tasking and Concurrency}
\label{ch:tasking}

Concurrency in Ada is part of the language, not a bolted-on RTOS API: \emph{tasks}
for independent threads of control, \emph{protected objects} for race-free sharing,
interrupt handlers as protected procedures, and a monotonic clock for real-time
work. This chapter shows the practical mechanics on the dual-core ESP32-S3 --
creating tasks and interrupt handlers, periodic scheduling, waiting with timeouts,
the recurring concurrency patterns, watchdogs, and putting both cores to work. It
draws on the encapsulation features of Chapter~\ref{ch:encap} and the drivers of
the HAL chapters.

\section{Tasks and interrupt handlers}
\label{sec:tasks}\index{tasks}

Concurrency in Ada is in the language, not an RTOS API: you declare \emph{tasks}
for independent threads of control and \emph{interrupt handlers} as protected
procedures. This section shows the practical mechanics; the Interrupts chapter has
the full story on how attachment works underneath.

\subsection{Creating and starting a task}

You do not ``start'' an Ada task with a call -- you \emph{declare} it, and it
activates automatically. A single task is declared with its body; it begins
running when its master (the enclosing scope) reaches the code after the
declarations. For a \emph{library-level} task that master is the program itself,
so the binder activates it before \texttt{Main} even begins.

\cnote{there is no \texttt{xTaskCreate}, \texttt{pthread\_create}, or
\texttt{vTaskStartScheduler}: you declare the task and it activates itself, and
the Ada runtime \emph{is} the scheduler -- already running when \texttt{Main}
begins.}

\begin{lstlisting}
   task Blinker;                       -- declaration: this creates the task
   task body Blinker is
   begin                              -- runs concurrently once activated
      loop
         ESP32S3.GPIO.Toggle (2);
         delay until Clock + Milliseconds (250);
      end loop;
   end Blinker;
\end{lstlisting}

Add a \texttt{pragma Priority} (or the \texttt{Priority} aspect) to place it in
the priority order, and on this dual-core target pin it to a core with
\texttt{CPU => 1} or \texttt{CPU => 2}:

\begin{lstlisting}
   task Sampler with Priority => 5, CPU => 2;   -- runs on the second core
\end{lstlisting}

For a \emph{family} of identical workers use a task \emph{type} and declare
objects of it. One restriction matters here: under the \texttt{light-tasking} and
\texttt{embedded} (Jorvik) profiles a task may not be nested inside a subprogram
(\texttt{No\_Task\_Hierarchy}), so application tasks live at \emph{library level}
-- in a package -- not inside \texttt{Main}.

\subsection{Creating and attaching an interrupt handler}

An interrupt handler is a \emph{protected procedure} bound to an interrupt. The
protected object's \texttt{Interrupt\_Priority} is the hardware level it runs at
(and the ceiling that serialises it), and \texttt{pragma Attach\_Handler} wires
it up at elaboration -- no run-time call (\S\ref{ch:registers} and the Interrupts
chapter):

\begin{lstlisting}
   protected Handler with Interrupt_Priority => Device_L3_Priority is
   private
      procedure ISR;
      pragma Attach_Handler (ISR, Device_L3_0);   -- bound to the L3 device slot
   end Handler;
\end{lstlisting}

For GPIO pins the HAL wraps that boilerplate: \texttt{ESP32S3.GPIO.Interrupts}
owns the one GPIO interrupt source, routes it to the level-3 slot, and demuxes by
pin to per-pin callbacks. You just register a callback with a trigger:

\begin{lstlisting}
   Enable (Pin => 0, On => Falling_Edge, Action => On_Edge'Access);
\end{lstlisting}

A handler -- protected \texttt{ISR} or HAL callback -- runs in \emph{interrupt
context} at the slot's ceiling. Keep it tiny, and do not call a lower-ceiling
protected object or block. The idiom is to record the event (bump an
\texttt{Atomic}, or \texttt{Set} a \texttt{Suspension\_Object}) and let an
ordinary task do the real work.

\subsection{Worked example: counting falling edges on GPIO0}

Putting both together: a counter that increments on every falling edge of GPIO0,
with a task that reports it once a second. The interrupt callback is the only
writer of \texttt{Edges} (the level-3 ISR cannot re-enter itself), so an
\texttt{Atomic} \texttt{Natural} written there and read by the task needs no lock.
Everything is library-level, so it is legal on the \texttt{embedded} profile.

\begin{lstlisting}
with Ada.Real_Time;           use Ada.Real_Time;
with Ada.Text_IO;             use Ada.Text_IO;
with ESP32S3.GPIO;            use ESP32S3.GPIO;
with ESP32S3.GPIO.Interrupts; use ESP32S3.GPIO.Interrupts;

package body Edge_Demo is

   Edges : Natural := 0 with Atomic;     -- written by the ISR, read by the task

   --  Runs in interrupt context (level-3 ceiling): do the minimum.
   procedure On_Edge is
   begin
      Edges := Edges + 1;                -- sole writer -> the inc is safe
   end On_Edge;

   --  A normal task, activated automatically at start-up.
   task Reporter;
   task body Reporter is
   begin
      loop
         delay until Clock + Seconds (1);
         Put_Line ("GPIO0 falling edges:" & Natural'Image (Edges));
      end loop;
   end Reporter;

begin
   --  Package elaboration: configure the pin and attach the interrupt.
   Configure (0, Mode => Input, Pull => Pull_Up);          -- pulled-up input
   Enable (Pin => 0, On => Falling_Edge, Action => On_Edge'Access);
end Edge_Demo;
\end{lstlisting}

\paragraph{Reading it line by line.} For a reader new to Ada, here is every
construct in the example:

\begin{description}[style=nextline]
\item[\texttt{with Ada.Real\_Time;}] names a unit this code \emph{depends on} and
  makes its declarations available -- like an \texttt{\#include}, but semantic:
  the compiler records the dependency rather than pasting text. The four
  \texttt{with}s pull in timing (\texttt{Ada.Real\_Time}), console output
  (\texttt{Ada.Text\_IO}), the GPIO driver, and its interrupt child.
\item[\texttt{use Ada.Real\_Time;}] makes that unit's names usable \emph{without}
  the package prefix (see below). \texttt{with} grants access; \texttt{use} grants
  short names.
\item[\texttt{package body Edge\_Demo is}] begins the implementation of a package
  called \texttt{Edge\_Demo}. A package groups related declarations; everything up
  to \texttt{end Edge\_Demo;} belongs to it. The task it declares is legal under the Jorvik profile because it is a top-level
(\emph{library-level}) package.
\item[\texttt{Edges : Natural := 0 with Atomic;}] declares a variable. The colon
  separates the name from its type; \texttt{Natural} is the predefined type of
  integers $\geq 0$; \texttt{:= 0} sets its initial value (\texttt{:=} is
  \emph{assignment}; plain \texttt{=} is the equality \emph{test}). \texttt{with
  Atomic} is an aspect making each read or write a single indivisible access.
\item[\texttt{-- \ldots}] a comment: \texttt{--} runs to the end of the line.
\item[\texttt{procedure On\_Edge is \ldots\ begin \ldots\ end On\_Edge;}] defines a
  \emph{procedure} -- a subprogram that returns no value -- taking no
  parameters. The statements between \texttt{begin} and \texttt{end} run when it
  is called. Repeating the name after \texttt{end} is optional but conventional.
\item[\texttt{Edges := Edges + 1;}] an assignment statement: evaluate
  \texttt{Edges + 1} and store it back. Statements end with a semicolon.
\item[\texttt{task Reporter;}] declares a \emph{task} -- an independent thread of
  control. This line creates it; it will start on its own (no ``start'' call).
\item[\texttt{task body Reporter is \ldots\ loop \ldots\ end loop; \ldots}] the
  task's code. \texttt{loop \ldots\ end loop;} repeats indefinitely.
\item[\texttt{delay until Clock + Seconds (1);}] suspends the task until an
  \emph{absolute} instant: the current time (\texttt{Clock}) plus one second
  (\texttt{Seconds (1)} is a time span). \texttt{delay until} waits to a point in
  time; plain \texttt{delay X} would wait a duration.
\item[\texttt{Put\_Line ("\ldots" \& Natural'Image (Edges));}] prints a line.
  \texttt{\&} concatenates strings, and \texttt{Natural'Image (Edges)} converts
  the number to text. The apostrophe introduces an \emph{attribute} -- a built-in
  query on a type or object; \texttt{'Image} is ``the textual image of this
  value.''
\item[\texttt{begin \ldots\ end Edge\_Demo;}] a package body may end with an
  executable part. It runs \emph{once, automatically}, when the package is
  elaborated at start-up (before \texttt{Main}). Here it brings up the hardware.
\item[\texttt{Configure (0, Mode => Input, Pull => Pull\_Up);}] calls the GPIO
  driver. \texttt{0} is a positional argument (the pin); \texttt{Mode => Input}
  and \texttt{Pull => Pull\_Up} are \emph{named} associations
  (\texttt{name => value}), so the call documents which parameter is which.
\item[\texttt{Enable (\ldots, Action => On\_Edge'Access);}] registers the
  interrupt. \texttt{On\_Edge'Access} is the \texttt{'Access} attribute -- a
  reference to the \texttt{On\_Edge} procedure -- which the HAL stores and calls
  from the ISR on each falling edge.
\end{description}

\paragraph{Why \texttt{use} makes it readable.} A \texttt{use} clause makes a
package's names directly visible, so you write \texttt{Configure},
\texttt{Falling\_Edge} and \texttt{Put\_Line} instead of their full paths. Without
it, the same two setup lines would read:

\begin{lstlisting}
   ESP32S3.GPIO.Configure (0, Mode => ESP32S3.GPIO.Input,
                              Pull => ESP32S3.GPIO.Pull_Up);
   ESP32S3.GPIO.Interrupts.Enable
     (Pin => 0, On => ESP32S3.GPIO.Interrupts.Falling_Edge,
      Action => On_Edge'Access);
\end{lstlisting}

The package path repeats on nearly every name and buries the intent. With
\texttt{use} the code reads in the language of the problem -- ``configure pin 0
as a pulled-up input; enable a falling-edge interrupt'' -- while the \texttt{with}
clauses still record exactly where those names come from. The trade-off is that a
name's origin is less explicit, and two \texttt{use}d packages could in principle
declare the same name; in a small, focused unit like this the clarity gain wins,
and you can always write a single name in full where it aids understanding.

The application's \texttt{Main} need only \texttt{with Edge\_Demo;} and stay
alive; \texttt{Reporter} and the interrupt are running before \texttt{Main}'s
first statement. Each time GPIO0 goes high$\rightarrow$low (release of a pull-up
button, say) the ISR bumps \texttt{Edges}, and the reporter prints the running
total -- a complete interrupt-driven counter in a couple of dozen lines, with no
RTOS, no manual ISR vectoring, and no explicit task start.

\section{Real-time tasks: periodic and time-driven}
\label{sec:realtime}\index{real-time tasks}

For periodic work use \texttt{Ada.Real\_Time} -- a \emph{monotonic} clock that
never jumps -- not \texttt{Ada.Calendar}. The canonical cyclic task computes its
\emph{next} release as an absolute instant and \texttt{delay until}s it, so timing
error does not accumulate:

\begin{lstlisting}
   task body Sampler is
      Period : constant Time_Span := Milliseconds (10);
      Next   : Time := Clock;             -- first release = now
   begin
      loop
         Read_And_Process;
         Next := Next + Period;           -- the next absolute instant
         delay until Next;                -- no drift, however long the body took
      end loop;
   end Sampler;
\end{lstlisting}

Because \texttt{Next} advances by exactly \texttt{Period} regardless of how long
\texttt{Read\_And\_Process} ran, the task holds its average rate with no cumulative
jitter -- the crucial difference from a naive \texttt{delay Period}, which slips
by the body's run time every cycle. Two more tools are available on all profiles:

\begin{itemize}
  \item \texttt{Time\_Span} arithmetic (\texttt{Milliseconds},
        \texttt{Microseconds}, \texttt{To\_Duration}) for expressing intervals
        exactly;
  \item \texttt{Ada.Real\_Time.Timing\_Events} -- register a procedure to run at
        a future instant without dedicating a task, handy for a one-shot timeout.
\end{itemize}

\texttt{Ada.Calendar} exists too, for wall-clock dates, but its absolute form is
unsuited as a scheduling base (see the Limitations chapter) and on the bare
runtime it is anchored only to boot. Schedule with \texttt{Ada.Real\_Time}.

\section{Waiting for an event -- and timing out}
\label{sec:waiting}\index{waiting for an event}

Two questions come up constantly in firmware: how do I \emph{wait} for something
to happen -- a GPIO pin to change, a byte to arrive on the UART -- and how do I
\emph{give up} after a while? There are two ways to wait (poll, or block until an
interrupt wakes you), and the clean way to bound the wait with a timeout differs
across the three profiles because the leaner two forbid the \texttt{select}
statement.

\subsection{Polling (any profile)}

The simplest wait re-checks a condition in a loop, backing off with
\texttt{delay until} so other tasks run; a deadline turns it into a timeout. To
wait for a UART byte:

\begin{lstlisting}
with Ada.Real_Time; use Ada.Real_Time;
with ESP32S3.UART;  use ESP32S3.UART;

function Wait_For_Byte (S : Session; Timeout : Time_Span) return Boolean is
   Deadline : constant Time := Clock + Timeout;
begin
   while Available (S) = 0 loop
      if Clock >= Deadline then
         return False;                       -- gave up
      end if;
      delay until Clock + Milliseconds (1);  -- back off; let other tasks run
   end loop;
   return True;                              -- a byte is waiting
end Wait_For_Byte;
\end{lstlisting}

The same shape waits on a GPIO level -- replace \texttt{Available (S) = 0} with
\texttt{not ESP32S3.GPIO.Read (Pin)}. Polling works on \emph{every} profile and
needs only \texttt{Ada.Real\_Time}; the cost is the CPU spent re-checking.

Rather than hand-roll the deadline loop each time, the HAL bundles it as
\texttt{ESP32S3.Wait.Until\_True} (lock-free, all profiles): give it a condition
function and a timeout and it returns \texttt{True} on success or \texttt{False}
on timeout.

\begin{lstlisting}
with ESP32S3.Wait;
function Byte_Ready return Boolean is (Available (S) > 0);
...
if ESP32S3.Wait.Until_True (Byte_Ready'Access, Milliseconds (200)) then
   Read (S, Buf, N);     -- a byte arrived in time
end if;
\end{lstlisting}

\subsection{Blocking until an interrupt signals you (any profile)}

Better, when the wait may be long or power matters: sleep until the hardware wakes
you. Attach an interrupt handler that signals a \texttt{Suspension\_Object}, and
suspend the task on it:

\begin{lstlisting}
with Ada.Synchronous_Task_Control;  use Ada.Synchronous_Task_Control;
with ESP32S3.GPIO.Interrupts;

Edge_Seen : Suspension_Object;

procedure On_Edge is               -- runs in interrupt context: keep it tiny
begin
   Set_True (Edge_Seen);
end On_Edge;
...
ESP32S3.GPIO.Interrupts.Enable
  (Pin => 0, On => ESP32S3.GPIO.Interrupts.Falling_Edge, Action => On_Edge'Access);
...
Suspend_Until_True (Edge_Seen);    -- the task sleeps here until the edge fires
\end{lstlisting}

This burns no CPU while waiting. But \texttt{Suspend\_Until\_True} has no timeout
of its own -- and \emph{that} is where the profiles diverge.

\subsection{Timing out: three profiles, three approaches}

\textbf{Jorvik (light-tasking).} No \texttt{select} and no exception propagation,
so the timeout is a \emph{value} and the wait blocks on a protected entry whose
barrier opens on either the event or a deadline. The deadline is an
\texttt{Ada.Real\_Time.Timing\_Events} handler; the task genuinely sleeps, and on
return a flag says which fired:

\begin{lstlisting}
with Ada.Real_Time.Timing_Events; use Ada.Real_Time.Timing_Events;

protected Gate is
   procedure Signal;                            -- from the interrupt callback
   procedure Expire (E : in out Timing_Event);  -- the deadline handler
   entry Wait (Timed_Out : out Boolean);
private
   Occurred, Expired : Boolean := False;
end Gate;

protected body Gate is
   procedure Signal is begin Occurred := True; end Signal;
   procedure Expire (E : in out Timing_Event) is begin Expired := True; end Expire;
   entry Wait (Timed_Out : out Boolean) when Occurred or Expired is
   begin
      Timed_Out := not Occurred;
      Occurred := False;  Expired := False;     -- re-arm for next time
   end Wait;
end Gate;

--  Use: the interrupt callback calls Gate.Signal; the caller arms a deadline.
declare
   Deadline : Timing_Event;
   Late     : Boolean;
begin
   Set_Handler (Deadline, Clock + Milliseconds (500), Gate.Expire'Access);
   Gate.Wait (Late);          -- sleeps until the edge OR 500 ms, whichever first
   if Late then ... else ... end if;
end;
\end{lstlisting}

\textbf{embedded.} The same Jorvik primitives (still no \texttt{select}), but
exceptions now propagate -- so a timeout reads naturally as a \emph{catchable
exception}, giving linear code, with any controlled handles released on the way
out:

\begin{lstlisting}
Timeout_Error : exception;

procedure Wait_For_Byte (S : Session; Timeout : Time_Span) is
   Deadline : constant Time := Clock + Timeout;
begin
   while Available (S) = 0 loop
      if Clock >= Deadline then
         raise Timeout_Error with "no byte within the deadline";
      end if;
      delay until Clock + Milliseconds (1);
   end loop;
end Wait_For_Byte;
...
begin
   Wait_For_Byte (S, Milliseconds (200));
   Read (S, Buf, N);                 -- a byte is here
exception
   when Timeout_Error =>             -- handle the timeout in one place
      ...
end;
\end{lstlisting}

(You can raise the same exception from the \texttt{Gate.Wait} barrier above
instead of polling, if you would rather sleep than spin.)

\textbf{full.} Full Ada gives the timeout directly, as a \emph{timed entry call}:
the task blocks on the entry and a \texttt{delay} races it -- no helper deadline,
no polling. The protected object only needs the event half:

\begin{lstlisting}
protected Gate is
   procedure Signal;          -- from the interrupt callback
   entry Wait;
private
   Occurred : Boolean := False;
end Gate;

protected body Gate is
   procedure Signal is begin Occurred := True; end Signal;
   entry Wait when Occurred is begin Occurred := False; end Wait;
end Gate;

--  Use: whichever happens first wins.
select
   Gate.Wait;                              -- the event arrived
   ...
or
   delay until Clock + Milliseconds (500); -- the timeout
   ...
end select;
\end{lstlisting}

(Use \texttt{delay until} -- an \emph{absolute} time -- here: on this port a
\emph{relative} \texttt{delay} as a \texttt{select} trigger has a known gap, see
the Limitations chapter.) \texttt{full} also offers \texttt{select \dots\ then
abort} to put a deadline around an \emph{arbitrary} computation, not just an entry
call -- abandoning it if it overruns.

\textbf{Choosing.} Poll for a short, infrequent wait or the leanest code;
interrupt-signalled blocking is better when the wait may be long or power matters.
For the timeout: a status value on Jorvik, a catchable exception on embedded, and
the \texttt{select}\,/\,\texttt{delay} timed entry call on full.

\section{Choosing how to wait: a spectrum, by timescale}
\label{sec:wait-spectrum}\index{busy-wait}\index{polling!spectrum}\index{spin loop}

The previous section drew the choice as binary -- poll, or block on an interrupt.
That is the right coarse split, but a driver author actually picks among
\emph{four} points on a spectrum, and the deciding factor is almost always the
same question: \emph{how long is the wait expected to be?} Match the strategy to
the timescale and the rest follows.

\subsection{A tight spin -- shorter than a context switch}

When you are waiting on a hardware register bit that clears in a microsecond or
two -- a command latching, a DMA-fed shift register draining -- the right thing
is to spin with \emph{no} back-off at all:

\begin{lstlisting}
B.Regs.CMD.USR := True;            -- kick the SPI transfer
while B.Regs.CMD.USR loop null;    -- ... and wait the few us for it to finish
end loop;
\end{lstlisting}

This looks like the cardinal sin of embedded code -- a busy-loop burning the CPU
-- but here it is correct, and a \texttt{delay until} would be \emph{worse}. The
wait is a handful of microseconds; the kernel tick is a millisecond, and even a
bare context switch costs more cycles than the bit takes to clear. Sleeping would
turn a 2\,\textmu s wait into a 1\,ms one and add scheduler churn for nothing. The
SPI engine (\texttt{CMD.USR}, \texttt{CMD.UPDATE}), the I2S engine
(\texttt{TX\_UPDATE}), and the SHA peripheral (\texttt{BUSY.STATE}) all spin like
this, deliberately. The rule: \emph{if the wait is shorter than the cost of
yielding, do not yield.}

\subsection{A bounded spin -- a tight loop that cannot hang}

The pure \texttt{loop null} spin has one flaw: if the bit \emph{never} clears -- a
wedged peripheral, a chip that is not there -- it hangs the core forever. Where
that risk is real, the same spin gets a guard counter, so a fault becomes a
bounded failure instead of a deadlock:

\begin{lstlisting}
Guard : Natural := 0;
while not Done (C, Dir) and then Guard < 2_000_000 loop   -- GDMA completion
   Guard := Guard + 1;
end loop;
--  fall through whether the DMA finished or the cap tripped; the caller checks Done
\end{lstlisting}

The cap is a crude worst-case-execution-time bound: still no \texttt{delay} (the
expected wait is short), but liveness is guaranteed. GDMA's \texttt{Wait} and the
W5500's command-accept (\texttt{Issue}, capped at 2000 tries) use this shape.

\subsection{A paced poll -- when the wait is milliseconds}

Once the expected wait reaches the millisecond range -- a reset settling, a slow
status bit, a byte trickling in over a UART -- yielding pays off, and you back off
with \texttt{delay until} so other tasks run meanwhile. This is the polling loop
from the previous section, and \texttt{ESP32S3.Wait.Until\_True} packages it.
Datasheet-mandated settle delays (the ST7789's 120\,ms power-on, the W5500's reset
hold) are the same idea with a \emph{fixed} delay rather than a re-checked
condition.

\subsection{Blocking on an interrupt -- long or unbounded}

When the wait could be long, is unpredictable, or power matters -- a network
packet that may arrive in 1\,ms or 10\,s, a key-press, a sensor's data-ready line
-- polling wastes both CPU and energy. Sleep instead: an interrupt handler
\texttt{Set\_True}s a \texttt{Suspension\_Object} and the task
\texttt{Suspend\_Until\_True}s on it, burning nothing until the hardware wakes it
(Chapter~\ref{ch:suspension}). This is the top of the spectrum: the most machinery
(a handler, a wired interrupt), but the only choice that costs \emph{zero} while
waiting.

\subsection{One driver, the whole spectrum}

The W5500 networking stack is a useful specimen because it spans nearly all of
these at once. Accepting a chip command spins, bounded (\texttt{Issue}); a short
register read spins; and the blocking waits -- for an incoming connection, for
received data, for room in the transmit buffer -- funnel through a single hook:

\begin{lstlisting}
procedure Wait_Event (S : Socket) is
begin
   if Waiter /= null then Waiter (S.Index);     -- sleep on INTn (Suspend_Until_True)
   else delay until Clock + Milliseconds (1);   -- else a 1 ms paced poll
   end if;
end Wait_Event;
\end{lstlisting}

By default the stack \emph{polls} (the 1\,ms fallback); calling
\texttt{Interrupts.Enable} registers a waiter and the very same waits become
interrupt-blocking, with no change to the socket code. That default carries a
throughput cost worth knowing: a bulk transfer paced at one decision per
millisecond is far slower than one that wakes the instant the buffer drains -- the
difference the 1\,MiB FTP upload of Chapter~\ref{ch:networking} ran into. For
sustained data, enable the interrupts.

\subsection{The rule of thumb}

\begin{center}
\begin{tabular}{l l l}
\toprule
\textbf{Expected wait} & \textbf{Strategy} & \textbf{Example} \\ \midrule
under a few \textmu s          & tight spin, no back-off        & SPI \texttt{CMD.USR}, SHA \texttt{BUSY} \\
short, but may never end       & bounded spin (guard count)     & GDMA \texttt{Wait}, W5500 \texttt{Issue} \\
milliseconds                   & paced poll (\texttt{delay until}) & reset/settle, \texttt{Wait.Until\_True} \\
long / unbounded / power       & block on an interrupt          & W5500 \texttt{Wait\_Data}, sensor DRDY \\
\bottomrule
\end{tabular}
\end{center}

Two failure modes bracket the table. Yield on a sub-microsecond wait and you pay a
thousandfold latency penalty for nothing; spin -- even bounded -- on a
multi-second wait and you burn a core and a battery. Picking the row that matches
the timescale is most of getting it right.

\section{Concurrency patterns}
\label{sec:concurrency}\index{concurrency patterns}

A few protected-object shapes recur often enough to be worth naming. All work on
every tasking profile; the Jorvik restrictions (one entry per protected object,
one caller queued at a time) shape them but do not prevent them.

\textbf{Producer/consumer over a bounded buffer.} The classic decoupling: one
task (or an interrupt) produces items, another consumes them, with a protected
ring buffer between. The consumer's \texttt{Get} is an \emph{entry} whose barrier
holds it until something is available -- no polling.

\begin{lstlisting}
protected Queue is
   procedure Put (X : Item);          -- producer side (also callable from an ISR)
   entry     Get (X : out Item);      -- consumer side: blocks while empty
private
   Buf        : Item_Array (0 .. 15);
   Count      : Natural := 0;
   Head, Tail : Natural := 0;
end Queue;

protected body Queue is
   procedure Put (X : Item) is
   begin
      if Count < Buf'Length then      -- (full -> drop; or count an overflow)
         Buf (Tail) := X;
         Tail := (Tail + 1) mod Buf'Length;
         Count := Count + 1;
      end if;
   end Put;

   entry Get (X : out Item) when Count > 0 is
   begin
      X := Buf (Head);
      Head := (Head + 1) mod Buf'Length;
      Count := Count - 1;
   end Get;
end Queue;
\end{lstlisting}

The consumer task just loops \texttt{Queue.Get (X)}; it sleeps on the barrier when
the queue is empty and wakes the instant \texttt{Put} adds an item. (Under Jorvik,
exactly one consumer task may queue on \texttt{Get} at a time -- which is the
usual one-producer/one-consumer arrangement.)

\textbf{A signal / mailbox.} The degenerate one-slot case is a handoff: an
interrupt deposits a value and a task collects it. It is the
``protected entry'' half of the waiting pattern (\S\ref{sec:waiting}) carrying a
payload -- a \texttt{procedure Post (V)} from the ISR and an
\texttt{entry Take (V : out \ldots)} barriered on ``something posted.''

\textbf{DMA double-buffering (ping-pong).} For a continuous stream -- audio in
over I2S, samples to a DAC -- you cannot stop to process between transfers. Use
\emph{two} buffers: while the GDMA fills (or drains) buffer~A, the CPU processes
buffer~B; when the transfer completes you swap, so the peripheral is never idle.
The completion interrupt signals the processing task (a suspension object or a
queue of ``buffer ready'' indices) and points the next transfer at the other
buffer. The protected object owns the ``which buffer is the hardware's'' state;
the task owns the other. This keeps a real-time stream gap-free with bounded,
analysable hand-off.

\begin{figure}[ht]
\centering
\begin{tikzpicture}[node distance=8mm and 30mm]
\node[blk, fill=green!9, text width=3.0cm]  (a1) {Buffer A -- DMA fills};
\node[blk, below=of a1, text width=3.0cm]   (b1) {Buffer B -- CPU works};
\node[blk, right=of a1, fill=green!9, text width=3.0cm] (b2) {Buffer B -- DMA fills};
\node[blk, below=of b2, text width=3.0cm]   (a2) {Buffer A -- CPU works};
\draw[flow] (a1.east) -- (a2.west);     % A: DMA -> CPU
\draw[flow] (b1.east) -- (b2.west);     % B: CPU -> DMA
\node[font=\scriptsize, fill=white, inner sep=1pt] at ($(b1)!0.5!(b2)$) {swap on DMA-done IRQ};
\node[font=\scriptsize] at ($(a1.north)!0.5!(b1.south)-(2.4cm,0)$) {phase 1};
\node[font=\scriptsize] at ($(b2.north)!0.5!(a2.south)+(2.4cm,0)$) {phase 2};
\end{tikzpicture}
\caption{DMA double-buffering: the peripheral always has one buffer while the CPU
works the other; the completion interrupt swaps their roles (the crossing), so the
stream never stalls.}
\label{fig:pingpong}
\end{figure}

\textbf{Why protected objects, not flags.} Each of these could be hacked with a
shared \texttt{Volatile} flag, but the protected object gives you the barrier (no
polling), mutual exclusion (no torn reads across cores), and a priority ceiling
(bounded blocking) for free -- the points the Encapsulation chapter
(Chapter~\ref{ch:encap}) makes about \texttt{protected}.

\section{Watchdogs and recovery}
\label{sec:watchdogs}\index{watchdog}

A watchdog is a hardware timer that resets the chip unless the software keeps
``feeding'' it -- the last line of defence against a hang that no exception or
check can catch (a wild loop, a deadlock, a stuck peripheral). The ESP32-S3 has
several: the two timer-group main watchdogs (TG0/TG1 MWDT), the RTC watchdog
(RWDT), and the RTC \emph{super}-watchdog (a hardware backstop).

\textbf{The bare boot leaves them off.} The ROM and the second-stage bootloader
arm watchdogs during boot, but this runtime's startup \emph{disables} them (so a
long elaboration or a debugger pause is not cut short). So by default a bare-metal
app runs with \emph{no} active watchdog -- a hang will sit forever rather than
recover. That is fine while developing (and is why a wedged example does not
reset on its own), but a deployed device usually wants protection.

\textbf{Arming the RTC watchdog.} To get auto-recovery, enable the RWDT with a
timeout and feed it from a place that only runs while the system is healthy. The
registers are key-locked; the sequence is unlock, configure, feed, re-lock:

\begin{lstlisting}
--  RTC watchdog registers (overlaid as plain words).
Key     : constant := 16#50D8_3AA1#;
Config0 : Unsigned_32 with Volatile, Address => To_Address (16#6000_8098#), Import;
Config1 : Unsigned_32 with Volatile, Address => To_Address (16#6000_809C#), Import;
Feed    : Unsigned_32 with Volatile, Address => To_Address (16#6000_80AC#), Import;
Wprot   : Unsigned_32 with Volatile, Address => To_Address (16#6000_80B0#), Import;

procedure Arm (Timeout_Cycles : Unsigned_32) is   -- RTC-slow-clk (~63 kHz) cycles
begin
   Wprot   := Key;                       -- unlock
   Config1 := Timeout_Cycles;
   Config0 := 2**31                      -- WDT_EN
            or (3 * 2**28)               -- stage 0 action = reset the main system
            or (7 * 2**16) or (7 * 2**13);  -- reset-pulse lengths
   Feed    := 2**31;
   Wprot   := 0;                         -- re-lock
end Arm;

procedure Pet is begin Wprot := Key; Feed := 2**31; Wprot := 0; end Pet;
\end{lstlisting}

\textbf{Feed on progress, not on a timer.} The robust idiom is to pet the dog from
points that only execute when work is actually advancing -- the main loop body, a
completed transaction -- not from a periodic task that would keep petting even if
the real work has stalled. The ACATS test harness in this project does exactly
this: it arms the RWDT (\,$\approx$19\,s) and feeds it on every progress line, so a
test that deadlocks stops feeding and is reset, while a slow-but-progressing test
survives.

\textbf{The super-watchdog} is a separate hardware backstop that a deep-sleep wake
can leave armed; \texttt{ESP32S3.RTC.Disable\_Super\_Watchdog} turns it off when an
app means to stay awake (the RTC driver also re-feeds it automatically after a
deep-sleep wake). Between an armed RWDT for hang-recovery, the precise
stack-overflow watchpoint (the Debugging chapter), and a last-chance handler for
the unrecoverable case, a deployed app has a layered safety net.

\section{Dual-core (SMP) programming}
\label{sec:smp}\index{SMP}

The ESP32-S3 has two cores and the runtime schedules across both. Two ordinary-Ada
facilities put them to use.

\textbf{Pin a task to a core} with the \texttt{CPU} aspect (\texttt{CPU => 1} is
the first core, \texttt{2} the second); omit it and the task may run on either:

\begin{lstlisting}
   task Beat0 with Priority => Priority'Last - 1, CPU => 1;
   task Beat1 with Priority => Priority'Last - 1, CPU => 2;   -- the other core
\end{lstlisting}

\textbf{Share data with a protected object} -- it works \emph{across} cores, not
just between tasks on one core. A protected \emph{entry} can even hand a caller
from one core to the other: in the \texttt{esp32s3\_smp} example a producer on
core 1 blocks in \texttt{entry Get when Full}, and serving the entry wakes it
through the inter-core interrupt. The locking is the same priority-ceiling monitor
(Chapter~\ref{ch:encap}); the runtime adds the cross-core wake-up underneath.

\begin{lstlisting}
   protected Mailbox is
      procedure Put (V : Integer);
      entry     Get (V : out Integer);   -- blocks until Full; may serve cross-core
   private
      Item : Integer;
      Full : Boolean := False;
   end Mailbox;
\end{lstlisting}

Two cautions specific to this target: time comes from a \emph{shared} hardware
timer (the SYSTIMER), not each core's own cycle counter, so \texttt{Clock} agrees
across cores; and the USB-serial-JTAG console is not robust for two tasks printing
at once -- funnel diagnostics through one task. With those in mind, an SMP
application is just ordinary tasks, pinned where you want them, sharing through
protected objects.


%  Short title: the full one overflows the running header.
\chapter[Mass Storage: SD Cards, SPI Flash, ext4]%
        {Mass Storage: SD Cards, SPI Flash, and the ext4 Filesystem}
\label{ch:storage}

The final layer is mass storage: block drivers for SD cards and for a raw SPI NOR
flash chip, and a complete pure-Ada ext2/3/4 filesystem (a reimplementation of
lwext4) that sits on top of them: mounting real \texttt{mkfs.ext4} cards,
\emph{and} wear-leveling and formatting a bare flash chip entirely on-device. The
smallest non-volatile stores --- serial EEPROM and its faster, wear-free cousin
FRAM, each driven by a generic over a whole catalogue of parts --- get their own
chapters next (Chapters~\ref{ch:eeprom} and~\ref{ch:fram}).

\section{SD\_SPI: an SD card over the SPI master}\index{SD card}\index{SD\_SPI}\index{SPI SD mode}

\texttt{ESP32S3.SD\_SPI} talks to an SD/SDHC card in SPI mode, layered on the
task-safe SPI driver. The chip-select is an ordinary GPIO so it can be held across
a whole command. Blocks are 512-byte logical sectors (LBA addressing); SDSC and
SDHC are detected and handled internally.

\begin{lstlisting}
type Block is array (0 .. 511) of Interfaces.Unsigned_8;
type Status is (OK, No_Card, Unusable, Init_Timeout, Read_Error, Write_Error);
procedure Setup (C : out Card; Host : ESP32S3.SPI.SPI_Host;
                 Sclk, Mosi, Miso, Cs : Pin_Id;
                 Init_Clock_Hz : Positive := 400_000;
                 Data_Clock_Hz : Positive := 8_000_000);
procedure Initialize (C : in out Card; Result : out Status);
procedure Read_Block  (C : in out Card; LBA : Block_Address;
                       Data : out Block; Result : out Status);
procedure Write_Block (C : in out Card; LBA : Block_Address;
                       Data : Block; Result : out Status);
\end{lstlisting}

\lstexample{Bring up a card and read sector 0.}
\begin{lstlisting}
with ESP32S3.SD_SPI; use ESP32S3.SD_SPI;
with ESP32S3.SPI;
procedure SD_Init is
   C   : Card;
   St  : Status;
   Sec : Block;
begin
   Setup (C, ESP32S3.SPI.SPI2, Sclk => 12, Mosi => 11, Miso => 13, Cs => 10);
   Initialize (C, St);
   if St = OK then
      Read_Block (C, LBA => 0, Data => Sec, Result => St);   --  the MBR sector
   end if;
end SD_Init;
\end{lstlisting}

\lstexample{A non-destructive read / write-back / re-read round-trip.}
\begin{lstlisting}
with ESP32S3.SD_SPI; use ESP32S3.SD_SPI;
with ESP32S3.SPI;
procedure SD_Roundtrip is
   C          : Card;
   St         : Status;
   Orig, Back : Block;
begin
   Setup (C, ESP32S3.SPI.SPI2, Sclk => 12, Mosi => 11, Miso => 13, Cs => 10);
   Initialize (C, St);
   Read_Block  (C, 16#2000#, Orig, St);     --  read a scratch sector
   Write_Block (C, 16#2000#, Orig, St);     --  write the same bytes back
   Read_Block  (C, 16#2000#, Back, St);     --  re-read; Back = Orig, no data lost
end SD_Roundtrip;
\end{lstlisting}

\section{SDMMC: the native SD/MMC host}\index{SDMMC}\index{SDHOST controller}

\texttt{ESP32S3.SDMMC} drives the dedicated SDHOST controller (clock + command +
1- or 4-bit data) in PIO mode, which is faster than SPI mode, and with no DMA or
finalization it works under every profile. The API mirrors \texttt{SD\_SPI}.

\begin{lstlisting}
procedure Setup (C : out Card; On : Slot; Clk, Cmd, D0 : Pin_Id;
                 D1, D2, D3 : Optional_Pin := No_Pin;
                 Width : Bus_Width := Width_1; ...);
procedure Initialize (C : in out Card; Result : out Status);
procedure Read_Block  (C : in out Card; LBA : Block_Address; Data : out Block;
                       Result : out Status);
\end{lstlisting}

\lstexample{4-bit bus, read a sector.}
\begin{lstlisting}
with ESP32S3.SDMMC; use ESP32S3.SDMMC;
procedure SDMMC_Read is
   C   : Card;
   St  : Status;
   Sec : Block;
begin
   Setup (C, On => Slot1, Clk => 14, Cmd => 15, D0 => 2, D1 => 4, D2 => 12,
          D3 => 13, Width => Width_4);
   Initialize (C, St);
   if St = OK then
      Read_Block (C, 0, Sec, St);
   end if;
end SDMMC_Read;
\end{lstlisting}

\section{The ext4 filesystem}\index{ext4 filesystem}\index{lwext4}\index{JBD2 journal}

\texttt{ESP32S3.Ext4} is a pure-Ada ext2/3/4 implementation with a JBD2 journal
(Chapter overview in \texttt{libs/esp32s3\_hal/src/ext4/README.md}). It reads
default \texttt{mkfs.ext4} cards (extents, 64-bit, HTree directories,
\texttt{metadata\_csum} verification), writes ext2/3/non-csum filesystems, and
recovers a dirty journal on mount. It needs the \texttt{embedded} or \texttt{full}
profile.

\subsection{Connecting a block device}\index{block device}

The filesystem talks to a small block-device abstraction (\texttt{ESP32S3.%
Block\_Dev.Device}). An adapter turns an initialised SD card into one with a
single call. The \texttt{Make} below appears in each complete example that
follows:

\begin{lstlisting}
with ESP32S3.Block_Dev;
with ESP32S3.Block_Dev.SD_SPI_Source;
--  ... given an initialised "Card : aliased ESP32S3.SD_SPI.Card" ...
   Dev : constant ESP32S3.Block_Dev.Device :=
           ESP32S3.Block_Dev.SD_SPI_Source.Make (Card'Access);
\end{lstlisting}

\subsection{Mounting and reading}\index{mount}\index{inode}\index{journal replay}

A \texttt{Mount} is a controlled handle: it flushes and closes on scope exit, and
\texttt{Open} automatically replays a dirty journal on a writable mount. The
complete example brings up the card, wraps it, mounts read-only, then reads a
file:

\begin{lstlisting}
with ESP32S3.SPI;
with ESP32S3.SD_SPI;
with ESP32S3.Block_Dev;
with ESP32S3.Block_Dev.SD_SPI_Source;
with ESP32S3.Ext4;
with ESP32S3.Ext4.FS;
with ESP32S3.Ext4.Inode;
procedure EXT4_Read is
   use type ESP32S3.SD_SPI.Status;
   Card : aliased ESP32S3.SD_SPI.Card;
   St   : ESP32S3.SD_SPI.Status;
begin
   ESP32S3.SD_SPI.Setup (Card, ESP32S3.SPI.SPI2,
                         Sclk => 12, Mosi => 11, Miso => 13, Cs => 10);
   ESP32S3.SD_SPI.Initialize (Card, St);
   if St /= ESP32S3.SD_SPI.OK then
      return;
   end if;
   declare
      Dev  : constant ESP32S3.Block_Dev.Device :=
               ESP32S3.Block_Dev.SD_SPI_Source.Make (Card'Access);
      M    : ESP32S3.Ext4.FS.Mount;
      I    : ESP32S3.Ext4.Inode.Info;
      Buf  : ESP32S3.Ext4.Byte_Array (0 .. 255);
      Last : Natural;
   begin
      M.Open (Dev, Read_Only => True);
      M.Stat (M.Lookup ("/etc/hello.txt"), I);  --  resolve a path -> its inode
      M.Read_File (I, Offset => 0, Into => Buf, Last => Last);
      --  Buf (0 .. Last - 1) holds the first Last bytes of the file
   end;   --  M finalizes: flush + close
end EXT4_Read;
\end{lstlisting}

\textbf{Listing a directory.} \texttt{Iterate} calls a visitor for each entry.
The card bring-up is identical to the example above; the body is:

\begin{lstlisting}
   declare
      Dev  : constant ESP32S3.Block_Dev.Device :=
               ESP32S3.Block_Dev.SD_SPI_Source.Make (Card'Access);
      M    : ESP32S3.Ext4.FS.Mount;
      Root : ESP32S3.Ext4.Inode.Info;
      procedure Show (Name : String;
                      Ino  : ESP32S3.Ext4.Inode_Number;
                      File_Type : ESP32S3.Ext4.U8) is
      begin
         Put_Line ("  " & Name);
      end Show;
   begin
      M.Open (Dev, Read_Only => True);
      M.Stat (M.Lookup ("/"), Root);
      M.Iterate (Root, Show'Access);
   end;
\end{lstlisting}

\subsection{Writing}

On a non-\texttt{metadata\_csum} filesystem (ext2/3, or
\texttt{mke2fs -O \textasciicircum metadata\_csum}) the filesystem can be modified.
The full POSIX-like set is available: \texttt{Create\_File}, \texttt{Write\_File},
\texttt{Truncate}, \texttt{Mkdir}, \texttt{Rmdir}, \texttt{Unlink},
\texttt{Rename} and hard \texttt{Link}.

\begin{lstlisting}
   declare
      Dev : constant ESP32S3.Block_Dev.Device :=
              ESP32S3.Block_Dev.SD_SPI_Source.Make (Card'Access);
      M   : ESP32S3.Ext4.FS.Mount;
      N   : ESP32S3.Ext4.Inode_Number;
   begin
      M.Open (Dev, Read_Only => False);
      N := M.Create_File ("/", "log.txt");
      M.Write_File (N, (Character'Pos ('h'), Character'Pos ('i'), 10));
      M.Mkdir ("/", "data");
      M.Close;          --  flush + close (or let M finalize)
   end;
\end{lstlisting}

\subsection{Streaming writes and file-size limits}\index{streaming write}\index{Append}\index{block map}\index{single-indirect block}\index{double-indirect block}

\texttt{Write\_File} sets a file's whole contents from one in-memory buffer.
That is fine for a small file, but you cannot hand an embedded part a gigabyte-sized
array. \texttt{Append} grows a file at its end one chunk at a time, so a large
file is built from small writes with no buffer to hold it all:

\begin{lstlisting}
   N := M.Create_File ("/", "big.log");
   for I in 1 .. 1000 loop          --  e.g. 1000 records, streamed
      M.Append (N, Chunk);          --  grows the file; no whole-file buffer
   end loop;
   M.Commit;
\end{lstlisting}

What bounds a single file is the \emph{block map}. The writer allocates classic
indirect blocks (12 direct, then single-indirect, then double-indirect), so
on a 4\,KiB-block volume one file reaches about 4\,MiB through the single-indirect
tier and about 4\,GiB through the double-indirect tier; the matching free and
truncate paths walk the same tiers, so a file that can be written can always be
deleted and shrunk. The \emph{read} path goes further: it follows the full
classic map (through triple-indirect) and ext4 extent trees, so a
multi-gigabyte file written by a host is readable on-device, bounded only by the
volume's free space.

\subsection{Directory capacity}\index{directory capacity}\index{dirent}\index{directory growth}

A directory is itself a file of variable-length entries (each 8 bytes plus the
name, rounded to 4), and creating a name must find room for one. When every
existing directory block is packed, \texttt{Add\_Entry} grows the directory:
it allocates a block, writes a single spanning unused entry (the shape
\texttt{mkfs} gives a fresh directory block), wires it into the inode's next
\emph{direct} map slot, and persists the grown inode before placing the name.
Twelve direct slots at 4\,KiB give 48\,KiB of entries -- on the order of
1{,}500--2{,}900 names depending on their length. Growth past the direct slots
raises \texttt{No\_Space}, and a directory a host created with an extent map
raises \texttt{Unsupported\_Feature} rather than being modified half-correctly.

That ceiling is not hypothetical bookkeeping. Before growth existed, a
directory was capped at its \emph{first} block -- roughly 140 mirror-length
names -- and one application built on this SDK hit it silently: it mirrors a
weather service's chart set at three files per product, filled the block
partway through the first full download, and every later create failed with
\texttt{No\_Space} from that day on. The products that happened to sort late in
its list simply never appeared on the card, which surfaced to the user as one
category of the UI ``only ever shows one image.'' A capacity limit in a
filesystem rarely announces itself where the limit lives; budget directory
sizes (files per directory, not just bytes per volume) when planning on-card
layouts, and prefer one subdirectory per product family once counts head into
the hundreds.

\subsection{Crash-safe transactions}\index{crash-safe transactions}\index{journaling}\index{commit barrier}

For atomicity, group an operation behind \texttt{Commit}: the dirty metadata and
the superblock are written to the journal (durably), the RECOVER flag is set
(the commit barrier), then the metadata is checkpointed to its final locations
and the flag cleared. A power loss after the barrier is recovered forward by the
next \texttt{Open} (or by any other ext4 implementation, including the Linux
kernel); a loss before it has no effect.

\begin{lstlisting}
   --  Payload : constant ESP32S3.Ext4.Byte_Array := (16#41#, 16#42#, 16#43#);
   N := M.Create_File ("/", "important.txt");
   M.Write_File (N, Payload);
   M.Commit;         --  one atomic, journaled, crash-safe transaction
\end{lstlisting}

A volume formatted \emph{without} a journal (\texttt{mke2fs -O \textasciicircum
has\_journal}) has no log to write through, so \texttt{Commit} instead flushes the
dirty metadata and superblock straight to their final locations. The filesystem
detects which kind it mounted and routes \texttt{Commit} automatically. A
no-journal volume trades the crash-safety barrier for smaller size and fewer
writes, the right choice on a small flash chip.\index{no-journal volume}

\subsection{Concurrency: one owner per mount}\index{block cache}\index{single-owner mount}

A \texttt{Mount} is a \emph{single-owner} handle, not a shared service. Every
operation takes it \texttt{in out} and mutates shared state, and the filesystem
holds no internal lock or protected object: drive one mount from one task, and let
its calls run to completion in sequence.

That rule is stricter than it first looks because even reads have side effects.
Underneath the mount sits a write-back LRU block cache, and \texttt{Read\_File}
fetches-on-miss and reorders it, so two tasks ``only reading'' the same mount
still race on the cache. Writes are stronger still: \texttt{Create\_File},
\texttt{Write\_File}, \texttt{Mkdir} and the rest accumulate a dirty-metadata
write-set that \texttt{Commit} flushes as \emph{one} journaled transaction (the
crash-safe commit above). Interleaving two tasks' writes scrambles that write-set
and breaks the crash-safety guarantee, which assumes a single writer building one
coherent transaction. The block drivers \emph{beneath} the filesystem
(\texttt{SD\_SPI}, \texttt{SDMMC}) are themselves task-safe, but that only
serialises the \emph{bus}; it does nothing for the filesystem state above it.

If more than one task must reach the filesystem, give the mount a single owner and
funnel every request through it: a dedicated task that owns the \texttt{Mount} and
performs the I/O, with other tasks handing it work through a \texttt{protected}
request queue (Ravenscar-compatible) or, on the \texttt{full} profile, a rendezvous
entry. The I/O itself must stay in the owning task, never inside a protected
action because the SD drivers can suspend and a protected operation must not block.

\begin{lstlisting}
--  One task owns the Mount and is the only code that touches it; other tasks
--  hand it work (a protected request queue under Ravenscar, or a rendezvous
--  entry on the full profile).  The blocking I/O stays in this task.
task body File_Server is
   M : ESP32S3.Ext4.FS.Mount;          --  the single owner
begin
   M.Open (Dev, Read_Only => False);
   loop
      --  take the next request, then service it on the one mount:
      --     M.Read_File (...)              --  a read request
      --     M.Write_File (...);  M.Commit; --  a write request, committed
      null;
   end loop;
end File_Server;
\end{lstlisting}

\subsection{Errors}

Filesystem operations raise: the standard \texttt{Ada.IO\_Exceptions} names
(\texttt{Name\_Error} for a missing path, \texttt{Use\_Error}, \texttt{Device\_%
Error}, \texttt{End\_Error}) plus filesystem-specific \texttt{Unsupported\_%
Feature}, \texttt{Bad\_Checksum}, \texttt{Corrupt}, \texttt{No\_Space},
\texttt{Not\_Empty} and \texttt{Read\_Only}. Wrap calls in a handler so a bad card
or a missing file never escapes a task:

\begin{lstlisting}
begin
   M.Stat (M.Lookup ("/maybe/missing"), I);
exception
   when ESP32S3.Ext4.Name_Error => Put_Line ("not found");
   when ESP32S3.Ext4.Bad_Checksum => Put_Line ("corrupt metadata");
end;
\end{lstlisting}

\subsection{Free-count integrity: the phantom-free tripwire}\index{phantom-free tripwire}\index{free-count drift}

A subtle class of corruption is the \emph{free-count drift}: the superblock and
block-group descriptors cache how many blocks and inodes are free, and if a
\texttt{Free} adjusts that cache for a bit that was \emph{already} clear (a
double-free, or a stale/incoherent bitmap read off a flaky SD card), the cached
count creeps above what the bitmap actually says. \texttt{e2fsck} then reports a
Pass-5 count mismatch on an otherwise-sound filesystem.

The fix makes \texttt{Free} \emph{idempotent} by construction. \texttt{Clear\_Bit}
now reports whether the bit was actually set, and the group/superblock counts are
bumped \emph{only} when it was, so a double or phantom free is a no-op rather
than a drift, and the count cannot diverge from the bitmap.

\begin{lstlisting}
--  Bitmap.Clear_Bit returns False if the bit was already clear:
if Bitmap.Clear_Bit (Group_BM, Index) then
   Dec_Used (Group);          --  bit really was set: adjust the counts
else
   Phantom_Frees := Phantom_Frees + 1;   --  already free: no drift, but flag it
end if;
\end{lstlisting}

Crucially the already-clear path is \emph{not} silent. That would mask a genuine
double-free bug. A \texttt{Phantom\_Free\_Count} tripwire counts the events (a plain
counter, not a \texttt{pragma Assert} because the target also builds with
\texttt{-gnata} and an assert would abort the filesystem operation). The host test
harness asserts the count is zero after every scenario, and the on-target
\texttt{esp32s3\_ext4\_write} example reports it as a battery check
(\texttt{[PASS] no phantom frees}; a flaky card flips it to \texttt{[FAIL]}).

\begin{lstlisting}
Bitmap.Reset_Phantom_Free_Count;            --  arm the tripwire
--  ... exercise allocate / free / truncate / unlink ...
Check ("no phantom frees (count consistent with bitmap)",
       Bitmap.Phantom_Free_Count = 0);
\end{lstlisting}

The original drift was ultimately traced to a \emph{worn} SD card: a brand-new card
runs the whole battery \texttt{e2fsck}-clean with zero phantom frees. The tripwire
is what turned an intermittent, card-dependent count mismatch into a single
yes/no signal, surfacing the fault instead of hiding it.

\subsection{How it is validated}

Because the filesystem is pure logic over the block interface, it is developed and
tested host-native (x86) against real \texttt{mke2fs} images, with the Linux
kernel and \texttt{e2fsck} as the oracle: every read is byte-compared to the
source file, every mutation is checked with \texttt{e2fsck -fn}, and journal
replay and commit are cross-checked against the kernel's own recovery. The same
code cross-compiles for the ESP32-S3 and runs over the SD drivers above. See
\texttt{tests/ext4\_host/} and \texttt{libs/esp32s3\_hal/src/ext4/REFERENCES.md}.

\subsection{Scope: deliberately smaller than Linux ext4}\index{ext4 scope}\index{HTree}\index{delayed allocation}\index{metadata\_csum}

It is worth being explicit that this filesystem does not aim for feature parity
with the Linux kernel's ext4, and why. The asymmetry is intentional: the
\emph{read} path is wide (extents, 64-bit, HTree-indexed directories,
\texttt{metadata\_csum} verification -- a default desktop card mounts and reads),
while the \emph{write} path is a narrow, fully-understood subset (classic
indirect maps, linear directories, non-csum volumes). Anything outside that
subset is \emph{refused} -- \texttt{Read\_Only}, \texttt{Unsupported\_Feature} --
rather than attempted, because the most dangerous state for an on-disk format
is ``implemented enough to write, not enough to write correctly'': that is
silent corruption a desktop mount discovers weeks later. Refusal fails loudly
at the call site.

The features left out earn nothing at this scale. HTree directory indexing
beats a linear scan somewhere in the thousands of entries; a directory here
tops out around two thousand entries in 48\,KiB that is hot in the block cache.
Delayed allocation and the multi-block allocator exist to optimise allocation
locality under multi-process write pressure that a single-owner embedded mount
does not have. Full \texttt{jbd2} parity would double metadata writes on a bus
measured in single-digit MB/s (and add SD wear) to strengthen guarantees the
simple journal-plus-\texttt{Commit} model already provides. And every one of
them wants cache and working state sized for a kernel page cache, not for a
32\,KiB block cache sharing PSRAM with an application's framebuffers.

What keeps the small subset honest is that it stays small enough to
\emph{verify}: the host harness runs whole scenario batteries with
\texttt{e2fsck} as the referee, as described above, and individual pieces are
within reach of SPARK proof. A port of Linux ext4's roughly sixty thousand lines
would be the largest, least-reviewable component in the SDK, with a test
matrix that grows combinatorially in feature interactions -- correctness that
upstream earned through two decades of fleet exposure no embedded project can
replicate. If a concrete interoperability need arises -- the strongest
candidate is \emph{writing} \texttt{metadata\_csum} volumes, so cards formatted
by a modern desktop \texttt{mkfs.ext4} need not be reformatted -- the right
move is to implement that one feature (checksum recompute on the metadata the
writer already touches), not to chase parity.

\section{Putting ext4 on raw SPI NOR flash}\index{SPI NOR flash}\index{wear leveling}\index{flash translation layer}

The same filesystem runs on a bare SPI NOR flash chip, no SD card in sight. Four
pure-Ada layers stack up, each a thin adapter over the one below:

\begin{center}
\texttt{ESP32S3.Ext4} $\rightarrow$ \texttt{Block\_Dev.WL} $\rightarrow$
\texttt{Block\_Dev.W25Q\_Source} $\rightarrow$ \texttt{ESP32S3.W25Q}
\end{center}

\noindent The flash driver speaks to the chip; \texttt{W25Q\_Source} presents it
as a 512-byte block device; \texttt{WL} is a wear-leveling flash translation
layer that spreads writes so no sector wears out first; and \texttt{Ext4} mounts
the result. The reference part is a Winbond \texttt{W25Q256FV} (32\,MB) that shares
the SPI2 bus with the W5500 Ethernet chip.

\subsection{The W25Q flash driver}\index{W25Q}\index{JEDEC ID}\index{4-byte addressing}\index{chip-select}

\texttt{ESP32S3.W25Q} reads the JEDEC id, then programs, erases and reads in
\emph{4-byte address mode} (a part larger than 16\,MB needs four address bytes).
Two design points stand out. The chip size is \emph{detected}, not hard-coded: the
JEDEC density byte is a power of two, so \texttt{Capacity\_Bytes} turns
\texttt{EF\,40\,19} into 32\,MB, and because each layer above sizes itself from the
one below, fitting a different W25Q part (8/16/32/64\,MB) needs no code change.
And the device record \emph{describes itself}: its host, its own bit clock, and its
chip select. The select is its own GPIO, not the SPI host's single hardware
\texttt{CS0}: the flash names its select pin and the SPI driver drives it
active-low, held across the whole command (so a streamed read can keep \texttt{CS}
asserted across many transfers). This is what lets the flash share the bus
with another device on \texttt{CS0}, each clocked its own way. Setting
\texttt{Clock\_Hz} and \texttt{CS\_Pin} is all it takes; a select that is not one
plain GPIO (a 3:8 decoder, an I/O-expander line) instead supplies a \texttt{CS\_CB}
callback. The host itself is brought up once with \texttt{ESP32S3.SPI.Setup} +
\texttt{Configure\_Pins} (the shared \texttt{Sclk}/\texttt{Mosi}/\texttt{Miso}).

\begin{lstlisting}
with ESP32S3.W25Q; use ESP32S3.W25Q;
--  Flash : ESP32S3.W25Q.Flash := (Host => ESP32S3.SPI.SPI2,
--    Clock_Hz => 8_000_000, CS_Pin => 21, others => <>);  --  8 MHz, CS on IO21
   ID      : JEDEC_ID;
   Mode_OK : Boolean;
begin
   Read_Identification (Flash, ID);          --  EF 40 19 on a W25Q256FV
   Initialize (Flash, Mode_OK);              --  enter 4-byte address mode
   if ID.Manufacturer = 16#EF# and then Capacity_Bytes (ID) /= 0 then
      null;                                  --  Capacity_Bytes (ID) = 32 MB
   end if;
\end{lstlisting}

\subsection{Block device and wear leveling}\index{W25Q\_Source}\index{block-erase}\index{ping-pong configuration}

NOR flash must be \emph{erased} (to all-ones, in 4\,KiB sectors) before it can be
re-programmed, so \texttt{W25Q\_Source} hides an erase-aware, write-through
512-byte sector interface over the chip; an optional block-erase capability lets
the layer above clear a whole 4\,KiB block in one operation.

\texttt{Block\_Dev.WL} is the wear-leveling FTL. It keeps an $O(1)$ state (a
single move counter) and maps a logical 4\,KiB block to a physical one by a
shifting offset, with one always-free ``hole''. Every few writes it advances the
hole by copying one block, so over time every logical block visits every physical
block and a hot block (an ext4 bitmap, say) never wears out one sector. The state
is written to two configuration blocks \emph{ping-pong} (alternating, each with a
higher sequence number and a CRC), so a power loss mid-update always leaves a
consistent earlier state; \texttt{Mount} picks the highest valid sequence.

That same ``highest sequence wins'' rule is why \texttt{Format} cannot just write a
fresh, low-sequence record and stop: a chip that was used before still holds an old
record with a far higher sequence in the \emph{other} ping-pong slot, and
\texttt{Mount} would then pick that stale slot and silently bring the volume back up
with its pre-format wear state, as though the format never happened.
\texttt{Format} therefore \emph{invalidates both} config slots first --- zeroing
them, so their magic and version no longer parse --- and only then writes the new
sequence-one record, so the fresh state is the only thing \texttt{Mount} can find.

The whole stack (erase-aware writes, the wear-level remap, the auto-detected size)
comes together in four calls:

\begin{lstlisting}
   BDW.Configure (Raw, Flash => Flash);      --  auto-size from the JEDEC id
   WL.Attach (Vol, BDW.Make (Raw'Access));   --  wear-leveling volume over it
   WL.Mount  (Vol, Formatted);               --  recover the saved wear state
   Dev := WL.Make (Vol'Access);              --  a 512-byte Block_Dev for ext4
\end{lstlisting}

\subsection{Formatting on-device: \texttt{Ext4.Mkfs}}\index{mkfs}\index{Ext4.Mkfs}\index{on-device format}

The external flash is not on the host flasher's path, and the filesystem reads and
writes but does not format, so to start from a blank chip the board makes the
filesystem itself. \texttt{ESP32S3.Ext4.Mkfs} is a minimal pure-Ada
\texttt{mkfs.ext4}: it writes a single-block-group ext4 (4\,KiB blocks, classic
block-mapped inodes, a root directory and \texttt{lost+found}), with an optional
4\,MiB JBD2 journal (\texttt{Journal => True}). It is the inverse of the read path,
validated against the host's \texttt{e2fsck}. From a blank chip to a mounted,
written filesystem is then:

\begin{lstlisting}
   --  ... build Dev as above, but WL.Format (Vol) for a fresh volume ...
   ESP32S3.Ext4.Mkfs.Format (Dev, Volume_Label => "FLASH", Journal => True);
   M.Open (Dev, Read_Only => False);
   N := M.Create_File ("/", "boot.txt");
   M.Write_File (N, (Character'Pos ('h'), Character'Pos ('i'), 10));
   M.Commit;                                 --  through the journal just created
\end{lstlisting}

\noindent The on-target examples \texttt{esp32s3\_w25q}, \texttt{esp32s3\_wl} and
\texttt{esp32s3\_ext4\_mkfs} run this stack end to end on the hardware.

\subsection{Writing in runs, not sectors}\index{Write\_Run}\index{erase amplification}

One more W25Q\_Source fact earns its space, because it is an 8$\times$
difference. The block layer's \texttt{Write\_Sector} is erase-aware but
per-sector: any change that is not purely clearing bits read-modify-writes the
whole 4\,KB erase block. Hand it the eight 512-byte sectors of one block, one
at a time, and it erases that block eight times --- eight times the wear, and at
$\sim$45\,ms an erase, eight times the wait. \texttt{Write\_Run} walks a run
of consecutive sectors one \emph{erase block} at a time instead: a block the
run covers completely is erased once and programmed; a block covered in part
takes one read-modify-write, not one per sector. Callers reach it through
\texttt{Block\_Dev.Write\_Sectors}, which uses the run primitive when the
device offers one and falls back to the per-sector loop when it does not.
Measured on a 32\,MB part behind a USB mass-storage function: a 64\,KiB
write fell from a per-sector 8\,kB/s crawl to wire speed.

\section{FAT16: the filesystem a PC can just mount}
\index{FAT16}\index{Fat16.Mkfs}\index{long filenames}

ext4 is the right home for the board's own data. It is exactly wrong for one
job: a volume a Windows PC must mount without drivers or questions.
\texttt{ESP32S3.Fat16} exists for that job --- a board exposes its flash over
USB mass storage, a PC drops files on it, the board reads them back. Interop
is the whole point, so the design follows what PCs expect rather than what
would be elegant.

The reader mounts an MBR-partitioned disk (the first FAT16 partition) or a
bare-boot-sector ``superfloppy'', walks long filenames as well as 8.3 names,
matches names case-insensitively because FAT's own rule does, and streams:
\texttt{Open} / \texttt{Read} / \texttt{Seek} over a caller-supplied buffer,
no heap. A \texttt{Read} that comes back short means end of file, never a
transient.

\texttt{Fat16.Mkfs.Format} lays down an empty volume so a blank chip appears
to the first PC as a usable drive rather than as ``you need to format this
disk''. Two facts drive its geometry. Windows expects a USB drive to carry a
partition table with the filesystem inside partition~1, so that is the
default. And the backing medium is NOR flash with a 4\,KB erase unit, so
clusters are 4\,KB and the data area starts on an erase-block boundary --- a
cluster never straddles two erase blocks, and a cluster write costs the block
layer one erase. Formatting is destructive and unconditional; deciding
\emph{whether} is the caller's job, and the honest test for ``blank'' is the
medium's own bytes (erased NOR reads all ones), not a missing signature ---
a volume the reader cannot parse looks exactly like no volume at all.

The host harness (\texttt{test/fat16\_host}) checks the pair three ways
against reference tooling: volumes written by \texttt{dosfstools}, volumes
written by Python, and volumes written by \texttt{Mkfs} itself, cross-read
and \texttt{fsck}-checked. One application built on this SDK runs the full
loop: it formats its own flash, serves it to a PC as a USB drive, and later
reads back the firmware images the PC left there.


\chapter{Conformance at Scale: A Standalone ACATS Sweep}
\label{ch:acats-sweep}

The runtime is only as trustworthy as the conformance evidence behind it. This
chapter describes how the project runs the official ACATS 4.2 suite \emph{one test
per image} across a pool of ESP32-S3 boards in parallel, and why that arrangement
recovers hundreds of passes that a bundled sweep loses.

\medskip
\noindent\fbox{\parbox{\dimexpr\linewidth-2\fboxsep-2\fboxrule\relax}{\small
\textbf{Note.} The ACATS test suite and its sweep harness
(\texttt{acats\_build.py}, \texttt{acats\_run.py}, the \texttt{ACATS/} tree) are
not shipped in this distribution; they belong to the development repository.
This chapter is kept as a record of the conformance methodology and results; the
commands below are shown to document \emph{how} the sweep was run, not as steps to
reproduce from this repository.}}
\medskip

\section{Two ways to run the suite}\index{ACATS}\index{conformance testing}\index{test harness}

The earlier sweep (\texttt{parallel\_sweep.py}) is \emph{bundled}: dozens of tests
share one flashed image and advance through it with a panic-resume cursor, so a
crash or hang resets the board and the next boot continues past the offending test.
That is fast and cheap on flash cycles, but it couples the tests. A neighbour's
heap churn or secondary-stack residue can exhaust memory mid-batch and reset a later
test \emph{before} it prints its grade: the test is recorded as ungraded even
though it passes in isolation. On the \texttt{full} profile the bundled run grades
914 PASS, with a large ungraded tail dominated by these batch artifacts.

The \emph{standalone} sweep removes the coupling completely: every test is compiled
as its own single-test program, flashed alone, run alone, and graded alone. There is
no shared cursor, no shared heap, no neighbour. The cost (roughly 1{,}500 separate
cross-compiles and flashes) is absorbed by running the work across many boards and
many build trees at once.

\section{Getting the grade off the chip}\index{USB-Serial-JTAG}\index{esp\_rom\_printf}

A subtlety decides whether the standalone approach is even possible. The bare runtime
sends its console (including the ACATS \texttt{Report} \texttt{PASSED}/\texttt{FAILED}
lines) out the ESP32-S3's \emph{built-in USB-Serial-JTAG}, not the UART pins. The
board pool here is wired with ordinary external USB-serial adapters (FTDI and CH343)
on UART0, which never see that JTAG console.

The fix is small. \texttt{ACATS/target/report.adb} routes each grade line through the
ROM's \texttt{esp\_rom\_printf} in addition to \texttt{Ada.Text\_IO}:

\begin{lstlisting}
procedure Line (Msg : String) is
   procedure Rom_Printf (Fmt, Arg : System.Address)
     with Import, Convention => C, External_Name => "esp_rom_printf";
   Fmt : constant String := "%s" & ASCII.LF & ASCII.NUL;
   Buf : String (1 .. Msg'Length + 1);
begin
   Buf (1 .. Msg'Length) := Msg;
   Buf (Buf'Last) := ASCII.NUL;
   Rom_Printf (Fmt'Address, Buf'Address);
end Line;
\end{lstlisting}

The ROM console reaches UART0, so the grades land on the external bridge regardless
of how the board is wired.

\section{The grader must not trust the test id}\index{grading}

The first version of the harness mis-scored most tests as hung. The cause was the
grading regex: it looked for the list name followed by \texttt{PASSED}, but the list
says \texttt{A22006} while \texttt{Report} prints the \emph{unit} name \texttt{A22006B}
-- the suffix letter breaks an exact-id match. Because each image contains exactly one
test, the grader instead matches the \texttt{Report} terminal line \emph{structurally},
ignoring the name:

\begin{lstlisting}[style=sh]
FAILED          ->  **** <anything> FAILED          (checked first)
NOT-APPLICABLE  ->  ++++ <anything> NOT-APPLICABLE
PASSED          ->  ==== <anything> PASSED
TENTATIVELY     ->  !!!! <anything> TENTATIVELY PASSED
\end{lstlisting}

A second early bug scored passing tests as crashes: the flash step reset the board to
run, then the capture step reset it again, producing two \texttt{ESP-ROM:} banners --
and the crash heuristic keyed on ``two or more boot banners.'' Flashing with
\texttt{-{}-no-reset} leaves the chip halted so the capture's own reset is the single
deterministic boot; a second banner then genuinely means a crash-reboot.

\section{Phase 1: parallel build}\index{parallel build}\index{mini-repo}

\texttt{acats\_build.py} compiles each test as a one-test image. The build is the
bottleneck (a single board sits idle for tens of seconds per cross-compile), so it is
parallelised. The complication is that the build performs several \emph{shared} writes:
it regenerates the runtime (\texttt{gen\_runtime.sh}), rebuilds
\texttt{common/bare/boot/obj}, and stages test sources under \texttt{ACATS/work/run},
so naive parallel builds in one tree corrupt each other.

Each builder therefore gets its own isolated ``mini-repo'': a copy of
\texttt{crates}, \texttt{libs}, \texttt{ACATS}, \texttt{ACATS42}, and the two needed
\texttt{examples} subtrees, laid out at the same depth so the build scripts'
\texttt{REPO=\$EX/../..} arithmetic still resolves. The copy carries the warm caches,
so per-test builds stay incremental: no cold runtime rebuild per worker.

\begin{lstlisting}[style=sh]
# OUT/<name>.bin       built OK (the flashable image)
# OUT/<name>.nobuild   could not build / all tests dropped (empty marker)
# OUT/BUILD_DONE        written once every test has been attempted
./acats_build.py ACATS/work/full_applicable.txt OUT 10
\end{lstlisting}

A test that drops to zero buildable units (it needs a library unit the bare runtime
omits) produces \texttt{tests in image: 0} and becomes a \texttt{.nobuild}. The
\texttt{.bin} is written atomically (a temp file renamed into place) so Phase 2 never
sees a partial image, and the phase is resumable: a test that already has a
\texttt{.bin} or \texttt{.nobuild} is skipped.

\section{Phase 2: parallel run and grade}\index{board pool}

\texttt{acats\_run.py} starts one worker per board. Each worker claims the next
un-run \texttt{.bin}, flashes its board with \texttt{-{}-no-reset}, pulses a reset,
captures the console until it sees a grade (or a timeout), classifies the run, and
writes a single marker file holding the raw log:

\begin{lstlisting}[style=sh]
# OUT/<name>.passed  .failed  .na  .crashed  .hung  .nooutput  .flasherr
./acats_run.py OUT 20        # 20-second capture window
\end{lstlisting}

The phases may overlap: Phase 2 drains \texttt{.bin} files as they appear and exits
once \texttt{BUILD\_DONE} exists and no un-run image remains. The whole sweep is
resumable at the file level: delete a marker and re-run, and only that test is
re-flashed. A handful of tests that briefly produced no console (a capture race during
parallel flashing) and several genuinely slow tasking tests were re-run with a 60-second
window simply by removing their markers.

\section{Results}

Full profile, \texttt{full\_applicable.txt} (1{,}518 tests), one image each:

\begin{longtable}{p{4.5cm} r p{6.5cm}}
\toprule
\textbf{Outcome} & \textbf{Count} & \textbf{Meaning} \\ \midrule
\endhead
PASS                & 1{,}286 & graded \texttt{PASSED} / \texttt{TENTATIVELY PASSED} \\
NOT-APPLICABLE      & 10      & test ruled itself out (correctly) \\
FAIL                & 0       & \\
HUNG                & 4       & documented limitations (see below) \\
CRASH               & 0       & \\
won't-build         & 218     & needs a library unit the bare runtime omits \\
\bottomrule
\end{longtable}

That is \textbf{1{,}296 conformant} (PASS plus N/A), \textbf{zero failures}, and
\textbf{zero crashes}, and \textbf{+372 PASS} over the bundled \texttt{full} run
(914), recovered purely by removing batch interaction. Building each test alone is what
turns a co-bundled ``ungraded'' tail back into passes.

\textbf{Three apparent FAIL were harness artifacts, now resolved.} \texttt{CXD1002} and
\texttt{CXD1007} check the \emph{main task's} \texttt{pragma Priority}, live only when the
test is built \emph{as the program main}; the harness now honours
\texttt{standalone\_tests.txt}'s \texttt{main} mode and they PASS. \texttt{CZ1101} is the
ACATS \emph{Report self-test}, which deliberately prints \texttt{**** FAIL\_TEST FAILED}
sub-test sentinels; the grader now matches the result line for the test's \emph{own} name
and takes its final verdict (\texttt{!!!! CZ1101A TENTATIVELY PASSED}).

\textbf{Five apparent HUNG were under-resourced, now passing.} \texttt{CXD1008} and
\texttt{CXD4007} are silent slow tests (\texttt{ImpDef.Switch\_To\_New\_Task} is 1.0\,s, so
they need minutes, not the default 60\,s window); \texttt{C433A01/02} and \texttt{C953002}
OOM'd at the 64\,KB secondary-stack default. The harness now honours two tuning lists:
\texttt{long\_wdt\_tests.txt} (a 240\,s capture window) and \texttt{small\_sec\_stack\_tests.txt}
(the 4\,KB default for many-task tests), and all five PASS.

\textbf{Two more were program-termination cases, now fixed.}
\texttt{C761001}'s grade is emitted from a \emph{library-level} controlled object's
\texttt{Finalize}, which runs only at program termination, but
\texttt{Configurable\_Run\_Time} makes gnatbind omit the \texttt{finalize\_library} chain,
so it never ran. Flipping that flag is not viable (it needs a full runtime rebuild and is
entangled with the runtime's codegen; the \texttt{.ali} flags are identical native-vs-target,
so it is purely the binder mode). Since the per-unit finalizers
(\texttt{<unit>\_\_finalize\_spec/body}) are still generated, the build synthesizes a
finalizer-runner that calls the user units' finalizers in reverse elaboration order (parsed
from the binder's \texttt{adainit}; weak + guarded), merges it into the image, and the bare
env body calls it after \texttt{adafinal}, reached only by a program-main image after its
main returns, so the wrapped sweep images are unaffected. \texttt{C761001} now PASSES.

\texttt{C94004} is a multi-file test whose real main is \texttt{C94004A}, a library task
``examined to see whether it terminates.'' As a wrapped batch test its non-terminating
library tasks pile up and hang; built as the program main, only \texttt{C94004A} runs and
prints \texttt{PASSED} from its own body. Its main unit (\texttt{c94004a}) differs from the
list key (\texttt{C94004}), so the build resolves the boot entry from
\texttt{batch\_runnable.txt} (\texttt{build\_ada.sh} writes \texttt{.noidf/ada\_main.sym},
which \texttt{bare\_build} honours instead of the guessed \texttt{\_ada\_c94004}).
\texttt{C94004} now PASSES.

\textbf{The remaining four HUNG are documented limitations.} \texttt{CXD1006} (an
\texttt{Interrupt\_Priority} task collides with the kernel's own Xtensa level 5),
\texttt{CB1010} (a large-frame recursive stack overflow the 64-byte watchpoint cannot
catch), \texttt{C480001} (unbounded recursion / GNAT codegen), and \texttt{CXC7004}
-- \emph{not} a deadlock but an environment-task \emph{termination} test (JTAG dual-core
post-mortem shows cpu0 spun past \texttt{ADA\_MAIN}, cpu1 idle) whose grade depends on
shutdown-path behavior the Configurable runtime does not provide.

\textbf{The ten N/A are correct.} \texttt{CXD2001-2004} and \texttt{CXD6001-6003} test
\emph{uniprocessor} dispatching (a single ready queue with deterministic ordering),
which does not apply to a genuine two-core SMP target, so each prints
\texttt{+ Multi-Processor Configuration} and grades NOT-APPLICABLE. \texttt{C45322},
\texttt{C45523}, and \texttt{C4A012} rule themselves out for their own reasons.

\subsection{A wider net: the maximal list across seven boards}\index{Ada.Interrupts}\index{Ada.Calendar}

A later run pushed the net wider on purpose. Rather than the curated
\texttt{full\_applicable.txt}, it swept the \emph{maximal union} (the applicable
list plus everything previously parked in \texttt{jorvik\_skipped.txt} and
\texttt{ungraded\_tests.txt}, 1{,}536 tests): one image each, across \emph{seven}
boards at once (the standalone harness is one worker per serial console, so seven
consoles run seven tests in parallel). The point of including the already-skipped
set is that capabilities added since the lists were drawn (the Calendar children,
and the interrupt work) might turn an old \texttt{nobuild} into a pass; running the
superset catches that automatically.

It did, and it also earned its keep by finding a real bug. The dynamic
\texttt{Ada.Interrupts} tests \texttt{CXC3007}/\texttt{CXC3008} \emph{failed},
which is how the restricted-\texttt{a-interr}-stub problem (every dynamic operation
\texttt{raise Program\_Error}; see the Interrupts chapter) came to light. With the
delegating \texttt{a-interr} body in place, \texttt{CXC3007} now \texttt{PASSES} and
\texttt{CXC3008} advances into its interactive ``generate the interrupt'' phase
(which the unattended rig cannot satisfy). The wider net found a second real bug
too: \texttt{C954025}/\texttt{C954026} (a requeue-with-abort whose timeout is a
\texttt{delay until} an \texttt{Ada.Calendar} time) expired at the wrong moment,
because the kernel's timed-sleep collapsed the \texttt{Absolute\_Calendar} and
\texttt{Absolute\_RT} delay modes and fed a Calendar-base deadline to the raw
counter. Shifting a Calendar deadline back into the raw base (adding the BB-time
\texttt{Epoch}) in \texttt{Timed\_Sleep} fixes both. They now \texttt{PASS}
(they are slow and silent, so they join the wide-watchdog list). After these
fixes the residue on the wider list is what you would expect: the documented
frontiers (\texttt{CXD1006}, \texttt{CB1010}, \texttt{C480001}, \texttt{CXC7004})
and the interactive interrupt and suspension tests that need a bench-generated
interrupt. No crashes, and no surprise failures in the conformant core.
(\texttt{C96004}'s post-2099 Calendar range was later closed too, by coarsening the
Julian-day base $5\times$ to reach 2399; see the Limitations chapter. A companion
fix set \texttt{System.Tick} to the true \texttt{1/240\,MHz} clock period (it had
shipped as \texttt{0.0}, telling clock-resolution tests the clock was infinitely
precise), which is what \texttt{C96001}/\texttt{C96004} read.)

\section{The embedded profile}\index{embedded profile}\index{Jorvik}\index{No\_Task\_Hierarchy}

The same harness runs the \texttt{embedded} (ZCX) profile against the
\texttt{jorvik\_hw\_runnable.txt} list (846 tests). It needed one runner fix first. For a
single-test image the generated driver adds a do-nothing task to force the environment
task's secondary stack; that task was declared \emph{nested inside} the \texttt{Main}
procedure, which is fine under \texttt{full} but violates the embedded / light-tasking
Jorvik restriction \texttt{No\_Task\_Hierarchy} (a task nested in a subprogram is illegal).
Every size-1 image therefore failed to build: the standalone sweep produced
all-no-build. The bundled \texttt{parallel\_sweep} never hit this because it pads
isolated images with filler tests to keep the batch larger than one, so the helper task
was never emitted. The fix is to emit the helper as a \emph{library-level} task instead,
which is legal under Jorvik and accepted by \texttt{full} alike.

\begin{longtable}{p{4.5cm} r p{6.5cm}}
\toprule
\textbf{Outcome} & \textbf{Count} & \textbf{Meaning} \\ \midrule
\endhead
PASS                & 840 & graded \texttt{PASSED} / \texttt{TENTATIVELY PASSED} \\
NOT-APPLICABLE      & 3   & test ruled itself out (correctly) \\
FAIL                & 1   & \texttt{CXC3004}, the interactive test (see below) \\
HUNG                & 0   & \\
CRASH               & 0   & \\
won't-build         & 2   & partition-wide build constraints (see below) \\
\bottomrule
\end{longtable}

\textbf{843 conformant, zero HUNG, zero CRASH, zero genuine FAIL}, above the bundled
embedded baseline (836). The one FAIL, \texttt{CXC3004}, is the \emph{interactive} test:
it waits for a manually-generated external interrupt, which nothing raises on the
unattended rig: a documented limitation, not a runtime bug.

The two won't-build are notable because neither is a missing Ada library unit --
each is a \emph{partition-wide} build constraint. \texttt{CXH1001} (Annex H) uses the
configuration pragma \texttt{Normalize\_Scalars}, which must be applied to every unit in
the partition including the runtime; the prebuilt embedded runtime is not compiled with it,
so the binder rejects the program (\emph{``some but not all files compiled with
Normalize\_Scalars''}). \texttt{CA5006} (a \texttt{pragma Elaborate} separate-compilation
test) binds to an \emph{elaboration circularity} under the embedded runtime's static
elaboration model. Both are correctly counted as won't-build rather than failures.

\section{Closing a won't-build category: the Ada.Calendar family}\index{calendar overlay}

Some of the ``won't-build'' tail is not a fundamental limitation but a simply-missing library
unit, and the sweep is what makes that visible. The bare runtime originally shipped only base
\texttt{Ada.Calendar} (plus \texttt{.Conversions}); the RM~9.6.1 \emph{child} units were
absent, so every Calendar formatting/arithmetic test dropped at compile. Those children are now
provided on the \texttt{embedded} and \texttt{full} profiles, from a shared
\texttt{calendar\_overlay/} applied by \texttt{gen\_runtime.sh}:

\begin{itemize}
  \item \texttt{Ada.Calendar.Arithmetic}: day arithmetic on the parent's
        Modified-Julian-Day representation (exact; the bare runtime has no leap-second table,
        so \texttt{Leap\_Seconds} is always zero).
  \item \texttt{Ada.Calendar.Formatting}: \texttt{Image}/\texttt{Value}/\texttt{Split}/
        \texttt{Time\_Of}/\texttt{Day\_Of\_Week}, written against the parent's public API with
        hand-rolled decimal formatting (the runtime excludes the \texttt{'Image}/\texttt{'Value}
        helper units, so the displayed fraction is truncated and \texttt{Value} raises on
        out-of-range fields, both as the language requires).
  \item \texttt{Ada.Calendar.Time\_Zones}: bare metal has no time-zone database, so the
        Calendar keeps UTC and \texttt{Local\_Time\_Offset} returns zero.
  \item \texttt{Ada.Calendar.Delays}: \texttt{delay until} on a \texttt{Calendar.Time},
        delegated to the SMP-hardened \texttt{Ada.Real\_Time.Delays}.
  \item the obsolescent J.1 top-level \texttt{Calendar} rename, so Ada-83 \texttt{with Calendar;}
        tests build.
\end{itemize}

Hardware-verified: \texttt{C961001} (Formatting \texttt{Image}/\texttt{Value} conformance) and
\texttt{C960002} (\texttt{Delay\_Until}) PASS on both profiles, and the \texttt{with Calendar;}
tests \texttt{C96005}/\texttt{C96006} PASS. \texttt{C96004}, which requires
\texttt{Year\_Number'Last = 2399} (Ada~2005), now PASSES as well: rather than the
full time-base rewrite first assumed, the Julian-day base is coarsened $5\times$
(sub-day scale \texttt{Duration'Small\,/\,17\_280.0}), trading 1\,ns Calendar
resolution for 5\,ns and buying the extra three centuries within signed 64 bits --
see the Limitations chapter. \texttt{light-tasking} is deliberately
\emph{excluded}: it is a \texttt{No\_Exception\_Propagation} profile, so every
\texttt{Time\_Error}/\texttt{Constraint\_Error} becomes a board reset. \texttt{Ada.Calendar}
is unusable and grades nothing there, which is exactly why bb-runtimes omits it from the Jorvik
profile. By contrast \texttt{Ada.Numerics} needed no work at all: the parent, elementary
functions, the random generators, complex types, real arrays and big integers already ship and
compile into both profiles (HW-verified via \texttt{CXG2003}, \texttt{CXG2009}, \texttt{CXG1004}).

\section{Why this matters}

The standalone sweep is slower per test, but it produces an unambiguous, file-per-test
record: every outcome is one marker file containing the exact console transcript that
justified it. Cross-checking the two methods is itself a result: the bundled run's
ungraded tail is explained, test by test, by the standalone run's passes, and the only
remaining non-passes are a short, named list of artifacts and documented frontiers.

%  Under-the-hood: how the runtime is built, how the chip boots, then the
%  scheduler internals it all leads into.
\chapter{Building and Porting the Runtime}
\label{ch:runtime-build}

The earlier \emph{Building and Running} chapter (\S\ref{ch:build}) showed the
runtime from the outside: pick a profile, run \texttt{./x}, get an image. This
chapter opens the box. It explains how the three runtimes are actually
\emph{constructed}: where the Ada runtime sources come from, what makes one
profile differ from another, how the experimental \texttt{full} profile is
synthesised, and how the register packages are generated. If you ever need to
change a runtime unit, add a peripheral, or port the whole stack to another
Xtensa chip, this is the map.

\section{The idea: a board on top of \texttt{bb-runtimes}}\index{bb-runtimes}

The runtime is not written from scratch. It is grown from AdaCore's
\texttt{bb-runtimes}, the same generator that produces the bareboard GNAT
runtimes for Cortex-M, LEON, RISC-V and the rest. \texttt{bb-runtimes} knows how
to assemble a runtime from a large pool of shared Ada units; what it needs from
\emph{you} is a \emph{board definition}: the target triple, the clock, the
interrupt model, and a handful of board-specific bodies (the console, the timer,
machine reset).

\texttt{bb-runtimes} lives in the repository as a git submodule: a fork
(\texttt{rowsail/bb-runtimes}, branch \texttt{xtensa-esp32s3-port}) that carries
the ESP32-S3 board on top of the upstream sources. Keeping it a submodule means
the board port is version-pinned and the upstream history stays visible.

\cnote{There is no equivalent step in an ESP-IDF project: the C toolchain ships a
fixed newlib and FreeRTOS. Here the language runtime itself is a build product
you can read, patch and regenerate: the scheduler, the tick, the secondary
stack and the exception machinery are all Ada source in the tree.}

\section{\texttt{gen\_runtime.sh}: from board to archive}\index{gen\_runtime.sh}

One script, \texttt{crates/esp32s3\_rts/gen\_runtime.sh}, turns the board
definition into a ready-to-link runtime. It is driven by a single environment
variable:

\begin{lstlisting}[style=sh]
export ESP32S3_RTS_PROFILE=embedded      # light-tasking (default) | embedded | full
crates/esp32s3_rts/gen_runtime.sh
\end{lstlisting}

The script does four things:

\begin{enumerate}
\item \textbf{Generate the source tree.} It runs \texttt{bb-runtimes}'
  \texttt{build\_rts.py} against the \texttt{esp32s3} board, which copies the
  selected pool units plus the board's own sources into a profile-specific
  directory:
\begin{lstlisting}[style=sh]
python3 build_rts.py -f -o "$HERE" \
   --rts-src-descriptor=gnat_rts_sources/lib/gnat/rts-sources.json esp32s3
\end{lstlisting}
  The output lands in \texttt{crates/esp32s3\_rts/<profile>-esp32s3/}: a
  complete runtime with \texttt{gnat/} (sequential units) and \texttt{gnarl/}
  (tasking units) subdirectories. These directories are \emph{git-ignored}: they
  are build products, regenerated on demand.
\item \textbf{(\texttt{full} only) Apply the overlay and patches} (the next
  section).
\item \textbf{Compile} the tree with \texttt{gprbuild} against a small
  \texttt{ravenscar\_build.gpr}.
\item \textbf{Archive} the objects into \texttt{adalib/libgnat.a} (sequential)
  and \texttt{adalib/libgnarl.a} (tasking) with \texttt{xtensa-esp32-elf-ar}.
  The \texttt{esp32s3\_rts.gpr} that applications \texttt{with} points its
  \texttt{Runtime} at \texttt{<profile>-esp32s3/}, so consuming the runtime is
  just naming the crate.
\end{enumerate}

Crucially the compile-and-archive step is guarded: it runs only when
\texttt{adalib/libgnat.a} is \emph{absent}. A second invocation with the tree
already built is therefore nearly instant, but it is also the source of a trap
(\S\ref{sec:stale-cache}).

\section{Three profiles, three \texttt{system.ads}}\index{runtime profile}

What actually distinguishes the profiles is mostly one file: \texttt{system.ads},
the package that declares the target's restrictions and capabilities. The board
class maps each profile name to a variant:

\begin{center}
\begin{tabular}{l l}
\toprule
\textbf{Profile} & \textbf{\texttt{system.ads} variant} \\ \midrule
\texttt{light-tasking} & \texttt{system-xi-xtensa-light-tasking.ads} \\
\texttt{embedded}      & \texttt{system-xi-xtensa-embedded.ads} \\
\texttt{full}          & synthesised locally (see below) \\
\bottomrule
\end{tabular}
\end{center}

The \texttt{light-tasking} variant carries the restrictive pragmas:
\texttt{No\_Exception\_Propagation}, \texttt{No\_Finalization},
\texttt{No\_Secondary\_Stack}, all under \texttt{pragma Profile (Jorvik)}. The
\texttt{embedded} variant is the \emph{same Jorvik profile} with those three
restrictions lifted: exceptions propagate (zero-cost, ZCX), controlled types
finalize, and a secondary stack exists. Changing profile is, at heart, swapping
this declaration. Everything downstream (which pool units \texttt{build\_rts}
selects, which code the compiler generates) follows from it.

The board's own bodies are shared across profiles and named with the
\texttt{\_\_esp32s3} convention that \texttt{bb-runtimes} uses for board files:

\begin{center}
\begin{tabular}{l l}
\toprule
\textbf{File} & \textbf{Role} \\ \midrule
\texttt{s-bbpara\_\_esp32s3.ads} & board parameters: 240\,MHz clock, interrupt range, stacks \\
\texttt{s-bbbosu\_\_esp32s3.adb} & board support: the tick (\texttt{CCOMPARE}), interrupt dispatch \\
\texttt{s-textio\_\_esp32s3.adb} & console \texttt{Put} via the ROM \texttt{esp\_rom\_printf} \\
\texttt{s-macres\_\_esp32s3.adb} & machine reset \\
\texttt{a-intnam\_\_esp32s3.ads}  & \texttt{Ada.Interrupts.Names} for the CPU's interrupts \\
\bottomrule
\end{tabular}
\end{center}

\section{The \texttt{full} profile: overlay plus patches}\index{full profile}

There is no upstream Jorvik-free Xtensa runtime, so \texttt{full} is built on top
of \texttt{embedded} rather than generated directly. \texttt{gen\_runtime.sh}
generates the embedded tree, then layers two things on it:

\begin{itemize}
\item an \textbf{overlay} (\texttt{full\_overlay/}): a de-restricted
  \texttt{system.ads} and a set of GNARL units (interrupt management, the OS
  primitives, the tasking primitives) that replace the Jorvik-limited versions
  with the full ones; and
\item a small set of \textbf{source patches}, kept as ordinary \texttt{.patch}
  files and applied with GNU \texttt{patch}:
\end{itemize}

\begin{lstlisting}[style=sh]
patch -p1 -d "$RTS" < full_overlay/patches/01-...-secondary-stacks.patch
\end{lstlisting}

Each patch fixes one concrete defect that only surfaces under full tasking, and
each is documented at its head:

\begin{center}
\begin{tabular}{p{4.6cm} p{8.6cm}}
\toprule
\textbf{Patch} & \textbf{What it fixes} \\ \midrule
secondary-stack reclaim & free a task's heap secondary stack when it terminates \\
trampoline IRAM alias & point nested-subprogram trampolines at the executable
  (IRAM) alias of SRAM, not the data view (otherwise the first call faults) \\
\texttt{To\_Time\_Span} rounding & a nonzero \texttt{Duration} must round to
  $\geq 1$ tick (the 4.17\,ns tick is coarser than \texttt{Duration'Small}) \\
dynamic-priority re-apply & re-apply a \texttt{Set\_Priority} after the lock is
  released, which the bareboard unlock would otherwise undo \\
FIFO-within-priority & implement RM~D.2.3 so a base-priority change moves a task
  to the tail of its ready queue \\
\bottomrule
\end{tabular}
\end{center}

Keeping these as diffs rather than wholesale copied files means each change is
small, reviewable, and obviously attributable to a specific runtime behaviour,
and it keeps \texttt{gen\_runtime.sh} readable.

\section{Register packages from the SVD}\index{svd2ada}\index{SVD}

The HAL never hard-codes a peripheral address. Every register is a typed Ada
record drawn from \texttt{ESP32S3\_Registers.*}, and those packages are
\emph{generated} from Espressif's machine-readable SVD (System View Description)
by \texttt{svd2ada}. The generator script,
\texttt{libs/esp32s3\_hal/regenerate.sh}:

\begin{enumerate}
\item fetches a pinned \texttt{esp32s3.svd} (a specific \texttt{espressif/svd}
  commit, overridable with \texttt{ESP32S3\_SVD=...});
\item builds \texttt{svd2ada} \emph{from AdaCore source} (the Alire-indexed
  build is too old and mishandles dimensioned register arrays);
\item runs it into \texttt{libs/esp32s3\_hal/svd/}:
\begin{lstlisting}[style=sh]
svd2ada esp32s3.svd -o svd -p ESP32S3_Registers --boolean
\end{lstlisting}
\end{enumerate}

The root package is named \texttt{ESP32S3\_Registers} on purpose: GNAT forbids
user-defined children of \texttt{Interfaces}, so the natural
\texttt{Interfaces.ESP32S3} is not allowed. One wrinkle needs a post-process:
the SVD has a peripheral literally called \texttt{SYSTEM}, which becomes
\texttt{ESP32S3\_Registers.SYSTEM}, and because Ada is case-insensitive, a
bare \texttt{System.\dots} reference inside that unit resolves to the wrong
package. A one-line \texttt{sed} rewrites those references to the fully qualified
\texttt{Standard.System.\dots}. Add a peripheral and the recipe is the same: add
\texttt{ESP32S3.<Peri>} to the HAL, backed by the generated register record.

\section{The stale-cache trap}\index{stale cache}
\label{sec:stale-cache}

Because the heavy compile/archive step runs only when
\texttt{adalib/libgnat.a} is missing, editing a runtime source \emph{in the
generated tree} does not, on its own, rebuild anything: the next build links
the old archive and your change silently does nothing. This has burned us more
than once, so it is worth stating plainly:

\begin{quote}
\textbf{After editing a generated runtime \texttt{.adb}, delete the archives
before rebuilding:}
\begin{lstlisting}[style=sh]
rm crates/esp32s3_rts/<profile>-esp32s3/adalib/libgnat.a \
   crates/esp32s3_rts/<profile>-esp32s3/adalib/libgnarl.a
\end{lstlisting}
\end{quote}

(Editing a source under \texttt{bb-runtimes/} or \texttt{full\_overlay/} instead
regenerates correctly because the whole tree is rebuilt; the trap is specifically
editing the \emph{regenerated copy}.) The same caution applies to verifying any
runtime fix: confirm the change is in the linked archive, not just on disk,
before drawing a conclusion.

\section{Porting to another board}

The structure makes the porting checklist short. To bring the runtime to another
Xtensa SoC you would: add a board class under \texttt{bb-runtimes/xtensa/} (target
triple, clock, interrupt range, SMP flag); supply the five \texttt{\_\_<board>}
bodies (parameters, board support, console, reset, interrupt names); write the
two \texttt{system-xi-xtensa-*.ads} variants (or reuse these if the restriction
set is the same); regenerate the register packages from that chip's SVD; and
point \texttt{gen\_runtime.sh} at the new board name. The boot path (the
2nd-stage loader, the clock and cache bring-up, PSRAM) is the subject of the
next chapter.

\chapter{The Boot Path: Bootloader and PSRAM}
\label{ch:bootloader}

When the chip comes out of reset, a great deal has to happen before the first
line of Ada in your \texttt{Main} runs: the application has to be read out of
flash and placed in RAM, the instruction and data caches have to be pointed at
the right flash windows, the CPU has to be sped up from its sleepy reset clock,
the register windows have to be made sane for Ada's calling convention, and --
if the board has it -- the external PSRAM has to be brought up and mapped. This
chapter follows that path. It is the part of the system that is \emph{not}
Espressif's: the project replaced the last vendored binary, the 2nd-stage
bootloader, with its own -- written, like almost everything else here, in Ada.

\section{Three things run before \texttt{Main}}

Boot happens in three hand-offs (the boot figure in
Chapter~\ref{ch:anatomy}, Fig.~\ref{fig:boot}, sketches the same sequence):

\begin{enumerate}
\item \textbf{The mask ROM} (silicon, unchangeable) wakes on reset, reads the
  image at flash offset \texttt{16\#0\#}, and loads it into SRAM. That image is our
  \emph{2nd-stage bootloader}.
\item \textbf{Our 2nd-stage bootloader} reads the \emph{application} image (at
  flash offset \texttt{16\#1\_0000\#}), copies its RAM segments into place, maps its
  flash-resident code and data through the cache MMU, brings up PSRAM, and jumps
  to the application's entry point.
\item \textbf{The application's \texttt{start.S}} normalises the CPU state,
  configures the instruction cache, switches to the full clock, and calls into
  the Ada world (eventually the environment task that runs \texttt{Main}).
\end{enumerate}

\section{The 2nd-stage bootloader}\index{bootloader}

The loader is ZFP-style Ada -- no runtime under it -- in
\texttt{examples/common/bare/bootloader/} (the core is \texttt{boot\_main.adb}).
It is deliberately small. Its job, in order:

\begin{itemize}
\item \textbf{Silence the watchdogs.} The boot ROM arms the timer-group and RTC
  watchdogs; the loader unlocks them (the magic key \texttt{16\#50D8\_3AA1\#}) and
  clears them so they cannot reset the chip mid-load.
\item \textbf{Read and validate the app header.} At flash offset
  \texttt{16\#1\_0000\#} sits the image header: a \texttt{16\#E9\#} magic byte, a segment
  count, and the entry-point address. (This is exactly the format the
  \texttt{esp\_elf2image} host tool emits -- \S\ref{ch:build}.)
\item \textbf{Place each segment by where it wants to live.} For every segment
  the loader looks at its load address:
\begin{center}
\begin{tabular}{l l l}
\toprule
\textbf{Window} & \textbf{Range} & \textbf{Action} \\ \midrule
IRAM   & \texttt{16\#4037\_0000\#..16\#403E\_0000\#} & copy from flash into SRAM \\
DRAM   & \texttt{16\#3FC8\_0000\#..16\#3FD0\_0000\#} & copy from flash into SRAM \\
IROM   & \texttt{16\#4200\_0000\#..16\#4400\_0000\#} & map (execute-in-place via cache) \\
DROM   & \texttt{16\#3C00\_0000\#..16\#3E00\_0000\#} & map (read-in-place via cache) \\
\bottomrule
\end{tabular}
\end{center}
  RAM segments are copied byte-for-byte with the ROM SPI-flash read routine; the
  large code and read-only-data segments are left \emph{in} flash and reached
  through the cache.
\item \textbf{Map the flash windows through the cache MMU.} The loader programs
  the instruction-bus and data-bus MMUs (\texttt{Cache\_Ibus\_MMU\_Set} /
  \texttt{Cache\_Dbus\_MMU\_Set}, 64\,KB pages) so the IROM/DROM virtual
  addresses resolve to the right flash pages, then clears the cache ``shut'' bits
  to switch the buses on.
\item \textbf{Bring up PSRAM} (next section) and \textbf{jump} to the entry point.
\end{itemize}

The loader's own code and stack are placed high in SRAM, above the region the
application's segments occupy, so copying the app never overwrites the running
loader.

\section{PSRAM: octal bring-up}\index{PSRAM}

External PSRAM is the trickiest thing the bootloader does: bringing up octal PSRAM
reconfigures the very MSPI timing that flash execute-in-place depends on, so it has
to run from RAM, carefully, at exactly the right moment -- and the 2nd-stage
bootloader is that moment. The bring-up (a small C shim, \texttt{psram\_boot.c}) is
now almost entirely \emph{from source}; only the lowest-level OPI and cache pokes
still call masked-ROM helpers (\texttt{esp\_rom\_opiflash\_*}, \texttt{Cache\_*}).
Chapter~\ref{ch:psram} tells the de-blobbing story in full; in outline it does:

\begin{enumerate}
\item \textbf{Wire the octal pins.} \texttt{esp\_rom\_opiflash\_pin\_config} and
  the MSPI drive-strength setup connect the extra data lines and DQS that octal
  mode needs.
\item \textbf{Enable and size the PSRAM} via the from-source octal driver
  (\texttt{psram\_impl\_enable\_src} / \texttt{psram\_impl\_get\_physical\_size\_src}),
  reporting the detected size.
\item \textbf{Tune the data-in timing.} At 80\,MHz OPI-DDR the SPI0 data-in delay
  must be calibrated; \texttt{psram\_tune\_din} sweeps the candidate modes, measures
  which pass, and centres on the passing window (replacing an earlier hardcoded
  poke -- see Chapter~\ref{ch:psram}).
\item \textbf{Map it into the address space.} \texttt{Cache\_Dbus\_MMU\_Set} with
  the external-RAM flag maps the PSRAM at virtual base \texttt{16\#3D00\_0000\#}; the
  number of 64\,KB pages comes from the project's \texttt{board.ads}
  (\texttt{PSRAM\_Size / 64\,KB}).
\end{enumerate}

Because the bootloader maps PSRAM \emph{before} the application's \texttt{start.S}
sets up its stack, the application can even place the environment task's primary
stack in PSRAM (the \texttt{ENV\_STACK\_PSRAM} build option), trading a little
latency for a great deal of headroom on deep elaboration or recursion -- which
is exactly what the full-profile ACATS sweep needs.

\cnote{In ESP-IDF, \texttt{CONFIG\_SPIRAM=y} hides all of this. Here it is a few
dozen lines you can read: the pins, the enable call, the one mode-register fix,
and the MMU map. The size comes from \texttt{board.ads}, so a different module is
a one-line change, not a menuconfig safari.}

\section{\texttt{start.S}: from reset state to Ada}\index{start.S}

Once the loader jumps in, the application's \texttt{start.S}
(\texttt{examples/common/bare/start.S}) turns the raw post-loader state into
something Ada can run on:

\begin{itemize}
\item \textbf{Vectors and processor state.} It sets \texttt{VECBASE} to our own
  vector table (so any fault from here uses our handlers) and sets \texttt{PS} to
  the windowed-ABI state the runtime expects.
\item \textbf{Sane register windows.} It rewrites \texttt{WINDOWSTART} to a single
  frame at the current \texttt{WINDOWBASE}, discarding the loader's stale call
  frames so the first window overflow/underflow does not fault. (Getting this
  ordering right was the cure for an early environment-task double fault.)
\item \textbf{FPU on.} It sets \texttt{CPENABLE} so the first floating-point
  instruction does not trap.
\item \textbf{Stack and \texttt{.bss}.} It loads the stack pointer from the
  linker symbol and zeroes \texttt{.bss}.
\item \textbf{The instruction cache.} The loader enabled caching but left the
  I-cache unconfigured; \texttt{start.S} configures a 16\,KB, 8-way, 32-byte-line
  I-cache and sets the IROM/DROM MMU split, so instruction fetches from the IROM
  window actually return code rather than zeros.
\item \textbf{The clock.} This is the big one. Out of reset the CPU runs at
  \emph{20\,MHz} (the crystal divided by two); the runtime's timer math assumes
  240\,MHz, so short delays would run twelve times too long. \texttt{start.S}
  selects the 480\,MHz PLL divided to 240\,MHz (registers
  \texttt{16\#600C\_0010\#} and \texttt{16\#600C\_0060\#}) and tells the ROM timing
  helpers about the new frequency, all from IRAM before the jump into flash code.
\end{itemize}

It then jumps (not calls) into the C/Ada glue, which sets up the environment task
via \texttt{\_\_gnat\_enter\_env} -- a small shim that makes the env-task body the
\emph{outermost} window frame, so the first context switch has nothing stale
above it to corrupt. On the second core, a parallel cold-start entry
(\texttt{core1\_start}) repeats the window/VECBASE/stack normalisation before the
SMP runtime releases it.

\section{The memory map}\index{memory map}

All of this is pinned down by the linker scripts under
\texttt{examples/common/bare/vendor/} (the same regions appear in the memory-map
Table~\ref{tab:memmap}). The ESP32-S3's internal SRAM is visible at two
addresses -- an instruction-bus (IRAM) view and a data-bus (DRAM) view of the
\emph{same} bytes -- which is why code and data segments land in different
windows of one physical memory:

\begin{center}
\begin{tabular}{l l l}
\toprule
\textbf{Region} & \textbf{Base} & \textbf{What lives there} \\ \midrule
IRAM code     & \texttt{16\#4037\_8000\#} & vectors + hot code copied to SRAM \\
IROM (flash)  & \texttt{16\#4200\_0000\#} & bulk code, execute-in-place via cache \\
DRAM          & \texttt{16\#3FC9\_0000\#} & data, \texttt{.bss}, heap \\
DROM (flash)  & \texttt{16\#3C00\_0000\#} & read-only data via cache \\
PSRAM         & \texttt{16\#3D00\_0000\#} & external RAM (mapped by the loader) \\
RTC slow mem  & \texttt{16\#5000\_0000\#} & retained across deep sleep \\
\bottomrule
\end{tabular}
\end{center}

Two practical consequences fall out of this map and recur elsewhere in the book.
First, executable code must live in (or be reached through) an instruction-bus
window: that is why the full profile relocates GCC's nested-subprogram
trampolines to the IRAM alias rather than leaving them on the data-view stack
(\S\ref{ch:runtime-build}). Second, PSRAM is reachable only after the loader maps
it, and it cannot be a DMA source -- which is why DMA-capable scratch buffers are
carved from internal SRAM even when the task stack itself lives in PSRAM (the SD
and filesystem drivers rely on this).

\chapter{Bringing Up the Octal PSRAM}
\label{ch:psram}

Many ESP32-S3 modules package \emph{external} pseudo-static RAM (PSRAM) alongside
the SoC -- but whether there is any, how much, and of what kind depends on the
module. Espressif's S3 family ships in several variants: no PSRAM at all, 2\,MB of
\emph{quad} (4-line) PSRAM, or up to 8\,MB of \emph{octal} (8-line) PSRAM, and from
more than one memory vendor. This chapter is about the octal variant -- the
high-bandwidth option -- which the boards used here happen to populate with an
8\,MB AP~Memory die. The bring-up does not hard-code that, though: it reads the
device's own mode registers and adapts to whichever octal part is fitted (vendor,
density, voltage, latency), so the 8\,MB and the exact constants below are what
\emph{these} parts reported, not universal facts.

Why bother? Internal SRAM is about 292\,KB -- generous for a microcontroller, but
nowhere near enough for the megabyte-scale buffers, heaps, and elaboration stacks
that a full Ada program wants. PSRAM closes that gap, but only if you can turn it
on, and turning it on is one of the most delicate sequences in the whole bring-up:
a high-speed double-data-rate (DDR) octal link that has to be configured register
by register, and then \emph{timed}, before a single byte will round-trip.

This chapter is the deep version of the PSRAM paragraph in the bootloader chapter
(\S\ref{ch:bootloader}). There we noted that the second-stage loader brings PSRAM
up before the application runs; here we tell the whole story -- the deterministic
register script, the empirical timing problem at 80\,MHz, and the de-blobbing
effort that replaced the vendored ESP-IDF binaries with readable from-source code.
It is a genuine war story, and the registers and values below are the real ones,
captured from working silicon.

\section{What the octal PSRAM is, and why it is hard}\index{PSRAM}\index{octal SPI}\index{MSPI}

The PSRAM is an external DRAM die that speaks an \emph{octal} (8-data-line) variant
of the memory-SPI (MSPI) protocol, clocked DDR -- data on both clock edges. The
S3's MSPI controller can map it transparently through the data cache, so once it is
up the chip appears as ordinary memory at \texttt{0x3D000000} and Ada code reads and
writes it with no special API at all (\S\ref{ch:psram-validate}). That transparency
is the whole point, and also the whole danger: a load instruction that targets a
mis-timed PSRAM word does not raise an exception, it \emph{stalls the bus}.

Physically the link uses dedicated MSPI pads: eight data lines (the four ordinary
SPI data pins \texttt{SPID}/\texttt{SPIQ}/\texttt{SPIWP}/\texttt{SPIHD} plus
\texttt{SPID4}--\texttt{SPID7}), a data strobe \texttt{DQS}, the shared
\texttt{SPICLK}, and a second chip select \texttt{CS1} on GPIO\,26. These are the
same pins the hardware chapter (\S\ref{ch:hardware}) flags as off-limits: GPIO\,26
through 37 are wired to the in-package flash and PSRAM, and driving them as general
I/O corrupts the very bus the program is executing from. The PSRAM bring-up is the
code that legitimately owns those pins.

The reason 8 octal DDR data lines at 80\,MHz is hard -- and the reason this chapter
exists -- is timing. At that rate the round-trip from the controller's clock edge
out to the chip and back as data is a meaningful fraction of one bit period, so the
controller cannot just sample data on a fixed edge; it has to \emph{learn} when the
returning bits are valid. That learning step is the only part of the whole sequence
that is not a fixed register script, and it is where almost all the difficulty lives.

\section{The de-blobbing goal}\index{de-blobbing}\index{ESP-IDF}

A working PSRAM bring-up already existed: the vendored ESP-IDF v5.4.4 binaries
\texttt{esp\_psram\_impl\_octal.c.obj} and the \texttt{mspi\_timing\_*.c.obj} set,
linked as opaque blobs. The goal of this effort was to replace those blobs with
readable from-source C (plus a handful of ROM calls), so that nothing in the boot
path is a black box. That was done in two stages, retiring five blobs:

\begin{itemize}
\item \textbf{Stage 1} retired the four timing/GPIO blobs
(\texttt{mspi\_timing\_by\_mspi\_delay}, \texttt{mspi\_timing\_config},
\texttt{mspi\_timing\_tuning}, \texttt{gpio\_periph}), replacing them with
\texttt{mspi\_timing\_src.c}: the clock dividers (core 160\,MHz $\rightarrow$ PSRAM
80\,MHz), the din application, and the pad drive strength.
\item \textbf{Stage 2} retired the last blob,
\texttt{esp\_psram\_impl\_octal.c.obj}, replacing it with
\texttt{psram\_impl\_src.c}: the chip-side mode-register programming and the SPI0
cache-phase configuration.
\end{itemize}

What was \emph{kept} is the ROM: \texttt{esp\_rom\_opiflash\_*} (the octal command
helper and pad config) and \texttt{Cache\_*} (cache disable/enable and the d-bus MMU
map). These are masked into the chip, not vendored blobs, so they cost nothing to
keep and de-ROM-ing them buys little. The net result is the few-dozen-line bring-up
the bootloader chapter advertised.

\cnote{In ESP-IDF this is \texttt{esp\_psram\_impl\_octal} -- thousands of lines of
generated, config-gated code you never see. Here it is two short C files you can
read top to bottom: the deterministic register script, and one honest timing sweep.}

\section{The deterministic part: a fixed register script}\index{register script}\index{mode registers}

Most of the bring-up is a fixed sequence of register writes and ROM calls that is
identical on every boot and every board. The driving function is
\texttt{psram\_bringup()} in \texttt{psram\_boot.c}, which wraps the from-source
enable, maps the chip through the cache, and then tunes the timing:

\begin{lstlisting}[style=sh]
esp_rom_opiflash_pin_config()          # ROM: wire the octal pads SPID4-7 + DQS
mspi_timing_set_pin_drive_strength()   # from-src: FUN_DRV=3 on the 9 octal pads
psram_impl_enable_src()                # from-src: the chip bring-up (below)
psram_impl_get_physical_size_src(&sz)  # from-src: report decoded density -> 8 MB
Cache_Disable_DCache()                 # ROM
Cache_Dbus_MMU_Set(.., 0x3D000000, ..) # ROM: map PSRAM pages x 64 KB
Cache_Enable_DCache(0)                 # ROM
psram_tune_din()                       # from-src: the 80 MHz din sweep (below)
\end{lstlisting}

The pad wiring is \emph{FIX 1}: the ROM flash setup leaves only four of the eight
data lines connected because flash is quad. Without \texttt{SPID4}--\texttt{SPID7}
and \texttt{DQS}, the chip's mode-register read comes back with half its bits
floating, so the density field decodes to \texttt{0x5} (16\,MB) instead of the real
\texttt{0x3} (8\,MB) and the chip is mis-configured from the start.

Inside \texttt{psram\_impl\_enable\_src()} the order is load-bearing: controller
setup, drop to a slow clock, do the chip transactions, raise the clock, then the
cache phases. The deterministic steps are:

\begin{longtable}{p{2.4cm} p{4.2cm} p{5.4cm}}
\toprule
\textbf{Step} & \textbf{Register / call} & \textbf{What it does} \\ \midrule
\endhead
init pins      & \texttt{0x6000906C}=\texttt{0xF00}        & route GPIO\,26 to \texttt{SPICS1}, \texttt{FUN\_DRV}=3 \\
clk pad drive  & \texttt{0x600033FC}=\texttt{0x0210105F}   & SMEM \texttt{SPICLK} pad drive \\
CS timing      & \texttt{0x600030DC}=\texttt{0x0400B18F}   & \texttt{SMEM\_AC}: CS hold 3, setup 3, hold-delay 2 \\
ECC off        & clear \texttt{0x60026058} bit 8           & \texttt{SRAM\_ACE0\_ATTR}: disable inline ECC \\
low speed      & \texttt{enter\_low\_speed\_mode}          & core 80\,MHz, flash+PSRAM /4 = 20\,MHz, clear din \\
var-dummy      & \texttt{0x600020E0} \verb#|=# 2           & SPI1 \texttt{SPI\_FMEM\_VAR\_DUMMY}; no DTR swap \\
write MR0      & \texttt{exec\_cmd} \texttt{0xC0C0}        & \texttt{lt}=1 (fixed), \texttt{read\_latency}=2, \texttt{drive\_str}=0 $\rightarrow$ \texttt{0x28} \\
connect probe  & \texttt{exec\_cmd} \texttt{0x8080}/\texttt{0x0000} & write+read-back \texttt{0x5a6b7c8d} at addr 0 \\
read MR0--8    & \texttt{exec\_cmd} \texttt{0x4040}        & decode vendor / density / Vcc / latency \\
\textbf{tune din} & \texttt{psram\_tune\_din}              & \textbf{empirical 80\,MHz sweep -- the non-deterministic step} \\
high speed     & \texttt{enter\_high\_speed\_mode}         & core 160\,MHz, /2 = 80\,MHz, apply din \\
cache phases   & 5 SPI0 regs (below)                       & program transparent OPI-DDR cache access \\
\bottomrule
\end{longtable}

A few of these deserve a closer look. The MR0 write is a read-modify-write of the
low six bits to \texttt{0x28}: \texttt{lt=1} selects \emph{fixed} read latency (so
the controller knows exactly how many dummy cycles to insert), and
\texttt{read\_latency=2} encodes a 10-cycle fixed latency. The connectivity probe is
the only ``is it there?'' test in the octal path -- there is no JEDEC \texttt{0x9F}
read-ID -- so it writes the witness word \texttt{0x5a6b7c8d} to address 0 and reads
it back; identity comes entirely from the mode registers. On a genuine
ESP32-S3-DevKitC the decode prints:

\begin{lstlisting}[style=sh]
PSRAM: AP Memory octal DDR @80MHz, 64 Mbit (8 MB), dev gen 3, Vcc 3.0V
  latency: read 10-cyc (fixed), write 5-cyc;  MR0=28 MR1=0d MR2=93 MR3=60 MR4=40 MR8=05
\end{lstlisting}

All the mode-register traffic -- the MR write, the probe, the MR reads -- goes
through the ROM helper \texttt{esp\_rom\_opiflash\_exec\_cmd} in OPI-DTR mode
(\texttt{mode 7}) with a 16-bit command, 32-bit address, and the PSRAM chip select
\texttt{CS=BIT(1)}. Read command is \texttt{0x4040}, write \texttt{0xC0C0}, sync read
\texttt{0x0000}, sync write \texttt{0x8080}; read dummy 18 half-cycles, write dummy 8.

\subsection{The golden cache-phase registers}\index{cache MMU}

The last deterministic step configures the SPI0 (cache) controller so that a plain
load from \texttt{0x3D000000} silently becomes an octal-DDR PSRAM read. These five
register values were captured from the working blob over JTAG -- the ``golden''
state -- and are written directly. They encode the same OPI-DDR geometry as the
manual transactions above, but for the transparent cache path:

\begin{longtable}{p{2.6cm} p{2.4cm} p{6.8cm}}
\toprule
\textbf{Register} & \textbf{Value} & \textbf{Meaning} \\ \midrule
\endhead
\texttt{SRAM\_CMD}  & \texttt{0x007C0000} & octal command/address/data lanes \\
\texttt{SRAM\_DRD\_CMD} & \texttt{0xF0000000} & cache read command \texttt{0x0000}, 16-bit \\
\texttt{SRAM\_DWR\_CMD} & \texttt{0xF0008080} & cache write command \texttt{0x8080}, 16-bit \\
\texttt{SMEM\_DDR}  & \texttt{0x00003023} & DDR enable + variable dummy \\
\texttt{CACHE\_SCTRL} & \texttt{0x01F7C479} & 32-bit addr, read dummy 18 / write 8, octal lanes \\
\bottomrule
\end{longtable}

\texttt{CACHE\_SCTRL} carries the lane bits that make every phase octal
(\texttt{SCMD\_OCT}, \texttt{SADDR\_OCT}, \texttt{SDOUT\_OCT}, \texttt{SDIN\_OCT},
\texttt{SRAM\_OCT}) and the \texttt{SCMD\_4BYTE} flag for the 32-bit address. After
these five writes, \texttt{Cache\_Resume\_DCache(0)} hands the cache back and the
chip is, deterministically, configured. Everything so far is reproducible from
source. What remains is the one step that is not.

\section{The non-deterministic part: timing tuning}\index{din timing}

At 80\,MHz DDR the controller must choose \emph{when} to sample the returning data.
The S3 exposes three per-data-line knobs for this, all in the SPI0 (cache) MSPI block
at \texttt{0x6000\_30xx}:

\begin{itemize}
\item \texttt{din\_mode} -- \texttt{SMEM\_DIN[0..7]\_MODE} (the register at
\texttt{0x600030C0}): the sampling edge/phase, 3 bits per lane, 8 settings.
\item \texttt{din\_num} -- \texttt{SMEM\_DIN[0..7]\_NUM} at \texttt{0x600030C4}: an
integer-cycle delay added to data-in.
\item \texttt{extra\_dummy} -- \texttt{SMEM\_TIMING\_CALI} at \texttt{0x600030BC}:
extra dummy cycles before the data window opens.
\end{itemize}

Of these, \texttt{din\_mode} is the one that matters here, and the one whose correct
value is not knowable from source -- it depends on the actual chip, controller, and
temperature. Picking it wrong does not fail gracefully: a load through the cache at a
bad \texttt{din\_mode} stalls the bus with no exception, the PC pins, and no interrupt
can preempt it. So this knob has to be chosen carefully, and chosen safely.

\section{The tuning war story}

\subsection{Why IDF's own tuning is a no-op}\index{din\_mode}

The natural assumption is that the vendor blob already calibrates \texttt{din\_mode}.
It does not. Single-stepping \texttt{mspi\_timing\_psram\_tuning} on live silicon
showed why: the IDF sweep drops the MSPI to \emph{20\,MHz} before testing and only
restores 80\,MHz afterwards. At 20\,MHz the data eye is enormous, so the sampling
phase is irrelevant -- every one of the 14 candidate configurations reads the
reference pattern back byte-perfectly. The DTR selection heuristic sees an
implausibly wide all-pass window (run length $\ge$ 6) and \emph{distrusts} it, falling
back to \texttt{default\_config\_id = 5}, which is \texttt{din\_mode = 4}.

\cnote{The captured ``calibration'' is really a vendor-default lookup. In the blob,
\texttt{din\_mode = 4} is dead-centre of the \emph{failing} region at 80\,MHz -- it
crashes the cache read. That is why the old shim carried a hand-coded ``FIX 2''
forcing \texttt{din\_mode = 1}: the only way to make the blob's chosen default usable
was to override it.}

So the working configuration was, historically, ``vendor default plus a magic
constant.'' De-blobbing meant either reproducing that magic constant honestly, or
measuring the right value for real.

\subsection{Why you cannot naively sweep the cache path}

The obvious fix is to sweep \texttt{din\_mode} yourself: set each mode, read a known
pattern back through the cache, keep the modes that pass. This was tried, and the
hardware fights back at every turn:

\begin{enumerate}
\item \textbf{A wrong cache din hangs the CPU.} A synchronous PSRAM load with a bad
din never completes -- the bus stalls, the PC pins, \texttt{exccause=0}, and no
interrupt can preempt a stalled load. You cannot ``read it and check if it is
garbage,'' because the read does not return.
\item \textbf{Async preload dodges the hang but poisons the controller.} Driving the
sweep with \texttt{Cache\_Start\_DCache\_Preload} plus a polled timeout lets the CPU
survive (it only polls a flag), but a stalled preload leaves the MSPI/cache
controller wedged: subsequent cache reads return the fault-fill pattern
\texttt{0xbad00bad}, and \emph{even \texttt{Cache\_Disable\_DCache} hangs}. Nothing
short of an MSPI peripheral reset clears it -- which also kills flash execute-in-place.
\item \textbf{SPI1 has its own din registers.} SPI1 (the manual-transaction path) keeps
separate din registers at \texttt{0x6000\_20xx}, distinct from the cache controller's
\texttt{0x6000\_30xx}. A bounded SPI1 read calibrates the wrong register -- it tells
you nothing directly about the cache path.
\end{enumerate}

The conclusion is uncomfortable but precise: the din-sensitive path \emph{is} the
cache path, and the cache path cannot be safely probed because any bad value risks an
unrecoverable wedge.

\subsection{The fix: a bounded SPI1 sweep at the real speed}

\texttt{psram\_tune\_din()} threads the needle. It sweeps \texttt{din\_mode} at the
real \emph{80\,MHz} operating clock, but it does the readback over a \emph{bounded SPI1
manual transaction} (\texttt{esp\_rom\_opiflash\_exec\_cmd}), not through the cache. A
wrong din over SPI1 returns garbage and the transaction \emph{still completes} -- so
the sweep can measure all eight modes without ever stalling the bus. SPI0 and SPI1 share
the same MSPI clock and pads, so the window measured over SPI1 is, to a known offset,
the window the cache wants. The core of it:

\begin{lstlisting}[style=sh]
# write a reference word (data-out: din-independent), then sweep the read din:
exec_cmd(spi1, OPI_DTR, 0x8080, addr, dummy=8,  &ref)      # reference write
for m in 0..7:
    REG(0x600030C0) = din_word(m)   # SPI0 (cache) din-mode, 3 bits x 9 lanes
    REG(0x600020C0) = din_word(m)   # SPI1 (manual) din-mode
    exec_cmd(spi1, OPI_DTR, 0x0000, addr, dummy=18, &rd)   # bounded read
    if rd == ref: pass |= 1 << m
\end{lstlisting}

The measured window is \texttt{passmask = 0xC3} -- modes \{6, 7, 0, 1\} pass, modes
\{2, 3, 4, 5\} fail. Crucially this is a \emph{circular} arc: \texttt{din\_mode} is a
sampling phase, so mode 7 is adjacent to mode 0, and the passing run wraps the
$7\rightarrow0$ seam. The selection logic walks the modes twice (two laps) to catch a
window that wraps, and picks the centre of the widest passing run:

\begin{lstlisting}[style=sh]
mode 0,1 PASS | mode 2,3 fail (0x..76) | mode 4,5 fail (0x..e2) | mode 6,7 PASS
                                                       passmask 0xC3  ->  centre = mode 0
\end{lstlisting}

Notice that IDF's default \texttt{din\_mode = 4} sits dead-centre in the \emph{failing}
region -- the 80\,MHz sweep correctly rejects exactly what the 20\,MHz tune handed
back. The measured centre (mode 0) has more margin than FIX 2's mode 1, which was a
valid but edge-of-eye point. If the sweep is degenerate (empty or all-pass) the code
falls back to mode 1. (One subtlety, recorded in the source: the SPI1-measured centre
can sit slightly off the cache path's true optimum, and that offset is invisible to a
quick checksum but can corrupt a heap-heavy app under sustained access; so when mode 1
-- the validated-robust cache point for this chip -- lies inside the measured window,
the code prefers it.)

Because din timing affects \emph{reads} only and the bootloader itself runs from
flash/IRAM rather than PSRAM, retuning the PSRAM din live, after the chip is mapped, is
safe -- nothing the bootloader is executing depends on it.

\subsection{Cross-validation on two boards}

The strongest evidence that this measures a real property and not a per-board fluke is
that two physically different boards land on the \emph{same} window. The whole effort
was developed on a Meshnology LoRa AIOT board; the finished from-source bring-up was
then run unchanged on a genuine Espressif ESP32-S3-DevKitC. Both independently measured
\texttt{passmask = 0xC3} $\rightarrow$ mode 0, on the same AP Memory 8\,MB octal die.
The din window tracks a stable chip-plus-controller property, which is exactly what you
want from a calibration: reproducible, not lucky.

\section{The 80\,MHz versus 40\,MHz trade}

It is worth being explicit about why this complexity exists at all because there is a
simpler road not taken. At 40\,MHz the MSPI does \emph{not} need tuning: below the
40\,MHz threshold the data eye is wide enough that the sampling phase is irrelevant,
din/num/dummy are all zero, and the entire bring-up collapses into a straight register
script with no empirical step -- every value a compile-time constant. That is the
clean, fully-deterministic option.

The cost is bandwidth: 40\,MHz DDR is half the throughput of 80\,MHz DDR. This build
chose 80\,MHz, and therefore inherited the tuning problem because the PSRAM here is
used for bandwidth-sensitive work, not just capacity. The 40\,MHz script remains a
documented fallback -- the right choice if a future board needs PSRAM only as a large,
slow backing store and would rather not carry a sweep at all.

\section{On-target validation}\index{external RAM}
\label{ch:psram-validate}

None of this matters unless an Ada program can actually use the result, so the
\texttt{esp32s3\_psram} example exists to prove it does. It places a 1\,MB static byte
array in external RAM -- no heap, no \texttt{malloc} -- using a linker section that
\texttt{psram.ld} maps onto the PSRAM region:

\begin{lstlisting}
Buffer : array (0 .. 1024 * 1024 - 1) of Unsigned_8
  with Linker_Section => ".ext_ram.bss";
\end{lstlisting}

\texttt{Big.Run} fills the buffer with a known pattern (byte $i$ gets $i \bmod 256$),
reads every byte back, and checksums the lot. Because the pattern is bytes
$0\ldots255$ repeating across exactly $2^{20}$ bytes, the sum is a fixed,
hand-computable constant: \texttt{checksum = 0x07F80000}. Any timing error anywhere in
the 1\,MB round-trip through the cache would perturb it. On hardware:

\begin{lstlisting}[style=sh]
[ada-free-boot] octal PSRAM up: rc=0  8 MB
[ada-free-boot] PSRAM mapped @0x3D000000 rc=0
[ada-free-boot] PSRAM din tuned @80MHz: passmask=0xc3 -> mode 0
[psram] buffer @ 0x3d000000  1048576 bytes  checksum=0x07f80000  (PSRAM)
\end{lstlisting}

The buffer's \texttt{0x3D000000} address confirms it really lives in the external-RAM
range, and the deterministic checksum -- byte-identical across resets and across both
test boards -- confirms the full megabyte round-trips correctly through a cache served
by a din the bootloader measured for itself. There is no FreeRTOS, no ESP-IDF, and no
remaining blob in the path: the entire octal-PSRAM bring-up is now readable source plus
a few ROM calls, which is exactly what de-blobbing set out to achieve.

One caveat the example notes for completeness: task stacks stay in internal SRAM. PSRAM
is unreachable whenever the cache is disabled -- for instance during a flash write --
so it is the right home for large data buffers and heaps, but not for a stack that an
interrupt might unwind while the cache is down (\S\ref{ch:bootloader}).

\chapter{The Heap: \texttt{malloc}/\texttt{free} in Ada, the TLSF Way}
\label{ch:heap}

Two of the runtime profiles give you a real dynamic heap, and that heap is now
written entirely in Ada. Both Ada's \texttt{new} and the C library's
\texttt{malloc} land in the same allocator, which is a \emph{Two-Level
Segregated Fit} (TLSF)\footnote{TLSF is the algorithm of M.~Masmano, I.~Ripoll,
A.~Crespo and J.~Real, ``TLSF: A New Dynamic Memory Allocator for Real-Time
Systems,'' \emph{16th Euromicro Conference on Real-Time Systems (ECRTS)}, 2004
(\url{http://www.gii.upv.es/tlsf/}). The implementation here is original Ada; the
design is theirs.} implementation: O(1) allocate, O(1) free-and-coalesce,
and bounded fragmentation. This chapter is the current truth about the heap. The
``first-fit'' descriptions in \S\ref{ch:profiles} and the appendix
(\S\ref{ch:appendix}) describe the allocator that TLSF replaced; what follows
supersedes them.

\section{Why a heap, and what lands here}\index{heap}\index{dynamic allocation}\index{secondary stack}

Ada programs reach the heap by two doors. The first is the language itself: every
\texttt{new T} (and every access-type allocator, and the storage pool behind
unconstrained returns when no other pool is named) ultimately calls the runtime's
allocation primitive, which is wired to \texttt{\_\_gnat\_malloc} / the C
\texttt{malloc} symbol. The second door is C code linked into the image (a
vendor blob, a bit of glue, anything that calls \texttt{malloc}, \texttt{calloc},
\texttt{realloc} or \texttt{free}). Both doors open onto the single allocator
described here.

What does \emph{not} come here is the secondary stack. Ada's secondary stack,
which carries function results whose size is not known to the caller, such as an
unconstrained \texttt{String} return, is a separate, contiguously-managed
region with its own mark/release discipline; it never touches \texttt{malloc}.
Keeping the two apart is what lets the leanest profile run with no heap at all.

That leanest profile is \texttt{light-tasking}: it is configured
\texttt{No\_Secondary\_Stack} and heap-less, with only a bump allocator that can
hand out memory but never reclaim it (see \S\ref{ch:profiles}). A program built
\texttt{light-tasking} that calls \texttt{new} in a loop will eventually run the
bump pointer off the end of the arena, by design because that profile is for
code that is statically sized and never frees. The \texttt{embedded} and
\texttt{full} profiles are the ones that link the real allocator, so on those
\texttt{free} actually returns memory to the pool and a long-running program can
allocate and release indefinitely.

\section{Why first-fit had to go}\index{first-fit allocator}\index{fragmentation}

The allocator TLSF replaced kept a single linked list of free blocks. To satisfy
a request it walked that list from the head looking for the first block large
enough: O(n) in the number of free blocks. Worse, to keep fragmentation in
check it coalesced on every \texttt{free} by scanning the whole list to find the
freed block's physical neighbours, again O(n). For a small internal-RAM heap
with a handful of live blocks, nobody notices. The problem is the motivating case
for this whole chapter: a \emph{large, slow} heap.

The ESP32-S3 can map up to 8\,MiB of external PSRAM, and that PSRAM is an order of
magnitude slower than on-chip SRAM: it runs as DDR with a wait state, so every
header you touch is an off-chip burst. Put a heap there (the
\texttt{HEAP\_PSRAM=1} mode of the on-target test does exactly this) and a list
that may hold thousands of free blocks turns each allocate \emph{and} each free
into a long off-chip walk. The cost is not just big; it is unpredictable, which is
exactly what you do not want on a real-time runtime. TLSF makes both operations
O(1) with a small, fixed number of memory touches, so the heap behaves the same
whether it lives in fast SRAM or slow PSRAM.

\cnote{ESP-IDF answers the slow-vs-fast question with \texttt{heap\_caps\_malloc}
and a set of capability flags (\texttt{MALLOC\_CAP\_INTERNAL},
\texttt{MALLOC\_CAP\_SPIRAM}, \ldots) routing to several heaps. Here there is one
allocator and one arena; \emph{which} arena (leftover DRAM or PSRAM) is chosen
once, at link time, in \texttt{bare\_build.sh}.}

\section{TLSF, taught from the code}\index{TLSF allocator}

The portable allocator lives in \texttt{tlsf\_core.ads}/\texttt{.adb}. It is
deliberately target-free: no lock, no assembly, no linker symbols. It takes an
arena from its caller and manages it. That purity is what lets the identical code
run on the host (\S\ref{sec:heap-validation}). Here is the whole interface:

\begin{lstlisting}
procedure Init (Base : System.Address; Size : Storage_Count);
function  Allocate (N : Storage_Count) return System.Address;
procedure Deallocate (P : System.Address);
function  Reallocate (P : System.Address; N : Storage_Count)
                      return System.Address;
function  Ready return Boolean;
function  Invariants_Hold return Boolean;   --  test-only
\end{lstlisting}

\subsection{Boundary tags: every block knows its neighbours}\index{boundary tags}

The arena is one contiguous chain of physically-adjacent blocks. Each block
carries a header before its payload:

\begin{lstlisting}
type Hdr is record
   Prev_Phys : System.Address;   --  physical previous block (Null = first)
   Size      : Storage_Count;    --  payload bytes (16-aligned)
   Free      : Boolean;
end record;
\end{lstlisting}

Because a block knows both its own \texttt{Size} and the address of its physical
predecessor (\texttt{Prev\_Phys}), the runtime can step forward to the next block
(\texttt{Next\_Phys = B + Hdr\_Sz + Size}) and back to the previous one in
constant time, without any list walk. That is the \emph{boundary tag} idea, and
it is what makes \texttt{free} O(1): on a free we look only at the block itself
and its two immediate physical neighbours, coalescing with whichever of them is
also free.

A header-size subtlety worth flagging because it bit this code: \texttt{Hdr\_Sz}
is a compile-time \emph{static} constant, not derived from
\texttt{Hdr'Object\_Size}. An object-size-based constant would be \emph{elaborated}
into \texttt{.bss}, and the ZFP boot has no \texttt{adainit} to run that
elaboration, so the constant would read as garbage and produce wild pointers.
For the same reason all the mutable state (\texttt{Heads}, the bitmaps,
\texttt{Pool\_Lo}, \texttt{Sentinel}) has zero initializers, which places it in
\texttt{.bss} (zeroed by \texttt{start.S}) rather than \texttt{.data}.

\subsection{Two-level segregated free lists}\index{segregated free lists}

Free blocks are bucketed by size into a two-dimensional grid of lists. The
\emph{first level} (FL) is a power-of-two class: roughly, $\mathrm{FL} =
\lfloor\log_2 \mathrm{size}\rfloor$. The \emph{second level} (SL) subdivides each
power-of-two range linearly into \texttt{SL\_Count} equal sub-ranges. Here that is
$2^5 = 32$ second-level lists per first-level class, across 25 first-level
classes:

\begin{lstlisting}
SL_Log2  : constant := 5;
SL_Count : constant := 2 ** SL_Log2;   --  32 second-level lists
FL_Count : constant := 25;             --  first-level classes
Heads : array (0 .. FL_Count - 1, 0 .. SL_Count - 1) of System.Address;
\end{lstlisting}

\texttt{Heads (FL, SL)} is the head of a doubly-linked list of free blocks whose
sizes all fall in that one bucket. A free block stores its list links in the first
two words of its \emph{payload} (it is free, so the payload is unused),
\texttt{Next\_Free} / \texttt{Prev\_Free}, which is why the minimum payload is
two words wide.

The \texttt{Mapping} function turns a size into its \texttt{(FL, SL)} indices.
Above the small-block threshold it finds the high bit with a find-last-set
(\texttt{FLS}), uses that for FL, and takes the next \texttt{SL\_Log2} bits below
it for SL:

\begin{lstlisting}
F  : constant Integer := FLS (S);                         --  high bit
SL := Integer (Shift_Right (S, F - SL_Log2) and (SL_Count - 1));
FL := F - FL_Shift + 1;
\end{lstlisting}

A worked example with \texttt{FL\_Shift = 9} (which is \texttt{SL\_Log2 + 4}): a
request rounding to 1536 bytes is \texttt{0b110\_0000\_0000}, whose top bit
(\texttt{FLS}) is at position 10. \texttt{SL} takes the five bits starting just
below that bit: $(1536 \gg (10-5)) \mathbin{\mathrm{and}} 31 = 48
\mathbin{\mathrm{and}} 31 = 16$, and $\mathrm{FL} = 10-9+1 = 2$. So 1536-byte
blocks live in bucket \texttt{Heads (2, 16)}. A 2048-byte block
(\texttt{0b1000\_0000\_0000}, top bit 11) lands in \texttt{Heads (3, 0)}, a
different first-level class, exactly as you want at a power-of-two boundary.

\subsection{Bitmaps + find-first-set: O(1) bucket lookup}

The grid would still need scanning if we had to search it. We do not because a
pair of bitmaps records which buckets are non-empty:

\begin{lstlisting}
FL_Bitmap : Unsigned_32;                       --  bit FL set if class FL has any free block
SL_Bitmap : array (0 .. FL_Count - 1) of Unsigned_32;  --  per-class SL occupancy
\end{lstlisting}

To allocate, \texttt{Mapping\_Search} first rounds the request \emph{up} so that
\emph{any} block in the chosen bucket is guaranteed big enough (this is what makes
it ``good-fit'' rather than risking a too-small first block). Then
\texttt{Find\_Suitable} masks off the buckets at or below the requested class and
uses find-first-set (\texttt{FFS}) to jump straight to the smallest non-empty
bucket that fits: first within the same FL class, then, if that class is
exhausted, by consulting \texttt{FL\_Bitmap} for the next non-empty class up. Two
masked bit-scans, no list walk. That is the O(1) \texttt{malloc}.

\texttt{FFS} and \texttt{FLS} here are written as bounded shift loops rather than
a hardware count-leading-zeros instruction. The loops run at most 32 (or 64)
iterations regardless of heap size, so they are O(1) in the heap and keep the core
free of target intrinsics.

\begin{figure}[ht]
\centering
\resizebox{0.92\linewidth}{!}{%
\begin{tikzpicture}[node distance=6mm]
\node[blk] (req) {request\\\scriptsize round up};
\node[blk, right=of req] (map) {\texttt{Mapping}\\\scriptsize size $\to$ (FL,SL)};
\node[blk, right=of map] (ffs) {bitmaps + FFS\\\scriptsize smallest fitting bucket};
\node[blk, right=of ffs] (pop) {pop head\\\scriptsize of \texttt{Heads(FL,SL)}};
\node[blk, right=of pop] (spl) {split?\\\scriptsize remainder $\to$ free};
\draw[flow] (req)--(map); \draw[flow] (map)--(ffs);
\draw[flow] (ffs)--(pop); \draw[flow] (pop)--(spl);
\end{tikzpicture}}
\caption{One \texttt{Allocate}: map the size, find the bucket by bit-scan, take
the head block, and split off any large remainder back into a free bucket. No
step depends on the number of free blocks.}
\label{fig:tlsf-alloc}
\end{figure}

\subsection{Allocate, free, split, coalesce}\index{coalescing}

\texttt{Allocate} rounds the request to a 16-byte multiple (and up to the
two-word minimum payload), finds a bucket, pops the head block off it, and,
if the block is substantially larger than needed, \texttt{Split}s it: the
front becomes the returned block and the tail is re-inserted as a smaller free
block (updating both its and its successor's \texttt{Prev\_Phys}). The leftover
is dropped back into whatever bucket its new size maps to.

\texttt{Deallocate} is the mirror image and the place boundary tags earn their
keep. It marks the block free, then coalesces \emph{forward} (if
\texttt{Next\_Phys} is free, remove it from its bucket and absorb it) and
\emph{backward} (if \texttt{Prev\_Phys} is free, absorb into it). The terminating
\texttt{Sentinel} block (a zero-size, permanently-\emph{used} block at the top
of the arena) is what makes the forward check safe without a bounds test: the
real top block always has a non-free successor to stop at. After coalescing, the
single resulting block is inserted into its bucket and the relevant bitmap bits
are set. Every one of these steps touches a fixed, small number of blocks.

\texttt{Reallocate} shrinks in place when the existing block is already big
enough, and otherwise allocates-copies-frees. \texttt{Init} lays the arena out as
exactly one big free block plus the sentinel.

\section{The target glue: \texttt{bare\_heap}}\index{bare\_heap}\index{malloc}\index{free}

\texttt{tlsf\_core} has no idea it is on a microcontroller. Everything
target-specific lives in \texttt{bare\_heap.ads}/\texttt{.adb}, which is the unit
that actually exports the C ABI. It replaces the old \texttt{bare\_heap.c}.

\begin{lstlisting}
function  Malloc  (N : Interfaces.C.size_t) return System.Address;
pragma Export (C, Malloc, "malloc");

procedure Free    (P : System.Address);
pragma Export (C, Free, "free");

function  Realloc (P : System.Address; N : Interfaces.C.size_t)
                   return System.Address;
pragma Export (C, Realloc, "realloc");

function  Calloc  (Nmemb, Size : Interfaces.C.size_t) return System.Address;
pragma Export (C, Calloc, "calloc");

procedure Task_Stack_Free (Stack, Thread : System.Address);
pragma Export (C, Task_Stack_Free, "__gnat_task_stack_free");
\end{lstlisting}

\subsection{The critical section}\index{critical section}

\texttt{tlsf\_core} is not internally locked, by design, so the host harness
can drive it single-threaded. On target, multiple agents reach the heap: the
environment task, finalizers running RAII cleanup, and any task calling
\texttt{new}. \texttt{bare\_heap} serialises them with a global critical section
built from the Xtensa \texttt{rsil} instruction at level 15 (raise interrupt
level to the maximum, returning the old PS) around each operation, restored with
\texttt{wsr.ps}\,/\,\texttt{rsync}:

\begin{lstlisting}
function Enter_Crit return Unsigned_32 is
   Ps : Unsigned_32;
begin
   Asm ("rsil %0, 15",
        Outputs  => Unsigned_32'Asm_Output ("=r", Ps),
        Volatile => True, Clobber => "memory");
   return Ps;
end Enter_Crit;
\end{lstlisting}

This matches what the old C heap did. \texttt{Calloc} is careful to zero the
returned block \emph{outside} the lock. It only needs the lock to allocate, not
to memset, which keeps the critical section short.

\subsection{An arena chosen at link time}

\texttt{bare\_heap} never names an address. It imports two linker symbols whose
\emph{addresses} are the arena bounds:

\begin{lstlisting}
Heap_Base : Storage_Element
  with Import, Convention => C, External_Name => "__bare_heap_base";
Heap_End  : Storage_Element
  with Import, Convention => C, External_Name => "__bare_heap_end";
\end{lstlisting}

\texttt{bare\_build.sh} supplies these via \texttt{-{}-defsym}, pointing them
either at the leftover internal DRAM after \texttt{.bss}, or (for the large
heap) at the PSRAM region. The Ada is identical either way; only the two
\texttt{-{}-defsym} values differ. The first \texttt{malloc} lazily calls
\texttt{T.Init} over whatever region those symbols delimit (a \texttt{Started}
flag, also \texttt{.bss}-resident, gates this).

\subsection{Freeing a task's own stack}\index{task stack reclamation}

The full profile allocates each task's primary stack on the heap, which raises a
nasty question: when a task terminates, who frees the stack it was \emph{running
on}, and how do we avoid pulling it out from under a context that has not yet
switched away, possibly on the other core? \texttt{Task\_Stack\_Free}
(exported as \texttt{\_\_gnat\_task\_stack\_free}) handles this with a bounded
spin against the runtime's per-core running-thread table:

\begin{lstlisting}
Running : array (0 .. 1) of System.Address
  with Import, Convention => C, External_Name => "__gnat_running_thread_table";

procedure Task_Stack_Free (Stack, Thread : System.Address) is
begin
   for I in 1 .. 2_000_000 loop
      if Running (0) /= Thread and then Running (1) /= Thread then
         Free (Stack);
         return;
      end if;
   end loop;
   --  timed out -- leak rather than risk a use-after-free
end Task_Stack_Free;
\end{lstlisting}

It frees only once the terminating thread is no longer the running thread on
either core. If the thread somehow never clears (the spin bound expires), it
chooses to \emph{leak} the stack rather than risk a use-after-free: a deliberate
``leak-on-timeout'' policy.

\section{Validation}
\label{sec:heap-validation}\index{heap validation}

\texttt{tlsf\_core} is tested off-chip as ordinary Ada because it is
target-free.
\texttt{heap\_stress.adb} is a generic that takes an allocator's six operations
and beats on them: an initial sanity pass (alloc/align/free, OOM returns null,
realloc-grow preserves data), then 300\,000 randomised alloc/free steps over a
64\,KiB arena, each allocation stamped with a unique byte pattern that is
re-verified before every free so any overlap or stale pointer is caught. After a
full drain it asserts the arena coalesced back to (nearly) one block: proof the
free path actually merges. \texttt{allocator\_test.adb} instantiates that harness
over \emph{both} the old first-fit core and \texttt{Tlsf\_Core}, so a behavioural
divergence between them would show up.

The structural check throughout is \texttt{Invariants\_Hold}, run after every
operation. It walks the physical chain from \texttt{Pool\_Lo} to the
\texttt{Sentinel} and verifies: each block is 16-aligned, each
\texttt{Prev\_Phys} points back to the actual previous block, sizes are at least
the minimum and stay within the arena, and (the key correctness property) no
two physically-adjacent blocks are both free (i.e.\ coalescing is complete). The
sentinel must be zero-size, used, and correctly back-linked.

On hardware, \texttt{examples/esp32s3\_heaptest} drives the live
\texttt{malloc}/\texttt{free} symbols directly over the real arena: 50\,000 random
operations, 32 live blocks, the same pattern-stamping idea, reporting
\texttt{PASS}/\texttt{FAIL} over \texttt{ESP32S3.Log}. Its \texttt{HEAP\_PSRAM=1}
build relocates the arena into the 8\,MiB PSRAM at \texttt{0x3D000000}: the
large, slow heap that motivated TLSF in the first place. See \S\ref{ch:hal} for
how PSRAM is brought up before the heap can use it.

\section{Which profiles get it}

\texttt{light-tasking} links neither \texttt{bare\_heap} nor \texttt{tlsf\_core};
it is heap-less and bump-only (\S\ref{ch:profiles}). \texttt{embedded} and
\texttt{full} link both, so \texttt{malloc}/\texttt{free}/\texttt{realloc}/%
\texttt{calloc} (and therefore Ada's \texttt{new}) are backed by TLSF; the
full profile additionally uses \texttt{Task\_Stack\_Free} for per-task primary-%
stack reclamation. And because the same \texttt{tlsf\_core} compiles and runs
host-native, the allocator the chip runs is the very allocator the host test
exercises, not a model of it.

\chapter{GNARL and GNULL: The Two-Layer Runtime}
\label{ch:gnarl-gnull}

When you write \texttt{task}, \texttt{protected}, \texttt{delay until}, or an entry
call in Ada, the compiler does not emit instructions to suspend a thread or to
read a hardware timer. It emits \emph{calls} into the GNAT tasking runtime, and
that runtime is what actually talks to the machine. This chapter is about the
internal architecture of that runtime: two stacked layers named \textbf{GNARL}
and \textbf{GNULL}, and how this bare-metal port wires the bottom of the stack
onto the ESP32-S3 instead of onto an operating system.

The user-facing side --- what the language constructs mean and how you use them
--- is covered in \S\ref{ch:tasking}. Here we look underneath: who is called,
what each layer is responsible for, and which real source units implement it.

\section{Two names, two layers}

The runtime that implements Ada tasking is split into two libraries with two
deliberately retro acronyms.

\begin{itemize}
  \item \textbf{GNARL} --- the \emph{GNat Ada Run-time Library}. This is the
        \emph{upper} layer. It knows Ada semantics: what a task control block is,
        how rendezvous and protected-object queues work, how a task is created,
        activated, and terminated, how priorities and the ceiling protocol
        behave. It is portable: it is written in terms of an abstract notion of
        ``a thread'' and ``the clock,'' and it does not name any specific
        hardware or OS.
  \item \textbf{GNULL} --- the \emph{GNu Low-Level (Library)}. This is the
        \emph{lower} layer. Its one job is to satisfy GNARL's abstract needs ---
        create a thread, find the current thread, set a priority, sleep, wake,
        delay until an absolute time, read a monotonic clock --- by mapping each
        of them onto the services of the underlying platform.
\end{itemize}

The split exists so that GNARL can be written once and ported many times: only
GNULL changes per target. On a hosted GNAT (Linux, Windows) the ``underlying
platform'' GNULL targets is the OS thread library --- pthreads. \textbf{The
single most important idea in this chapter is that here there is no OS, so the
platform GNULL targets is the bare-board kernel} shipped with the bb-runtimes:
the \texttt{System.BB.*} packages. Every place a hosted runtime would call
\texttt{pthread\_create} or \texttt{clock\_gettime}, this port calls into
\texttt{System.BB.Threads} or \texttt{System.BB.Time} instead.

\cnote{On a hosted target the ``OS'' under GNULL is pthreads: \texttt{Self}
becomes \texttt{pthread\_self}, \texttt{Sleep} becomes a condition-variable
wait, the clock is \texttt{clock\_gettime(CLOCK\_MONOTONIC)}. On the ESP32-S3
there is no pthreads and (after FreeRTOS is removed) no scheduler at all but the
one we ship. The bb-runtimes \texttt{System.BB} kernel \emph{is} the OS:
\texttt{System.BB.Threads} provides the ready queue and context switch,
\texttt{System.BB.Time} provides the alarm timer and monotonic clock. GNULL is
the thin Ada shim that names those as ``the OS.''}

\section{The layer cake}

From the source line you write down to the silicon, a tasking operation passes
through five strata. Reading top to bottom:

\begin{enumerate}
  \item \textbf{Ada tasking source} --- \texttt{task}, \texttt{protected},
        \texttt{delay until}, entry calls. The compiler expands these into calls
        into GNARL via \texttt{Rtsfind}.
  \item \textbf{GNARL (upper)} --- \texttt{System.Tasking} and its children. The
        Task Control Block (ATCB), task stages, rendezvous, protected entries.
  \item \textbf{GNULL} --- \texttt{System.Task\_Primitives.Operations} plus the
        abstraction boundary \texttt{System.OS\_Interface}. The primitive set
        (\texttt{Create\_Task}, \texttt{Self}, \texttt{Sleep},
        \texttt{Delay\_Until}, \texttt{Monotonic\_Clock}, \dots).
  \item \textbf{Bare-board kernel} --- \texttt{System.BB.Threads},
        \texttt{System.BB.Time}, \texttt{System.BB.Parameters}, and the Xtensa
        register-level context switch.
  \item \textbf{Hardware} --- the two Xtensa LX7 cores, the system timer that
        drives alarms, and the windowed register file the context switch saves
        and restores.
\end{enumerate}

Table~\ref{tab:gnarl-units} maps the layers to the actual units you will find in
\texttt{crates/esp32s3\_rts/light-tasking-esp32s3/gnarl/}.

\begin{longtable}{@{}p{0.16\linewidth}p{0.34\linewidth}p{0.42\linewidth}@{}}
\caption{The tasking stack, by layer and real source unit}
\label{tab:gnarl-units}\\
\hline
\textbf{Layer} & \textbf{Unit (file)} & \textbf{Role} \\
\hline
\endfirsthead
\hline
\textbf{Layer} & \textbf{Unit (file)} & \textbf{Role} \\
\hline
\endhead
GNARL & \texttt{System.Tasking} (\texttt{s-taskin}) & The ATCB type and task-state model \\
GNARL & \texttt{System.Tasking.Restricted.Stages} (\texttt{s-tarest}) & Task creation/activation under Ravenscar \\
GNARL & \texttt{System.Tasking.Protected\_Objects.Single\_Entry} (\texttt{s-tposen}) & Single-entry protected objects (Ravenscar) \\
GNARL & \texttt{System.Tasking.Stages} (\texttt{s-tassta}), \texttt{...Rendezvous} (\texttt{s-tasren}) & Full-profile task stages and rendezvous (full only) \\
GNULL & \texttt{System.Task\_Primitives.Operations} (\texttt{s-taprop}) & The primitive set GNARL calls \\
GNULL & \texttt{System.OS\_Interface} (\texttt{s-osinte}) & The ``this is the OS'' renaming boundary \\
Kernel & \texttt{System.BB.Threads} (\texttt{s-bbthre}) & Threads, ready queue, scheduling \\
Kernel & \texttt{System.BB.Time} (\texttt{s-bbtime}) & Monotonic clock + alarm queue \\
Kernel & \texttt{System.BB.Parameters} (\texttt{s-bbpara}) & Board constants (\texttt{Ticks\_Per\_Second}, \dots) \\
Kernel & \texttt{context\_switch.S} & Xtensa windowed-register switch (\S\ref{ch:context-switch}) \\
\hline
\end{longtable}

The header comment at the top of \texttt{s-taprop.ads} is the canonical
statement of where the boundary sits. It reads ``GNAT RUN-TIME LIBRARY (GNARL)
COMPONENTS'' and then, in the body of the comment, ``This package contains all
the GNULL primitives that interface directly with the underlying OS.'' That one
file \emph{is} the GNULL spec; everything below it in the table is ``the OS.''

\section{GNARL's upper layer: the ATCB}

The central data structure of GNARL is the \emph{Ada Task Control Block}. In
\texttt{System.Tasking} (\texttt{s-taskin.ads}) it is the type
\texttt{Ada\_Task\_Control\_Block}, and a \texttt{Task\_Id} is simply an access
to one:

\begin{lstlisting}
type Ada_Task_Control_Block;
type Task_Id is access all Ada_Task_Control_Block;
Null_Task : constant Task_Id := null;
\end{lstlisting}

The ATCB holds the Ada-level facts about a task: its state
(\texttt{Unactivated}, \texttt{Runnable}, and so on), its priority, its entry
queues, its activation chain. This is the bare-board ``Jorvik/HI-E version'' of
the package, as its header says --- a restricted, static cut of the full ATCB
with no fields for features the profile forbids. GNARL operates entirely on
\texttt{Task\_Id} values; it never touches a register or a timer directly.
When it needs the machine to actually \emph{do} something --- run this task,
suspend that one --- it calls down into GNULL.

\section{The GNULL primitive set}

\texttt{System.Task\_Primitives.Operations} is a small, fixed menu. Walking the
real spec of the Ravenscar version (\texttt{s-taprop.ads}), the operations GNARL
relies on are:

\begin{lstlisting}
procedure Initialize (Environment_Task : ST.Task_Id);
procedure Create_Task
  (T          : ST.Task_Id;
   Wrapper    : System.Address;
   Stack_Size : System.Parameters.Size_Type;
   Priority   : ST.Extended_Priority;
   Base_CPU   : System.Multiprocessors.CPU_Range;
   Succeeded  : out Boolean);
procedure Enter_Task (Self_ID : ST.Task_Id);
function  Self return ST.Task_Id;
procedure Set_Priority (T : ST.Task_Id; Prio : ST.Extended_Priority);
function  Get_Priority (T : ST.Task_Id) return ST.Extended_Priority;
procedure Sleep (Self_ID : ST.Task_Id; Reason : ST.Task_States);
procedure Wakeup (T : ST.Task_Id; Reason : ST.Task_States);
procedure Delay_Until (Abs_Time : Time);
function  Monotonic_Clock return Time;
function  Is_Task_Context return Boolean;
\end{lstlisting}

Each one is a thin adapter onto the bare kernel. The mapping is the whole point
of the port:

\begin{itemize}
  \item \texttt{Initialize} --- one-time bring-up of the environment task.
        Renames through to \texttt{System.BB.Threads.Initialize}, which sets up
        the ready queue, the alarm timer, and the FPU, and installs the
        environment thread as the first runnable thread.
  \item \texttt{Create\_Task} --- allocate a kernel thread for a new Ada task. It
        calls \texttt{System.BB.Threads.Thread\_Create} with the task's entry
        wrapper, priority, CPU affinity, and a preallocated stack. (Ravenscar
        forbids task termination, so there is deliberately no destroy
        primitive.)
  \item \texttt{Self} --- ``which task is running right now?'' Implemented as
        \texttt{To\_Task\_Id (System.OS\_Interface.Get\_ATCB)} --- the kernel
        stores the ATCB pointer inside each thread descriptor, and
        \texttt{Get\_ATCB} returns the one for the running thread. This is the
        bare-board equivalent of \texttt{pthread\_self} plus thread-local
        storage, but it is a single pointer load.
  \item \texttt{Set\_Priority} / \texttt{Get\_Priority} --- rename onto
        \texttt{System.BB.Threads.Set\_Priority} / \texttt{Get\_Priority}. A
        priority \emph{decrease} (leaving a protected action) is a scheduling
        point; that is where the bare kernel may pick a different thread.
  \item \texttt{Sleep} --- suspend the caller. The body asserts
        \texttt{Self\_ID = Self} (a task can only suspend itself) and then calls
        \texttt{System.OS\_Interface.Sleep}, which renames
        \texttt{System.BB.Threads.Sleep}. The kernel marks the thread
        \texttt{Suspended} and switches away (honoring the \texttt{Wakeup\_Signaled}
        token so a wakeup racing the sleep is not lost).
  \item \texttt{Wakeup} --- make a suspended task runnable again, via
        \texttt{System.BB.Threads.Wakeup}. Note this is \emph{not} itself a
        scheduling point under Ravenscar: the woken task enters the ready queue,
        but the switch happens later when the waker lowers its own priority.
  \item \texttt{Delay\_Until} --- the engine of \texttt{delay until}. It sets the
        caller's state to \texttt{Delay\_Sleep} and arms the kernel alarm queue
        in \texttt{System.BB.Time} for the absolute wake time. When the system
        timer fires at that time, the kernel makes the thread runnable.
  \item \texttt{Monotonic\_Clock} --- the time source behind \texttt{Ada.Real\_Time.Clock}.
        \texttt{System.OS\_Interface} renames \texttt{Clock} onto
        \texttt{System.BB.Time.Clock}, which reads the hardware timer;
        \texttt{Ticks\_Per\_Second} comes from \texttt{System.BB.Parameters}.
  \item \texttt{Is\_Task\_Context} --- ``are we in a task, or in an interrupt
        handler?'' True for task context, False inside an ISR. The kernel tracks
        this with the \texttt{In\_Interrupt} flag on the thread descriptor; the
        restricted runtime uses it to keep tasking operations out of interrupt
        context.
\end{itemize}

The pivot file is \texttt{System.OS\_Interface} (\texttt{s-osinte.ads}). It is
almost entirely \texttt{renames} declarations --- it does no work itself, it just
\emph{relabels} the bare kernel as ``the OS.'' For example:

\begin{lstlisting}
function  Clock return Time renames System.BB.Time.Clock;
procedure Sleep            renames System.BB.Threads.Sleep;
procedure Wakeup (Id : Thread_Id) renames System.BB.Threads.Wakeup;
function  Thread_Self return Thread_Id
                           renames System.BB.Threads.Thread_Self;
function  Get_ATCB return System.Address
                           renames System.BB.Threads.Get_ATCB;
\end{lstlisting}

This is the literal seam between the two worlds: above it, GNARL and GNULL speak
of \texttt{Time} and \texttt{Thread\_Id}; below it, everything is
\texttt{System.BB}. Swapping \texttt{s-osinte} for a pthreads version is, in
principle, all it takes to retarget the runtime to a hosted OS --- which is
exactly how AdaCore ships GNAT on Linux.

\subsection{Where the context switch lives}

The very bottom GNULL primitive --- the moment one task stops running and another
starts --- is not Ada at all. \texttt{Sleep}, \texttt{Delay\_Until}, and priority
changes ultimately reach the kernel's scheduler, which performs the actual
register save/restore in \texttt{context\_switch.S}. On the Xtensa LX7 that means
flushing and restoring the \emph{windowed} register file, not a flat set of
callee-saved registers. That mechanism is intricate enough to have its own
chapter; see \S\ref{ch:context-switch}. From GNULL's point of view the context
switch is simply ``the OS does this''; from the silicon's point of view it is the
single most ESP32-S3-specific line of the whole runtime.

\section{The thread descriptor: GNULL's view of a task}

GNARL has the ATCB; the kernel has its own, leaner record, the
\texttt{Thread\_Descriptor} in \texttt{System.BB.Threads} (\texttt{s-bbthre.ads}).
The two are linked: each \texttt{Thread\_Descriptor} carries an \texttt{ATCB}
field (a \texttt{System.Address}) pointing back at the Ada task, which is what
makes \texttt{Self}/\texttt{Get\_ATCB} a one-pointer lookup. The descriptor holds
exactly what the kernel needs to schedule the thread and nothing more:

\begin{lstlisting}
type Thread_Descriptor is record
   Context         : aliased System.BB.CPU_Primitives.Context_Buffer;
   ATCB            : System.Address;       --  back-pointer to the Ada ATCB
   Base_CPU        : System.Multiprocessors.CPU_Range;
   Base_Priority   : Integer;
   Active_Priority : Integer;              --  raised during protected actions
   Top_Of_Stack    : System.Address;
   Bottom_Of_Stack : System.Address;
   Next            : Thread_Id;            --  ready-queue link
   Alarm_Time      : System.BB.Time.Time;  --  for delay until
   Next_Alarm      : Thread_Id;            --  alarm-queue link
   State           : Thread_States;        --  Runnable | Suspended | Delayed
   In_Interrupt    : Boolean;
   Wakeup_Signaled : Boolean;
   ...
end record;
\end{lstlisting}

Note the deliberate ordering: the comment in the spec insists that
\texttt{Context} sit at the very start of the record, because the assembly
context switch assumes the saved-register buffer is at offset zero of the thread
descriptor. The three-valued \texttt{State} type --- \texttt{Runnable},
\texttt{Suspended}, \texttt{Delayed} --- is the entire scheduling vocabulary of
the restricted profile: a thread is either eligible to run, blocked on an entry,
or waiting out a delay. There is no ``zombie,'' no ``joining,'' no dynamic
priority band, because the profile forbids the features that would need them.

\section{Three profiles, one boundary}

This repository builds three runtime profiles, and the GNARL/GNULL split is the
hinge on which they differ. The selection happens in
\texttt{crates/esp32s3\_rts/gen\_runtime.sh} via the \texttt{ESP32S3\_RTS\_PROFILE}
environment variable; \S\ref{ch:runtime-build} covers the build mechanics in
full.

\begin{itemize}
  \item \textbf{light-tasking} (the default) and \textbf{embedded} both use the
        \emph{restricted} GNARL+GNULL --- the Ravenscar/Jorvik runtime described
        so far. The ATCB is the small \texttt{s-taskin} cut; task creation goes
        through \texttt{s-tarest} (\texttt{Tasking.Restricted.Stages}); protected
        objects are single-entry (\texttt{s-tposen}). There is no rendezvous, no
        \texttt{abort}, no dynamic priority change beyond ceiling locking. The
        two differ only in exception/finalization support
        (\texttt{No\_Exception\_Propagation} + \texttt{No\_Finalization} for
        light-tasking; full ZCX exceptions and finalization for embedded) --- not
        in the tasking model.
  \item \textbf{full} swaps in the \emph{complete} GNARL and a different GNULL.
        \texttt{gen\_runtime.sh} synthesizes it by cloning the embedded source
        tree and overlaying \texttt{crates/esp32s3\_rts/full\_overlay/}: it
        imports the full GNARL units (\texttt{s-tassta}, \texttt{s-tasren} for
        rendezvous, dynamic priorities, abort/ATC machinery) from a native GNAT,
        deletes the restricted stages (\texttt{rm -f s-tarest.*}), and drops in
        \texttt{full\_overlay/gnarl/s-taprop} --- a from-scratch GNULL spec/body
        written against the same \texttt{System.BB} kernel.
\end{itemize}

The full GNULL spec is visibly larger than the Ravenscar one. Its own header
calls it ``the STANDARD Task\_Primitives.Operations interface (as the full GNARL
--- s-tassta / s-tasren --- expects it), implemented over the bare-board BB
kernel.'' Beyond the restricted set it must provide the full lock family
(\texttt{Write\_Lock}/\texttt{Read\_Lock}/\texttt{Unlock} on \texttt{Lock} and
\texttt{RTS\_Lock}), \texttt{Timed\_Sleep}/\texttt{Timed\_Delay},
suspension-object operations, dynamic \texttt{Set\_Priority} with inheritance,
SMP affinity (\texttt{Set\_Task\_Affinity}, \texttt{Get\_CPU}), and the abort/ATC
family (\texttt{Abort\_Task}, \texttt{Suspend\_Task}, \texttt{Stop\_All\_Tasks},
\dots). The crucial design fact: \emph{the bottom of the stack did not change.}
Both the restricted and the full GNULL sit on the same \texttt{System.BB}
kernel. Going to the full profile means a richer upper layer and a fatter
adapter, but the ``OS'' underneath --- threads, clock, alarm queue, context
switch --- is the same bare-board executive in both cases.

\section{Following one operation all the way down}

To make the layering concrete, trace a single \texttt{delay until Next} for a
periodic task:

\begin{enumerate}
  \item \textbf{Source / compiler.} The compiler turns \texttt{delay until Next}
        into a call to a GNARL delay routine, passing the absolute wake time.
  \item \textbf{GNARL.} GNARL converts the \texttt{Ada.Real\_Time.Time} to the
        runtime's internal \texttt{Time} and calls the GNULL primitive
        \texttt{System.Task\_Primitives.Operations.Delay\_Until}.
  \item \textbf{GNULL.} \texttt{Delay\_Until} records
        \texttt{Self\_ID.Common.State := Delay\_Sleep} on the ATCB, then calls
        through \texttt{System.OS\_Interface.Delay\_Until}, which renames
        \texttt{System.BB.Time.Delay\_Until}.
  \item \textbf{Kernel.} \texttt{System.BB.Time} writes the wake time into the
        thread descriptor's \texttt{Alarm\_Time}, links it into the alarm queue
        ordered by expiry, marks the thread \texttt{Delayed}, and programs the
        hardware timer for the nearest alarm. It then invokes the scheduler.
  \item \textbf{Hardware + context switch.} With the caller no longer runnable,
        the scheduler picks the next ready thread and \texttt{context\_switch.S}
        restores its windowed registers. Later, the timer interrupt fires at
        \texttt{Next}, the kernel moves the thread back to \texttt{Runnable}, and
        at the next scheduling point it resumes --- on the line after the
        \texttt{delay until}.
\end{enumerate}

Five layers, one statement. That is the whole reason the runtime is split this
way: GNARL wrote step~2 once, for every Ada target in the world, and this port
only had to supply steps~3--5 by pointing \texttt{System.OS\_Interface} at
\texttt{System.BB} instead of at pthreads. The user view of these same
constructs is in \S\ref{ch:tasking}; the interrupt that ends the delay is in
\S\ref{ch:interrupts}; and the register-level switch in step~5 is in
\S\ref{ch:context-switch}.

\chapter{The Context Switch}
\label{ch:context-switch}

Every preemptive runtime turns on one small, brutal piece of code: the routine
that freezes one task's CPU state, thaws another's, and jumps. On a flat
register machine that routine is tedious but obvious. On the Xtensa LX7, with its
\emph{windowed} register file and a hardware FPU bolted on the side, it is neither
\textemdash{} and getting it wrong produces failures that look like anything but a
broken context switch. This chapter walks through the mechanism end to end: where
it lives, the state it must preserve, the two completely different ways it is
entered, and how the trickier parts were verified on hardware.

\section{Where it lives}\index{context switch}

The switch spans three files, two languages, and the boundary between the generic
GNARL kernel and this CPU port:

\begin{itemize}
  \item \texttt{crates/bb-runtimes/xtensa/context\_switch.S} \textemdash{} the
        register-level switch itself, in hand-written Xtensa assembly. It is
        copied verbatim into each profile's runtime tree
        (\texttt{<profile>-esp32s3/gnarl/context\_switch.S}) when
        \texttt{gen\_runtime.sh} builds the RTS, so editing the bb-runtimes copy
        and regenerating is what changes the running behaviour.
  \item \texttt{System.BB.CPU\_Primitives} (\texttt{s-bbcppr.adb}) \textemdash{}
        the Ada entry point \texttt{Context\_Switch}, a thin wrapper that picks
        the per-CPU thread-table slots and calls the assembly. This is the seam
        the portable kernel sees; every CPU port provides its own.
  \item \texttt{System.BB.CPU\_Specific} (\texttt{s-bbcpsp.ads}) \textemdash{} the
        \texttt{Context\_Buffer} record: the per-task saved CPU state, whose field
        offsets the assembly hard-codes.
\end{itemize}

Above all of this sits the portable scheduler, \texttt{System.BB.Threads.Queues}
(\texttt{s-bbthqu.adb}). It never touches a register; it only maintains the ready
queue and answers one question \textemdash{} \texttt{Context\_Switch\_Needed},
defined simply as \texttt{First\_Thread /= Running\_Thread}. When the highest
ready task is not the one on the CPU, the kernel calls
\texttt{CPU\_Primitives.Context\_Switch} and the machinery below takes over.

\section{The data structures}\index{context buffer}\index{stack pointer}\index{thread table}

Two structures carry the state.

\textbf{The \texttt{Context\_Buffer}}, one per task, is deliberately tiny:

\begin{longtable}{p{3cm} p{1.2cm} p{8cm}}
\toprule
\textbf{Field} & \textbf{Off} & \textbf{Meaning} \\ \midrule
\endhead
\texttt{SP}         & 0  & task stack pointer (top of its spilled frame) \\
\texttt{PC}         & 4  & resume target: a \texttt{retw} stub, or the new-task trampoline \\
\texttt{PS}         & 8  & processor state (INTLEVEL, window control, \dots) \\
\texttt{A0}         & 12 & windowed return address \\
\texttt{THREADPTR}  & 16 & thread-local-storage base \\
\texttt{CP\_State}  & 20 & pointer to the FPU/coprocessor save area \\
\texttt{SAR}        & 24 & shift-amount register \\
\texttt{LBEG/LEND/LCOUNT} & 28..36 & zero-overhead loop registers \\
\texttt{Frame\_Kind} & 40 & 0 = solicited, 1 = interrupt (see below) \\
\bottomrule
\end{longtable}

What is conspicuously \emph{absent} is the general register file. That is the
whole point of the windowed design, and the next section explains why.

\textbf{The thread tables}, \texttt{Running\_Thread\_Table} and
\texttt{First\_Thread\_Table}, are indexed by CPU (this is a genuine two-core SMP
target). \texttt{Running} is who is on the core now; \texttt{First} is who the
scheduler wants there. \texttt{Context\_Switch} passes the \emph{addresses} of
these two slots to the assembly, which dereferences them to the descriptors whose
first field is the \texttt{Context\_Buffer}. The switch updates
\texttt{Running := First} as part of resuming, so the table always names the live
task.

\section{The windowed-ABI premise}\index{register window}\index{window overflow}

On the Xtensa windowed ABI a \texttt{call} does not push registers; it
\emph{rotates} a 64-entry physical register file so the callee sees a fresh
window. Live data from outer frames stays in the physical registers until the
file fills, at which point the hardware raises window-overflow exceptions that
spill the oldest frames to their stacks; returning raises window-underflow
exceptions that reload them. The implication for a context switch is liberating:
to save a task, you do not copy sixteen or thirty-two registers \textemdash{} you
issue \texttt{SPILL\_ALL\_WINDOWS}, which flushes every live frame to that task's
own stack, and then almost the entire register state lives on the stack. All you
must record separately is the stack pointer and the handful of special registers
the window machinery does not cover (\texttt{SAR}, the loop registers, the FPU,
\texttt{THREADPTR}). Resuming is the mirror image: set the stack pointer, execute
one \texttt{retw}, and let window-underflow lazily fault the frames back in as the
task returns up its own call chain. This is why the \texttt{Context\_Buffer} is so
small, and why this port is closest in spirit to the SPARC (LEON) windowed port
rather than the flat ARM and RISC-V ones.

\section{Two ways in}\index{preemption}\index{interrupt frame}

A task can lose the CPU two very different ways, and the switch handles them with
two different entry points.

\textbf{Cooperative} \textemdash{} the running task blocks (a \texttt{delay
until}, a protected entry that cannot proceed) or a higher task becomes ready at
a point where the kernel is already running. The kernel calls
\texttt{Context\_Switch}, which calls \texttt{\_\_gnat\_context\_switch}. This path
runs in the task's own context, so it can do the full \texttt{SPILL\_ALL\_WINDOWS}
and save the outgoing task in \emph{solicited} form.

\textbf{Preemptive} \textemdash{} a level-5 timer or IPI interrupt fires while a
task is running. The L5 vector (\texttt{highint5.S}) saves the interruptee with
\texttt{\_xt\_context\_save} into a full \texttt{XT\_STK} interrupt frame, services
the tick, and \textemdash{} if that made a higher task ready \textemdash{}
tail-jumps \texttt{\_\_gnat\_preempt\_dispatch} to switch away. This path saves the
outgoing task in \emph{interrupt} form.

The two must not be mixed, and that is the subtle part. A cooperative
\texttt{SPILL\_ALL\_WINDOWS} executed from inside an interrupt's borrowed window
chain corrupts the register windows. So when the kernel decides mid-interrupt that
a switch is needed, \texttt{Context\_Switch} does \emph{not} switch: it sets a
per-CPU \texttt{Switch\_Pending} flag and returns. The interrupt epilogue, running
in the vector's clean single-window context, notices the flag and performs the
dispatch itself via \texttt{\_\_gnat\_preempt\_dispatch}. (Tracking down a bug here
\textemdash{} the cooperative spill running in interrupt context, marching the
stack pointer down one frame per tick until a window spill double-faulted
\textemdash{} is the story behind several of the comments in \texttt{highint5.S}
and \texttt{context\_switch.S}, and behind ACATS \texttt{CXD8002}.)

\section{Dual-format dispatch}\index{WINDOWBASE}\index{PS register}

Because a task can be saved either way, each \texttt{Context\_Buffer} records its
\texttt{Frame\_Kind}, and the shared resume code reads it to resume correctly:

\begin{itemize}
  \item \textbf{Solicited} (\texttt{Frame\_Kind = 0}): the live windows were
        spilled to the task's stack. Resume by restoring \texttt{SP}, \texttt{SAR},
        the loop registers and \texttt{PS}, then a single \texttt{retw} \textemdash{}
        window-underflow reloads the rest.
  \item \textbf{Interrupt} (\texttt{Frame\_Kind = 1}): the full register state is in
        an \texttt{XT\_STK} frame at \texttt{SP}. Resume by calling
        \texttt{\_xt\_context\_restore} and returning with \texttt{rfi}.
\end{itemize}

Both entry points converge on a single label, \texttt{.Lresume}: it makes
\texttt{First} the running task, restores the format-independent extra state (FPU,
\texttt{THREADPTR}), then branches on \texttt{Frame\_Kind} to one of the two
resume tails. The interrupt tail carries hard-won detail \textemdash{} it forces a
single clean window (\texttt{WINDOWSTART = 1 << WINDOWBASE}) so a freshly scheduled
task does not inherit the dispatcher's stale window mask, reloads \texttt{SAR}
(clobbered by the window-start arithmetic), and returns with \texttt{rfi 5} so
\texttt{INTLEVEL} stays high until the very last instruction, closing the race
where a fresh tick preempts a half-restored task.

\section{Walkthrough: a cooperative switch}

A task calls \texttt{delay until}. The kernel marks it blocked, finds a different
\texttt{First\_Thread}, and calls \texttt{\_\_gnat\_context\_switch} with the two
slot addresses:

\begin{enumerate}
  \item \texttt{SPILL\_ALL\_WINDOWS} flushes the task's live frames to its stack.
  \item Save into the outgoing buffer: \texttt{SP}, \texttt{A0}, \texttt{PS},
        \texttt{SAR}, the loop registers, \texttt{THREADPTR}, and (if the task owns
        an FPU save area) \texttt{F0..F15} + \texttt{FCR} + \texttt{FSR}.
        \texttt{Frame\_Kind := 0}.
  \item Fall into \texttt{.Lresume}: set \texttt{Running := First}, restore the
        incoming task's FPU and \texttt{THREADPTR}, and \textemdash{} since the
        incoming task was itself last saved solicited \textemdash{} restore its
        \texttt{SP}/\texttt{SAR}/loops/\texttt{PS} and \texttt{retw} into it.
\end{enumerate}

\section{Walkthrough: a preemptive switch}

The L5 timer fires while a compute task runs:

\begin{enumerate}
  \item The vector saves the interruptee with \texttt{\_xt\_context\_save} into an
        \texttt{XT\_STK} frame and services the tick.
  \item If a higher task is now ready, the epilogue tail-jumps
        \texttt{\_\_gnat\_preempt\_dispatch} from the vector's clean single window
        (no spill here \textemdash{} the full state is already in \texttt{XT\_STK}).
  \item \texttt{\_\_gnat\_preempt\_dispatch} records the interruptee:
        \texttt{SP} (= the \texttt{XT\_STK} frame), \texttt{THREADPTR}, the FPU,
        \texttt{Frame\_Kind := 1}.
  \item \texttt{.Lresume} restores the incoming task's FPU/\texttt{THREADPTR} and
        resumes it in its own format \textemdash{} solicited via \texttt{retw}, or
        interrupt via \texttt{\_xt\_context\_restore} + \texttt{rfi}.
\end{enumerate}

\section{Starting a task that has never run}\index{task stack}\index{trampoline}

A brand-new task has no spilled frames to reload, so its first dispatch cannot use
\texttt{retw}. \texttt{Initialize\_Context} sets the new buffer's \texttt{PC} field
to the trampoline \texttt{\_\_gnat\_start\_thread} and stashes the task's entry
point and argument at the top of its fresh stack. The first time the scheduler
resumes the task, control lands in the trampoline, which loads the argument and
\texttt{callx} es the task wrapper \textemdash{} which never returns. The buffer's
loop registers are zero-initialised (a non-zero \texttt{LCOUNT} would make the
hardware spuriously loop on first entry) and \texttt{THREADPTR} starts at zero,
i.e.\ ``no thread-local storage'' until the task sets it.

\section{What state is saved \textemdash{} and what is not}\index{coprocessor state}\index{THREADPTR}

\begin{longtable}{p{5.5cm} p{2.2cm} p{5.5cm}}
\toprule
\textbf{State} & \textbf{Saved?} & \textbf{Notes} \\ \midrule
\endhead
\texttt{SP}/\texttt{PS}/\texttt{A0}, register windows & always & spill + frame \\
\texttt{SAR}, \texttt{LBEG}/\texttt{LEND}/\texttt{LCOUNT} & always & loops: \texttt{XCHAL\_HAVE\_LOOPS}=1 \\
\texttt{THREADPTR} & always & TLS base; per task \\
\texttt{F0..F15} + \texttt{FCR} + \texttt{FSR} & conditional & only when the task owns a \texttt{CP\_State} area \\
MAC16 \texttt{ACCLO}/\texttt{ACCHI}/\texttt{Mn} & no & \texttt{XCHAL\_HAVE\_MAC16}=0 \textemdash{} the unit is absent \\
\bottomrule
\end{longtable}

The FPU is handled \emph{eagerly per task}: the run-time enables the FPU bit
(\texttt{CP0}) of \texttt{CPENABLE} for every task, gives each a \texttt{CP\_State}
save area, and the switch copies all sixteen \texttt{F} registers plus
\texttt{FCR}/\texttt{FSR} whenever the buffer's \texttt{CP\_State} pointer is
non-null. Note this covers \emph{only} the FPU: the PIE/SIMD coprocessor
(\texttt{CP3}) is not part of the per-task enable, so a task using
\texttt{ESP32S3.SIMD} must switch it on itself (\S\ref{ch:asm}) and keep its vector
work single-task, since the \texttt{q}-registers are not in the save area. \texttt{THREADPTR} was the last addition to
the list; for some time the field existed in the record but the assembly never
touched it, leaving all tasks sharing one value \textemdash{} harmless only because
nothing yet used per-task TLS.

\section{Verifying the hard parts}

The FPU and \texttt{THREADPTR} paths share a verification hazard: a test that
exercises them na\"ively passes whether or not the switch actually preserves them,
because the value under test was sitting safely in memory the whole time. A green
test means nothing unless it is \emph{sensitive} \textemdash{} demonstrably red
when the property is broken. Both were therefore verified two-sided on the
\texttt{full} profile: a deliberately broken build must fail, and only then does
the intact build's success count.

\textbf{FPU.} A low-priority checker holds six floating-point accumulators live in
registers across a two-million-step loop (each step the exact identity
\texttt{X := X * Lm * Li}, with \texttt{Lm}=2.0 and \texttt{Li}=0.5 read once from
\texttt{Volatile} cells so the compiler can neither fold the pair away nor spill
the accumulators \textemdash{} disassembly confirmed zero memory operations in the
loop). A higher-priority hog wakes every millisecond and tramples all of
\texttt{F0..F15}. With the preemptive FPU save disabled the checker logs corruption
continuously; with it intact, zero corruption across hundreds of preemptions per
batch.

\textbf{THREADPTR.} A checker task sets \texttt{THREADPTR} to one sentinel; a
higher-priority hog sets a different one and preempts it (both confirmed on the
same core). The clean discriminator turned out to be the \emph{hog's} value on
re-entry, not the checker's \textemdash{} the checker, as the background task,
reads its own value either way. With the restore disabled the hog comes back
seeing the checker's sentinel (its own was lost); with the restore intact it keeps
its own. That asymmetry is the proof the per-task save/restore works.

Two methodology traps recurred and are worth stating once. First, optimisation:
a value only tests the switch if the compiler keeps it \emph{live in a register}
across the preemption window \textemdash{} verify that by disassembling the loop,
not by reading the source. Second, staleness: editing
\texttt{bb-runtimes/xtensa/context\_switch.S} does not take effect until the
generated \emph{base} runtime tree is deleted and regenerated, and the only
trustworthy confirmation that a change is in the running image is to disassemble
the final linked ELF. Both traps produce a green run from a binary that does not
contain your change \textemdash{} the most dangerous kind of pass.

\chapter{Interrupts}
\label{ch:interrupts}

The context switch in the previous chapter does not run by magic; something
has to \emph{call} the scheduler. On this port that something is almost always an
interrupt: the systimer tick, or a device interrupt whose handler makes a
higher-priority task ready. This chapter follows an interrupt from the silicon
vector all the way to an Ada handler, and explains how \texttt{pragma
Attach\_Handler} is supported across all three profiles, including the
bare-board full \texttt{System.Interrupts} the \texttt{full} profile needs,
which no stock bare-board runtime provides.

\section{The Xtensa interrupt levels}\index{interrupts}\index{interrupt levels}\index{xt\_highint5}

The ESP32-S3's Xtensa LX7 has six maskable interrupt levels plus a non-maskable
one, and two of those levels are special:

\begin{longtable}{p{2.2cm} p{10.5cm}}
\toprule
\textbf{Level} & \textbf{Role on this port} \\ \midrule
\endhead
1     & The general \emph{exception} vector (synchronous faults). No async
        interrupt is assigned here. \\
2--3  & \textbf{The OS band} (\texttt{\_\_gnat\_level2\_vector} /
        \texttt{\_\_gnat\_level3\_vector}): every interrupt that talks to
        the scheduler. The device interrupts live here, and so do the systimer
        tick and the cross-core IPI ``poke'' (both share CPU interrupt 21, at
        level 2). \\
4--5  & \emph{Free for OS-oblivious handlers}: lean assembly that never
        touches the scheduler. Nothing is installed today; the old level-5
        vector (\texttt{xt\_highint5}) survives only as a park stub. \\
6     & \texttt{DEBUGLEVEL} (debug exceptions: breakpoints and the
        \texttt{DBREAK} data watchpoints, used for stack-overflow detection). \\
7     & \texttt{NMILEVEL} (non-maskable; not used). \\
\bottomrule
\end{longtable}

One architectural constant shapes everything below.
\texttt{XCHAL\_EXCM\_LEVEL} is 3, and it draws the line the whole design obeys:
\emph{every interrupt that saves context or can trigger a task switch sits at
or below it}. \texttt{PS.EXCM} -- which the hardware sets during every window
over/underflow exception -- masks exactly levels 1--3. So does the raised
\texttt{INTLEVEL} across every register-window spill. Keep all the
scheduler-facing interrupts inside that band and no interrupt can ever observe
a half-spilled register window. This is the FreeRTOS Xtensa port's rule, and
this port adopted it deliberately (verified against the esp-idf source). The
original design ran the tick at level 5, above the line -- and a level-5 tick
can land in the middle of window-critical code and corrupt the register state.
Levels above \texttt{EXCM\_LEVEL} remain available, but only for handlers
that never touch the scheduler.

\section{Priorities and the masking model}\index{interrupt priorities}\index{interrupt masking}

Ada interrupt priorities map onto Xtensa levels in
\texttt{System.BB.Board\_Support.Priority\_Of\_Interrupt}: a hardware interrupt's
fixed level becomes the Ada priority its handler runs at, and conversely a task's
ceiling becomes the \texttt{INTLEVEL} it raises while it holds a lock
(\texttt{Enable\_Interrupts} in \texttt{s-bbcppr}). \texttt{Ada.Interrupts.Names}
gives symbolic IDs for the slots that are free for application use:
\texttt{Device\_L2\_0/1} (CPU interrupts 19/20) and \texttt{Device\_L3\_0/1}
(23/27), with ceilings \texttt{Device\_L2\_Priority} /
\texttt{Device\_L3\_Priority}. CPU interrupt 21 (\texttt{Device\_L2\_2}) is
kernel-reserved: the systimer tick and the cross-core poke share it, and the
kernel attaches its alarm handler there at start-up. Do not attach to it.

\section{What the silicon does}\index{EPC}\index{EPS}\index{EXCSAVE}\index{VECBASE}\index{rfi}

Before any of our code runs, the CPU itself performs a fixed sequence. It is
small, and everything above depends on it, so here it is in full. When an
enabled interrupt of level $n$ (2--7) is the highest one pending:

\begin{enumerate}
  \item The hardware copies the current \texttt{PC} into \texttt{EPC$n$} and
        the current \texttt{PS} into \texttt{EPS$n$}. Each level has its own
        pair, so a level-3 interrupt cannot clobber the state a level-2 entry
        is still using.
  \item It sets \texttt{PS.INTLEVEL} to $n$, masking level $n$ and below.
  \item It jumps to the level's fixed slot in the vector table:
        \texttt{VECBASE} plus a per-level offset. Nothing is pushed on any
        stack --- there is no hardware stacking on Xtensa. The vector gets the
        CPU exactly as the interrupted code left it, registers and all.
\end{enumerate}

That last point drives the vector's first moves. The slot has \emph{no} free
register, so its first instruction stashes \texttt{a0} in the level's
\texttt{EXCSAVE$n$} scratch register; only then can it start saving state
properly. From C this is the part the NVIC did for you on a Cortex-M ---
here the eight-instruction prologue that builds the frame is ours, in
\texttt{highint5.S}.

The return trip is one instruction: \texttt{rfi $n$} restores
\texttt{PS} from \texttt{EPS$n$} and jumps to \texttt{EPC$n$}
\emph{atomically}. That atomicity is load-bearing: \texttt{INTLEVEL} stays
at $n$ until the very last instruction, so there is no window in which a
fresh interrupt could preempt a half-restored context.

Two registers complete the picture. \texttt{INTERRUPT} (read-only) shows
which of the 32 CPU interrupts are pending right now; the dispatch reads it
to find the sources to service. \texttt{INTENABLE} gates them per core; the
runtime sets a bit when a handler is attached. Most sources are
\emph{level-triggered}: the CPU-side bit stays asserted until the handler
clears the \emph{peripheral's} status register (usually write-1-to-clear),
which is why every HAL handler acks the device, not the CPU.

\section{The interrupt matrix}\index{interrupt matrix}\index{INT\_MAP}

If you come from a Cortex-M, one ESP32 idea needs installing first: peripheral
interrupts do not have fixed CPU positions. The chip has roughly a hundred
peripheral \emph{sources} (TWAI, each GDMA channel, the systimer comparators,
GPIO, \dots) and each core has just 32 CPU interrupts, each with a
\emph{fixed, hardwired} level and trigger type. Between the two sits the
\emph{interrupt matrix}: one map register per source, per core
(\texttt{INTERRUPT\_CORE0\_*\_INT\_MAP}), holding the number of the CPU
interrupt that source should raise.

Three consequences matter in practice:

\begin{itemize}
  \item \textbf{Routing is a design decision.} A driver picks its
        interrupt's level by picking which CPU interrupt to map its source
        to. That is how the LCD refill moved from level 2 to level 3, and how
        the tick moved from level 5 to level 2 --- one store to a map
        register each.
  \item \textbf{Sources can share.} Several sources may map to one CPU
        interrupt; the handler then reads the peripherals to tell them apart.
        The kernel uses this deliberately: the systimer alarm and the
        cross-core poke share CPU interrupt 21, and the dispatch checks the
        \texttt{FROM\_CPU} register to distinguish them.
  \item \textbf{Not every CPU interrupt is usable.} The 32 slots are a mixed
        bag: some are level-triggered, some edge-triggered, some are internal
        (software) interrupts the matrix cannot drive at all. The usable
        level-triggered OS-band slots are few --- 19/20/21 at level 2, 23/27
        at level 3 --- which is why the allocation in
        \texttt{Ada.Interrupts.Names} is as tight as it is. (Two findings
        from the field: CPU interrupt 22 is edge-triggered and will simply
        never fire for a level source, and 29 is an internal software
        interrupt the matrix cannot reach. Both look free on paper.)
\end{itemize}

\section{From silicon to handler}\index{interrupt vector}\index{interrupt wrapper}

\texttt{start.S} points \texttt{VECBASE} at our vendored vector table before
anything else runs, so every vector is ours. The level-2 and level-3 vectors
are two instantiations of one assembly macro (\texttt{gnat\_os\_vector} in
\texttt{highint5.S}), and they mirror the FreeRTOS Xtensa port's
\texttt{\_frxt\_int\_enter}/\texttt{\_frxt\_int\_exit} structure. When an
OS-band interrupt fires (say the tick, on CPU interrupt 21), this is the path:

\begin{enumerate}
  \item \textbf{Vector entry.} The CPU jumps to the level-2 slot. The
        vector's first instruction masks to \texttt{DEBUGLEVEL}
        (\texttt{rsil 6}): the whole prologue is now atomic against the OS
        band, and the frame build below cannot trip the per-task
        stack-overflow watchpoint.
  \item \textbf{Save.} It builds an \texttt{XT\_STK} frame on the
        interrupted stack and calls \texttt{\_xt\_context\_save} to spill the
        full register state into it. The spill leaves
        \texttt{WINDOWSTART = 1 << WINDOWBASE} by construction, so a plain
        resume needs no window bookkeeping at all.
  \item \textbf{Enter.} Still masked, it increments the per-CPU nesting
        count (\texttt{\_\_gnat\_int\_nest}, owned entirely by the vector).
        On the outermost entry it saves the task-side SP and hops onto the
        per-CPU \emph{interrupt stack}, so handler frames and their spills
        never land on task stacks.
  \item \textbf{Service.} It sets \texttt{INTLEVEL} to
        \texttt{XCHAL\_EXCM\_LEVEL} and calls the Ada dispatch routine
        (\texttt{\_\_gnat\_level2\_dispatch} or
        \texttt{\_\_gnat\_level3\_dispatch}). The dispatch reads the Xtensa
        \texttt{INTERRUPT} register and, for each pending source it owns,
        calls \texttt{System.BB.Interrupts.Interrupt\_Wrapper (Id)}. Because
        every handler runs at \texttt{EXCM\_LEVEL}, OS-band handlers never
        preempt one another; the interrupt \emph{level} decides which pending
        source is taken first and which interrupts preempt task code.
  \item \textbf{The wrapper} (\texttt{s-bbinte.adb}) is the GNARL bridge: it
        records the interrupt being handled, raises the running task's active
        priority to the interrupt's priority (so the ceiling protocol holds),
        runs the \emph{registered handler}, then restores priority. The
        handler is an \texttt{access procedure (Id : Interrupt\_ID)}.
  \item \textbf{Epilogue: the one switch site.} Back in the vector, masked
        again, the nesting count is decremented. Only the \emph{outermost}
        exit may switch, and it decides by itself: it compares
        \texttt{First\_Thread} with \texttt{Running\_Thread} and, if they
        differ, tail-jumps \texttt{\_\_gnat\_preempt\_dispatch} from the
        vector's clean single-window context. One guard applies: if the
        interrupted PC lies inside vector text, the switch is deferred to the
        next interrupt exit -- a half-entered vector's \texttt{EPC/EPS/%
EXCSAVE} registers are still live in hardware and must not be stranded.
  \item \textbf{Resume.} Otherwise \texttt{\_xt\_context\_restore} +
        \texttt{rfi} returns to the interruptee atomically:
        \texttt{INTLEVEL} stays high until the very last instruction, so
        there is no re-entrancy window.
\end{enumerate}

The level-3 vector is the same macro with \texttt{EPS\_3/EPC\_3} and
\texttt{rfi 3}. Note what is \emph{not} here: no per-vector deferral flags,
no software \texttt{WINDOWSTART} repair on the plain-resume path, and no
special tick vector -- the tick is just another level-2 source, dispatched
like any device interrupt.

\section{Attaching a handler: two axes}\index{Attach\_Handler}\index{protected handler}

It helps to keep two \emph{independent} distinctions apart. The first is
\emph{which API layer} you attach through; the second (which lives entirely
inside the high-level layer) is \emph{static versus dynamic} attachment.

\subsection{Axis 1: API layer (high level vs low level)}\index{Ada.Interrupts}

\textbf{High level: \texttt{Ada.Interrupts} / \texttt{pragma Attach\_Handler}.}
The portable Ada way: the handler is a protected procedure, the enclosing
protected object's \texttt{Interrupt\_Priority} is its ceiling, and the compiler
generates the registration and ceiling locking. This is what the GPIO HAL and the
\texttt{esp32s3\_intr\_levels} example use. The same counting handler as the
low-level version below, written this way, is:

\begin{lstlisting}
with Ada.Interrupts.Names; use Ada.Interrupts.Names;

package body High is
   protected PO with Interrupt_Priority => Device_L2_Priority is
      function Count return Natural;
   private
      procedure Handle;
      pragma Attach_Handler (Handle, Device_L2_0);  -- installed at elaboration
      N : Natural := 0;
   end PO;

   protected body PO is
      procedure Handle is           -- runs at the PO ceiling; locking is implicit
      begin
         Clear_Source;              -- ack the (level-triggered) source
         N := N + 1;
      end Handle;
      function Count return Natural is (N);
   end PO;
end High;
\end{lstlisting}

There is no \texttt{Install} call: \texttt{pragma Attach\_Handler} wires the
handler up when \texttt{PO} elaborates. The handler is a protected procedure (no
\texttt{Id} parameter, since it is bound to one interrupt), the count \texttt{N} is a
protected component rather than a bare \texttt{Volatile}, and the ceiling protocol
serialises it against other accesses to \texttt{PO} for free. Contrast the
hand-rolled low-level form:

\textbf{Low level: \texttt{System.BB.Interrupts.Attach\_Handler}.} The
bareboard kernel's own primitive: register a plain \texttt{access procedure (Id)}
for an \texttt{Interrupt\_ID}, with an explicit priority. This is what
\texttt{Interrupt\_Wrapper} ultimately calls, it is available on \emph{all}
profiles, and it needs no protected object, but also gives you no
compiler-enforced ceiling. It is the escape hatch when the portable layer is
unavailable or unwanted:

\begin{lstlisting}
with System; use System;
with System.BB.Interrupts; use System.BB.Interrupts;

package body Low is
   Count : Natural := 0 with Volatile;

   procedure Handle (Id : Interrupt_ID) is
   begin
      Clear_Source (Id);          --  ack the (level-triggered) source
      Count := Count + 1;
   end Handle;

   procedure Install is
   begin
      --  plain procedure, explicit Id + priority, no protected object
      Attach_Handler (Handle'Access, Id => 19,
                      Prio => Interrupt_Priority'Last);
   end Install;
end Low;
\end{lstlisting}

\subsection{Axis 2: static vs dynamic (within the high-level layer)}\index{dynamic attachment}

Static versus dynamic is the dividing line of the \texttt{No\_Dynamic\_Attachment}
restriction; it only concerns the high-level \texttt{Ada.Interrupts} layer.

\textbf{Static: \texttt{pragma Attach\_Handler}.} The handler is fixed at
compile time and installed once, automatically, when the protected object is
elaborated. There is no run-time attach call; the compiler lowers it straight to
\texttt{System.Interrupts.Install\_Handlers}, \emph{bypassing}
\texttt{Ada.Interrupts}. This is the form to reach for by default:

\begin{lstlisting}
with Ada.Interrupts.Names; use Ada.Interrupts.Names;

package body Blink is
   protected L2_PO with Interrupt_Priority => Device_L2_Priority is
   private
      procedure Handle;
      pragma Attach_Handler (Handle, Device_L2_0);  --  installed at elaboration
      N : Integer := 0;
   end L2_PO;

   protected body L2_PO is
      procedure Handle is
      begin
         Clear_L2;                 --  ack the source
         N := N + 1;
      end Handle;
   end L2_PO;
end Blink;
\end{lstlisting}

\textbf{Dynamic: the \texttt{Ada.Interrupts} operations.} Here the handler is
attached, swapped, or detached \emph{at run time}. The protected procedure must
first be made eligible with \texttt{pragma Interrupt\_Handler}; then
\texttt{Attach\_Handler}, \texttt{Detach\_Handler}, \texttt{Exchange\_Handler} and
\texttt{Current\_Handler} act on it during execution. These calls go \emph{through}
the \texttt{Ada.Interrupts} body: the path the bare runtime initially stubbed
out (every operation raised \texttt{Program\_Error}) until a delegating
\texttt{a-interr} body was supplied (see \emph{Why the profiles differ} below):

\begin{lstlisting}
with Ada.Interrupts;       use Ada.Interrupts;
with Ada.Interrupts.Names; use Ada.Interrupts.Names;

package body Dyn is
   protected PO is
      procedure Handle;
      pragma Interrupt_Handler (Handle);   --  eligible for dynamic attach
   end PO;
   protected body PO is
      procedure Handle is begin null; end Handle;
   end PO;

   procedure Run is
   begin
      Attach_Handler (PO.Handle'Access, Device_L2_0);  --  attach now
      --  ... later, restore the default treatment ...
      Detach_Handler (Device_L2_0);
   end Run;
end Dyn;
\end{lstlisting}

So the two axes compose: a high-level handler is either static (the \texttt{pragma},
installed at elaboration) or dynamic (the run-time \texttt{Ada.Interrupts} calls),
while the low-level \texttt{System.BB.Interrupts} primitive sits beneath both,
itself a run-time call, but \emph{not} the Ada \texttt{No\_Dynamic\_Attachment}
notion of ``dynamic''.

\section{Why the profiles differ}\index{runtime profile}\index{System.Interrupts}

The high-level form compiles differently depending on the tasking profile, and
that is the crux:

\begin{itemize}
  \item \texttt{embedded} and \texttt{light-tasking} set \texttt{pragma Profile
        (Jorvik)}. Under a Ravenscar-family profile the compiler lowers
        \texttt{pragma Attach\_Handler} to the \emph{restricted} installation path,
        \texttt{System.Interrupts.Install\_Restricted\_Handlers}, which the
        bareboard \texttt{s-interr} provides. So it works.
  \item \texttt{full} sets \emph{no} profile pragma (its whole point is to lift
        the Jorvik restrictions). The compiler then lowers interrupt-handler
        protected objects to the \emph{full dynamic} machinery:
        \texttt{Register\_Interrupt\_Handler}, \texttt{Install\_Handlers}, and the
        \texttt{Static\_Interrupt\_Protection} / \texttt{Dynamic\_Interrupt\_-
        Protection} protected types.
\end{itemize}

(Note this is the \emph{profile} distinction, not \texttt{Configurable\_Run\_Time},
which is \texttt{True} in both.) The catch is that AdaCore ships a full
\texttt{System.Interrupts} only for hosted (POSIX) targets, where it is built on
a signal-server task and is unportable to bare metal; \emph{no} bare-board
runtime (not even AdaCore's own \texttt{ravenscar-full}) implements the
full surface because Ravenscar forbids dynamic attachment in the first place.
For a long time, then, a \texttt{full}-profile build of an \texttt{Ada.Interrupts}
handler simply failed with \texttt{construct not allowed in this configuration}
and \texttt{System.Interrupts.Register\_Interrupt\_Handler not defined}.

\subsection{The bare-board full \texttt{System.Interrupts}}\index{umbrella handler}

That gap is now closed. \texttt{full\_overlay/gnarl/s-interr}
is a bare-board re-implementation of the full \texttt{System.Interrupts}
surface (\texttt{Register\_Interrupt\_Handler}, the \texttt{Static\_-} and
\texttt{Dynamic\_Interrupt\_Protection} protected bases, \texttt{Install\_-
Handlers} and the dynamic \texttt{Attach\_}/\texttt{Detach\_}/\texttt{Exchange\_-
Handler} operations) but layered on the machinery that already exists rather
than on a signal server. It only had to add the GNARL-facing surface; the hard
parts were already done and proven (the previous sections):

\begin{itemize}
  \item It installs a single \emph{umbrella} handler per interrupt through the
        low-level \texttt{System.OS\_Interface.Attach\_Handler}
        (= \texttt{System.BB.Interrupts}), then redirects it by re-pointing a
        handler table. The bare-board primitive accepts only \emph{one} attach
        per source, so the umbrella is installed once and the table does the rest.
  \item The interrupt's run priority is fixed by its Xtensa level, so the
        controller-programming step ignores the priority argument
        (\texttt{Install\_Interrupt\_Handler} marks it \texttt{Unreferenced}).
        That is what makes the dynamic \texttt{Attach\_Handler} path, which
        carries no priority, safe.
  \item Interrupt \emph{entries} and the POSIX.5 signal services have no
        bare-board meaning and raise \texttt{Program\_Error}.
\end{itemize}

The overlay file replaces the restricted \texttt{s-interr} for the \texttt{full}
profile \emph{only}; \texttt{embedded} and \texttt{light-tasking} keep the
restricted one they need for the Jorvik path. So the same protected-handler
source now builds on all three profiles.

There is one wrinkle, and it is subtler than it first looks. The full
\texttt{System.Interrupts} alone is necessary but not \emph{sufficient}: the
\emph{static} \texttt{pragma Attach\_Handler} path reaches it directly (the
compiler lowers the handler to \texttt{Install\_Handlers}), but the \emph{dynamic}
operations go through \texttt{Ada.Interrupts}, and the bare runtime inherited the
\emph{restricted} (\texttt{No\_Dynamic\_Attachment}) \texttt{Ada.Interrupts} body
from the Jorvik base, in which every operation is simply \texttt{raise
Program\_Error}. So for a while static attach worked while
\texttt{Ada.Interrupts.Attach\_Handler}, \texttt{Detach\_Handler},
\texttt{Exchange\_Handler}, \texttt{Current\_Handler}, \texttt{Is\_Attached} and
\texttt{Reference} all raised \texttt{Program\_Error} (the full ACATS sweep caught
this: \texttt{CXC3007}). The overlay now also supplies the \emph{delegating}
\texttt{a-interr.adb} that forwards each operation to \texttt{System.Interrupts},
so the full dynamic API works (\texttt{CXC3007} HW-verified \texttt{PASSED}).

The genuine residual limitation is narrow: because \texttt{Configurable\_Run\_Time}
is \texttt{True}, the front end does not generate library-level finalization, so a
\emph{library-level} interrupt protected object's \texttt{Finalize} (which would
restore the previously installed handlers at program shutdown) does not run
automatically. Explicit \texttt{Detach\_Handler} works; only the implicit
restore-on-finalization of a library-level handler PO is absent. That is harmless
on an embedded target that never exits, and the low-level
\texttt{System.BB.Interrupts.Attach\_Handler} primitive remains available on all
profiles as the no-finalization escape hatch.

\section{Exercising the vectors}\index{CPU\_INT}

Because a bug in the vectors (a corrupted context, a dropped interrupt) is exactly
the kind of thing that hides until the worst moment, the
\texttt{esp32s3\_intr\_levels} example is a standalone regression test for them. It
fires the level-2 and level-3 device vectors with \emph{no external wiring} by
routing the \texttt{FROM\_CPU} interrupt-matrix sources (the same mechanism as the
cross-core poke) to \texttt{Device\_L2\_0} and \texttt{Device\_L3\_0} and asserting
them from software, while a low-priority victim task holds floating-point
accumulators and a \texttt{THREADPTR} sentinel live in registers, continually
preempted by those two vectors \emph{and} the timer tick. If any vector failed
to preserve the interrupted context, an accumulator or the sentinel comes back
wrong. On hardware the handler counts climb and the corruption count stays zero,
confirming all three vectors both dispatch correctly and preserve context.
(Level~4 has no vector to test; level~1 carries no asynchronous interrupt on this
port.)

The companion \texttt{esp32s3\_full\_intr} example is the same test built under
the \texttt{full} profile. Its job is narrower: to prove that the identical
\texttt{pragma Attach\_Handler} protected objects now build, bind, link \emph{and
run} against the bare-board full \texttt{System.Interrupts} described above:
the exact source that used to fail with \texttt{Register\_Interrupt\_Handler not
defined} on \texttt{full}. On hardware it boots clean on both cores and the L2
and L3 handler counts climb together with the clean-batch counter, with zero
context-loss and zero faults. So on \texttt{full}, too, the portable
\texttt{Ada.Interrupts} protected-handler form dispatches correctly and preserves
the preempted context.

\chapter{Full-Profile Limitations (and Why They Rarely Bite)}

The \texttt{full} profile is a complete GNARL tasking runtime rebuilt on bare
metal: no Ravenscar restrictions, with rendezvous, protected objects with
entries, dynamic and nested tasks, dynamic priorities, task attributes,
exception propagation, \texttt{abort}, and \texttt{Ada.Interrupts} handlers all
running on hardware. A handful of rough edges remain. This chapter lists them
honestly and, for each, explains why it seldom matters for a real embedded
program. The short version: the gaps cluster at the margins of the language
(constructs a startup-time, fixed-task-set application almost never reaches), and
each has a clean workaround or is explicitly permitted by the Reference Manual.

\section{Toolchain-bound: asynchronous \texttt{select} with a relative delay}\index{asynchronous transfer of control}\index{select statement}

A \texttt{select \ldots\ delay D; then abort \ldots\ end select} (the relative,
\texttt{Duration}-triggered form of asynchronous transfer of control, ACATS
\texttt{C974*}) does not compile under \texttt{full}. This is \emph{not} a runtime
gap: under \texttt{Configurable\_Run\_Time} the GNAT front end emits a call that
passes a \texttt{Duration} to \texttt{Enqueue\_RT (T : Real\_Time.Time)}, which,
crossed with this port's \texttt{type Time}, does not type-check. Fixing it is a
compiler change or a deep \texttt{Real\_Time.Time} redesign, not an RTS unit we
can add.

\textbf{Why it rarely bites.} The \emph{absolute} form (\texttt{select \ldots\
delay until T}) and ATC triggered by an \emph{entry call} both work. A relative
timeout is sugar for ``compute \texttt{Clock + D}, then wait until that instant'',
which the absolute form expresses directly. Timed operations are also commonly
written with a guard plus \texttt{Ada.Real\_Time} arithmetic, untouched by this.

\section{RM-permitted: aborting a pure CPU loop}\index{abort statement}

\texttt{abort} of a task blocked at any \emph{abort-completion point}
(a \texttt{delay}, an entry call, a rendezvous) works and is HW-verified
(\texttt{'Terminated} flips, the task unwinds cleanly). The one case that does
\emph{not} take effect is a task spinning in a tight computational loop with no
completion point at all: it keeps running until it next reaches one.

\textbf{Why it rarely bites.} This is explicitly allowed: ARM 9.8 only
requires abort to take effect at completion points, and a loop with no such point
is permitted to run on. Delivering it anyway would mean injecting an exception
into a preempted task from the timer ISR, which (as the context-switch and
interrupts chapters describe) is deep, windowed-Xtensa, shared-kernel surgery
for a behavior the standard does not mandate. In practice any loop that
does I/O, delays, or calls an entry is abortable; a genuinely tight numeric
kernel that must be cancellable can poll a \texttt{Volatile} flag.

\section{\texttt{Configurable\_Run\_Time} finalization}\index{finalization}\index{Configurable\_Run\_Time}\index{Ada.Interrupts}

The front end does not generate \emph{library-level} finalization because
\texttt{system.ads} sets \texttt{Configurable\_Run\_Time => True}. Two consequences:

\begin{itemize}
  \item A \emph{library-level} controlled object's \texttt{Finalize} runs only if
        the object is finalized at scope exit, not as part of program shutdown.
  \item Consequently, the \emph{implicit} restore-on-finalization of a
        \emph{library-level} interrupt-handler protected object does not run
        automatically at shutdown. (Both the static \texttt{pragma Attach\_Handler}
        path \emph{and} the full dynamic \texttt{Ada.Interrupts} operations
        (\texttt{Attach\_Handler}, \texttt{Detach\_Handler}, \texttt{Exchange\_-
        Handler}, \texttt{Current\_Handler}) are now supported and
        HW-verified; see the Interrupts chapter. Only this one implicit-restore
        case is outstanding.)
\end{itemize}

\textbf{Why it rarely bites.} Scope-level finalization -- the overwhelmingly
common case (a controlled object in a subprogram or block) -- works normally.
Library-level objects live for the entire life of the program anyway, so on an
embedded target that never exits there is nothing to finalize. And interrupt
handlers are wired up once at elaboration and stay for the life of the system, so
the missing implicit restore-at-shutdown never matters in practice; explicit
\texttt{Detach\_Handler} works if you do need it.

\section{\texttt{Ada.Calendar} resolution: 5\,ns, with the full range to 2399}\index{Ada.Calendar}

The bare \texttt{Ada.Calendar} originally declared \texttt{Year\_Number} as
\texttt{1901\,..\,2099} (the Ada~95 bound), whereas Ada~2005 onward (and the
modern ACATS \texttt{C96004}) require \texttt{2399}. The obstacle was the time
representation: the Calendar uses a \emph{Julian-day fixed-point} base sized so the
span fits in signed 64 bits at sub-second precision, and stretching the range to
2399 (roughly five centuries) overflows that base at 1\,ns granularity, which is
why mainline GNAT eventually abandoned the Julian-day Calendar for an
integer-nanosecond one.

Rather than swap in a new time base (a risky rewrite of a unit shared by the
\texttt{embedded} and \texttt{full} profiles), the range is bought by
\emph{coarsening} the existing base by exactly $5\times$: the sub-day fixed-point
scale changes from \texttt{Duration'Small\,/\,86\_400.0} to
\texttt{/\,17\_280.0}\,($=86\_400/5$), and the Modified-Julian-Day and Julian-Day
type ranges widen to cover 2399-12-31 (\texttt{Time'Last} compile-checks to
\texttt{197\_641}). The 1858 epoch and the day-arithmetic algorithm are untouched.
The trade is precision for range: the Calendar's resolution becomes $5$\,ns instead
of $1$\,ns, and in exchange \texttt{Year\_Number} now runs the full
\texttt{1901\,..\,2399}. \texttt{C96004} \textbf{PASSES} (HW-verified, both
profiles), and the change is applied by \texttt{gen\_runtime.sh} when it
materialises the runtime.

\textbf{Why the 5\,ns rarely bites.} \texttt{Ada.Calendar} is a wall-clock: its job
is civil date-and-time stamping, where a five-nanosecond quantum on a value that
counts whole seconds is meaningless. Any code that genuinely needs sub-microsecond
timing uses \texttt{Ada.Real\_Time} instead, whose \texttt{Time\_Span} keeps its
full \texttt{Duration'Small} (1\,ns) precision and is untouched by this change.
Arithmetic, formatting, \texttt{Split}/\texttt{Time\_Of}, and \texttt{delay until}
on a \texttt{Calendar.Time} remain HW-verified across the whole 1901--2399 span.

\section{Per-task heap reclamation under task churn}\index{task churn}

A terminated task's primary stack is returned to the heap, but a program that
\emph{churns} tasks (creating and destroying them in a tight loop) can still
exhaust the heap: the reclamation of the per-task resources is not yet fully
race-free against the terminating thread, and the bare allocator fragments. The
underlying cause is that the kernel has no explicit ``retire this thread'' step;
the design for fixing it is recorded in the project notes.

\textbf{Why it rarely bites.} Idiomatic embedded Ada creates a \emph{fixed} set
of tasks at startup and runs them for the life of the system. The fixed-set
case never churns and never leaks. Dynamic create/destroy in a hot loop is an
unusual pattern on a microcontroller; where it is needed, a task pool (allocate
once, reuse) sidesteps the issue entirely.

\section{Minor robustness edges}

\begin{itemize}
  \item \textbf{Lowering the running task's priority below a ready peer} can hang
        the caller (a bare-kernel reschedule edge). Raising priority, or staying
        the highest, works. \emph{Why it rarely bites:} priority is normally set
        once at task creation; the dynamic-lowering-below-peers pattern is rare,
        and priority \emph{ceilings} (the common dynamic case) go up, not down.
  \item \textbf{\texttt{Ada.Text\_IO} over the USB-Serial-JTAG console} is not
        robust for two tasks printing concurrently. \emph{Why it rarely bites:}
        this is a console-transport detail, not the tasking kernel. Route
        diagnostic output through one task, or use a real UART (the UART driver
        is task-safe).
\end{itemize}

\section{Architectural: \texttt{Interrupt\_Priority'Last} and the kernel tick}\index{interrupt priority}

A task declared at \texttt{Interrupt\_Priority'Last} masks the kernel's own
Xtensa level-5 tick and cross-core IPI, deadlocking its activation (ACATS
\texttt{CXD1006}; see the ACATS chapter). The kernel owns the top hardware level
and shares it with the highest Ada interrupt priority.

\textbf{Why it rarely bites.} Only the single \emph{topmost} interrupt priority
collides. Any priority below it (including every priority the peripheral HAL
drivers use for their device interrupts) is unaffected, and ordinary tasks use
\texttt{Priority}, not \texttt{Interrupt\_Priority}, at all. A handler simply
should not claim the one ceiling reserved for the scheduler.

\section{The common case is clean}

None of the above touches the path a typical full-tasking program actually walks:
declare a fixed set of tasks and protected objects at startup; rendezvous and
share data through protected entries; use absolute delays and timing events;
attach interrupt handlers statically; finalize controlled objects at scope exit;
abort at completion points. That path is implemented and hardware-verified. The
limitations are real, but they live at the edges (toolchain corners, an
RM-permitted abort case, and patterns that idiomatic embedded code seldom uses:
library-level finalization, task churn, dynamic priority lowering), and every
one has either a clean workaround or standing permission from the language.


\appendix
\chapter{Quick Reference}
\label{ch:appendix}

\section{Usable GPIO pins}
\label{app:pins}

The HAL's \texttt{Pin\_Id} subtype (\S\ref{sec:subtypes}) admits only the pads the
ESP32-S3 actually exposes. Everything else is a compile-time error.

\begin{center}
\begin{tabular}{l l}
\textbf{Pads} & \textbf{Status} \\ \hline
\texttt{0 .. 21}   & usable (general-purpose) \\
\texttt{22 .. 25}  & absent on the package \\
\texttt{26 .. 32}  & SPI flash --- driving them hangs the chip \\
\texttt{33 .. 37}  & octal PSRAM --- driving them hangs the chip \\
\texttt{38 .. 48}  & usable (general-purpose) \\
\end{tabular}
\end{center}

So \texttt{Pin\_Id} is \texttt{0\,..\,21 | 38\,..\,48}, and \texttt{Optional\_Pin}
adds the sentinel \texttt{No\_Pin} ($=-1$) for ``leave this line unrouted.''

\section{Runtime profiles at a glance}
\label{app:profiles}

\begin{center}
\begin{tabular}{l l l l}
\textbf{Feature} & \textbf{light-tasking} & \textbf{embedded} & \textbf{full} \\ \hline
Tasks, protected objects, \texttt{delay until} & yes & yes & yes \\
Exception propagation (ZCX)        & no  & yes & yes \\
Controlled types / finalization (RAII) & no & yes & yes \\
Secondary stack                    & no  & yes & yes \\
Heap                               & bump only & TLSF (O(1)) & TLSF (O(1)) \\
\texttt{select} / rendezvous / \texttt{abort} & no & no & yes \\
RAII-handle HAL drivers (SPI, I2S, \ldots) & no & yes & yes \\
Smallest footprint                 & yes & --- & --- \\
\end{tabular}
\end{center}

Choose \texttt{light-tasking} for the smallest, most analysable build,
\texttt{embedded} for most applications (the default), and \texttt{full} only when
you need dynamic tasking, rendezvous or \texttt{abort}
(Chapter~\ref{ch:profiles}).

\section{Example index}
\label{app:examples}

Every example builds and runs with \texttt{./x run <name> -p /dev/ttyACM0}; the
profile column is its default.

\begin{center}\small
\begin{longtable}{@{}l l l@{}}
\textbf{Name} & \textbf{Profile} & \textbf{Shows} \\ \hline
\endhead
\texttt{gpio0\_blink}    & light-tasking & GPIO output; the minimal app \\
\texttt{heartbeat}       & light-tasking & a periodic task \\
\texttt{psram}           & light-tasking & external PSRAM bring-up \\
\texttt{smp}             & light-tasking & dual-core scheduling \\
\texttt{embedded}        & embedded & tagged dispatch, finalization, exceptions \\
\texttt{exceptions}      & embedded & local / propagated / unhandled exceptions \\
\texttt{crypto}          & embedded & SHA + AES accelerators \\
\texttt{adc\_read}       & embedded & SAR ADC one-shot \\
\texttt{touch\_read}     & embedded & capacitive touch \\
\texttt{i2c\_loopback}   & embedded & I2C master self-test \\
\texttt{uart\_loopback}  & embedded & UART internal loopback \\
\texttt{i2s\_loopback}   & embedded & I2S full-duplex DMA loopback \\
\texttt{i2s\_pdm}        & embedded & I2S PDM (PCM\,$\leftrightarrow$\,PDM) \\
\texttt{gdma\_copy}      & embedded & memory-to-memory DMA \\
\texttt{ledc\_pwm}, \texttt{mcpwm\_pwm} & embedded & PWM outputs \\
\texttt{rmt\_loopback}   & embedded & RMT pulse TX/RX \\
\texttt{pcnt\_count}     & embedded & pulse / edge counter \\
\texttt{sdm\_output}     & embedded & sigma-delta density output \\
\texttt{twai\_loopback}  & embedded & CAN (standard, extended, RTR) \\
\texttt{timer\_count}    & embedded & general-purpose timers \\
\texttt{lcd\_i8080}      & embedded & 8-bit i80 parallel LCD \\
\texttt{rtc\_sleep}, \texttt{rtcio\_hold} & embedded & deep sleep, RTC pad hold \\
\texttt{sd\_spi}, \texttt{sdmmc}, \texttt{ext4} & embedded & SD cards + ext4 filesystem \\
\texttt{intr\_levels}    & embedded & interrupt levels/priorities \\
\texttt{full\_tasking}, \texttt{rendezvous}, \texttt{full\_intr} & full & rendezvous, dynamic tasking \\
\end{longtable}
\end{center}

\chapter{Internal Peripheral Drivers}
\label{app:internal}

This appendix is a per-driver reference for the SoC's \emph{on-chip} peripherals
--- the blocks built into the ESP32-S3 itself (GPIO, the serial buses, the
timers and PWM units, DMA, the analog and audio paths, the low-power domain, and
the crypto accelerators). The companion appendix~\ref{app:external} covers the
drivers for \emph{external} parts wired to these peripherals. The narrative HAL
chapters (\S\ref{ch:hal} onward) introduce the conventions and the design
rationale; this appendix collects, for each driver, a short description, a digest
of its public API, and a worked example drawn from the \texttt{examples/} tree.

A few patterns recur and are worth stating once. The bus drivers (UART, SPI,
I2C, I2S, LCD, TWAI) seal their unsynchronised register access in a private
\texttt{.Engine} child and hand the application an acquired, auto-releasing
\texttt{Session}; the resource drivers (LEDC, MCPWM, Timer, PCNT, SDM, RMT, GDMA,
ADC) hand out a \emph{limited, controlled} handle that claims hardware on
\texttt{Claim}/\texttt{Acquire} and frees it on scope exit, even during exception
unwinding. Both idioms depend on finalization, so those drivers require the
\texttt{embedded} or \texttt{full} profile; the register-poke drivers (GPIO,
RNG, RTC, RTC-IO, Touch, SHA, AES, SDMMC) work under every profile, and each
section notes which case applies. See the driver-design chapter
(\S\ref{ch:driver-design}) for why the ownership gateway is built this way.

\section{GPIO --- General-Purpose I/O}\index{GPIO driver}
\label{app:int:gpio}

The ESP32-S3 routes every digital pad through a \emph{GPIO matrix}: a pad can be
a software-driven GPIO or be cross-connected to any peripheral's input/output
signal. The Ada driver \texttt{ESP32S3.GPIO} sits on top of the generated
register layer (\texttt{Interfaces.ESP32S3.GPIO} plus \texttt{IO\_MUX}) and
exposes the plain-GPIO case: configure a pad's direction, pull and drive
strength, then \texttt{Set} / \texttt{Clear} / \texttt{Toggle} / \texttt{Write}
/ \texttt{Read} it. \texttt{Configure} always selects the matrix path
(\texttt{MCU\_SEL = 1}, output index 256); routing a pad to a \emph{peripheral}
signal instead is the job of that peripheral's own \texttt{Configure\_Pins}.

The pad space is modelled by \texttt{Pad\_Number} (\texttt{-1\,..\,48}), with
\texttt{-1} the \texttt{No\_Pin} sentinel for an optional unrouted line. The
subtype \texttt{Pin\_Id} carries a \emph{static} predicate that excludes the
pads the silicon does not bring out (22\,..\,25) and those bonded to the
in-package flash and octal PSRAM (26\,..\,37): naming one of those as a static
value is a compile-time error, and a dynamic value is caught at run time under
\texttt{-gnata}. \texttt{Set} / \texttt{Clear} / \texttt{Write} use the
hardware-atomic W1TS/W1TC banks and \texttt{Read} is a pure load, so those are
concurrency-safe as written; the read-modify-write paths (\texttt{Configure},
\texttt{Toggle}) are serialised through a protected object, so the package needs
a tasking runtime. A companion child, \texttt{ESP32S3.GPIO.Interrupts}, owns the
single GPIO interrupt source, demultiplexes it by latched status, and dispatches
a per-pin \texttt{Callback} that runs in interrupt context under the level-3
ceiling.

\begin{lstlisting}
package ESP32S3.GPIO is
   type Pad_Number is range -1 .. 48;
   No_Pin : constant Pad_Number := -1;

   subtype Pin_Id is Pad_Number range 0 .. 48 with
     Static_Predicate => Pin_Id in 0 .. 21 | 38 .. 48;
   subtype Optional_Pin is Pad_Number with
     Static_Predicate => Optional_Pin in -1 | 0 .. 21 | 38 .. 48;

   type Pin_Mode       is (Input, Output);
   type Pull_Mode      is (Floating, Pull_Up, Pull_Down);
   type Drive_Strength is (Drive_Weak, Drive_Medium, Drive_Strong, Drive_Strongest);

   procedure Configure (Pin : Pin_Id; Mode : Pin_Mode;
                        Pull  : Pull_Mode      := Floating;
                        Drive : Drive_Strength := Drive_Medium);
   procedure Set    (Pin : Pin_Id);   --  atomic W1TS
   procedure Clear  (Pin : Pin_Id);   --  atomic W1TC
   procedure Toggle (Pin : Pin_Id);
   procedure Write  (Pin : Pin_Id; On : Boolean);
   function  Read   (Pin : Pin_Id) return Boolean;
end ESP32S3.GPIO;

--  Pin interrupts (ESP32S3.GPIO.Interrupts):
type Trigger  is (Rising_Edge, Falling_Edge, Any_Edge, Low_Level, High_Level);
type Callback is access procedure;
procedure Enable  (Pin : Pin_Id; On : Trigger; Action : Callback);
procedure Disable (Pin : Pin_Id);
\end{lstlisting}

The \texttt{gpio0\_blink} example drives GPIO0 from a library-level task ---
configure once as a push-pull output, then toggle on a fixed period:

\begin{lstlisting}
with ESP32S3.GPIO;
with Ada.Real_Time; use Ada.Real_Time;

Pin : constant ESP32S3.GPIO.Pin_Id := 0;
task Blinker;
task body Blinker is
   Period : constant Time_Span := Milliseconds (250);   --  ~2 Hz
   Next   : Time;
   High   : Boolean := False;
begin
   ESP32S3.GPIO.Configure
     (Pin, ESP32S3.GPIO.Output, Drive => ESP32S3.GPIO.Drive_Strong);
   Next := Clock + Period;
   loop
      delay until Next;
      High := not High;
      ESP32S3.GPIO.Write (Pin, High);
      Next := Next + Period;
   end loop;
end Blinker;
\end{lstlisting}

\cnote{ESP-IDF wraps the same matrix in \texttt{gpio\_config()} +
\texttt{gpio\_set\_level()}, taking a raw \texttt{int} pin number checked only at
run time. The Ada \texttt{Pin\_Id} predicate moves that check to compile time for
the common static-pin case --- a flash/PSRAM pad never even links.}

\noindent\textbf{Example:} \texttt{examples/esp32s3\_gpio0\_blink}.

\section{UART --- Asynchronous Serial}\index{UART driver}
\label{app:int:uart}

The chip has three general-purpose UART controllers; the bare runtime uses
USB-Serial-JTAG for its console, so even UART0's pads are free to repurpose.
\texttt{ESP32S3.UART} is the only interface the application sees: the raw
register driver lives in a private child, \texttt{ESP32S3.UART.Engine}, that
cannot be \texttt{with}-ed from outside the subtree, so the unsynchronised
primitives can never be reached by accident.

Configuration is split in two. \texttt{Setup} runs once per port at startup
(single-threaded) and fixes baud, frame format and pin routing --- all four
lines (\texttt{Tx}, \texttt{Rx}, \texttt{Rts}, \texttt{Cts}) are
\texttt{Optional\_Pin}, so a link routes only what it uses. Everything after
that goes through an acquired \texttt{Session}: \texttt{Acquire} hands out a
limited, non-copyable token that owns the port exclusively (other tasks suspend
until it is released), and every transfer \emph{and} every later
reconfiguration (\texttt{Set\_Baud}, \texttt{Configure\_Pins},
\texttt{Set\_Inversion}, \ldots) is a method on the held \texttt{Session} --- so
a reconfigure can never race another task's transfer. The \texttt{Session} is a
controlled type that auto-releases on scope exit (including during exception
unwinding), so this driver, like the other task-safe HAL drivers, targets the
\texttt{embedded}/\texttt{full} profile (full exceptions and finalization).
\texttt{Enable\_Loopback} ties TX back to RX inside the controller for a
no-wiring self-test.

\begin{lstlisting}
package ESP32S3.UART is
   type UART_Port  is (UART0, UART1, UART2);
   subtype Baud_Rate is Positive range 300 .. 5_000_000;
   type Data_Bits   is range 5 .. 8;
   type Parity_Mode is (None, Even, Odd);
   type Stop_Bits   is (One, Two);
   type Byte        is new Interfaces.Unsigned_8;
   type Byte_Array  is array (Natural range <>) of Byte;
   type Session     is limited private;            --  auto-releasing

   procedure Setup (Port : UART_Port;
                    Baud : Baud_Rate   := 115_200;
                    Bits : Data_Bits   := 8;
                    Parity : Parity_Mode := None;
                    Stop   : Stop_Bits   := One;
                    Tx, Rx, Rts, Cts : ESP32S3.GPIO.Optional_Pin
                                       := ESP32S3.GPIO.No_Pin;
                    Rx_Flow_Threshold : Natural := 100);

   procedure Acquire (S : in out Session; Port : UART_Port);
   procedure Write   (S : Session; Data : Byte_Array);
   procedure Read    (S : Session; Data : out Byte_Array; Count : out Natural);
   function  Available (S : Session) return Natural;
   procedure Enable_Loopback (S : Session; On : Boolean := True);
   procedure Release (S : in out Session);
   Not_Initialized, Not_Owned : exception;
end ESP32S3.UART;
\end{lstlisting}

The \texttt{uart\_loopback} self-test brings up UART1, ties TX to RX internally,
and round-trips a known buffer --- no wiring needed:

\begin{lstlisting}
with ESP32S3.UART; use ESP32S3.UART;

Port : constant UART_Port := UART1;
Tx   : constant Byte_Array := (16#55#, 16#AA#, 16#00#, 16#FF#, 16#12#, 16#34#);
S    : Session;
Rx   : Byte_Array (Tx'Range);
Got  : Natural;
begin
   Setup (Port, Baud => 115_200);     --  8-N-1 defaults
   Acquire (S, Port);
   Enable_Loopback (S);               --  internal TX->RX
   Write (S, Tx);
   Read  (S, Rx, Got);
   Release (S);
   --  Got = Tx'Length and Rx = Tx  on a healthy data path
\end{lstlisting}

\noindent\textbf{Example:} \texttt{examples/esp32s3\_uart\_loopback}.

\section{SPI --- Serial Peripheral Interface}\index{SPI driver}
\label{app:int:spi}

\texttt{ESP32S3.SPI} is a task-safe, full-duplex \emph{master} over the two
general-purpose hosts, \texttt{SPI2} and \texttt{SPI3} (\texttt{SPI0}/\texttt{SPI1}
drive the in-package flash and PSRAM and are deliberately not offered). As with
UART, the unsynchronised register driver is sealed in a private child,
\texttt{ESP32S3.SPI.Engine}, and reached only through an acquired
\texttt{Session}.

\texttt{Setup} brings a host up at one of the four clock modes (CPOL/CPHA 0\,..\,3)
and a bit clock from roughly 80\,kHz to 80\,MHz, and claims a GDMA channel;
\texttt{Set\_Clock} re-tunes the rate without re-claiming DMA, which lets a
device driver init slowly then run fast (an SD card negotiates at
$\le$400\,kHz, then switches to several MHz). \texttt{Configure\_Pins} routes
\texttt{Sclk}/\texttt{Mosi}/\texttt{Miso}/\texttt{Cs} to pads (each an
\texttt{Optional\_Pin}), and \texttt{Enable\_Loopback} ties MOSI to MISO through
one pad for a self-test. The single transfer primitive is a blocking,
full-duplex DMA \texttt{Transfer} of 1\,..\,4095 bytes between two SRAM buffers,
addressed by \texttt{System.Address}. A device that brings its own chip-select ---
to share the bus, or behind a 3:8 decoder or an I/O-expander --- registers a
library-level \texttt{CS\_Select} callback at \texttt{Acquire} and brackets its
transaction with \texttt{Select\_Device}; the driver then suppresses the hardware
\texttt{CS0} for that session. The controlled \texttt{Session} auto-releases on
scope exit, so this driver targets \texttt{embedded}/\texttt{full}.

\begin{lstlisting}
package ESP32S3.SPI is
   type SPI_Host is (SPI2, SPI3);
   subtype SPI_Mode is Natural range 0 .. 3;
   type Session is limited private;            --  auto-releasing

   procedure Setup       (Host : SPI_Host; Mode : SPI_Mode := 0;
                          Clock_Hz : Positive := 1_000_000);
   procedure Set_Clock   (Host : SPI_Host; Hz : Positive);
   procedure Enable_Loopback (Host : SPI_Host; Pad : ESP32S3.GPIO.Pin_Id);
   type CS_Select is access procedure            --  app chip select: library-
     (Ctx : System.Address; Active : Boolean);   --  level, closure-free

   procedure Configure_Pins (Host : SPI_Host;
                             Sclk, Mosi, Miso : ESP32S3.GPIO.Optional_Pin;
                             Cs   : ESP32S3.GPIO.Optional_Pin := No_Pin);

   procedure Acquire  (S : in out Session; Host : SPI_Host;
                       Select_CB : CS_Select := null;
                       Ctx : System.Address := System.Null_Address);
   procedure Select_Device (S : in out Session; On : Boolean);
   procedure Transfer (S : Session; Tx, Rx : System.Address; Length : Natural);
   procedure Release  (S : in out Session);
   Not_Initialized, Not_Owned : exception;
end ESP32S3.SPI;
\end{lstlisting}

A no-wiring round-trip: bring up SPI2, loop MOSI to MISO on a free pad, and DMA
a buffer through it --- what comes back on \texttt{Rx} equals what went out:

\begin{lstlisting}
with ESP32S3.SPI; use ESP32S3.SPI;
with Interfaces;

Tx : aliased array (0 .. 7) of Interfaces.Unsigned_8
       := (16#01#, 16#23#, 16#45#, 16#67#, 16#89#, 16#AB#, 16#CD#, 16#EF#);
Rx : aliased array (Tx'Range) of Interfaces.Unsigned_8 := (others => 0);
S  : Session;
begin
   Setup (SPI2, Mode => 0, Clock_Hz => 1_000_000);
   Enable_Loopback (SPI2, Pad => 8);          --  MOSI -> MISO on GPIO8
   Acquire (S, SPI2);
   Transfer (S, Tx'Address, Rx'Address, Tx'Length);
   Release (S);                                --  Rx now mirrors Tx
\end{lstlisting}

The same driver carries real traffic in \texttt{esp32s3\_sd\_spi} (an SD card on
\texttt{SPI2}, init at 400\,kHz then data at 8\,MHz) and behind the ST7789
display driver.

\noindent\textbf{Example:} \texttt{examples/esp32s3\_sd\_spi} (and the ST7789
display) exercise this driver against real hardware.

\section{I2C --- Two-Wire Master}\index{I2C driver}
\label{app:int:i2c}

\texttt{ESP32S3.I2C} is a polled, single-segment master over the two
controllers \texttt{I2C0} and \texttt{I2C1}. Each \texttt{Write} or
\texttt{Read} is one complete \texttt{START}\,\ldots\,\texttt{STOP} transaction
of up to \texttt{Max\_Transfer} (32) bytes --- the controller's FIFO depth.
\texttt{Setup} fixes the SCL bit clock (100\,kHz standard / 400\,kHz fast are
typical), and \texttt{Configure\_Pins} routes \texttt{Scl}/\texttt{Sda} as
open-drain lines with the internal pull-ups on, enough for bring-up (add real
resistors for a production bus). The same sealed-\texttt{Engine},
\texttt{Acquire}-a-\texttt{Session} structure as the other buses applies, and the
controlled \texttt{Session} auto-releases, so the package targets
\texttt{embedded}/\texttt{full}.

A 7-bit \texttt{Slave\_Address} is given without the R/W bit (the driver adds
it). \texttt{Write} reports \texttt{Success} = \texttt{True} only if the slave
ACKed the address and every byte; a zero-length \texttt{Write} is an
address-only probe, which makes bus scanning trivial, and \texttt{Check\_Ack =>
False} clocks the frame out regardless of ACK. Unlike UART/SPI there is no
on-chip loopback: I2C's SDA is a bidirectional wired-AND node and the matrix
gives each pad one output source, so verifying the read/ACK direction needs a
real device or a jumper.

\begin{lstlisting}
package ESP32S3.I2C is
   type I2C_Host is (I2C0, I2C1);
   subtype Slave_Address is Natural range 0 .. 16#7F#;
   type Byte       is new Interfaces.Unsigned_8;
   type Byte_Array is array (Natural range <>) of Byte;
   Max_Transfer : constant := 32;
   type Session  is limited private;            --  auto-releasing

   procedure Setup          (Host : I2C_Host; Clock_Hz : Positive := 100_000);
   procedure Configure_Pins (Host : I2C_Host;
                             Scl, Sda : ESP32S3.GPIO.Pin_Id);

   procedure Acquire (S : in out Session; Host : I2C_Host);
   procedure Write   (S : Session; Addr : Slave_Address; Data : Byte_Array;
                      Success : out Boolean; Check_Ack : Boolean := True);
   procedure Read    (S : Session; Addr : Slave_Address; Data : out Byte_Array;
                      Success : out Boolean);
   procedure Release (S : in out Session);
   Not_Initialized, Not_Owned : exception;
end ESP32S3.I2C;
\end{lstlisting}

The \texttt{i2c\_loopback} self-test brings up \texttt{I2C0} and probes an empty
bus: an ACK-checked write to an absent address must come back NACKed
(\texttt{Success = False}), proving START, addressing and ACK detection on
silicon:

\begin{lstlisting}
with ESP32S3.I2C; use ESP32S3.I2C;

Host   : constant I2C_Host      := I2C0;
Absent : constant Slave_Address := 16#55#;   --  nothing answers here
S      : Session;
Ok     : Boolean;
begin
   Setup (Host, Clock_Hz => 100_000);
   Configure_Pins (Host, Scl => 6, Sda => 4);
   Acquire (S, Host);
   Write (S, Absent, (1 => 16#00#), Ok);     --  address-only-ish probe
   Release (S);
   --  PASS when (not Ok): the master saw the NACK from the empty bus
\end{lstlisting}

\noindent\textbf{Example:} \texttt{examples/esp32s3\_i2c\_loopback}.

\section{RNG --- Hardware Random Number Generator}\index{RNG driver}
\label{app:int:rng}

\texttt{ESP32S3.RNG} reads the on-chip random number generator: each load of its
data register yields a fresh 32-bit word, drawn from analog and clock-jitter
noise. The package is deliberately the simplest in the HAL --- \texttt{Read} is a
single atomic register load and \texttt{Fill} writes only the caller's buffer,
so there is no shared mutable state and \emph{no} protected object. It is marked
\texttt{Preelaborate} and stays usable from a ZFP (no-runtime) context as well
as the tasking profiles; concurrent readers simply get independent words.

Two caveats carry over from the spec. First, this is the RNG \emph{peripheral}
register, not ESP-IDF's \texttt{WDEV\_RND\_REG} --- that one only yields entropy
with the RF clock domain up (Wi-Fi/BT), which the bare runtime never starts, so
it would read a constant here. Second, without an active hardware entropy source
(RF, or SAR-ADC bootloader entropy) the output is fine for dithering, non-secret
IDs, test data or seeding a software PRNG, but should \emph{not} be treated as a
CSPRNG. Don't poll faster than the hardware refreshes, or successive words may
correlate.

\begin{lstlisting}
package ESP32S3.RNG with Preelaborate is
   subtype Word is ESP32S3_Registers.UInt32;
   function Read return Word with Inline;       --  one fresh 32-bit value

   type Byte       is mod 2 ** 8 with Size => 8;
   type Byte_Array is array (Natural range <>) of Byte with Pack;
   procedure Fill (Buffer : out Byte_Array);    --  any length, word at a time
end ESP32S3.RNG;
\end{lstlisting}

Reading a word and filling an arbitrary-length buffer:

\begin{lstlisting}
with ESP32S3.RNG;

W   : ESP32S3.RNG.Word;
Key : ESP32S3.RNG.Byte_Array (0 .. 11);    --  12 random bytes (3 words + 0 tail)
begin
   W := ESP32S3.RNG.Read;        --  e.g. seed a software PRNG with W
   ESP32S3.RNG.Fill (Key);       --  fill the whole buffer with random bytes
\end{lstlisting}

\noindent\textbf{Example:} no dedicated example; the driver is a one-call read
(\texttt{ESP32S3.RNG.Read}) usable directly from any profile, ZFP included.

\section{Temperature --- On-Chip Die Sensor}\index{Temperature driver}
\label{app:int:temperature}

\texttt{ESP32S3.Temperature} reads the SoC's internal temperature sensor. It
reports \emph{die} temperature, not ambient air: the chip self-heats under load,
so an idle board reads a few degrees above room temperature, with roughly
$\pm$1\,\textdegree C accuracy after the factory trim and best near the middle of
the selected range. Bring-up follows ESP-IDF's \texttt{temperature\_sensor\_ll}
sequence --- gate and pulse-reset the SAR clock, open the analog REGI2C bus,
program the DAC range (a ROM call), then power the sensor up.

\texttt{Initialize} is optional: the first \texttt{Read\_*} brings the sensor up
with the default range if you skip it, and you may call \texttt{Initialize}
again to switch range. Five \texttt{Measure\_Range} values bracket different
spans; pick the one that contains your expected die temperature for best
accuracy. Readings come as signed centi-degrees
(\texttt{Read\_Centi\_Celsius}, e.g.\ 1807 = 18.07\,\textdegree C), whole degrees
(\texttt{Read\_Celsius}), or the raw 8-bit code (\texttt{Read\_Raw}) if you
prefer your own curve. The sensor is owned by a protected object so concurrent
reads are serialised (each busy-waits microseconds for a fresh conversion under
the lock), which means the package requires a tasking runtime.

\begin{lstlisting}
package ESP32S3.Temperature is
   type Measure_Range is
     (Range_Minus10_80,    --  -10 ..  80 C  (default)
      Range_20_100,        --   20 .. 100 C
      Range_50_125,        --   50 .. 125 C
      Range_Minus30_50,    --  -30 ..  50 C
      Range_Minus40_20);   --  -40 ..  20 C

   procedure Initialize (Span : Measure_Range := Range_Minus10_80);
   function  Read_Centi_Celsius return Integer;    --  1807 = 18.07 C
   function  Read_Celsius       return Integer;    --  truncated toward zero
   function  Read_Raw           return ESP32S3_Registers.Byte;
end ESP32S3.Temperature;
\end{lstlisting}

Initialise for a room-temperature board, then read the die temperature:

\begin{lstlisting}
with ESP32S3.Temperature;

Centi : Integer;
Whole : Integer;
begin
   ESP32S3.Temperature.Initialize (ESP32S3.Temperature.Range_Minus10_80);
   Centi := ESP32S3.Temperature.Read_Centi_Celsius;   --  e.g. 4231 = 42.31 C
   Whole := ESP32S3.Temperature.Read_Celsius;         --  e.g. 42
\end{lstlisting}

\noindent\textbf{Example:} no dedicated example; the driver is a direct read
(\texttt{ESP32S3.Temperature.Read\_Celsius}) after the optional
\texttt{Initialize}.

\section{LEDC --- LED Controller (PWM)}\index{LEDC driver}
\label{app:int:ledc}

The ESP32-S3's LEDC block is the simple ``dim an LED / generate a clean PWM''
peripheral: eight low-speed channels fed by four timers. A channel selects a
timer (which fixes its frequency and duty resolution) and drives a GPIO with a
duty cycle changeable at run time. Unlike MCPWM, it has no dead-time, fault or
capture logic.

The Ada driver in \texttt{ESP32S3.LEDC} hands a channel out as a \texttt{Channel}
handle that is \emph{limited} (non-copyable, so two tasks cannot drive one
channel) and \emph{controlled} (it stops the output and returns the channel to
the free pool on scope exit, including on an exception). Because you exclusively
own a claimed channel, \texttt{Set\_Duty} needs no lock. The finalization
dependency targets the embedded/full profile, not the light-tasking build.

One chip caveat surfaces in the API: a channel uses timer \texttt{(Index mod 4)},
so channels whose indices differ by 4 (e.g.\ 0 and 4) share one frequency and
resolution. Use channels 0..3 for four independent frequencies. Achievable
frequency depends on the duty bits: $f_{\max} = 80\,\mathrm{MHz} / 2^{\texttt{Bits}}$.

\begin{lstlisting}
type Channel_Index is range 0 .. 7;
subtype Resolution    is Positive range 1 .. 14;   --  duty-cycle bits
subtype Duty_Percent  is Float range 0.0 .. 100.0;

type Channel is limited private;                   --  non-copyable, auto-releasing

procedure Claim     (C : in out Channel; Index : Channel_Index);
function  Is_Valid  (C : Channel) return Boolean;
procedure Release   (C : in out Channel);
procedure Configure (C    : in out Channel;
                     Freq : Positive;
                     Pin  : ESP32S3.GPIO.Pin_Id;
                     Bits : Resolution := 10);
procedure Set_Duty  (C : Channel; Percent : Duty_Percent);  --  loads next period
procedure Stop      (C : Channel);                          --  force output low
\end{lstlisting}

\begin{lstlisting}
declare
   Ch0    : Channel;
   Duties : constant array (1 .. 2) of Duty_Percent := (25.0, 75.0);
begin
   Claim (Ch0, 0);
   Configure (Ch0, Freq => 5_000, Pin => 4, Bits => 10);   --  5 kHz, 10-bit
   for I in Duties'Range loop
      Set_Duty (Ch0, Duties (I));
      delay until Clock + Milliseconds (5);                --  let the duty latch
      --  (the self-test now GPIO-samples pin 4 to verify duty + frequency)
   end loop;
end;   --  Ch0 finalizes: output stopped, channel freed
\end{lstlisting}

\noindent\textbf{Example:} \texttt{examples/esp32s3\_ledc\_pwm}.

\section{MCPWM --- Motor-Control PWM}\index{MCPWM driver}
\label{app:int:mcpwm}

MCPWM is the heavyweight PWM block: two units (\texttt{MCPWM0}, \texttt{MCPWM1}),
each with three generator channels and three capture channels. A generator
produces edge-aligned PWM on output A (high at period start, low when the
up-counter reaches the duty comparator). It can also drive a \emph{complementary}
pair --- an inverted B output with programmable dead-time inserted on each edge so
A and B are never high together (the half-bridge / H-bridge motor case) --- plus
carrier (chopper) modulation, fault/trip-zone shutdown, and edge capture.

The Ada driver in \texttt{ESP32S3.MCPWM} models both a generator channel
(\texttt{Channel}) and a capture channel (\texttt{Capture}) as limited, controlled
handles. A unit's clock must be brought up once with \texttt{Setup} (PWM clock =
160\,MHz) before any of its channels can be claimed; claiming a channel of an
un-set-up unit raises \texttt{Not\_Initialized}. Capture timestamps are in
\texttt{Capture\_Clock\_Hz} (80\,MHz) ticks. A claimed handle's operations,
including \texttt{Set\_Duty}, need no further locking.

\cnote{ESP-IDF's \texttt{mcpwm\_*} component spreads a generator across timer,
operator, comparator and generator objects you wire together. Here one
\texttt{Configure\_Channel} call (frequency, pin, optional complement pin,
dead-time) does the equivalent, and ownership is enforced by the type system
rather than by convention.}

\begin{lstlisting}
type MCPWM_Unit    is (MCPWM0, MCPWM1);
type Channel_Index is (Ch0, Ch1, Ch2);
type Channel is limited private;     type Capture is limited private;
Not_Initialized : exception;

procedure Setup   (Unit : MCPWM_Unit);                       --  once per unit
procedure Claim   (C : in out Channel; Unit : MCPWM_Unit; Index : Channel_Index);
procedure Configure_Channel
  (C : Channel; Freq : Positive; Pin : ESP32S3.GPIO.Pin_Id;
   Complement_Pin : ESP32S3.GPIO.Optional_Pin := ESP32S3.GPIO.No_Pin;
   Dead_Time_Ns   : Natural := 0);
procedure Start (C : Channel);   procedure Stop (C : Channel);
procedure Set_Duty (C : Channel; Percent : Duty_Percent);
procedure Set_Carrier (C : Channel; Enable : Boolean := True; ...);
--  Fault / trip-zone:
procedure Configure_Fault (Unit : MCPWM_Unit; Input : Fault_Input;
                           Pin : ESP32S3.GPIO.Pin_Id; Active_High : Boolean := True);
procedure Protect_Channel (C : Channel; Input : Fault_Input;
                           Mode : Fault_Mode := One_Shot; Action : Trip_Action := Force_Low);
procedure Clear_Fault (C : Channel);   function Faulted (C : Channel) return Boolean;
--  Capture:
Capture_Clock_Hz : constant := 80_000_000;
procedure Configure_Capture (Cap : Capture; Pin : ESP32S3.GPIO.Pin_Id;
                             Edge : Cap_Edge := Both_Edges);
function  Capture_Pending (Cap : Capture) return Boolean;
procedure Read_Capture (Cap : Capture; Value : out Interfaces.Unsigned_32;
                        Falling : out Boolean);
\end{lstlisting}

\begin{lstlisting}
Setup (MCPWM0);
Claim (Gen0, MCPWM0, Ch0);
Configure_Channel (Gen0, Freq => 20_000, Pin => 4);      --  20 kHz, edge-aligned
Start (Gen0);
Set_Duty (Gen0, 75.0);

--  Complementary half-bridge pair with 1 us dead-time on channel 1:
Claim (Gen1, MCPWM0, Ch1);
Configure_Channel (Gen1, Freq => 20_000, Pin => 6,
                   Complement_Pin => 7, Dead_Time_Ns => 1_000);
Start (Gen1);  Set_Duty (Gen1, 50.0);   --  A and B each ~50 %, never overlapping

--  Over-current trip: force channel 0 low when Fault0 asserts, latched:
Configure_Fault (MCPWM0, Fault0, Pin => 8, Active_High => True);
Protect_Channel (Gen0, Fault0, One_Shot, Force_Low);
--  ... later, once the fault clears:  Clear_Fault (Gen0);
\end{lstlisting}

\noindent\textbf{Example:} \texttt{examples/esp32s3\_mcpwm\_pwm}.

\section{Timer --- General-Purpose Timers}\index{Timer driver}
\label{app:int:timer}

The S3's two timer groups, \texttt{TIMG0} and \texttt{TIMG1}, each contain one
54-bit up/down counter clocked from the APB clock through a 16-bit prescaler,
with a programmable alarm. (The per-group watchdogs are a separate block and are
left untouched.) The Ada driver in \texttt{ESP32S3.Timer} exposes one
general-purpose timer per group as a limited, controlled \texttt{Timer} handle,
released --- and the counter stopped --- automatically on scope exit.

\texttt{Configure} sets the tick rate in ticks per second (default 1\,MHz =
1\,\textmu{}s per tick; the range is roughly \texttt{80\_000\_000} down to
\texttt{80\_000\_000/65536}) and leaves the counter stopped at 0.
\texttt{Set\_Alarm} arms a one-shot: \texttt{Alarm\_Fired} reports the latched
interrupt-status flag, which stays set until \texttt{Clear\_Alarm}. Counts are
read back as the 64-bit \texttt{Ticks} type. As a finalization-using driver it
targets the embedded/full profile.

\begin{lstlisting}
type Timer_Index is range 0 .. 1;       --  0 = TIMG0, 1 = TIMG1
type Ticks       is new Interfaces.Unsigned_64;
type Timer is limited private;

procedure Claim     (T : in out Timer; Index : Timer_Index);
function  Is_Valid  (T : Timer) return Boolean;
procedure Release   (T : in out Timer);
procedure Configure (T : in out Timer; Tick_Hz : Positive := 1_000_000);
procedure Start (T : Timer);   procedure Stop (T : Timer);   procedure Reset (T : Timer);
function  Value (T : Timer) return Ticks;
procedure Set_Alarm  (T : Timer; At_Ticks : Ticks);
function  Alarm_Fired (T : Timer) return Boolean;
procedure Clear_Alarm (T : Timer);
\end{lstlisting}

\begin{lstlisting}
declare
   T : Timer;
begin
   Claim (T, 0);
   Configure (T, Tick_Hz => 1_000_000);          --  1 tick = 1 us
   Reset (T);  Start (T);
   delay until Clock + Milliseconds (50);
   pragma Assert (Value (T) in 49_000 .. 51_000); --  ~50_000 ticks, cross-checked

   Reset (T);
   Set_Alarm (T, 30_000);                         --  fire at 30 ms
   Start (T);
   while not Alarm_Fired (T) loop null; end loop;
   Clear_Alarm (T);  Stop (T);
end;   --  T finalizes: stopped and released
\end{lstlisting}

\noindent\textbf{Example:} \texttt{examples/esp32s3\_timer\_count}.

\section{PCNT --- Pulse Counter}\index{PCNT driver}
\label{app:int:pcnt}

PCNT counts edges on an input signal. The S3 has four counter units, each
accumulating into a \emph{signed 16-bit} counter that wraps at $\pm 32768$; the
classic uses are quadrature decoding, tachometers and flow-meters. The Ada driver
in \texttt{ESP32S3.PCNT} exposes the common ``count edges on a pin'' case --- the
per-unit direction-control input and the threshold-event comparators are left at
their pass-through defaults.

A unit is handed out as a limited, controlled \texttt{Unit} handle (auto-released
on scope exit). \texttt{Configure} routes a pin into the counter and starts it; by
default each rising edge increments, or set \texttt{Both\_Edges} to count both.
\texttt{Count} returns the signed value as an \texttt{Integer}, and
\texttt{Pause}/\texttt{Resume} hold and continue counting while retaining the
count.

\begin{lstlisting}
type Unit_Index is range 0 .. 3;        --  the four counter units
type Unit is limited private;

procedure Claim     (U : in out Unit; Index : Unit_Index);
function  Is_Valid  (U : Unit) return Boolean;
procedure Release   (U : in out Unit);
procedure Configure (U : in out Unit; Pin : ESP32S3.GPIO.Pin_Id;
                     Both_Edges : Boolean := False);
function  Count  (U : Unit) return Integer;   --  signed, wraps at +/- 32768
procedure Clear  (U : Unit);
procedure Pause  (U : Unit);
procedure Resume (U : Unit);
\end{lstlisting}

\begin{lstlisting}
declare
   U : Unit;
begin
   Claim (U, 0);
   Configure (U, Pin => 4);             --  count rising edges on GPIO4
   for I in 1 .. 100 loop               --  the self-test drives the same pad ...
      ESP32S3.GPIO.Set (4);   delay until Clock + Microseconds (20);
      ESP32S3.GPIO.Clear (4); delay until Clock + Microseconds (20);
   end loop;
   pragma Assert (Count (U) = 100);     --  ... and counts the edges back
end;   --  U finalizes: paused and released
\end{lstlisting}

\noindent\textbf{Example:} \texttt{examples/esp32s3\_pcnt\_count}.

\section{SDM --- Sigma-Delta Modulator}\index{SDM driver}
\label{app:int:sdm}

The SDM block (part of the GPIO sigma-delta unit, \texttt{GPIO\_SD}) provides
eight 1-bit density-modulated outputs. Each channel emits a high-frequency pulse
stream whose average density is set by a signed 8-bit value; passed through an
external RC low-pass filter it becomes a cheap analog output for LED dimming, a
bias voltage, or simple audio. The Ada driver in \texttt{ESP32S3.SDM} exposes a
channel as a limited, controlled \texttt{Channel} handle.

\texttt{Configure} routes the output to a pin and starts it at 0\,\% density.
\texttt{Carrier\_Hz} requests the pulse-stream (carrier) frequency: the modulator
runs at the APB clock ($\sim$80\,MHz) divided by an integer 1..256, so the
achieved rate is the nearest \texttt{APB/N}, roughly 312\,500\,Hz to 80\,MHz. A
higher carrier is easier to smooth; the exact value rarely matters and is rounded
to the nearest available divider. \texttt{Set\_Density} is a single register
write.

\begin{lstlisting}
type Channel_Index   is range 0 .. 7;      --  the eight SDM channels
subtype Density_Percent is Float range 0.0 .. 100.0;
type Channel is limited private;

procedure Claim     (C : in out Channel; Index : Channel_Index);
function  Is_Valid  (C : Channel) return Boolean;
procedure Release   (C : in out Channel);
procedure Configure (C : in out Channel; Pin : ESP32S3.GPIO.Pin_Id;
                     Carrier_Hz : Positive := 1_000_000);
procedure Set_Density (C : Channel; Percent : Density_Percent);
\end{lstlisting}

\begin{lstlisting}
declare
   Ch        : Channel;
   Densities : constant array (1 .. 3) of Density_Percent := (25.0, 50.0, 75.0);
begin
   Claim (Ch, 0);
   --  A lower carrier oversamples cleanly when the pad is GPIO-sampled back:
   Configure (Ch, Pin => 4, Carrier_Hz => 400_000);
   for I in Densities'Range loop
      Set_Density (Ch, Densities (I));
      delay until Clock + Milliseconds (5);
      --  averaged pad reading should track the set density within a few percent
   end loop;
end;   --  Ch finalizes: output low, channel released
\end{lstlisting}

\noindent\textbf{Example:} \texttt{examples/esp32s3\_sdm\_output}.

\section{RMT --- Remote Control Transceiver}\index{RMT driver}
\label{app:int:rmt}

RMT is the ``infinitely-flexible pulse generator'': it transmits and receives
sequences of \{level, duration\} pulses, making it the workhorse for IR remotes,
WS2812 (``NeoPixel'') LEDs, 1-Wire, and arbitrary bit-banged timing. The S3 has
eight channels --- 0..3 transmit, 4..7 receive --- each backed by a 48-symbol RAM
block. One \texttt{RMT\_Symbol} packs two \{level, duration\} pulses into 32 bits
(durations are 15-bit channel ticks, where the tick length is
$1/\texttt{Resolution\_Hz}$).

The Ada driver in \texttt{ESP32S3.RMT} uses \emph{distinct} limited, controlled
handle types --- \texttt{TX\_Channel} and \texttt{RX\_Channel} --- so transmit and
receive channels can't be confused. \texttt{Transmit} blocks until the burst
finishes; bursts larger than the channel's RAM are streamed by re-filling the
symbol RAM in halves as it drains, so \texttt{Symbols} may be any length. The
\texttt{Blocks} parameter (1..4) borrows the RAM of higher-numbered TX channels
to raise the one-shot ceiling. On the receive side, \texttt{Start} arms the
channel and \texttt{Receive} blocks until the line stays idle for
\texttt{Idle\_Ticks}. Because \texttt{Transmit} busy-polls its re-fill, keep
higher-priority interrupts short.

\begin{lstlisting}
type TX_Index    is range 0 .. 3;   type RX_Index is range 0 .. 3;
type Tick_Count  is range 0 .. 32_767;     --  duration in channel ticks (15-bit)

type RMT_Symbol is record
   Level0 : Boolean := False;   Duration0 : Tick_Count := 0;
   Level1 : Boolean := False;   Duration1 : Tick_Count := 0;
end record;
for RMT_Symbol'Size use 32;
type Symbol_Array is array (Natural range <>) of RMT_Symbol;

type TX_Channel is limited private;   type RX_Channel is limited private;

--  Transmit:
procedure Claim     (C : in out TX_Channel; Index : TX_Index);
procedure Configure (C : in out TX_Channel; Resolution_Hz : Positive;
                     Pin : ESP32S3.GPIO.Pin_Id; Blocks : Positive := 1);
procedure Transmit  (C : TX_Channel; Symbols : Symbol_Array);
--  Receive:
procedure Claim     (C : in out RX_Channel; Index : RX_Index);
procedure Configure (C : in out RX_Channel; Resolution_Hz : Positive;
                     Pin : ESP32S3.GPIO.Pin_Id; Idle_Ticks : Tick_Count := 2_000);
procedure Start     (C : RX_Channel);
procedure Receive   (C : RX_Channel; Into : out Symbol_Array; Count : out Natural);
\end{lstlisting}

\begin{lstlisting}
declare
   Tx   : TX_Channel;   Rx : RX_Channel;
   Sent : constant Symbol_Array :=
     ((Level0 => True, Duration0 =>  50, Level1 => False, Duration1 =>  60),
      (Level0 => True, Duration0 =>  80, Level1 => False, Duration1 =>  90),
      (Level0 => True, Duration0 => 120, Level1 => False, Duration1 => 130));
   Got   : Symbol_Array (0 .. 15);
   Count : Natural;
begin
   Claim (Tx, 0);   Claim (Rx, 0);
   Configure (Tx, Resolution_Hz => 1_000_000, Pin => 4);            --  1 us / tick
   Configure (Rx, Resolution_Hz => 1_000_000, Pin => 4, Idle_Ticks => 1_000);
   Start    (Rx);              --  arm the receiver first
   Transmit (Tx, Sent);        --  drive the burst onto the pad
   Receive  (Rx, Got, Count);  --  block until idle, read it back (single-pad loopback)
end;   --  Tx, Rx finalize: released
\end{lstlisting}

\noindent\textbf{Example:} \texttt{examples/esp32s3\_rmt\_loopback}.

\section{GDMA --- General DMA}\index{GDMA driver}
\label{app:int:gdma}

The ESP32-S3 has a single AHB GDMA block with five channel pairs (\texttt{0..4}).
Each pair carries an independent transmit (OUT) and receive (IN) path, and either
can be wired to any peripheral through the GDMA crossbar --- a channel is a
runtime-assigned resource, not fixed per peripheral. The Ada driver
(\texttt{ESP32S3.GDMA}) models that directly: \texttt{Claim} hands out a free
\texttt{Channel} bound to a \texttt{Peripheral} (\texttt{Mem2Mem}, \texttt{SPI2},
\texttt{I2S0}, \texttt{LCD\_CAM}, \texttt{ADC\_DAC}, \ldots); peripheral drivers
such as I2S and LCD claim a channel and drive their own descriptors over it.

A \texttt{Channel} is \texttt{limited private} --- non-copyable, so two tasks
cannot alias one channel --- and \emph{controlled}: claiming goes through a
protected allocator, and the handle auto-releases on scope exit (including during
exception unwinding) so a channel cannot leak or be reused through a stale copy.
Because finalization is involved, this driver targets the embedded/full profile
(excluded from the light-tasking build). Transfers are bounded by
\texttt{Max\_Transfer = 4095} (the 12-bit buffer-size field); the implemented
memory-to-memory \texttt{Copy} loops one channel's OUT path into its own IN path,
and source, destination and descriptors must all live in internal SRAM. Beyond
\texttt{Copy}, the peripheral-bound \texttt{Start}/\texttt{Start\_Loop}/%
\texttt{Done}/\texttt{Wait} API arms a single buffer (or a self-looping gapless
one) in a \texttt{Direction} and lets the peripheral pace the move.

\begin{lstlisting}
type Channel_Id is mod 5;
type Peripheral is
  (Mem2Mem, SPI2, SPI3, UHCI0, I2S0, I2S1, LCD_CAM, AES, SHA, ADC_DAC, RMT);
type Channel is limited private;            --  non-copyable, auto-released
type Direction is (Mem_To_Periph, Periph_To_Mem);
Max_Transfer : constant := 4095;

procedure Claim   (C : in out Channel; Peri : Peripheral);
function  Is_Valid (C : Channel) return Boolean;
procedure Release (C : in out Channel);
procedure Copy    (C : Channel; Dst, Src : System.Address; Length : Natural);
procedure Start   (C : Channel; Dir : Direction;
                   Buffer : System.Address; Length : Natural);
procedure Start_Loop (C : Channel; Buffer : System.Address; Length : Natural);
function  Done    (C : Channel; Dir : Direction) return Boolean;
procedure Wait    (C : Channel; Dir : Direction);
\end{lstlisting}

\begin{lstlisting}
declare
   C  : Channel;                            --  limited: cannot be aliased
   Ok : Boolean := False;
begin
   Claim (C, Mem2Mem);
   if Is_Valid (C) then
      Copy (C, Dst'Address, Src'Address, Buffer'Length);
      Ok := (for all I in Buffer'Range => Dst (I) = Src (I));
   end if;
   Copy_Result (Ok);
end;                                        --  C finalizes -> channel released
\end{lstlisting}

\cnote{In ESP-IDF a channel is a \texttt{gdma\_channel\_handle\_t} you
\texttt{gdma\_new\_channel} / \texttt{gdma\_del\_channel} by hand, with descriptor
lifetimes left to you. Here the limited-controlled \texttt{Channel} \emph{is} the
allocation: leaving scope frees it, and the type system forbids the aliasing that
leaks one.}

\noindent\textbf{Example:} \texttt{examples/esp32s3\_gdma\_copy}.

\section{I2S --- Digital Audio Serial Bus}\index{I2S driver}
\label{app:int:i2s}

The S3 has two I2S controllers (\texttt{I2S0}, \texttt{I2S1}), each moving PCM
samples only through the GDMA crossbar --- the block has no CPU FIFO, so data
flows exclusively via DMA. The Ada driver (\texttt{ESP32S3.I2S}) brings a port up
as a stereo master that generates BCLK and WS and shifts a sample buffer out (TX)
and/or captures one in (RX), in either \texttt{Standard} (verbatim I2S/TDM) or
\texttt{PDM} mode. In PDM mode the on-chip sigma-delta converters sit between the
buffer and the wire: TX turns each PCM sample into a 1-bit pulse-density stream,
RX decimates a 1-bit PDM input back to PCM, while \texttt{Write}/\texttt{Read}/%
\texttt{Transfer} stay unchanged (the converters high-pass filter, so a constant
level does not survive a PDM round trip).

Mirroring \texttt{ESP32S3.SPI}, the raw register access lives in the private child
\texttt{ESP32S3.I2S.Engine} (which claims the GDMA channel); the application sees
only the parent. \texttt{Setup} configures a port once at startup; thereafter
\texttt{Acquire} hands out a \texttt{Session} --- limited, non-copyable, and
controlled --- that owns the port exclusively and auto-releases on scope exit, so a
fault between \texttt{Acquire} and \texttt{Release} cannot leak the lock.
Transfers (\texttt{Write}, \texttt{Read}, \texttt{Transfer}, plus the gapless
\texttt{Start\_Continuous}/\texttt{Capture}/\texttt{Stop} streaming path) run
outside the lock and raise \texttt{Not\_Owned} if the session does not hold a
port, or \texttt{Not\_Initialized} if a port was never \texttt{Setup}. The driver
needs a tasking runtime (Jorvik or richer) and, via finalization, targets the
embedded/full profile.

\begin{lstlisting}
type I2S_Port    is (I2S0, I2S1);
type Sample_Bits is (Bits_8, Bits_16, Bits_24, Bits_32);
type I2S_Mode    is (Standard, PDM);
type Session     is limited private;        --  non-copyable, auto-released
Not_Initialized, Not_Owned : exception;

procedure Setup (Port : I2S_Port; Sample_Rate : Positive := 16_000;
                 Bits : Sample_Bits := Bits_16; Mode : I2S_Mode := Standard);
procedure Enable_Loopback (Port : I2S_Port; Pad : ESP32S3.GPIO.Pin_Id);
procedure Configure_Pins  (Port : I2S_Port;
                           Bclk, Ws, Dout, Din, Mclk : Optional_Pin := No_Pin);
procedure Acquire  (S : in out Session; Port : I2S_Port);
procedure Write    (S : Session; Tx : System.Address; Length : Natural);
procedure Read     (S : Session; Rx : System.Address; Length : Natural);
procedure Transfer (S : Session; Tx, Rx : System.Address; Length : Natural);
procedure Start_Continuous (S : Session; Tx : System.Address; Length : Natural);
procedure Capture  (S : Session; Rx : System.Address; Length : Natural);
procedure Stop     (S : Session);
procedure Release  (S : in out Session);
\end{lstlisting}

\begin{lstlisting}
Setup (I2S0, Sample_Rate => 16_000, Bits => Bits_16);
Enable_Loopback (I2S0, Pad => Data_Pin);    --  out->in on one pad, no wiring
declare
   S  : Session;                            --  limited: cannot be copied/shared
   Ok : Boolean := False;
begin
   Acquire (S, I2S0);                       --  suspends until the port is free
   Transfer (S, Tx'Address, Rx'Address, Samples'Length * 2);
   Ok := (for all I in Samples'Range => Rx (I) = Tx (I));
end;                                        --  S finalizes -> port released
\end{lstlisting}

\noindent\textbf{Example:} \texttt{examples/esp32s3\_i2s\_loopback} (Standard
full-duplex) and \texttt{examples/esp32s3\_i2s\_pdm} (PDM microphone capture).

\section{LCD --- i8080 Parallel Display}\index{LCD driver}
\label{app:int:lcd}

The LCD half of the S3's LCD\_CAM controller drives an 8-bit Intel-8080
(``i80''/MCU) parallel display: a byte buffer is streamed out the data bus, one
byte per pixel clock (PCLK), over the GDMA crossbar. The Ada driver
(\texttt{ESP32S3.LCD}) covers exactly this 8-bit i80 transmit path (the
camera-receive half and the RGB/16-bit modes are not modelled). \texttt{Setup}
brings the controller up at (about) \texttt{Pclk\_Hz} --- quantised as
\texttt{20\,MHz / round(20\,MHz / Pclk\_Hz)} --- and claims its GDMA channel.

The single controller is guarded by a protected object; \texttt{Acquire} hands out
a limited, controlled \texttt{Session} that owns it exclusively and releases on
scope exit. Notably, both transfers \emph{and} post-\texttt{Setup} reconfiguration
go through the held session: \texttt{Configure\_Pins} (route the eight data lines
\texttt{Data\_Pins} and PCLK) and \texttt{Enable\_Clock\_Out} (free-run the clock
with no data) require ownership, so changing a setting can never race another
task. \texttt{Transmit} streams \texttt{1..4095} bytes from internal SRAM and
reports completion in \texttt{Ok}. Every operation raises \texttt{Not\_Owned}
without the lock, or \texttt{Not\_Initialized} if \texttt{Setup} did not succeed
(including a failed channel claim). All register access lives in the private
\texttt{ESP32S3.LCD.Engine} child; via finalization the driver targets the
embedded/full profile.

\begin{lstlisting}
No_Pin : constant ESP32S3.GPIO.Pad_Number := ESP32S3.GPIO.No_Pin;
type Data_Pins is array (0 .. 7) of ESP32S3.GPIO.Optional_Pin;  -- D0..D7
type Session   is limited private;          --  non-copyable, auto-released
Not_Initialized, Not_Owned : exception;

procedure Setup (Pclk_Hz : Positive := 1_000_000);
procedure Acquire (S : in out Session);
procedure Configure_Pins   (S : Session; Data : Data_Pins;
                            Pclk : ESP32S3.GPIO.Optional_Pin);
procedure Enable_Clock_Out (S : Session; Pclk_Pad : ESP32S3.GPIO.Pin_Id);
procedure Transmit (S : Session; Tx : System.Address; Length : Natural;
                    Ok : out Boolean);
procedure Release  (S : in out Session);
\end{lstlisting}

\begin{lstlisting}
Setup (Pclk_Hz => 200 * 1000);              --  200 kHz pixel clock
declare
   S  : Session;
   Ok : Boolean;
begin
   Acquire (S);
   Configure_Pins (S, D_Pins, Pclk => Pclk_Pin);  --  route pads on held ctrl
   Transmit (S, Buf'Address, Buffer'Length, Ok);  --  blocks until done
   Release (S);
   Put_Line (if Ok then "PASS" else "FAIL");
end;
\end{lstlisting}

\noindent\textbf{Example:} \texttt{examples/esp32s3\_lcd\_i8080}.

\section{ADC --- SAR Analog-to-Digital Converter}\index{ADC driver}
\label{app:int:adc}

The S3 has two SAR ADC units, each with up to ten 12-bit channels on fixed GPIOs:
ADC1 channel \texttt{n} maps to \texttt{GPIO(n+1)} (ch0\,=\,GPIO1\ldots ch9\,=\,GPIO10),
ADC2 channel \texttt{n} to \texttt{GPIO(n+11)}. The Ada driver
(\texttt{ESP32S3.ADC}) does software-triggered single (one-shot) conversions
through the RTC controller; each returns a raw \texttt{0..4095} code, with
per-channel \texttt{Attenuation} (\texttt{Db\_0} $\approx$ 1.1\,V full-scale
\ldots \texttt{Db\_12} $\approx$ 3.3\,V) setting the input range.

A unit is a shared resource handed out as a \texttt{Reader} --- limited
(non-copyable) and controlled (released automatically on scope exit) --- so
claiming a unit cannot leak. After \texttt{Claim}, check \texttt{Is\_Valid};
\texttt{Read} samples a channel (blocking, a few microseconds), returning
\texttt{0} on an invalid handle. \texttt{Channel\_Pin} reports the GPIO a channel
is wired to (routing/documentation), and the diagnostics \texttt{Cal\_Code}
(self-calibrated initial code) and \texttt{Last\_Done} (the SAR DONE flag) help
confirm a conversion actually ran. Via finalization the driver targets the
embedded/full profile.

\begin{lstlisting}
type ADC_Unit      is (ADC1, ADC2);
type Channel_Index is range 0 .. 9;
type Attenuation   is (Db_0, Db_2_5, Db_6, Db_12);
subtype Raw_Value  is Natural range 0 .. 4095;        --  12-bit result
type Reader        is limited private;                --  auto-released

procedure Claim   (R : in out Reader; Unit : ADC_Unit);
function  Is_Valid (R : Reader) return Boolean;
procedure Release (R : in out Reader);
function  Read    (R : Reader; Ch : Channel_Index;
                   Atten : Attenuation := Db_12) return Raw_Value;
function  Channel_Pin (Unit : ADC_Unit; Ch : Channel_Index)
                       return ESP32S3.GPIO.Pin_Id;
function  Cal_Code  (Unit : ADC_Unit) return Natural;
function  Last_Done return Boolean;
\end{lstlisting}

\begin{lstlisting}
declare
   R : Reader;
   High, Low : Natural := 0;
begin
   Claim (R, ADC1);                         --  ADC1 ch0 -> GPIO1
   ESP32S3.GPIO.Configure (Pin, ESP32S3.GPIO.Output);
   ESP32S3.GPIO.Set (Pin);                  --  drive the sensed pad high
   delay until Clock + Milliseconds (2);
   High := Read (R, Ch);                    --  expect near full scale
   ESP32S3.GPIO.Clear (Pin);
   delay until Clock + Milliseconds (2);
   Low := Read (R, Ch);                     --  expect near zero
end;                                        --  R finalizes -> unit released
\end{lstlisting}

\noindent\textbf{Example:} \texttt{examples/esp32s3\_adc\_read}.

\section{TWAI --- CAN Controller}\index{TWAI driver}
\label{app:int:twai}

The S3's TWAI block (Two-Wire Automotive Interface) is a single
SJA1000-compatible CAN 2.0 controller. The Ada driver (\texttt{ESP32S3.TWAI})
sends and receives both standard (11-bit \texttt{Standard\_Id}) and extended
(29-bit \texttt{Extended\_Id}, CAN 2.0B) frames, data or remote-request (RTR).
Each frame type carries its own identifier subtype, so the id is range-checked to
its standard and \texttt{Send} is overloaded on the frame type --- a 29-bit id
cannot end up in a standard frame. A real bus needs an external transceiver; for a
wiring-free self-test, \texttt{Self\_Test} mode transmits and self-receives a
frame (no second node to acknowledge), with TX looped back to RX through one pad
via \texttt{Enable\_Loopback}.

The single controller is guarded by a protected object; \texttt{Acquire} hands out
a limited, controlled \texttt{Session} owning it exclusively, auto-released on
scope exit. \texttt{Setup} (mode + bit rate, accepting all ids) runs once;
transfers and post-\texttt{Setup} reconfiguration (\texttt{Configure\_Pins},
\texttt{Enable\_Loopback}) all go through the held session and raise
\texttt{Not\_Owned} otherwise (or \texttt{Not\_Initialized} before \texttt{Setup}).
Because a received frame may be either width, receiving is a two-step idiom:
\texttt{Available} reports a waiting frame, \texttt{Is\_Extended} its width, then
the matching \texttt{Receive} overload reads it (returning \texttt{Got => False}
if the waiting frame is the other width, leaving it buffered). All register access
lives in the private \texttt{ESP32S3.TWAI.Engine} child; via finalization the
driver targets the embedded/full profile.

\begin{lstlisting}
type Bus_Mode    is (Normal, Listen_Only, Self_Test);
subtype Data_Length is Natural range 0 .. 8;
type Data_Bytes  is array (0 .. 7) of Interfaces.Unsigned_8;
subtype Standard_Id is Interfaces.Unsigned_32 range 0 .. 16#7FF#;        -- 11-bit
subtype Extended_Id is Interfaces.Unsigned_32 range 0 .. 16#1FFF_FFFF#;  -- 29-bit
type Standard_Frame is record
   Id : Standard_Id := 0; Remote : Boolean := False;
   Length : Data_Length := 0; Data : Data_Bytes := (others => 0);
end record;
type Extended_Frame is record ... end record;   --  same shape, Extended_Id
type Session is limited private;

procedure Setup (Mode : Bus_Mode := Normal; Bit_Rate : Positive := 125_000);
procedure Acquire (S : in out Session);
procedure Enable_Loopback (S : Session; Pad : ESP32S3.GPIO.Pin_Id);
procedure Send (S : Session; F : Standard_Frame);   --  overloaded on frame type
procedure Send (S : Session; F : Extended_Frame);
function  Available   (S : Session) return Boolean;
function  Is_Extended (S : Session) return Boolean;
procedure Receive (S : Session; F : out Standard_Frame; Got : out Boolean);
procedure Receive (S : Session; F : out Extended_Frame; Got : out Boolean);
\end{lstlisting}

\begin{lstlisting}
Setup (Mode => Self_Test, Bit_Rate => 125_000);
declare
   S   : Session;
   RS  : Standard_Frame;
   Got : Boolean;
begin
   Acquire (S);
   Enable_Loopback (S, Pad);                --  out->in on one pad, no wiring
   Send (S, Std);                           --  picks the Standard_Frame overload
   Got := Available (S) and then not Is_Extended (S);
   if Got then
      Receive (S, RS, Got);                 --  matching width -> Got => True
   end if;
end;                                        --  S finalizes -> controller released
\end{lstlisting}

\noindent\textbf{Example:} \texttt{examples/esp32s3\_twai\_loopback}.

\section{RTC --- Real-Time Clock and Deep Sleep}\index{RTC driver}
\label{app:int:rtc}

The ESP32-S3's RTC (low-power) domain stays alive while the digital core is
powered down. In deep sleep, CPU and main RAM lose power; only the RTC domain
survives, so on wake the chip \emph{resets} and re-runs from the start --- but
anything kept in the 8\,KB of RTC slow memory (at \texttt{0x5000\_0000}) persists
across the reset. \texttt{ESP32S3.RTC} exposes that retained region, the cause of
the current boot (power-on versus a deep-sleep wake), and deep-sleep entry with a
timer or an RTC-GPIO (EXT1) wake source.

The deep-sleep timer runs on the uncalibrated RTC slow clock
(\textasciitilde136\,kHz), so \texttt{Deep\_Sleep\_For} is approximate. The entry
procedures do not return on success: the chip powers its digital core down and
resets on wake. Retained memory can be reached word-by-word with
\texttt{Read}/\texttt{Write}, by overlaying a variable on the \texttt{Slow\_Memory}
address, or through the generic \texttt{Retained} package. The operations are plain
register pokes with no tasking or finalization, so the driver works under every
runtime profile.

\begin{lstlisting}
Slow_Memory      : constant System.Address := ...;   --  0x5000_0000
Slow_Memory_Size : constant := 8 * 1024;
subtype Word_Index is Natural range 0 .. Slow_Memory_Size / 4 - 1;

function  Read  (Index : Word_Index) return Interfaces.Unsigned_32;
procedure Write (Index : Word_Index; Value : Interfaces.Unsigned_32);

generic                                  --  typed retained object
   type Item is private;
   Offset : Natural := 0;
package Retained is
   Object : Item with Import, Volatile, Address => ...;
end Retained;

type Wake_Cause is (Power_On, Deep_Sleep_Timer, Deep_Sleep_GPIO, Other_Reset);
function Last_Wake return Wake_Cause;

procedure Deep_Sleep_For   (Wake_After : Duration);            --  timer wake
procedure Deep_Sleep_Until (Pin : ESP32S3.GPIO.Pin_Id;        --  EXT1 wake
                            High : Boolean := True);
procedure Disable_Super_Watchdog;
\end{lstlisting}

A boot counter kept in retained word 0 survives the deep-sleep reset; the wake
cause turns into \texttt{Deep\_Sleep\_Timer} on each timer wake.

\begin{lstlisting}
use ESP32S3.RTC;
WC         : constant Wake_Cause := Last_Wake;
Boot_Count : Unsigned_32 := Read (0);          --  retained across deep sleep
begin
   Disable_Super_Watchdog;                     --  a wake can leave it armed
   if WC = Deep_Sleep_Timer or else WC = Deep_Sleep_GPIO then
      Boot_Count := Boot_Count + 1;
   else
      Boot_Count := 1;                          --  power-on / flash reset
   end if;
   Write (0, Boot_Count);                       --  persist it

   if Boot_Count < 4 then
      Deep_Sleep_For (2.0);                      --  ~2 s; DOES NOT RETURN
   end if;                                       --  chip resets on wake
\end{lstlisting}

\noindent\textbf{Example:} \texttt{examples/esp32s3\_rtc\_sleep}.

\section{RTC-IO --- Low-Power GPIO and Pad Hold}\index{RTC\_IO driver}
\label{app:int:rtcio}

GPIO0 through GPIO21 are RTC-capable: they belong to the low-power IO domain. The
headline feature is pad \emph{hold} --- latch a pad at its current output level so
it keeps driving while the rest of the chip changes, and, crucially, while the
chip is in deep sleep (the digital core powers down, but a held RTC pad stays put)
and across the reset a deep-sleep wake causes. Use it to keep a load enabled or a
reset line asserted while you sleep. A held pad ignores ordinary GPIO writes until
you \texttt{Release} it.

\texttt{ESP32S3.RTC\_IO} also exposes the RTC-domain pulls, which are distinct
from the digital IO\_MUX pulls configured through \texttt{ESP32S3.GPIO}: route a
pad into the RTC domain with \texttt{Enable\_RTC\_Input}, then \texttt{Set\_Pull}
gives it a defined level that holds (including in deep sleep). Register pokes only,
so RTC-IO works under every runtime profile.

\begin{lstlisting}
subtype RTC_Pin is ESP32S3.GPIO.Pin_Id range 0 .. 21;

procedure Hold    (Pin : RTC_Pin);    --  latch pad at its current level
procedure Release (Pin : RTC_Pin);    --  follow the GPIO register again
function  Is_Held (Pin : RTC_Pin) return Boolean;

procedure Enable_RTC_Input (Pin : RTC_Pin);   --  route into RTC domain
type Pull_Mode is (No_Pull, Up, Down);
procedure Set_Pull (Pin : RTC_Pin; Mode : Pull_Mode);
\end{lstlisting}

Driving a pad high and holding it makes a subsequent GPIO clear a no-op until
release:

\begin{lstlisting}
Pin : constant ESP32S3.GPIO.Pin_Id := 5;     --  RTC-capable pad
begin
   ESP32S3.GPIO.Configure (Pin, ESP32S3.GPIO.Output);
   ESP32S3.GPIO.Set (Pin);                    --  drive high

   Hold (Pin);                                --  latch it high
   ESP32S3.GPIO.Clear (Pin);                  --  IGNORED while held
   --  Read (Pin) still reports high here

   Release (Pin);                             --  unlatch
   ESP32S3.GPIO.Clear (Pin);                  --  now takes effect -> low
\end{lstlisting}

\noindent\textbf{Example:} \texttt{examples/esp32s3\_rtcio\_hold}.

\section{Touch --- Capacitive Touch Sensor}\index{Touch driver}
\label{app:int:touch}

The S3's capacitive touch sensor (v2) provides 14 channels, channel $n$ wired to
GPIO$n$ ($n = 1\ldots14$). Each channel measures the self-capacitance of its pad
by counting charge/discharge cycles in a fixed window; a finger near the pad
raises the capacitance and changes the count. A hardware FSM scans the enabled
channels continuously on the RTC timer, and \texttt{Read} returns the latest raw
count.

\texttt{ESP32S3.Touch} keeps the model deliberately simple: bring the controller
up once with \texttt{Setup}, \texttt{Enable} the channels you want, then
\texttt{Read} their raw counts. \texttt{Touched} is a software comparison against a
caller-captured \texttt{Reference} with a \texttt{Margin} --- deterministic, with
no on-chip threshold registers in play. It lives in the RTC/SENS domain and is
register pokes only, so it works under every runtime profile.

\begin{lstlisting}
type Channel is range 1 .. 14;               --  channel n -> GPIO n
function Pad (Ch : Channel) return ESP32S3.GPIO.Pin_Id;

procedure Setup;                             --  start the scanning FSM (once)
procedure Enable (Ch : Channel);            --  add a channel to the scan
function  Read   (Ch : Channel) return Natural;   --  latest raw count

function Touched (Ch        : Channel;
                  Reference : Natural;
                  Margin    : Natural := 20_000) return Boolean;
\end{lstlisting}

A bare pad reads a stable non-zero baseline; two different pads read different
values, proving the FSM runs on silicon:

\begin{lstlisting}
A : constant Channel := 1;     --  GPIO1
B : constant Channel := 3;     --  GPIO3
begin
   Setup;
   Enable (A);
   Enable (B);
   delay until Clock + Milliseconds (50);     --  let the FSM scan

   declare
      Ra : constant Natural := Read (A);
      Rb : constant Natural := Read (B);
   begin
      --  Ra, Rb both non-zero and distinct => capacitance FSM is live
      null;
   end;

   --  Touch detection against a captured baseline:
   Baseline : constant Natural := Read (A);
   Hit      : constant Boolean := Touched (A, Baseline, Margin => 50_000);
\end{lstlisting}

\noindent\textbf{Example:} \texttt{examples/esp32s3\_touch\_read}.

\section{SHA --- Hardware Hash Accelerator}\index{SHA driver}
\label{app:int:sha}

The SHA accelerator computes SHA-1, SHA-224 and SHA-256 of a byte message in
hardware. All three variants share the same 512-bit block and padding, differing
only in the hardware mode and digest length. \texttt{ESP32S3.SHA} exposes them as
three one-shot functions that take a \texttt{Byte\_Array} of any length (padded
internally) and return the fixed-size digest.

The accelerator is a single shared resource, so a protected object serializes the
message-load / start / read handshake, making concurrent \texttt{Hash} calls from
different tasks safe. There is no finalization --- just register pokes --- so the
driver works under every runtime profile. (The block hardware also does
SHA-384/512 on a 1024-bit block; those are not exposed here.)

\begin{lstlisting}
type Byte_Array is array (Natural range <>) of Interfaces.Unsigned_8;

subtype SHA1_Digest   is Byte_Array (0 .. 19);   --  160-bit
subtype SHA224_Digest is Byte_Array (0 .. 27);   --  224-bit
subtype SHA256_Digest is Byte_Array (0 .. 31);   --  256-bit

function Hash_1   (Data : Byte_Array) return SHA1_Digest;
function Hash_224 (Data : Byte_Array) return SHA224_Digest;
function Hash_256 (Data : Byte_Array) return SHA256_Digest;
\end{lstlisting}

Hashing \texttt{"abc"} and comparing against the FIPS-180 vector:

\begin{lstlisting}
use type ESP32S3.SHA.Byte_Array;

Abc : constant ESP32S3.SHA.Byte_Array := (16#61#, 16#62#, 16#63#);  --  "abc"

Sha_Expect : constant ESP32S3.SHA.SHA256_Digest :=
  (16#BA#, 16#78#, 16#16#, 16#BF#, 16#8F#, 16#01#, 16#CF#, 16#EA#,
   16#41#, 16#41#, 16#40#, 16#DE#, 16#5D#, 16#AE#, 16#22#, 16#23#,
   16#B0#, 16#03#, 16#61#, 16#A3#, 16#96#, 16#17#, 16#7A#, 16#9C#,
   16#B4#, 16#10#, 16#FF#, 16#61#, 16#F2#, 16#00#, 16#15#, 16#AD#);
begin
   --  PASS when the hardware digest matches the published vector:
   if ESP32S3.SHA.Hash_256 (Abc) = Sha_Expect then ... end if;
\end{lstlisting}

\noindent\textbf{Example:} \texttt{examples/esp32s3\_crypto}.

\section{AES --- Hardware AES Accelerator}\index{AES driver}
\label{app:int:aes}

The AES accelerator performs single-block (128-bit) AES in hardware.
\texttt{ESP32S3.AES} exposes the ECB building block --- one-block encrypt and
decrypt --- for 128- and 256-bit keys. As with SHA, the engine is a single shared
resource guarded by a protected object, so calls from different tasks are safe,
and the register-poke implementation needs no finalization and works under every
runtime profile.

\cnote{The ESP32-S3 AES hardware supports only 128- and 256-bit keys; there is no
192-bit mode on this silicon (selecting ``AES-192'' silently falls back to AES-128
on the first 16 key bytes). The \texttt{Supported\_Key} precondition turns an
unsupported key length into a contract violation instead of a silent fallback ---
callers that pass the \texttt{Key\_128} / \texttt{Key\_256} subtypes satisfy it
statically, at no run-time cost.}

\begin{lstlisting}
type Block is array (0 .. 15) of Interfaces.Unsigned_8;   --  128-bit block

type Key_Bytes is array (Natural range <>) of Interfaces.Unsigned_8;
subtype Key_128 is Key_Bytes (0 .. 15);
subtype Key_256 is Key_Bytes (0 .. 31);

function Supported_Key (Key : Key_Bytes) return Boolean is
  (Key'Length = 16 or else Key'Length = 32);

function Encrypt_ECB (Key : Key_Bytes; Plain  : Block) return Block
  with Pre => Supported_Key (Key);
function Decrypt_ECB (Key : Key_Bytes; Cipher : Block) return Block
  with Pre => Supported_Key (Key);
\end{lstlisting}

Encrypting the FIPS-197 example block and round-tripping it back:

\begin{lstlisting}
use type ESP32S3.AES.Block;

Aes_Key : constant ESP32S3.AES.Key_128 :=
  (16#00#, 16#01#, 16#02#, 16#03#, 16#04#, 16#05#, 16#06#, 16#07#,
   16#08#, 16#09#, 16#0A#, 16#0B#, 16#0C#, 16#0D#, 16#0E#, 16#0F#);
Aes_Plain : constant ESP32S3.AES.Block :=
  (16#00#, 16#11#, 16#22#, 16#33#, 16#44#, 16#55#, 16#66#, 16#77#,
   16#88#, 16#99#, 16#AA#, 16#BB#, 16#CC#, 16#DD#, 16#EE#, 16#FF#);
Aes_Cipher : constant ESP32S3.AES.Block :=
  (16#69#, 16#C4#, 16#E0#, 16#D8#, 16#6A#, 16#7B#, 16#04#, 16#30#,
   16#D8#, 16#CD#, 16#B7#, 16#80#, 16#70#, 16#B4#, 16#C5#, 16#5A#);
begin
   --  encrypt matches the vector, and decrypt recovers the plaintext:
   if ESP32S3.AES.Encrypt_ECB (Aes_Key, Aes_Plain)  = Aes_Cipher and then
      ESP32S3.AES.Decrypt_ECB (Aes_Key, Aes_Cipher) = Aes_Plain
   then ... end if;
\end{lstlisting}

\noindent\textbf{Example:} \texttt{examples/esp32s3\_crypto}.

\section{SDMMC --- Native SD/MMC Host Controller}\index{SDMMC driver}
\label{app:int:sdmmc}

The dedicated SDHOST controller drives an SD card on the real SD bus --- a clock,
a bidirectional command line, and 1 or 4 data lines --- with its own
command/response state machine. Unlike the SPI-mode SD path, this is how you reach
an SDHC/SDXC card at speed (up to 50\,MHz in High-Speed mode).
\texttt{ESP32S3.SDMMC} moves data in PIO/FIFO mode (the CPU pushes/pops each
512-byte block through the controller FIFO), so there are no GDMA descriptors and
no finalization: a library-level protected object serializes the single shared
controller, and it works under every runtime profile, including light tasking. The
API is always 512-byte logical blocks (LBA); the SDHC/SDXC versus SDSC addressing
difference is resolved internally. (Covered in full in the Mass Storage chapter,
\S\ref{ch:storage}.)

\begin{lstlisting}
type Slot      is (Slot1, Slot2);
type Bus_Width is (Width_1, Width_4);
type Block         is array (0 .. 511) of Interfaces.Unsigned_8;
type Block_Address is new Interfaces.Unsigned_32;
type Card_Kind is (Unknown, SDSC, SDHC);             --  SDHC also covers SDXC
type Status    is (OK, No_Card, Unusable, Init_Timeout, Cmd_Timeout,
                   Cmd_CRC, Read_Error, Write_Error);
type Card is limited private;

procedure Setup (C : out Card; On : Slot;
                 Clk, Cmd, D0 : ESP32S3.GPIO.Pin_Id;
                 D1, D2, D3   : ESP32S3.GPIO.Optional_Pin := ESP32S3.GPIO.No_Pin;
                 Width        : Bus_Width := Width_1;
                 Init_Clock_Hz, Data_Clock_Hz : Positive;
                 High_Speed   : Boolean := False);
procedure Initialize  (C : in out Card; Result : out Status);
function  Kind        (C : Card) return Card_Kind;
procedure Read_Block  (C : in out Card; LBA : Block_Address;
                       Data : out Block; Result : out Status);
procedure Write_Block (C : in out Card; LBA : Block_Address;
                       Data : Block; Result : out Status);
\end{lstlisting}

Bringing up Slot1 in 4-bit mode and doing a non-destructive sector round-trip:

\begin{lstlisting}
use type ESP32S3.SDMMC.Status;
Test_LBA : constant ESP32S3.SDMMC.Block_Address := 16#2000#;   --  sector 8192
C    : ESP32S3.SDMMC.Card;
St   : ESP32S3.SDMMC.Status;
Orig : ESP32S3.SDMMC.Block;
begin
   ESP32S3.SDMMC.Setup
     (C, On => ESP32S3.SDMMC.Slot1,
      Clk => 14, Cmd => 15, D0 => 2, D1 => 4, D2 => 12, D3 => 13,
      Width => ESP32S3.SDMMC.Width_4,
      Init_Clock_Hz => 400_000, Data_Clock_Hz => 20_000_000);

   ESP32S3.SDMMC.Initialize (C, St);
   if St = ESP32S3.SDMMC.OK then
      ESP32S3.SDMMC.Read_Block  (C, Test_LBA, Orig, St);   --  read scratch sector
      ESP32S3.SDMMC.Write_Block (C, Test_LBA, Orig, St);   --  write it back unchanged
   end if;
\end{lstlisting}

\noindent\textbf{Example:} \texttt{examples/esp32s3\_sdmmc}.

\chapter{External Device Drivers}
\label{app:external}

This appendix is the companion to appendix~\ref{app:internal}: where that one
covers the SoC's on-chip peripherals, this one covers the drivers for
\emph{external} parts -- the I/O expanders, sensors, clocks, codec, display,
GPS module, SD card, LED string and Ethernet controller that hang off the
ESP32-S3's buses. Each driver layers on one (or two) of the on-chip drivers: the
I2C devices ride \texttt{ESP32S3.I2C}, the display rides \texttt{ESP32S3.SPI}, the
codec uses \texttt{ESP32S3.I2C} for control and \texttt{ESP32S3.I2S} for audio, the
GPS service reads \texttt{ESP32S3.UART}, and the LED string is clocked out over
\texttt{ESP32S3.RMT}. The W5500 Ethernet controller is the largest of these -- a
whole networking stack -- so it has its own chapter
(\ref{ch:networking}); this appendix gives only its bus and a pointer.

The same ownership discipline recurs. Most of these drivers expose a
\texttt{Device} configured once by \texttt{Setup} and, for the multi-step ones, a
\texttt{Session} (or \texttt{Output}/\texttt{Strip}) handle taken with
\texttt{Acquire} that holds the part exclusively across a sequence of operations
and auto-releases on scope exit, while the underlying bus is locked only per
transaction -- so other devices on the same bus stay reachable between your
calls. Each entry gives the device's bus, its I2C address where relevant, a digest
of the public API, and a worked example from the \texttt{examples/} tree. The
controlled handles mean these drivers require the \texttt{embedded} or
\texttt{full} profile, as noted per section.

\section{CH422G -- I2C I/O Expander}\index{CH422G driver}
\label{app:ext:ch422g}

The WCH CH422G is an I2C I/O expander offering eight bidirectional pins
(\texttt{IO0..IO7}) plus four output-only pins (\texttt{OC0..OC3}). Unlike a
conventional register-pointer expander, it is driven entirely through \emph{fixed,
function-specific} I2C command addresses -- \texttt{0x24} (write system config),
\texttt{0x23} (write OC outputs), \texttt{0x38} (write IO outputs) and \texttt{0x26}
(read IO inputs) -- so there are no address straps and \texttt{Setup} takes no
address: one CH422G per bus.

\texttt{ESP32S3.CH422G} sits on \texttt{ESP32S3.I2C} and follows the two-level
locking shape of the other expander/RTC drivers: \texttt{Acquire} a
\texttt{Session} to hold the chip across many operations, while the I2C host is
locked only per transaction. Direction is \emph{global} -- one
\texttt{IO\_Direction} bit makes all of \texttt{IO0..IO7} inputs or outputs;
\texttt{OC0..OC3} are output-only with a global \texttt{Push\_Pull}/%
\texttt{Open\_Drain} drive. Because the config/OC/IO-output registers cannot be
read back, the driver keeps a shadow initialised to the datasheet power-on
defaults (IO = inputs, OC = high, push-pull). The chip has no interrupt output,
hence no \texttt{.Interrupts} child. The controlled \texttt{Session} restricts this
driver to the embedded and full profiles.

\begin{lstlisting}
type IO_Value is mod 2 ** 8;          --  bit i = IOi
type OC_Value is mod 2 ** 4;          --  low nibble = OC0..OC3
type IO_Pin is range 0 .. 7;  type OC_Pin is range 0 .. 3;
type Pin_State    is (Low, High);
type IO_Direction is (Inputs, Outputs);        --  all of IO0..IO7
type OC_Drive     is (Push_Pull, Open_Drain);  --  all of OC0..OC3
type Status  is (OK, Bus_Error);
type Device  is limited private;
type Session is limited private;

procedure Setup (Dev : out Device; Sda, Scl : ESP32S3.GPIO.Pin_Id;
                 Host : ESP32S3.I2C.I2C_Host := ESP32S3.I2C.I2C0;
                 Clock_Hz : Positive := 400_000);
procedure Acquire (S : in out Session; Dev : Device);
procedure Release (S : in out Session);
function  Present (S : Session) return Boolean;   --  ACK probe of 0x24
procedure Configure (S : Session; IO_Dir : IO_Direction := Inputs;
                     OC_Mode : OC_Drive := Push_Pull; Result : out Status);
procedure Write_IO  (S : Session; Value : IO_Value; Result : out Status);
procedure Read_IO   (S : Session; Value : out IO_Value; Result : out Status);
procedure Write_OC  (S : Session; Value : OC_Value; Result : out Status);
\end{lstlisting}

The worked example sets the chip up on I2C0 (SDA=IO8, SCL=IO9), probes the config
command address, then reads the input pins once a second. It is deliberately
read-only and never drives a pin.

\begin{lstlisting}
package CH renames ESP32S3.CH422G;
use type CH.Status;
Dev : CH.Device;  S : CH.Session;  V : CH.IO_Value;  St : CH.Status;
begin
   CH.Setup (Dev, Sda => 8, Scl => 9);   --  I2C0, 400 kHz, no address
   CH.Acquire (S, Dev);                  --  held for the whole run

   Put_Line (if CH.Present (S) then "ACK (present)" else "no ACK");

   loop
      delay until Clock + Seconds (1);
      CH.Read_IO (S, V, St);             --  RD-IO (0x26)
      if St = CH.OK then
         Put ("IO inputs = 0x");  Put_Hex (Unsigned_32 (V), 2);  New_Line;
      end if;
   end loop;
\end{lstlisting}

\noindent\textbf{Example:} \texttt{examples/esp32s3\_ch422g}.

\section{TCA9555 -- I2C 16-bit I/O Expander}\index{TCA9555 driver}
\label{app:ext:tca9555}

The Texas Instruments TCA9555 is a 16-bit I2C GPIO expander: two 8-bit ports
(\texttt{P0} = pins 0..7, \texttt{P1} = pins 8..15), each pin independently an
input or output, with an optional input-polarity-inversion register and an
active-low open-drain \texttt{INT} output. The 7-bit address is
\texttt{Base\_Address} (\texttt{0x20}) plus the \texttt{A2/A1/A0} strap value
(0..7), so up to eight parts share one bus.

\texttt{ESP32S3.TCA9555} layers on \texttt{ESP32S3.I2C} and is the canonical
two-level-locking driver: a \texttt{Session} is an exclusive RAII hold on one
expander (so a per-pin read-modify-write is safe), while the I2C host is locked
only inside each transaction -- leaving the bus free for other devices between
your operations. The per-device locks are a fixed library-level array keyed by
(host, strap). The optional \texttt{INT} pin is wired through \texttt{Setup}'s
\texttt{Int\_Pin}, and the \texttt{ESP32S3.TCA9555.Interrupts} child arms it as a
pulled-up falling-edge input. The controlled \texttt{Session} restricts this
driver to the embedded and full profiles.

\begin{lstlisting}
Base_Address : constant ESP32S3.I2C.Slave_Address := 16#20#;
subtype Hardware_Address is Natural range 0 .. 7;
type Pin_Number is range 0 .. 15;     --  0..7 = Port 0, 8..15 = Port 1
type Port_Value is mod 2 ** 16;       --  bit i = Pin_Number i
type Direction is (Output, Input);    type Pin_State is (Low, High);
type Status is (OK, Bus_Error);
type Device is limited private;  type Session is limited private;

procedure Setup (Dev : out Device; Addr : Hardware_Address;
                 Sda, Scl : ESP32S3.GPIO.Pin_Id;
                 Int_Pin  : ESP32S3.GPIO.Optional_Pin := ESP32S3.GPIO.No_Pin;
                 Host : ESP32S3.I2C.I2C_Host := ESP32S3.I2C.I2C0;
                 Clock_Hz : Positive := 400_000);
procedure Acquire (S : in out Session; Dev : Device);
procedure Release (S : in out Session);
procedure Set_Directions (S : Session; Inputs : Port_Value; Result : out Status);
procedure Write_Port (S : Session; Value : Port_Value; Result : out Status);
procedure Write_Pin  (S : Session; Pin : Pin_Number; State : Pin_State;
                      Result : out Status);          --  read-modify-write
procedure Read_Port    (S : Session; Value : out Port_Value; Result : out Status);
procedure Read_Outputs (S : Session; Value : out Port_Value; Result : out Status);
procedure Set_Polarity (S : Session; Inverted : Port_Value; Result : out Status);
--  child: ESP32S3.TCA9555.Interrupts.Attach (Dev, Action) / Detach (Dev)
\end{lstlisting}

The example holds one \texttt{Session} for the whole test, forces every pin to an
input (so nothing fights the external wiring), then exercises the input port, the
output register round-trip, a per-pin RMW, and the polarity register.

\begin{lstlisting}
package GPX renames ESP32S3.TCA9555;
use type GPX.Status;  use type GPX.Port_Value;
Dev : GPX.Device;  S : GPX.Session;  St : GPX.Status;  V, Orig : GPX.Port_Value;
begin
   GPX.Setup (Dev, Addr => 0, Sda => 8, Scl => 7);
   GPX.Acquire (S, Dev);                           --  hold the expander

   GPX.Set_Directions (S, Inputs => 16#FFFF#, Result => St);  --  all inputs
   GPX.Read_Port (S, Orig, St);                    --  comms check

   GPX.Write_Port (S, 16#A55A#, St);               --  output register only
   GPX.Read_Outputs (S, V, St);                    --  V should read 0xA55A

   GPX.Write_Pin (S, 5, GPX.High, St);             --  per-pin RMW (Session-safe)
   GPX.Read_Outputs (S, V, St);

   GPX.Set_Polarity (S, 16#A55A#, St);
   GPX.Set_Polarity (S, 0, St);                    --  restore
   GPX.Release (S);
\end{lstlisting}

\noindent\textbf{Example:} \texttt{examples/esp32s3\_tca9555}.

\section{PCF85063A -- I2C Real-Time Clock}\index{PCF85063A driver}
\label{app:ext:pcf85063a}

The NXP PCF85063A is a small fixed-address (\texttt{Bus\_Address} = \texttt{0x51})
I2C real-time clock and calendar: seconds through years (stored BCD on the wire,
binary here), a programmable alarm matching any subset of
second/minute/hour/day/weekday, an oscillator-stop flag reporting whether time
survived a power loss, and an active-low open-drain \texttt{INT} line the alarm
asserts. Up to 400 kHz I2C.

\texttt{ESP32S3.PCF85063A} sits on \texttt{ESP32S3.I2C}. Unlike the expander
drivers it exposes no \texttt{Session} type to the caller -- each operation opens
its own short-lived controlled I2C session for one complete transaction and
auto-releases the host on scope exit, so concurrent callers serialise. Register
reads use the chip's auto-incrementing pointer (a 1-byte write sets it, a following
read streams from it). \texttt{Get\_Time} returns a \texttt{Valid} flag (clock
integrity); \texttt{Set\_Time} brackets its write with stop/start and clears the OS
flag. The optional \texttt{INT} pin is wired through \texttt{Setup} and armed by
\texttt{ESP32S3.PCF85063A.Interrupts.Attach} as a pulled-up falling-edge input. The
controlled session restricts this driver to the embedded and full profiles.

\begin{lstlisting}
Bus_Address : constant ESP32S3.I2C.Slave_Address := 16#51#;
subtype Year_Number is Natural range 2000 .. 2099;   --  also Month/Day/Hour/...
type Weekday is (Sunday, Monday, Tuesday, Wednesday, Thursday, Friday, Saturday);
type Time is record
   Year : Year_Number; Month : Month_Number; Day : Day_Number;
   Day_Of_Week : Weekday; Hour : Hour_Number;
   Minute : Minute_Number; Second : Second_Number;
end record;
type Alarm is record   --  any subset of fields participates
   Use_Second : Boolean; Second : Second_Number;  Use_Minute  : Boolean; ...
   Use_Hour, Use_Day, Use_Weekday : Boolean; ...
end record;
type Status is (OK, Bus_Error);   type Device is limited private;

procedure Setup (Dev : out Device; Sda, Scl : ESP32S3.GPIO.Pin_Id;
                 Int_Pin : ESP32S3.GPIO.Optional_Pin := ESP32S3.GPIO.No_Pin;
                 Host : ESP32S3.I2C.I2C_Host := ESP32S3.I2C.I2C0;
                 Clock_Hz : Positive := 400_000);
procedure Get_Time (Dev : Device; T : out Time; Valid : out Boolean;
                    Result : out Status);
procedure Set_Time (Dev : Device; T : Time; Result : out Status);
procedure Reset (Dev : Device; Result : out Status);
procedure Set_Alarm (Dev : Device; A : Alarm; Result : out Status);
procedure Alarm_Triggered (Dev : Device; Fired : out Boolean; Result : out Status);
procedure Acknowledge_Alarm (Dev : Device; Result : out Status);   --  releases INT
--  child: ESP32S3.PCF85063A.Interrupts.Attach (Dev, Action) / Detach (Dev)
\end{lstlisting}

The example probes, resets, loads a known calendar, arms a seconds-match alarm five
seconds out, then watches the clock tick. With no \texttt{INT} line wired
(\texttt{No\_Pin}) the alarm is found by polling \texttt{AF} over I2C;
\texttt{Attach} is then a no-op.

\begin{lstlisting}
package RTC renames ESP32S3.PCF85063A;
use type RTC.Status;
Dev : RTC.Device;  T : RTC.Time;  Valid : Boolean;  St : RTC.Status;
Initial : constant RTC.Time :=
  (Year => 2026, Month => 6, Day => 22, Day_Of_Week => RTC.Monday,
   Hour => 14, Minute => 30, Second => 0);
begin
   RTC.Setup (Dev, Sda => 8, Scl => 7, Int_Pin => ESP32S3.GPIO.No_Pin);
   RTC.Interrupts.Attach (Dev, Alarm_IRQ.Handler'Access);   --  no-op: No_Pin

   RTC.Get_Time (Dev, T, Valid, St);          --  probe (ACK?)
   RTC.Reset (Dev, St);
   RTC.Set_Time (Dev, Initial, St);           --  clears OS -> Valid = True
   RTC.Set_Alarm (Dev, (Use_Second => True, Second => 5, others => <>), St);

   for Tick in 1 .. 10 loop
      delay until Clock + Seconds (1);
      RTC.Get_Time (Dev, T, Valid, St);
      declare Fired : Boolean; begin
         RTC.Alarm_Triggered (Dev, Fired, St);
         if St = RTC.OK and then Fired then
            RTC.Acknowledge_Alarm (Dev, St);  --  releases INT
            exit;
         end if;
      end;
   end loop;
\end{lstlisting}

\noindent\textbf{Example:} \texttt{examples/esp32s3\_pcf85063a}.

\section{SHT41 -- I2C Temperature / Humidity Sensor}\index{SHT41 driver}
\label{app:ext:sht41}

The Sensirion SHT41 (SHT4x family) is a command-based I2C temperature and humidity
sensor with no registers: a measurement is ``write a 1-byte command, wait the
conversion time, read 6 bytes'' -- two CRC-8-protected 16-bit words (temperature
then humidity). The SHT41-AD1B answers at \texttt{Default\_Address} (\texttt{0x44});
other variants use \texttt{0x45}/\texttt{0x46}. There is no interrupt -- it is read
on request.

\texttt{ESP32S3.SHT41} sits on \texttt{ESP32S3.I2C} and shares the PCF85063A shape:
no caller-visible session, each operation opens a short-lived controlled I2C
session and auto-releases the host. \texttt{Measure} blocks for the conversion (per
\texttt{Precision}) while holding the bus and returns integer milli-units --
\texttt{Temperature} in milli-degrees Celsius, \texttt{Humidity} in milli-percent
RH -- so no float library is required; CRC failures surface as \texttt{CRC\_Error}.
The controlled session restricts this driver to the embedded and full profiles.

\begin{lstlisting}
Default_Address : constant ESP32S3.I2C.Slave_Address := 16#44#;
type Measurement is record
   Temperature : Integer := 0;   --  milli-degrees Celsius
   Humidity    : Integer := 0;   --  milli-percent RH (0 .. 100_000)
end record;
type Precision is (Low, Medium, High);          --  High ~8 ms, Low ~2 ms
type Status is (OK, Bus_Error, CRC_Error);
type Device is limited private;

procedure Setup (Dev : out Device; Sda, Scl : ESP32S3.GPIO.Pin_Id;
                 Address : ESP32S3.I2C.Slave_Address := Default_Address;
                 Host : ESP32S3.I2C.I2C_Host := ESP32S3.I2C.I2C0;
                 Clock_Hz : Positive := 400_000);
procedure Measure (Dev : Device; Value : out Measurement; Result : out Status;
                   Repeatability : Precision := High);
procedure Read_Serial_Number
  (Dev : Device; Serial : out Interfaces.Unsigned_32; Result : out Status);
procedure Reset (Dev : Device; Result : out Status);
\end{lstlisting}

The example reads the 32-bit serial number as a presence check, then takes a
high-precision measurement once a second.

\begin{lstlisting}
package SHT renames ESP32S3.SHT41;
use type SHT.Status;
Dev : SHT.Device;  St : SHT.Status;  SN : Interfaces.Unsigned_32;
begin
   SHT.Setup (Dev, Sda => 8, Scl => 7);

   SHT.Read_Serial_Number (Dev, SN, St);   --  doubles as a comms check
   Put ("serial : 0x");  Put_Hex (SN, 8);  New_Line;

   for Tick in 1 .. 15 loop
      delay until Clock + Seconds (1);
      declare M : SHT.Measurement; begin
         SHT.Measure (Dev, M, St);          --  blocks ~8 ms (High)
         exit when St /= SHT.OK;
         Put ("T=");  Put_Fixed (M.Temperature, 1000, 2);   --  m'C -> C
         Put (" C  RH=");  Put_Fixed (M.Humidity, 1000, 2);  Put_Line (" %");
      end;
   end loop;
\end{lstlisting}

\cnote{No vendor SDK is linked. The SHT4x's CRC-8 and the milli-unit fixed-point
conversion are done in Ada, so the whole sensor path is float-free -- unlike the
typical ESP-IDF driver that returns \texttt{float} degrees and pulls in soft-float.}

\noindent\textbf{Example:} \texttt{examples/esp32s3\_sht41}.

\section{QMI8658C -- I2C 6-Axis IMU}\index{QMI8658C driver}
\label{app:ext:qmi8658c}

The QST QMI8658C is a 6-axis IMU (3-axis accelerometer + 3-axis gyroscope) on I2C.
The \texttt{SA0}/\texttt{SDO} pin selects the 7-bit address:
\texttt{Address\_SA0\_High} (\texttt{0x6A}, the floating/power-on default via the
internal pull-up) or \texttt{Address\_SA0\_Low} (\texttt{0x6B}).
\texttt{WHO\_AM\_I} returns \texttt{Who\_Am\_I\_Value} (\texttt{0x05}). A
programmable \texttt{INT} line (INT1/INT2) is push-pull active-high by default.

\texttt{ESP32S3.QMI8658C} sits on \texttt{ESP32S3.I2C} with the PCF85063A shape (no
caller-visible session; short-lived auto-released I2C session per operation;
auto-incrementing register pointer set by \texttt{Configure}, run little-endian).
\texttt{Configure} sets the accelerometer and gyroscope full scale and a shared
output rate and enables both sensors; \texttt{Read\_Accelerometer}/%
\texttt{Read\_Gyroscope} return raw signed 16-bit counts in \texttt{Axes}, which
the \texttt{Accel\_LSB\_Per\_G}/\texttt{Gyro\_LSB\_Per\_DPS} accessors scale to
physical units. The optional \texttt{INT} pin is wired through \texttt{Setup} and
armed by \texttt{ESP32S3.QMI8658C.Interrupts.Attach} as a pulled-down rising-edge
input. The controlled session restricts this driver to the embedded and full
profiles.

\begin{lstlisting}
Address_SA0_High : constant ESP32S3.I2C.Slave_Address := 16#6A#;  --  default
Address_SA0_Low  : constant ESP32S3.I2C.Slave_Address := 16#6B#;
Who_Am_I_Value   : constant Interfaces.Unsigned_8     := 16#05#;
type Accel_Range is (Range_2G, Range_4G, Range_8G, Range_16G);
type Gyro_Range  is (Range_16DPS, ... Range_2048DPS);
type Output_Rate is (ODR_7520_Hz, ... ODR_235_Hz, ... ODR_29_Hz);
type Axes is record X, Y, Z : Interfaces.Integer_16 := 0; end record;
type Status is (OK, Bus_Error);   type Device is limited private;

procedure Setup (Dev : out Device; Sda, Scl : ESP32S3.GPIO.Pin_Id;
                 Int_Pin : ESP32S3.GPIO.Optional_Pin := ESP32S3.GPIO.No_Pin;
                 Address : ESP32S3.I2C.Slave_Address := Address_SA0_Low;
                 Host : ESP32S3.I2C.I2C_Host := ESP32S3.I2C.I2C0;
                 Clock_Hz : Positive := 400_000);
procedure Read_Who_Am_I (Dev : Device; Id : out Interfaces.Unsigned_8;
                         Result : out Status);
procedure Reset (Dev : Device; Result : out Status);     --  then wait ~15 ms
procedure Configure (Dev : in out Device; Accel : Accel_Range := Range_8G;
                     Gyro : Gyro_Range := Range_512DPS;
                     Rate : Output_Rate := ODR_235_Hz; Result : out Status);
procedure Read_Accelerometer (Dev : Device; A : out Axes; Result : out Status);
procedure Read_Gyroscope     (Dev : Device; G : out Axes; Result : out Status);
procedure Read_Temperature
  (Dev : Device; Raw : out Interfaces.Integer_16; Result : out Status);
function  Accel_LSB_Per_G  (Dev : Device) return Positive;
function  Gyro_LSB_Per_DPS (Dev : Device) return Positive;
\end{lstlisting}

The example probes both SA0 addresses for \texttt{WHO\_AM\_I} = \texttt{0x05},
resets and configures the part, then streams accel and gyro in milli-units using
the sensitivity accessors and integer math.

\begin{lstlisting}
package IMU renames ESP32S3.QMI8658C;
use type IMU.Status;
Dev : IMU.Device;  St : IMU.Status;  Id : Unsigned_8;  Found : Boolean := False;
Candidates : constant array (1 .. 2) of ESP32S3.I2C.Slave_Address :=
  (IMU.Address_SA0_Low, IMU.Address_SA0_High);
begin
   for I in Candidates'Range loop
      IMU.Setup (Dev, Sda => 8, Scl => 7, Address => Candidates (I));
      IMU.Read_Who_Am_I (Dev, Id, St);
      if St = IMU.OK and then Id = IMU.Who_Am_I_Value then
         Found := True;  exit;
      end if;
   end loop;

   IMU.Reset (Dev, St);
   delay until Clock + Milliseconds (15);          --  reset settle
   IMU.Configure (Dev, Accel => IMU.Range_8G, Gyro => IMU.Range_512DPS,
                  Rate => IMU.ODR_235_Hz, Result => St);

   declare
      A, G  : IMU.Axes;
      A_Lsb : constant Positive := IMU.Accel_LSB_Per_G (Dev);
   begin
      IMU.Read_Accelerometer (Dev, A, St);
      IMU.Read_Gyroscope (Dev, G, St);
      --  milli-g = Integer (A.X) * 1000 / A_Lsb
   end;
\end{lstlisting}

\noindent\textbf{Example:} \texttt{examples/esp32s3\_qmi8658c}.

\section{ES8311 -- Audio Codec}\index{ES8311 driver}
\label{app:ext:es8311}

The ES8311 is an Everest low-power \emph{mono} audio codec: a register-configured
DAC/ADC pair that the SoC controls over I2C while carrying PCM audio over I2S. The
driver in \texttt{ESP32S3.ES8311} brings it up for 16-bit output, running the
codec-vendor register-init sequence on I2C and configuring an I2S port as the
master that generates MCLK, BCLK (SCLK) and LRCK and shifts PCM out on the codec's
DSDIN line. The codec runs as an I2S slave clocked from the ESP's MCLK $= 256
\times$ sample-rate (4.096\,MHz at 16\,kHz); the register coefficients are
rate-independent for that $256\times$/16-bit configuration.

It sits on two on-chip drivers at once -- \texttt{ESP32S3.I2C} for control and
\texttt{ESP32S3.I2S} for data. \texttt{Setup} is a one-time, single-threaded
bring-up; pass an \texttt{Asdout} pin to also enable the ADC (microphone) capture
path. After setup, audio is owned through a limited, controlled \texttt{Output}
handle that holds the I2S port exclusively and releases it on scope exit, so it
requires a tasking runtime (embedded/full profile). \texttt{Play\_Continuous}
starts a self-looping DMA for gapless, click-free, zero-CPU playback;
\texttt{Capture} samples the ADC underneath it without disturbing playback or the
shared master clock.

\begin{lstlisting}
package ESP32S3.ES8311 is
   subtype Address is ESP32S3.I2C.Slave_Address;
   Default_Address : constant Address := 16#18#;

   procedure Setup
     (I2C_Bus     : ESP32S3.I2C.I2C_Host;
      Sda, Scl    : ESP32S3.GPIO.Pin_Id;
      Port        : ESP32S3.I2S.I2S_Port;
      Mclk, Sclk, Lrck, Dsdin : ESP32S3.GPIO.Pin_Id;
      Asdout      : ESP32S3.GPIO.Optional_Pin := ESP32S3.GPIO.No_Pin;
      Sample_Rate : Positive := 16_000;
      Volume      : Natural  := 70;     --  DAC volume, 0 .. 100 %
      Mic_Gain_Db : Natural  := 24;     --  ADC PGA gain, 0 .. 42 dB
      Addr        : Address  := Default_Address;
      Ok          : out Boolean);

   procedure Set_Volume (Percent : Natural; Ok : out Boolean);

   type Output is limited private;          --  exclusive, auto-releasing
   Not_Ready : exception;

   procedure Acquire (O : in out Output);
   procedure Play            (O : Output; Samples : System.Address; Length : Natural);
   procedure Play_Continuous (O : Output; Samples : System.Address; Length : Natural);
   procedure Capture         (O : Output; Samples : System.Address; Length : Natural);
   procedure Stop    (O : Output);
   procedure Release (O : in out Output);
end ESP32S3.ES8311;
\end{lstlisting}

The example plays a gapless 440\,Hz sine on the DAC and simultaneously captures the
ADC to estimate the tone heard back -- a loopback test. Control is on I2C (SDA=IO8,
SCL=IO7); audio on I2S (MCLK=IO1, SCLK=IO2, LRCK=IO4, DSDIN=IO5 out, ASDOUT=IO3 in).

\begin{lstlisting}
declare
   Ok        : Boolean;
   Audio     : ESP32S3.ES8311.Output;
begin
   ESP32S3.ES8311.Setup
     (I2C_Bus => ESP32S3.I2C.I2C0,
      Sda     => 8, Scl => 7,
      Port    => ESP32S3.I2S.I2S0,
      Mclk    => 1, Sclk => 2, Lrck => 4, Dsdin => 5,
      Asdout  => 3,                       --  codec ADC out -> our data in
      Sample_Rate => 16_000, Volume => 75, Mic_Gain_Db => 24, Ok => Ok);

   ESP32S3.ES8311.Acquire (Audio);        --  hold the audio port
   --  Self-looping DMA: replays Buf forever, click-free, zero CPU.
   ESP32S3.ES8311.Play_Continuous (Audio, Buf'Address, Buf_Bytes);
   loop                                   --  capture under the running tone
      ESP32S3.ES8311.Capture (Audio, Cap'Address, Cap_Bytes);
      Analyse (Peak, Est);                --  should read ~440 Hz
      delay until Clock + Milliseconds (1000);
   end loop;
end;
\end{lstlisting}

\noindent\textbf{Example:} \texttt{examples/esp32s3\_es8311}.

\section{ST7789 -- SPI TFT Display Controller}\index{ST7789 driver}
\label{app:ext:st7789}

The ST7789 (ST77xx family) drives a small colour TFT panel over 4-wire SPI. The
link is \emph{write-only}: CLK + MOSI, no MISO, so the controller cannot be read
back and there is no probe or status query. Three GPIOs are driven directly -- DC
(data/command), CS (chip select), and an optional RST (software reset is used if it
is \texttt{No\_Pin}). Pixels are 16-bit RGB565, MSB-first. \texttt{ESP32S3.ST7789}
sits on \texttt{ESP32S3.SPI} (mode 0, full-duplex master, default 40\,MHz); CS is
driven as a plain GPIO rather than routed to the peripheral.

Locking is two-level, like the on-chip expander drivers. A \texttt{Session} is an
exclusive, RAII hold on one display, taken with \texttt{Acquire} and auto-released
on scope exit; hold it across a whole sequence so no other task corrupts the
controller mid-stream. The SPI host itself is locked only inside each operation
(assert CS, transfer, deassert CS, release), so the bus is free between operations.
The per-display guards are a library-level array keyed by the CS pin, and the
controlled \texttt{Session} means embedded/full profiles only.

A companion child \texttt{ESP32S3.ST7789.Text} layers a 5$\times$7 bitmap font on
the parent's \texttt{Draw\_Bitmap} primitive: \texttt{Draw\_Char}/\texttt{Draw\_Text}
blit opaque 6$\times$8 cells (the panel is write-only, so there is no read-back for a
transparent overlay) at an integer \texttt{Scale}, with LF handling for newlines.

\begin{lstlisting}
package ESP32S3.ST7789 is
   type Color is mod 2 ** 16;                        --  RGB565, MSB-first
   function RGB (R, G, B : Natural) return Color;    --  each 0 .. 255
   Black, White, Red, Green, Blue : constant Color;  --  named colours
   type Color_Array is array (Natural range <>) of Color;
   type Rotation is (Rot_0, Rot_90, Rot_180, Rot_270);

   type Device  is limited private;
   type Session is limited private;
   Not_Initialized, Not_Owned : exception;

   procedure Setup
     (Dev                : out Device;
      Sclk, Mosi, DC, CS : ESP32S3.GPIO.Pin_Id;
      Width, Height      : Positive := 240;
      RST                : ESP32S3.GPIO.Optional_Pin := ESP32S3.GPIO.No_Pin;
      Clock_Hz           : Positive := 40_000_000);

   procedure Acquire (S : in out Session; Dev : Device);
   procedure Release (S : in out Session);

   procedure Init         (S : Session);             --  reset + power-on sequence
   procedure Set_Rotation (S : Session; Rot : Rotation);
   procedure Fill         (S : Session; C : Color);
   procedure Fill_Rect    (S : Session; X, Y, W, H : Natural; C : Color);
   procedure Set_Pixel    (S : Session; X, Y : Natural; C : Color);
   procedure Draw_Bitmap  (S : Session; X, Y, W, H : Natural; Pixels : Color_Array);
end ESP32S3.ST7789;
\end{lstlisting}

The example drives a 240$\times$240 panel on SPI2 (SCLK=IO12, MOSI=IO13, DC=IO16,
CS=IO10; backlight on IO6 is driven by the example, not the driver). One
\texttt{Session} is held for the whole demo while each fill locks the bus only for
its own transfer.

\begin{lstlisting}
declare
   Dev : ESP32S3.ST7789.Device;
   S   : ESP32S3.ST7789.Session;
begin
   ESP32S3.ST7789.Setup (Dev, Sclk => 12, Mosi => 13, DC => 16, CS => 10);
   ESP32S3.ST7789.Acquire (S, Dev);            --  protect the display
   ESP32S3.ST7789.Init (S);
   ESP32S3.ST7789.Fill (S, ESP32S3.ST7789.Red);
   ESP32S3.ST7789.Fill_Rect (S, X => 70, Y => 70, W => 100, H => 100,
                             C => ESP32S3.ST7789.White);
   ESP32S3.ST7789.Text.Draw_Text
     (S, X => 6, Y => 16, Str => "ESP32-S3 + Ada",
      FG => ESP32S3.ST7789.White, BG => ESP32S3.ST7789.RGB (0, 0, 32));
   ESP32S3.ST7789.Release (S);
end;                                            --  Session auto-released
\end{lstlisting}

\noindent\textbf{Example:} \texttt{examples/esp32s3\_st7789}.

\section{TLV2556 -- SPI 12-bit ADC}\index{TLV2556 driver}\index{ADC driver}
\label{app:ext:tlv2556}

The TLV2556 is a Texas Instruments 12-bit, 11-channel, 200-kSPS
successive-approximation ADC with an internal reference, on an SPI-compatible
serial interface (mode 0). It is the off-chip complement to the SoC's own SAR ADC
(in the HAL chapter): more channels, a stable internal reference, and -- because
it lives on the SPI bus -- usable even when the internal ADC's pins are committed
elsewhere.

Two things make it a useful worked example. First, it \emph{shares the SPI2 bus
with the W25Q flash}, through the SPI driver's built-in software chip-select: you
name the select GPIO as \texttt{CS\_Pin} and the driver drives it -- active-low,
held across the whole conversion -- and suppresses the hardware CS0 for the hold,
so it cannot disturb the flash. (Contrast the ST7789 above, which toggles its CS
as a plain GPIO by hand.) Second, the converter is \emph{pipelined}: each 16-clock
frame shifts an 8-bit command in (channel select plus configuration) while the
\emph{previous} conversion's result shifts out, MSB-first. \texttt{Read} hides
that -- it primes the conversion, waits out the (max $\sim$5.5\,\textmu s)
conversion time, then reads the result back.

\begin{lstlisting}
package ESP32S3.TLV2556 is
   type Sample       is range 0 .. 4095;              --  12-bit unipolar result
   type Analog_Input is (AIN0, AIN1, AIN2, AIN3, AIN4, AIN5, AIN6, AIN7,
                         AIN8, AIN9, AIN10,            --  the 11 analog inputs ...
                         Test_Half, Test_Zero, Test_Full);  --  ... and 3 self-tests
   type Reference    is (Internal_4096mV, Internal_2048mV, External);

   type Device is record                              --  describes itself on the bus
      Host     : ESP32S3.SPI.SPI_Host;
      Clock_Hz : Positive                  := 8_000_000;
      CS_Pin   : ESP32S3.GPIO.Optional_Pin := ESP32S3.GPIO.No_Pin;  -- driver-driven CS
      CS_CB    : ESP32S3.SPI.CS_Select     := null;   --  or a callback (decoder/expander)
      Ctx      : System.Address            := System.Null_Address;
   end record;

   procedure Initialize (Dev : Device; Ref : Reference := External);
   function  Read       (Dev : Device; Input : Analog_Input) return Sample;
   function  Millivolts (S : Sample; Ref : Reference) return Natural;
end ESP32S3.TLV2556;
\end{lstlisting}

The standout feature for bring-up is the \emph{self-test inputs}: three of the
\texttt{Analog\_Input} values are not pins at all but internal voltages
\emph{ratiometric to the reference rails}, so \texttt{Test\_Zero} reads 0,
\texttt{Test\_Half} reads 2048, and \texttt{Test\_Full} reads 4095 -- regardless
of the reference voltage or any analog wiring. They are the ADC analog of the
W25Q's JEDEC-ID check: a reference- and wiring-independent proof that the SPI link,
the framing, and the converter all work, before you trust a real reading.

\begin{lstlisting}
declare
   package ADC renames ESP32S3.TLV2556;
   Dev : ADC.Device :=                                --  shares SPI2; CS on IO12
     (Host => ESP32S3.SPI.SPI2, Clock_Hz => 8_000_000, CS_Pin => 12, others => <>);
begin
   ESP32S3.SPI.Setup (ESP32S3.SPI.SPI2);              --  bus shared with the flash
   ESP32S3.SPI.Configure_Pins (ESP32S3.SPI.SPI2, Sclk => 1, Mosi => 4, Miso => 45);
   ADC.Initialize (Dev, Ref => ADC.External);

   --  Bring-up proof: the ratiometric self-tests, independent of reference/wiring.
   Put_Line ("zero=" & ADC.Read (Dev, ADC.Test_Zero)'Image    -- ~ 0
           & " half=" & ADC.Read (Dev, ADC.Test_Half)'Image   -- ~ 2048
           & " full=" & ADC.Read (Dev, ADC.Test_Full)'Image); -- ~ 4095

   Value := ADC.Read (Dev, ADC.AIN0);                 --  then a real channel
end;
\end{lstlisting}

On the hardware (the self-tests carry a couple of LSB of analog slack):

\begin{lstlisting}[style=sh]
[tlv2556] self-test: zero=0 half=2047 full=4095   PASS
[tlv2556] AIN0 = 195 / 4095
\end{lstlisting}

\noindent\textbf{Example:} \texttt{examples/esp32s3\_tlv2556}.

\section{GPS -- NMEA-0183 Receiver}\index{GPS driver}
\label{app:ext:gps}

The GPS driver targets a UART GNSS module (developed against the Quectel L76K) that
streams NMEA-0183 sentences. Unlike the passive I2C device drivers you poll through
a handle, \texttt{ESP32S3.GPS} is a singleton background \emph{service}: a
library-level task owns one UART, continuously reads and decodes the sentence
stream, and publishes results into a protected store. The application has no
\texttt{Device} handle -- it just reads the store. It sits on \texttt{ESP32S3.UART};
only \texttt{Rx} is required (wired to the module's TXD), with optional \texttt{Tx}
(for configuration) and \texttt{Pps} (a 1PPS GPIO interrupt for time alignment).

Every published value is written and read under the store's lock, so a reader never
sees a half-updated value; in particular \texttt{Latitude}/\texttt{Longitude} are
one \texttt{Position} record updated atomically. Each reading carries the
\texttt{Updated\_At} time at which a \emph{valid} sentence last refreshed it --
compare \texttt{Age} against your tolerance to detect staleness. The driver decodes
GGA, RMC, GLL, VTG, GSV, GSA and ZDA into fix, position, time, date, velocity,
signal and satellite-in-view readings. The task plus protected objects mean
embedded/full profiles only.

\cnote{Latitude and longitude follow the u-blox convention: signed integers in
$10^{-7}$ degrees (\texttt{Integer\_32}), not floats -- exact, atomically paired,
and free of any soft-float dependency in the bare runtime.}

A child \texttt{ESP32S3.GPS.L76K} adds the Quectel PCAS vendor commands
(constellation selection, baud, update rate, restart) sent verbatim through
\texttt{Send}; only \texttt{Set\_Constellation} (PCAS04) is hardware-tested.
\texttt{Inject} feeds one framed sentence straight into the decoder for
host/on-target self-test with no live receiver.

\begin{lstlisting}
package ESP32S3.GPS is
   type Position is record                       --  1e-7 degrees, +N/-S +E/-W
      Latitude  : Interfaces.Integer_32 := 0;
      Longitude : Interfaces.Integer_32 := 0;
   end record;
   type Fix_Quality is (No_Fix, GPS_Fix, DGPS_Fix);
   type Fix_Type    is (Fix_None, Fix_2D, Fix_3D);

   type Position_Reading is record               --  value + freshness + valid
      Value : Position; Updated_At : Ada.Real_Time.Time; Valid : Boolean := False;
   end record;
   --  Fix_Reading, Time_Reading, Date_Reading, Velocity_Reading, Signal_Reading ...

   procedure Setup
     (Port : ESP32S3.UART.UART_Port;
      Rx   : ESP32S3.GPIO.Optional_Pin;
      Tx   : ESP32S3.GPIO.Optional_Pin := ESP32S3.GPIO.No_Pin;
      Pps  : ESP32S3.GPIO.Optional_Pin := ESP32S3.GPIO.No_Pin;
      Baud : ESP32S3.UART.Baud_Rate    := 9600);

   function Current_Position return Position_Reading;
   function Current_Fix      return Fix_Reading;
   function Current_Time     return Time_Reading;
   function Age (Updated_At : Ada.Real_Time.Time) return Ada.Real_Time.Time_Span;
   procedure Inject (Sentence : String; Accepted : out Boolean);
end ESP32S3.GPS;
\end{lstlisting}

The example self-tests the decoder by injecting canned sentences (including a
bad-checksum reject) before \texttt{Setup}, then brings up UART0 (Rx=IO44, Tx=IO43,
9600 baud) and prints the latest fix and its age once a second.

\begin{lstlisting}
GPS.Inject ("$GPGGA,123519,4807.038,N,01131.000,E,1,08,0.9,545.4,M,46.9,M,,*47", Ok);
declare
   P : constant GPS.Position_Reading := GPS.Current_Position;
begin
   pragma Assert (P.Valid and then P.Value.Latitude = 481_173_000);
end;

GPS.Setup (Port => ESP32S3.UART.UART0, Rx => 44, Tx => 43, Baud => 9_600);
loop
   delay until Clock + Seconds (1);
   declare
      P : constant GPS.Position_Reading := GPS.Current_Position;
      Fresh : constant Boolean :=
        P.Valid and then To_Duration (GPS.Age (P.Updated_At)) < 3.0;
   begin
      null;   --  print P.Value when Fresh, else show acquisition progress
   end;
end loop;
\end{lstlisting}

\noindent\textbf{Example:} \texttt{examples/esp32s3\_gps}.

\section{W25Q -- SPI NOR Flash}\index{W25Q driver}\index{SPI flash}
\label{app:ext:w25q}

The W25Q256FV is a 32\,MB (256\,Mbit) SPI NOR flash -- the on-board non-volatile
store the local filesystem of Chapter~\ref{ch:storage} lives on. It sits on SPI2,
sharing the bus with the W5500 (and, in the ADC example, the TLV2556) through the
SPI driver's software chip-select. The driver is deliberately small: identify, read,
erase by sector, program by page -- the primitives a wear-levelling layer and a
filesystem build on.

\begin{lstlisting}
package ESP32S3.W25Q is
   Page_Size   : constant := 256;        --  program granularity (one Page-Program)
   Sector_Size : constant := 4096;       --  smallest erase unit

   type JEDEC_ID is record
      Manufacturer : Unsigned_8;         --  0xEF for Winbond
      Memory_Type  : Unsigned_8;
      Capacity     : Unsigned_8;         --  log2 of the size in bytes
   end record;

   type Flash is record                  --  describes itself on the bus
      Host     : ESP32S3.SPI.SPI_Host;
      Clock_Hz : Positive                  := 8_000_000;
      CS_Pin   : ESP32S3.GPIO.Optional_Pin := ESP32S3.GPIO.No_Pin;  -- driver-driven CS
      CS_CB    : ESP32S3.SPI.CS_Select     := null;   --  or a callback
      Ctx      : System.Address            := System.Null_Address;
   end record;

   procedure Read_Identification (Dev : Flash; ID : out JEDEC_ID);
   function  Capacity_Bytes      (ID : JEDEC_ID) return Address;
   procedure Initialize          (Dev : Flash; OK : out Boolean);

   procedure Read         (Dev : Flash; Addr : Address; Data : out Byte_Array);
   procedure Erase_Sector (Dev : Flash; Addr : Address);
   procedure Program_Page (Dev : Flash; Addr : Address; Data : Byte_Array)
     with Pre => Data'Length in 1 .. Page_Size
                 and then Natural (Addr mod Page_Size) + Data'Length <= Page_Size;
end ESP32S3.W25Q;
\end{lstlisting}

The bring-up proof is the \emph{JEDEC-ID check}: \texttt{Read\_Identification}
returns \texttt{EF 40 19} on a healthy W25Q256 (Winbond, 256\,Mbit) -- a
wiring-independent confirmation that the SPI link and framing work before any read
or write is trusted, the flash analogue of the TLV2556's self-test inputs
(\S\ref{app:ext:tlv2556}). The page/sector asymmetry is the usual NOR contract: you
may program any aligned 256-byte page but can only erase a whole 4\,KB sector, and
the \texttt{Program\_Page} precondition enforces the page rule at the call site.
Above this driver, \texttt{Block\_Dev.WL} adds wear levelling and \texttt{Ext4.FS} a
filesystem (Chapter~\ref{ch:storage}); from the LISP REPL, \texttt{(spi-xfer ...)}
reads the same JEDEC id by hand (Chapter~\ref{ch:lisp}).

\noindent\textbf{Example:} \texttt{examples/esp32s3\_w25q}.

\section{SD Card (SD\_SPI) -- Removable Storage}\index{SD card driver}
\label{app:ext:sdspi}

\texttt{ESP32S3.SD\_SPI} talks to an SD/SDHC card in SPI mode, layering the SD
command protocol (CMD0/8/58, ACMD41, CMD17/24, CRC7) on top of the task-safe
\texttt{ESP32S3.SPI} master. Four lines plus power are routed (MOSI/MISO/SCLK and a
CS driven as a plain GPIO, so it can stay asserted across a whole
command/response/data sequence); a 10\,k pull-up on MISO is recommended. Cards
initialise at $\leq$400\,kHz then run at \texttt{Data\_Clock\_Hz}. The API is always
512-byte LBA blocks; SDHC/SDXC block-addressing vs.\ older byte-addressing is
handled internally. Every operation takes the SPI Session for the whole transaction
(so concurrent callers serialise), which means embedded/full profiles only.

\emph{(Covered in full in the Mass Storage chapter, \S\ref{ch:storage}.)} This
entry is just the bus-level driver beneath that chapter's filesystem stack.

\begin{lstlisting}
package ESP32S3.SD_SPI is
   type Block         is array (0 .. 511) of Interfaces.Unsigned_8;
   type Block_Address is new Interfaces.Unsigned_32;
   type Card_Kind is (Unknown, SD_V1, SD_V2_SC, SD_V2_HC);
   type Status    is (OK, No_Card, Unusable, Init_Timeout, Read_Error, Write_Error);
   type Card is limited private;

   procedure Setup (C : out Card; Host : ESP32S3.SPI.SPI_Host;
                    Sclk, Mosi, Miso, Cs : ESP32S3.GPIO.Pin_Id;
                    Init_Clock_Hz : Positive := 400_000;
                    Data_Clock_Hz : Positive := 8_000_000);
   procedure Initialize  (C : in out Card; Result : out Status);
   function  Kind        (C : Card) return Card_Kind;
   procedure Read_Block  (C : in out Card; LBA : Block_Address;
                          Data : out Block; Result : out Status);
   procedure Write_Block (C : in out Card; LBA : Block_Address;
                          Data : Block; Result : out Status);
end ESP32S3.SD_SPI;
\end{lstlisting}

The example does a non-destructive round-trip on a scratch sector on SPI2
(SCLK=IO12, MOSI=IO11, MISO=IO13, CS=IO10): read it, write the same bytes back,
re-read, and verify they match.

\begin{lstlisting}
ESP32S3.SD_SPI.Setup (C, Host => ESP32S3.SPI.SPI2,
                      Sclk => 12, Mosi => 11, Miso => 13, Cs => 10);
ESP32S3.SD_SPI.Initialize (C, St);            --  Kind (C) now reports the card type
if St = ESP32S3.SD_SPI.OK then
   ESP32S3.SD_SPI.Read_Block  (C, Test_LBA, Orig, St);
   ESP32S3.SD_SPI.Write_Block (C, Test_LBA, Orig, St);   --  same bytes back
   ESP32S3.SD_SPI.Read_Block  (C, Test_LBA, Back, St);
   pragma Assert (St = ESP32S3.SD_SPI.OK and then Orig = Back);
end if;
\end{lstlisting}

\noindent\textbf{Example:} \texttt{examples/esp32s3\_sd\_spi}.

\section{TX1812 -- Addressable RGB LED String}\index{TX1812 driver}
\label{app:ext:tx1812}

The TX1812 is a single-wire, daisy-chainable RGB LED in the WS2812 (``NeoPixel'')
family: 24 bits of colour are clocked in MSB-first as a precisely timed pulse train
(a `1' is long-high/short-low, a `0' short-high/long-low), and a $>$80\,\textmu s
low ``reset'' latches the chain. \texttt{ESP32S3.TX1812} generates that waveform
with the \texttt{ESP32S3.RMT} peripheral, one RMT symbol per data bit. A
\texttt{Strip} is a claimed handle that takes an RMT transmit channel on
\texttt{Acquire} and releases it on scope exit (its RMT channel component is
controlled).

A \texttt{Strip} is statically sized by its \texttt{Count} discriminant: it carries
both the \texttt{Count}-pixel colour buffer and the \texttt{Count}$\times$24
RMT-symbol frame buffer, so declaring e.g.\ \texttt{Panel : Strip (64)} reserves all
of it at elaboration with no heap, and the linker verifies it fits. Set colours into
the buffer with \texttt{Set}/\texttt{Set\_All}, then \texttt{Show} clocks them out;
longer strings are streamed via RMT wrap re-fill, so the default of one RMT block
handles any length.

\begin{lstlisting}
package ESP32S3.TX1812 is
   type Color is record R, G, B : Unsigned_8 := 0; end record;
   Off, Red, Green, Blue, White : constant Color;

   --  Statically sized: holds the pixel buffer AND the Count*24 RMT frame.
   type Strip (Count : Positive) is limited private;

   procedure Acquire (S       : in out Strip;
                      Pin     : ESP32S3.GPIO.Pin_Id;
                      Channel : ESP32S3.RMT.TX_Index := 0;
                      Blocks  : Positive := 1);
   function  Is_Valid (S : Strip) return Boolean;
   procedure Release  (S : in out Strip);
   procedure Set      (S : in out Strip; Index : Positive; C : Color);
   procedure Set_All  (S : in out Strip; C : Color);
   procedure Show     (S : in out Strip);   --  clock the buffer out and latch
end ESP32S3.TX1812;
\end{lstlisting}

The example drives a 64-LED string on IO48 (RMT channel 0), declared once at
library level so its $\sim$6.4\,KiB of buffers are reserved at elaboration. With no
physical string attached, the board's on-board LED on IO48 is pixel 1 of the chain,
so it cycles red\,$\rightarrow$\,green\,$\rightarrow$\,blue\,$\rightarrow$\,white\,$\rightarrow$\,off.

\begin{lstlisting}
--  LED_Panel.Panel : ESP32S3.TX1812.Strip (64);   (library-level)
LED.Acquire (LED_Panel.Panel, Pin => 48, Channel => 0);
if not LED.Is_Valid (LED_Panel.Panel) then
   loop delay until Clock + Seconds (3600); end loop;   --  channel busy
end if;
loop
   for I in Colors'Range loop
      LED.Set_All (LED_Panel.Panel, Colors (I));   --  all 64 pixels
      LED.Show    (LED_Panel.Panel);               --  stream the whole frame
      delay until Clock + Milliseconds (600);
   end loop;
end loop;
\end{lstlisting}

\noindent\textbf{Example:} \texttt{examples/esp32s3\_tx1812}.

\section{W5500 -- SPI Ethernet Controller}\index{W5500 driver}
\label{app:ext:w5500}

The WIZnet W5500 is an Ethernet controller with a hardwired TCP/IP stack (ARP,
ICMP, IP, TCP, UDP in silicon), eight hardware sockets and 32\,KB of buffer RAM,
driven over SPI. On this board it sits on \texttt{ESP32S3.SPI} (SPI2) with
\texttt{MISO=IO45}, \texttt{MOSI=IO4}, \texttt{SCLK=IO1}, \texttt{SCSn=IO39},
\texttt{INTn=IO3} and \texttt{RSTn=IO11}; the chip-select is an ordinary GPIO held
low across each variable-length SPI frame.

Unlike the other external parts, the W5500 is a whole subsystem rather than a single
driver: a transport layer (\texttt{ESP32S3.W5500}), a socket engine
(\texttt{ESP32S3.W5500.Sockets}), optional interrupts, a DHCP client, and an adapter
(\texttt{ESP32S3.W5500.Net\_Device}) that presents the chip as a standard
\texttt{GNAT.Sockets} interface. Because that is far more than an API digest, it has
its own chapter -- see chapter~\ref{ch:networking}, ``Networking: GNAT.Sockets over
the W5500,'' for the driver layers, the chip-neutral \texttt{Net\_Devices}
interface, the DNS and DHCP helpers, and the worked examples (echo server, HTTP,
NTP, DNS, DHCP, and a weather client). Like the other controlled-handle drivers it
requires the embedded or full profile.

\noindent\textbf{Examples:} \texttt{examples/esp32s3\_w5500*}.

\chapter{Glossary}
\label{ch:glossary}

A quick reference to the acronyms and terms this book leans on. Where a term has
a chapter that treats it properly, the entry points there.

\begin{description}[style=nextline, leftmargin=0pt]

\item[ACATS] Ada Conformity Assessment Test Suite -- the official conformance
  tests for an Ada implementation. This project runs them on real silicon
  (Chapter~\ref{ch:acats-sweep}).

\item[ACK polling] Probing an I2C device's address until it answers. An EEPROM stops acknowledging while it programs a page, so the master polls instead of sleeping the worst case (Chapter~\ref{ch:eeprom}).

\item[ADU] \emph{Application Data Unit} -- a complete Modbus message as it travels: the \texttt{MBAP} header (TCP) or address+CRC (RTU) wrapping a \texttt{PDU} (Chapter~\ref{ch:modbus}).

\item[Alire] The Ada package manager (\texttt{alr}). It supplies the
  cross-toolchain and pins the runtime crate; no other SDK is needed.

\item[APB] \emph{Advanced Peripheral Bus} -- the lower-speed on-chip bus many ESP32-S3 peripherals are clocked from ($\sim$80\,MHz), divided down per peripheral.

\item[ARP] \emph{Address Resolution Protocol} -- resolves an IPv4 address to a \texttt{MAC} address on a local link; part of the W5500's hardwired stack.

\item[\texttt{bb-runtimes}] AdaCore's ``bareboard runtimes'' generator. The Ada
  runtime here is a board port grown from it (Chapter~\ref{ch:runtime-build}).

\item[BCD] \emph{Binary-Coded Decimal} -- each decimal digit in its own 4-bit nibble; the on-the-wire format of some RTC and sensor registers.

\item[\texttt{.bss}] The zero-initialised data segment; cleared by \texttt{start.S}
  at boot.

\item[CAN] \emph{Controller Area Network} -- a multi-drop automotive/industrial serial bus. The ESP32-S3's controller for it is the \texttt{TWAI} peripheral.

\item[CCOUNT / CCOMPARE] Xtensa cycle-counter and compare registers. The runtime
  uses them to drive the clock tick and timed delays.

\item[ceiling priority] The locking policy for protected objects: a caller is
  raised to the object's ceiling priority for the duration, bounding blocking and
  preventing priority inversion.

\item[closure] A subprogram bundled with the enclosing-scope variables it uses but
  does not declare -- the state it ``closes over.'' A nested subprogram that
  captures a local needs a stack \emph{trampoline}, which faults on this chip, so
  the HAL builds under \texttt{No\_Implicit\_Dynamic\_Code} and uses closure-free
  callbacks or tagged dispatch instead (\S\ref{driver:plugin}).

\item[context switch] Saving one task's CPU state and restoring another's. The
  Xtensa version (windows, the coprocessor, \texttt{THREADPTR}) is
  Chapter~\ref{ch:context-switch}.

\item[CRL] \emph{Certificate Revocation List} -- a CA-signed list of certificates revoked before expiry; an alternative to \texttt{OCSP} (Chapter~\ref{ch:tls}).

\item[CSPRNG] \emph{Cryptographically-Secure Pseudo-Random Number Generator} -- a \texttt{PRNG} whose output is unpredictable enough for keys and nonces.

\item[DAG] \emph{Directed Acyclic Graph} -- a graph with no cycles; e.g.\ the recursion-free call graph the stack analyser walks for a worst-case bound (Chapter~\ref{ch:static-mem}).

\item[DDR] \emph{Double Data Rate} -- transferring on both clock edges to double throughput, as in octal (\texttt{OPI}) PSRAM/flash timing.

\item[din tuning] Calibrating the read-data (\emph{data-in}) sampling phase --- the \texttt{din\_mode} knob --- that an octal PSRAM controller needs at 80\,MHz DDR. Its correct value is not knowable from source (it varies by chip, controller and temperature) and a wrong one hangs the CPU on a cache read with no recovery, so the bootloader measures it at boot with a bounded SPI1 sweep and takes the centre of the passing window (Chapter~\ref{ch:psram}).

\item[DRAM / IRAM] The data-bus and instruction-bus \emph{views} of the same
  internal SRAM. Code runs from the IRAM view; data lives in the DRAM view
  (\S\ref{ch:bootloader}).

\item[DROM / IROM] The data and instruction \emph{flash} windows, reached
  execute-/read-in-place through the cache MMU.

\item[ECC] \emph{Elliptic-Curve Cryptography} -- public-key cryptography over elliptic curves; smaller keys than RSA for equal strength (Chapter~\ref{ch:tls}).

\item[ECDH] \emph{Elliptic-Curve Diffie--Hellman} -- ECC key agreement; \texttt{ECDHE} is its ephemeral, forward-secret form (Chapter~\ref{ch:tls}).

\item[EEPROM] \emph{Electrically Erasable Programmable Read-Only Memory} -- small, byte-addressable non-volatile memory needing no erase before a write. The 24C family speaks I2C (Chapter~\ref{ch:eeprom}).

\item[ELF] \emph{Executable and Linkable Format} -- the object/executable file format GNAT emits, converted to a flash image for the board.

\item[FFS / FLS] \emph{Find-First-Set / Find-Last-Set} -- bit-scan operations the \texttt{TLSF} allocator uses to pick a size class in $O(1)$ (Chapter~\ref{ch:heap}).

\item[FIPS] \emph{Federal Information Processing Standards} -- US government security standards; some hardened servers are FIPS-configured and reject non-NIST curves (Chapter~\ref{ch:tls}).

\item[FSM] \emph{Finite-State Machine} -- a controller with a fixed set of states and transitions; several peripherals embed one in hardware.

\item[FTL] \emph{Flash Translation Layer} -- the layer that hides NOR flash's erase-before-write and limited endurance behind a plain read/write block interface, with wear levelling (Chapter~\ref{ch:storage}).

\item[GHASH] The keyed Galois-field hash that computes the authentication tag in AES-GCM (Galois/Counter Mode); the Ada half of the hybrid record cipher (Chapter~\ref{ch:tls}).

\item[GNARL] The GNU Ada Runtime Library -- the tasking half of the runtime
  (tasks, protected objects, the scheduler), archived as \texttt{libgnarl.a}.

\item[GNAT] The GNU Ada compiler. Its sequential runtime is \texttt{libgnat.a}.

\item[GNSS] \emph{Global Navigation Satellite System} -- the umbrella term for GPS, Galileo, BeiDou and the like.

\item[GPR / \texttt{gprbuild}] GNAT project files and their builder. Each crate
  (\texttt{esp32s3\_rts}, \texttt{esp32s3\_hal}) is a GPR project.

\item[HAL] Hardware Abstraction Layer -- the \texttt{ESP32S3.*} peripheral
  drivers (Chapters~\ref{ch:hal}--\ref{ch:hal3}).

\item[HKDF] \emph{HMAC-based Key Derivation Function} -- expands a shared secret into keying material; the heart of the TLS\,1.3 key schedule (Chapter~\ref{ch:tls}).

\item[ICMP] \emph{Internet Control Message Protocol} -- IP's diagnostic/error sublayer (ping lives here); part of the W5500's hardwired stack.

\item[IDF] \emph{(ESP-IDF) Espressif IoT Development Framework} -- the vendor's C SDK and HAL; this project deliberately does without it (Chapter~\ref{ch:encap}).

\item[IMU] \emph{Inertial Measurement Unit} -- an accelerometer plus gyroscope (sometimes magnetometer) motion sensor.

\item[IPI] Inter-Processor Interrupt -- the poke one core sends another to wake a
  task or rebalance, central to the SMP scheduler.

\item[ISRG] \emph{Internet Security Research Group} -- the nonprofit behind Let's Encrypt; its \emph{ISRG Root X1} CA is pinned by the HTTPS example (Chapter~\ref{ch:tls}).

\item[JEDEC] The semiconductor standards body; here, the JEDEC identifier every SPI-NOR flash returns (manufacturer + device id), read by the W25Q driver (Chapter~\ref{ch:storage}).

\item[Jorvik] A tasking profile (Ada RM~D.13) -- a less restrictive successor to
  Ravenscar that still bans the unbounded constructs. Both \texttt{light-tasking}
  and \texttt{embedded} build under it.

\item[LBA] \emph{Logical Block Addressing} -- a flat sector index into a block device; the unit the filesystem reads and writes (Chapter~\ref{ch:storage}).

\item[LRU] \emph{Least Recently Used} -- the eviction policy of the ext4 write-back block cache (Chapter~\ref{ch:storage}).

\item[LSB / MSB] \emph{Least-} / \emph{Most-Significant Bit} (or \emph{Byte}) -- which end of a value or transfer comes first; e.g.\ RGB565 pixels are sent MSB-first.

\item[MBAP] \emph{Modbus Application Protocol} header -- the seven-byte TCP framing (transaction, protocol, length, unit id) prefixing each \texttt{PDU} (Chapter~\ref{ch:modbus}).

\item[MCU] \emph{Microcontroller Unit} -- a CPU integrated with memory and peripherals; the ESP32-S3 itself.

\item[MMU] Memory Management Unit. Here it is the cache MMU that maps 64\,KB flash
  pages (and PSRAM) into the CPU's address windows.

\item[MPI] \emph{Multi-Precision Integer} -- an arbitrary-width integer; the workhorse of RSA and elliptic-curve big-number arithmetic (Chapter~\ref{ch:tls}).

\item[MSPI] The memory SPI controller that talks to flash and PSRAM; its timing
  is what makes PSRAM bring-up delicate (\S\ref{ch:bootloader}).

\item[NAT] \emph{Network Address Translation} -- rewriting addresses/ports at a gateway so many hosts share one public IP; the connect-out patterns here are NAT-friendly.

\item[NIST] \emph{National Institute of Standards and Technology} -- publishes the curves (P-256) and cipher/hash standards the TLS stack implements (Chapter~\ref{ch:tls}).

\item[NMEA] \emph{National Marine Electronics Association} -- its 0183 sentence format (e.g.\ \texttt{RMC}, \texttt{GGA}) is how GPS modules report a fix.

\item[OCSP] \emph{Online Certificate Status Protocol} -- a live query for whether a certificate has been revoked; an alternative to a \texttt{CRL} (Chapter~\ref{ch:tls}).

\item[OUI] \emph{Organizationally Unique Identifier} -- the top 24 bits of a \texttt{MAC} address, assigned per manufacturer.

\item[overloading] Several subprograms (or operators) sharing one name, told apart
  by their parameter and result \emph{profile}; the compiler resolves each call to
  the one whose profile matches. No keyword is involved -- distinct from
  \emph{overriding}. (E.g.\ \texttt{Send} overloaded on the CAN frame type,
  Chapter~\ref{ch:hal2}.)

\item[overriding] Replacing a tagged type's inherited primitive operation with a
  new body. The optional \texttt{overriding} keyword makes the compiler confirm the
  declaration really does replace an inherited operation -- catching a typo that
  would otherwise add a new one instead. The basis of dispatching, and the shape of
  the book's pluggable backends (override the handlers you support; the framework
  calls them by dispatch). Explained in Chapter~\ref{ch:ada-primer}; cf.\
  \emph{overloading}.

\item[PDM] Pulse-Density Modulation -- the 1-bit serial format some microphones
  and ADCs emit; the I2S driver converts it to/from PCM.

\item[PDU] \emph{Protocol Data Unit} -- the protocol-specific payload; for Modbus, a function code plus its data, carried inside an \texttt{ADU} (Chapter~\ref{ch:modbus}).

\item[PHY] The \emph{physical-layer} transceiver -- the front-end that turns link signalling into bits, e.g.\ the Ethernet PHY inside the W5500.

\item[PKCS] \emph{Public-Key Cryptography Standards} -- the RSA-era format family, e.g.\ PKCS\#1 signatures and PKCS\#8 keys (Chapter~\ref{ch:tls}).

\item[PLL] \emph{Phase-Locked Loop} -- the on-chip multiplier that synthesises the high CPU/peripheral clocks (e.g.\ 240\,MHz) from the crystal.

\item[PRNG] \emph{Pseudo-Random Number Generator} -- a deterministic generator of random-looking values; cf.\ \texttt{CSPRNG}.

\item[protected object] An Ada monitor: data plus operations with built-in mutual
  exclusion and condition synchronisation (\S\ref{ch:encap}).

\item[PSRAM] Pseudo-static external RAM, mapped at \texttt{16\#3D00\_0000\#} by the
  bootloader. Not DMA-capable.

\item[PSS] \emph{Probabilistic Signature Scheme} -- the modern randomised RSA signature padding (RSASSA-PSS), used for the TLS CertificateVerify (Chapter~\ref{ch:tls}).

\item[RAII] ``Resource Acquisition Is Initialization'' -- tying a resource's
  lifetime to an object's scope. In Ada it is a \texttt{controlled} type whose
  \texttt{Finalize} releases the resource (\S\ref{ch:encap}).

\item[Ravenscar] The original restrictive tasking profile for high-integrity
  real-time systems; Jorvik's stricter predecessor.

\item[rendezvous] Direct task-to-task synchronisation via \texttt{entry}/
  \texttt{accept}. Available only in the \texttt{full} profile.

\item[restriction profile] A named bundle of \texttt{pragma Restrictions} (e.g.
  Jorvik) that the compiler enforces, trading language features for analysability
  and footprint (Chapter~\ref{ch:profiles}).

\item[repeated START] A second I2C START issued without an intervening STOP, so the slave sees one command rather than two. It is how a master writes a register address and then reads that register (\texttt{ESP32S3.I2C.Write\_Read}, Section~\ref{driver:end-opcode}).

\item[RFC] \emph{Request For Comments} -- the IETF document series specifying the internet protocols (TCP, TLS, DHCP, FTP\,\dots).

\item[RTS] Run-Time System -- the compiled runtime an application links against;
  here, one of the three profile archives.

\item[SAN] \emph{Subject Alternative Name} -- the X.509 extension listing the hostnames a certificate is valid for; the leaf's SAN must cover the host (Chapter~\ref{ch:tls}).

\item[secondary stack] Where Ada returns unconstrained results (e.g. a
  \texttt{String} of run-time-determined length). Absent in \texttt{light-tasking}.

\item[SMP] Symmetric Multi-Processing -- both ESP32-S3 cores under one scheduler,
  with tasks pinned per core and woken across cores by IPI.

\item[SVD / \texttt{svd2ada}] A machine-readable register description (System View
  Description) and the AdaCore tool that turns it into typed Ada register packages
  (\S\ref{ch:runtime-build}).

\item[TDM] \emph{Time-Division Multiplexing} -- interleaving several channels on one line by time slot; the multi-slot I2S audio mode (Chapter~\ref{ch:hal2}).

\item[\texttt{THREADPTR}] An Xtensa per-thread pointer register, saved/restored on
  context switch so per-task thread-local data works.

\item[trampoline] A scrap of executable code the compiler writes onto the stack to
  carry a closure's captured state. The stack is non-executable here, so a
  trampoline faults; \texttt{No\_Implicit\_Dynamic\_Code} forbids generating one
  (\S\ref{driver:plugin}, Chapter~\ref{ch:debug}).

\item[VECBASE] The Xtensa register holding the base address of the exception/
  interrupt vector table; \texttt{start.S} points it at our vectors.

\item[watchdog (WDT)] A timer that resets the chip unless fed. The bootloader
  disables the boot watchdogs; the RTC super-watchdog is re-disarmed after a
  deep-sleep wake.

\item[window (register window)] The Xtensa calling convention rotates a window of
  physical registers per call instead of spilling to the stack; mismanaged
  windows are a recurring source of boot and context-switch bugs.

\item[XIP] Execute-In-Place -- running code directly from flash through the cache
  rather than copying it to RAM.

\item[ZCX] Zero-Cost Exceptions -- the exception model in \texttt{embedded} and
  \texttt{full}: no overhead unless an exception is actually raised.

\item[ZFP] ``Zero FootPrint'' -- a runtime (or coding style) with no tasking,
  exceptions or secondary stack; the 2nd-stage bootloader is written this way.

\end{description}

\chapter{Further Reading}
\label{ch:references}

Pointers for going deeper, grouped by topic. Online resources were current at the
time of writing.

\section{The Ada language and its profiles}

\begin{itemize}
\item \textbf{Ada Reference Manual} -- the standard itself.
  \url{http://ada-auth.org/standards/ada22.html}. The most relevant corners for
  this book: \emph{D.13} (the Ravenscar and Jorvik profiles), \emph{D.2.3}
  (priority and dispatching policies), \emph{C.6} (\texttt{Volatile} and
  \texttt{Atomic}), \emph{13.3}/\emph{13.5} (address and record-representation
  clauses), and \emph{7.6} / \texttt{Ada.Finalization} (controlled types).
\item \textbf{John Barnes, \emph{Programming in Ada 2012}} -- the standard
  textbook for the language as a whole.
\item \textbf{AdaCore learn.adacore.com} -- free courses, including
  \emph{Introduction to Embedded Systems Programming} and the
  \emph{Ada for the C++/Java developer} guide. \url{https://learn.adacore.com}.
\end{itemize}

\section{The runtime and the toolchain}

\begin{itemize}
\item \textbf{\texttt{bb-runtimes}} -- AdaCore's bareboard runtime generator that
  this project's runtime is ported from.
  \url{https://github.com/AdaCore/bb-runtimes}. The ESP32-S3 board lives in the
  fork \url{https://github.com/rowsail/bb-runtimes} (branch
  \texttt{xtensa-esp32s3-port}).
\item \textbf{\texttt{svd2ada}} -- generates the typed register packages from an
  SVD. \url{https://github.com/AdaCore/svd2ada}.
\item \textbf{Alire} -- the Ada package manager that supplies the toolchain.
  \url{https://alire.ada.dev}.
\item \textbf{GNAT GPL / GNAT Pro documentation} -- for \texttt{pragma
  Restrictions}, the bareboard runtime structure, and \texttt{gprbuild}.
\end{itemize}

\section{The silicon}

\begin{itemize}
\item \textbf{ESP32-S3 Technical Reference Manual} -- the authoritative register
  reference; every peripheral chapter in this book traces back to it.
  \url{https://www.espressif.com/sites/default/files/documentation/esp32-s3_technical_reference_manual_en.pdf}.
\item \textbf{ESP32-S3 Datasheet} -- pin-out, electrical characteristics, the
  GPIO matrix.
\item \textbf{Espressif SVD files} -- the register descriptions consumed by
  \texttt{svd2ada}. \url{https://github.com/espressif/svd}.
\item \textbf{Xtensa Instruction Set Architecture (LX7)} and the
  \emph{Windowed Register Option} -- background for the context-switch and boot
  chapters (Cadence/Tensilica documentation).
\end{itemize}

\section{Cryptography and TLS}

\begin{itemize}
\item \textbf{SPARKNaCl} -- Roderick Chapman's SPARK\,2014 port of TweetNaCl,
  formally proved free of run-time errors. It supplies the public-key, hashing and
  HMAC/HKDF primitives the pure-Ada TLS client stands on (X25519, the SHA family,
  Ed25519); the entire HTTPS chapter (Chapter~\ref{ch:tls}) rests on it. BSD-3-Clause
  (\textcopyright~2021 Protean Code Limited), vendored under \texttt{crates/sparknacl}.
  \url{https://github.com/rod-chapman/SPARKNaCl}. \emph{Our thanks to Rod Chapman and
  the SPARKNaCl project.}
\item \textbf{NaCl / TweetNaCl} -- the originals SPARKNaCl follows: Daniel~J.
  Bernstein, Tanja Lange, Peter Schwabe and colleagues' compact, misuse-resistant
  cryptography library. \url{https://nacl.cr.yp.to} and
  \url{https://tweetnacl.cr.yp.to}.
\item \textbf{RFC 8446} -- \emph{The Transport Layer Security (TLS) Protocol Version
  1.3}, the specification the TLS chapter implements.
  \url{https://www.rfc-editor.org/rfc/rfc8446}.
\end{itemize}

\section{Conformance and high-integrity Ada}

\begin{itemize}
\item \textbf{ACATS} -- the Ada Conformity Assessment Test Suite.
  \url{http://www.ada-auth.org/acats.html}.
\item \textbf{The Ravenscar Profile for Real-Time Systems} (Burns, Dobbing,
  Vardanega) -- the rationale behind the restricted tasking profiles.
\end{itemize}

\section{This project}

\begin{itemize}
\item \textbf{Source repository} --
  \url{https://github.com/rowsail/ada_esp32s3}. The
  \texttt{QUICKSTART.md} and \texttt{TOOLING.md} there cover setup and the
  developer tooling; the \LaTeX{} source of this book is under \texttt{book/}.
\end{itemize}


\printindex

\end{document}
